% !TEX program = xelatex
\documentclass[12pt,a4paper,fontset=fandol]{ctexbook}

% ==================== 页面 ====================
\usepackage[top=2.5cm,bottom=2.5cm,left=2.5cm,right=2.5cm]{geometry}

% ==================== 字体 ====================
\usepackage{fontspec}
\usepackage{iffont}
% CJK 钉 fandol（两侧都取自 TeX 树里的同一份 FandolSong-Regular.otf）；
% mono 钉 DejaVu Sans Mono：CI 由 fontconfig 命中同名，本机没有该系统字体，
% 就用 MiKTeX dejavu 包里的同一版本 TTF（kpathsea 按文件名命中）。
\IfFontExistsTF{DejaVu Sans Mono}
  {\setmonofont{DejaVu Sans Mono}[Scale=MatchLowercase]}
  {\setmonofont[Extension=.ttf,
      UprightFont=DejaVuSansMono,
      BoldFont=DejaVuSansMono-Bold,
      ItalicFont=DejaVuSansMono-Oblique,
      BoldItalicFont=DejaVuSansMono-BoldOblique]
     {DejaVuSansMono}[Scale=MatchLowercase]}

% ==================== 颜色 ====================
\usepackage[dvipsnames,svgnames,x11names]{xcolor}
\definecolor{codebg}{HTML}{F5F5F5}
\definecolor{codeframe}{HTML}{E0E0E0}
\definecolor{cppkeyword}{HTML}{1A1AFF}
\definecolor{cppcomment}{HTML}{2E8B2E}
\definecolor{cppstring}{HTML}{B22222}
\definecolor{linenumgray}{HTML}{999999}
\definecolor{clheadbg}{HTML}{1E2761}
\definecolor{elfblue}{HTML}{2F6DB5}
\definecolor{pegolden}{HTML}{8A6D1F}
\definecolor{machogray}{HTML}{4A5568}
\definecolor{okgreen}{HTML}{1A7F37}
\definecolor{errred}{HTML}{B62324}

% ==================== listings ====================
\usepackage{listings}

% ---- 所有代码块共用的排版（先定义，各语言 style 组合它）----
\lstdefinestyle{codebase}{
  basicstyle=\small\ttfamily,
  keywordstyle=\color{cppkeyword}\bfseries,
  commentstyle=\color{cppcomment}\itshape,
  stringstyle=\color{cppstring},
  numberstyle=\tiny\color{linenumgray},
  numbers=left,
  frame=single,
  rulecolor=\color{codeframe},
  framesep=6pt,
  breaklines=true,
  tabsize=4,
  columns=flexible,
  keepspaces=true,
  backgroundcolor=\color{codebg},
  showstringspaces=false,
  escapeinside={(*@}{@*)},
}

\lstdefinestyle{cppstyle}{
  style=codebase,
  language=C++,
  morekeywords={constexpr,consteval,constinit,decltype,auto,
    concept,requires,co_await,co_yield,co_return,
    noexcept,override,final,nullptr,static_assert,
    template,typename,namespace,using,mutable,volatile,
    explicit,operator,friend,inline,virtual,
    thread_local,alignas,alignof,extern},
  deletekeywords={int,char,bool,float,double,long,short,unsigned,signed,void,wchar_t},
  emph={size_t,int32_t,uint32_t,int64_t,uint64_t,ssize_t,ptrdiff_t,
    uintptr_t,intptr_t,wchar_t,string,vector,multimap,
    unique_ptr,shared_ptr,optional,string_view,span},
  emphstyle=\color{teal!80!black},
}
\lstset{style=cppstyle}

% ---- C：ABI 边界的类型、属性与内建函数 ----
\lstdefinestyle{cstyle}{
  style=codebase,
  language=C,
  morekeywords={bool,FILE,restrict,inline,thread_local},
  emph={size_t,int32_t,uint32_t,int64_t,uint64_t,ssize_t,ptrdiff_t,
    uintptr_t,intptr_t,wchar_t,va_list,__attribute__,__declspec,
    __builtin_expect,_dl_sym,DL_FUNC},
  emphstyle=\color{teal!80!black},
}

% ---- 本册自建 shell 方言：编译器、汇编器、链接器与二进制工具命令行 ----
% 注意：邻册的 ShellPM / ShellLD / ShellUB 只在各自文件内定义，不能借用。
\lstdefinelanguage{ShellCL}{
  morekeywords={if,then,else,elif,fi,for,in,do,done,while,case,esac,
    function,return,local,export,source,echo,cd,set,exec,trap,read,
    printf,test,exit,shift,eval,alias},
  sensitive=true,
  morecomment=[l]{\#},
  morestring=[b]",
  morestring=[d]',
}
\lstdefinestyle{shellstyle}{
  style=codebase,
  language=ShellCL,
  keywordstyle=\color{cppkeyword}\bfseries,
  emph={gcc,g++,clang,clang++,cl,link,ld,lld,lld-link,gold,ar,ranlib,
    lib,mt,dlltool,windres,cmake,ninja,make,ccache,
    as,gas,nasm,yasm,llvm-mc,llvm-objdump,llvm-readelf,llvm-nm,
    objdump,nm,readelf,strip,objcopy,size,strings,addr2line,cxxfilt,
    dumpbin,editbin,mt,link,dumpbin,otool,vtool,lipo,dyldinfo,dyld,
    vmmap,codesign,dsymutil,plutil,gdb,lldb,valgrind,
    grep,sed,awk,xargs,jq,python3,pip,pwsh,bash,git,tar,mkdir,touch,
    ls,cat,wc,diff,find,sort,uniq,tee,rm,cp,mv,head,tail,cut,hexdump,
    xxd,od,dd,file,ldd,ulimit,taskset,time},
  emphstyle=\color{teal!80!black}\bfseries,
}

% ---- 本册自建 CMake 方言：门禁拒收裸 language=CMake ----
\lstdefinelanguage{CMakeCL}{
  morekeywords={cmake_minimum_required,project,add_executable,
    add_library,target_link_libraries,target_include_directories,
    target_compile_definitions,target_compile_options,set,if,elseif,
    else,endif,find_package,option,message,install,file,
    target_sources,add_subdirectory,enable_language,set_target_properties,
    check_cxx_compiler_flag,target_link_options,list,string},
  sensitive=false,
  morecomment=[l]{\#},
  morestring=[b]",
}
\lstdefinestyle{cmakestyle}{
  style=codebase,
  language=CMakeCL,
  keywordstyle=\color{cppkeyword}\bfseries,
}

% ---- 汇编与反汇编：AT&T、MASM 与 ARM64 三种写法共用一套方言 ----
% listings 的关键词默认只含字母，.text / .L 这类伪指令必须以 alsoletter={.} 放行。
\lstdefinelanguage{CLAsm}{
  alsoletter={.},
  morekeywords={.text,.data,.rodata,.bss,.section,.globl,.global,.local,
    .type,.size,.align,.p2align,.byte,.word,.long,.quad,.string,.asciz,
    .file,.ident,.weak,.hidden,.set,.equ,.macro,.endm,.if,.else,.endif,
    .rept,.endr,.cfi_startproc,.cfi_endproc,.cfi_def_cfa_offset,
    .extern,.comm,.lcomm,.symname,.tag,.endstruct,
    mov,movq,movl,movzx,movsx,lea,add,sub,and,or,xor,not,shl,shr,sar,
    cmp,test,jmp,je,jne,jl,jg,jle,jge,ja,jb,call,ret,push,pop,
    setne,setl,cmovne,cs,rep,bswap,cqto,cltd,endbr64,leave,nop,
    ldp,stp,bl,br,cbz,cbnz,adrp,adr,mul,sdiv,ldr,str,sxtw,uxtb,rev,
    movn,msr,mrs,dcbz,svc},
  sensitive=true,
  morecomment=[l]{\#},
  morecomment=[l]{;},
  morecomment=[n]{/*}{*/},
  morestring=[b]",
}
\lstdefinestyle{asmstyle}{
  style=codebase,
  language=CLAsm,
  keywordstyle=\color{elfblue}\bfseries,
}

% ---- 工具输出：readelf、objdump、dumpbin、otool 的报告（无语言，纯文本）----
\lstdefinestyle{outstyle}{
  style=codebase,
  language={},
  numbers=none,
  frame=single,
  morecomment=[l][\color{machogray}\bfseries]{\#},
  emph={ELF,PE,PE32,PE32plus,Mach-O,LOAD,TEXT,DATA,
    UND,GLOBAL,WEAK,LOCAL,BIND,RELA,text,data,bss,rodata,
    FOO,bar,main,_ZN,start,SECTION,HEADER},
  emphstyle=\color{elfblue}\bfseries,
}

% ==================== tcolorbox ====================
\usepackage[most]{tcolorbox}

\newtcolorbox{guideline}[1][]{
  enhanced,breakable,
  colback=blue!5!white,colframe=blue!75!black,
  fonttitle=\bfseries,title=要点,#1,
  boxed title style={colback=blue!75!black},
  attach boxed title to top left={yshift=-2mm,xshift=2mm},
  arc=2mm,boxrule=0.8mm,
}
\newtcolorbox{compare}[1][]{
  enhanced,breakable,
  colback=green!3!white,colframe=green!50!black,
  fonttitle=\bfseries,title=对比,#1,
  boxed title style={colback=green!50!black},
  attach boxed title to top left={yshift=-2mm,xshift=2mm},
  arc=2mm,boxrule=0.8mm,
}
\newtcolorbox{warning}[1][]{
  enhanced,breakable,
  colback=red!5!white,colframe=red!75!black,
  fonttitle=\bfseries,title=常见陷阱,#1,
  boxed title style={colback=red!75!black},
  attach boxed title to top left={yshift=-2mm,xshift=2mm},
  arc=2mm,boxrule=0.8mm,
}
\newtcolorbox{note}[1][]{
  enhanced,breakable,
  colback=orange!5!white,colframe=orange!80!black,
  fonttitle=\bfseries,title=注意,#1,
  boxed title style={colback=orange!80!black},
  attach boxed title to top left={yshift=-2mm,xshift=2mm},
  arc=2mm,boxrule=0.8mm,
}
\newtcolorbox{bestpractice}[1][]{
  enhanced,breakable,
  colback=purple!5!white,colframe=purple!60!black,
  fonttitle=\bfseries,title=最佳实践,#1,
  boxed title style={colback=purple!60!black},
  attach boxed title to top left={yshift=-2mm,xshift=2mm},
  arc=2mm,boxrule=0.8mm,
}
% 本册专用：流水线阶段框与平台差异框
\newtcolorbox{flow}[1][]{
  enhanced,breakable,
  colback=elfblue!4!white,colframe=elfblue,
  fonttitle=\bfseries,title=流水线阶段,#1,
  boxed title style={colback=elfblue},
  attach boxed title to top left={yshift=-2mm,xshift=2mm},
  arc=2mm,boxrule=0.8mm,
}
\newtcolorbox{abibox}[1][]{
  enhanced,breakable,
  colback=pegolden!5!white,colframe=pegolden,
  fonttitle=\bfseries,title=平台差异,#1,
  boxed title style={colback=pegolden},
  attach boxed title to top left={yshift=-2mm,xshift=2mm},
  arc=2mm,boxrule=0.8mm,
}

% ==================== 章节标题 ====================
\usepackage{titlesec}
\usepackage{fancyhdr}
\usepackage{graphicx}
\usepackage{amssymb}
\usepackage{amsmath}
\usepackage{tabularx}
\usepackage{booktabs}
\usepackage{longtable}
\usepackage{array}
\usepackage{multirow}
\usepackage{enumitem}
\usepackage{seqsplit}
\usepackage{float}

\titleformat{\chapter}[display]
  {\normalfont\huge\bfseries\color{clheadbg}}
  {\CTEXthechapter}{20pt}{\Huge}
\titleformat{\section}
  {\normalfont\Large\bfseries\color{blue!70!black}}
  {\thesection}{1em}{}
\titleformat{\subsection}
  {\normalfont\large\bfseries\color{blue!60!black}}
  {\thesubsection}{1em}{}

\pagestyle{fancy}
\fancyhf{}
\fancyhead[LE,RO]{\thepage}
\fancyhead[RE]{\leftmark}
\fancyhead[LO]{\rightmark}
\setlength{\headheight}{14.5pt}

\setlist[itemize]{leftmargin=2em,itemsep=2pt}
\setlist[enumerate]{leftmargin=2em,itemsep=2pt}
\emergencystretch=6em

% ==================== 超链接 ====================
\usepackage{hyperref}
\hypersetup{
  colorlinks=true,
  linkcolor=blue!70!black,
  urlcolor=blue!60!black,
  bookmarks=true,
  pdfauthor={C++ Reference},
  pdftitle={C/C++ 编译器与链接器参考手册：从源码到可执行文件与平台 ABI},
}

% ==================== 自定义命令 ====================
\newcommand{\cpp}[1]{C++#1}
\newcommand{\Fn}[1]{\texttt{#1()}}
\newcommand{\Tp}[1]{\texttt{#1}}
\newcommand{\Opt}[1]{\texttt{-\/-#1}}
\newcommand{\Bk}{\discretionary{}{}{}}
\newcommand{\Lang}[1]{\textsf{#1}}
\newcommand{\CL}{\Lang{C}}
\newcommand{\CPP}{\Lang{C++}}
% 工具、开关值与产物：本册统一这四个写法，\Tool 是工具名，\Sw 是单个开关
\newcommand{\Tool}[1]{\texttt{#1}}
\newcommand{\Sw}[1]{\texttt{#1}}
\newcommand{\Sect}[1]{\texttt{.#1}}
\newcommand{\Sym}[1]{\texttt{#1}}
\newcommand{\Mx}[1]{\textsf{#1}}
% 三种目标文件格式与两类平台代号：只用文字标签，避免正文里同义多形
\newcommand{\Elf}{\textbf{\textsf{ELF}}}
\newcommand{\PeFmt}{\textbf{\textsf{PE}}}
\newcommand{\Macho}{\textbf{\textsf{Mach-O}}}
\newcommand{\Nix}{\textsf{Linux}}
\newcommand{\Win}{\textsf{Windows}}
\newcommand{\Mac}{\textsf{macOS}}
\newcommand{\xSixtyFour}{\textsf{x86-64}}
\newcommand{\ArmSixtyFour}{\textsf{AArch64}}

% ====================================================================
% book 双页模式默认 \flushbottom，会靠拉伸行距把每页撑满；本仓库框体与长表多，
% 撑页产生的 Underfull \vbox 是排版噪声而非缺陷，这里统一改用小底部对齐。
\raggedbottom

\begin{document}

\frontmatter

\begin{titlepage}
\centering
\vspace*{4cm}
{\Huge\bfseries C/C++ 编译器与链接器参考手册\par}
\vspace{1cm}
{\LARGE 从源码到可执行文件\par}
\vspace{0.4cm}
{\Large 前端、中间表示与优化、目标文件与链接、\\ 加载执行与平台 ABI\par}
\vfill
{\large \today\par}
\end{titlepage}

\tableofcontents

\mainmatter

\cleardoublepage

% ====================================================================
% 骨架冒烟章：把本册全部 listing 方言、框体与自定义宏各用一次，
% 保证 preamble 里的每一项都在编译路径上。正式章节合并时整章替换掉。
% ====================================================================% 章节正文由 frag_chNN.tex 合并插入到下面这行之前（merge_frags.py）
\part{前端}
\chapter{流水线形状与实现的载体}\label{ch:cl01-pipeline}

书名里有“编译器”，也有“链接器”，而两者之间那条缝——一份源码究竟经过几个阶段才变成镜像——
是本书其余章节的地基。本章先给出一条确定的形状：从源文件到可执行镜像，至少经过预处理、
编译、汇编、链接四个可分离的步骤，之后还有加载与运行期初始化；再处理一件更根本的事：
为什么几乎所有语言的实现都坐在 \CL{} 或 \CPP{} 之上。

先把结论摆在前面。语言层面的“自举”与实现层面的“载体”是两件事：一门语言完全可以由它自己
写出编译器，而那份编译器源码仍然要用别人的 \CL{} 编译器编译、仍然要链接一个 \CL{} 写的运行时。
把这两件事混在一起，是讨论工具链时最常见的语义事故。

\section{阶段图：从源文件到镜像}\label{sec:cl01-stages}

一条典型的 \CPP{} 流水线不是“一个程序”，而是若干程序串起来的序列，每一段都可以被单独调用，
每一段都留下一个可读、可写、可比较的产物。形状如下。

\begin{flow}[title={从源文件到可执行镜像的阶段序列}]
\begin{enumerate}[leftmargin=2.2em]
  \item 入口：一个 \CPP{} 翻译单元的起点（源文件与其包含闭包）；进入下一阶段前，
    它必须已经被展开成一份完整文本。
  \item 预处理：把源文件的包含闭包展开成一份完整文本，同时求值条件编译、做宏替换。
    机制本身在第 \ref{ch:cl02-preprocessing} 章展开。
  \item 编译（前端）：词法与语法分析、语义检查、名字查找、模板实例化、名字修饰，产出语法树。
  \item 编译（中端）：把语法树转换成中间表示，跑与目标机器无关的优化。
    中间表示的形状与优化序列见第 \ref{ch:cl04-ir} 与第 \ref{ch:cl05-passes} 章。
  \item 编译（后端）：指令选择、寻址模式匹配、寄存器分配、调度，产出目标机汇编文本。
    这一段的取舍见第 \ref{ch:cl06-codegen} 章。
  \item 汇编：汇编文本换成目标文件，符号引用此时仍是未解决的占位。
  \item 链接：符号解析与重定位，把目标文件与库拼成一个镜像，见第 \ref{ch:cl03-resolution}、
    第 \ref{ch:cl07-object} 与第 \ref{ch:cl08-static-link} 章。
  \item 加载与运行期初始化：加载器解析依赖、做重定位、铺开线程局部存储、跑全局构造，
    最后把控制交给语言入口。动态这一半在第 \ref{ch:cl09-dynamic} 章。
\end{enumerate}
\end{flow}

这条序列里最重要的不是阶段的名字，而是下面这句话：\emph{每个阶段都有可查产物}。
预处理后的文本、汇编文本、目标文件、静态库、镜像，全都是磁盘上的字节，能被打印、被比较、
被单独重做。所以编译器的行为从来不是黑箱——当你怀疑某一步时，把那一步的产物落到磁盘上、
用只读工具看它，这比读源码或猜默认值可靠得多。本章最后一节给出现成的命令形状；
想把这些动作放进编辑器里做，见第 \ref{ch:cl14-ide} 章与附录 \ref{ap:clA-cli}。

把这张表当诊断坐标系用，四条最常见症状各自定位到不同阶段。
\begin{itemize}
  \item 症状是“符号找不到”：问题在链接期，去看目标文件的符号表，而不是去改前端代码。
  \item 症状是“布局不对”：问题在类型布局与名字修饰，去看 \CPP{} 的修饰规则与平台 ABI。
  \item 症状是“只在 \Win{} 上崩”：问题在阶段八（加载与初始化），去看镜像的依赖与节布局。
  \item 症状是“同样的源码，换个优化档位结果不同”：问题在中端优化序列或后端指令选择，
    先落盘汇编文本再比较。
\end{itemize}

\section{一次驱动与分步驱动}\label{sec:cl01-drivers}

日常构建用的是“一次编译驱动”：把源文件交给一个命令，让它把整条流水线跑到底，中间产物
留在实现自己的临时目录里，用完即弃。它的优点是简洁，缺点是你没有介入点：临时文件的命名、
存活时间、格式都由实现决定，出错时看不到任何东西。

“分步驱动”把同一条流水线的每一段变成显式产物：预处理单独一步、编译单独一步、汇编单独一步、
链接单独一步。它的代价是构建脚本变长、要自己维护产物之间的依赖关系；它的收益是每一段都能
被存档、被比较、被送进另一个工具读。本书讨论机制时一律用分步视角，因为机制只存在于缝里。

\begin{compare}[title={两种驱动方式的实际差别}]
\begin{itemize}
  \item 产物位置：一次驱动把中间产物丢给临时目录；分步驱动每个产物都有一个你能命名的路径。
  \item 可复现性：分步驱动能把一个可疑的预处理文本存档，多年后仍可用同一份输入复现。
  \item 并行度：分步驱动的每一步都是天然的任务边界，构建系统按文件粒度并行；
    一次驱动只能按翻译单元粒度并行。
  \item 工具复用：想拿汇编文本喂给另一个汇编器、或把目标文件塞进归档，都需要分步产物。
  \item 风险：分步驱动里“产物属于哪种语言”要靠扩展名或显式声明告诉实现，
    说错了就会走错前端；一次驱动则由源文件名自己推断。
\end{itemize}
\end{compare}

把阶段切开的“停止在某一步”的能力，\Tool{gcc} 与 \Tool{clang} 这一族驱动都有对应开关：
只预处理、只出汇编文本、只出目标文件、以及最后一步交给链接器。具体写法见下例；
\Tool{cl} 与 \Tool{link} 这一族的等价能力，开关名以其文档为准，本书不凭印象给。

\begin{lstlisting}[style=shellstyle,language=ShellCL,caption={典型分步流程：每步留一个显式产物（示意）}]
# 0) 源文件 widget.cpp；下面的阶段切分以 gcc/clang 族的驱动为例
# 1) 预处理：得到展开后的翻译单元文本
g++ -std=c++20 -E widget.cpp -o widget.ii

# 2) 编译：把预处理文本降到汇编文本（中间表示此时不落盘）
g++ -std=c++20 -S widget.ii -o widget.s

# 3) 汇编：汇编文本换成目标文件
g++ -std=c++20 -c widget.s -o widget.o

# 4) 归档：把目标文件收进静态库，供链接期按需抽取成员
ar rcs libwidget.a widget.o

# 5) 链接：解析符号、做重定位，产出镜像
g++ -std=c++20 main.o -L. -lwidget -o app

# 6) 共享库：同一批目标文件也能产出一个被加载器解析的动态库
#    位置无关代码的开关随实现不同，此处不写，名称以该实现文档为准
g++ -std=c++20 -shared widget.o -o libwidget.so
\end{lstlisting}

\begin{warning}[title={分步之后，“源文件名”这个信息就丢了}]
一旦你手写各步骤，实现只能靠输入文件的扩展名判断该走哪条前端。把 \CPP{} 的预处理产物命名成
\Sect{i}（那是 \CL{} 的预处理产物约定），会得到一个按 \CL{} 规则解析的翻译单元：
模板关键字被拒、名字修饰不会发生、闭包里那些只合法于 \CPP{} 的内容全部报错——
此时源文件的包含展开早已完成，报错点与真正的错配点不在同一处。
这类“语言模式不匹配”的错误长得像语法错误，很容易被误读成代码写坏了。
如果确实需要显式声明输入的语言种类，该实现提供相应开关，名称以其文档为准。
\end{warning}

\section{每个阶段的可查产物}\label{sec:cl01-products}

产物扩展名是工具链之间约定的接口之一，但它的约定强度远低于直觉：同一语义在三个平台上
常有三套写法。下表只列本书会反复用到的、且在文档里有明确说法的那些；不确定的一律
写成“以该实现文档为准”，不凑数。

\begin{longtable}{@{}>{\raggedright\arraybackslash}p{0.16\textwidth} >{\raggedright\arraybackslash}p{0.14\textwidth} >{\raggedright\arraybackslash}p{0.20\textwidth} >{\raggedright\arraybackslash}p{0.30\textwidth}@{}}
\caption{流水线各阶段留下的产物与其常见扩展名}\label{tab:cl01-products}\\
\toprule
阶段 & 产物 & 常见扩展名 & 谁会读它 \\
\midrule
\endfirsthead
\toprule
阶段 & 产物 & 常见扩展名 & 谁会读它 \\
\midrule
\endhead
\bottomrule
\endfoot
预处理 & 展开后的翻译单元 & \Sect{i}（\CL）、\Sect{ii}（\CPP） & 前端；源文件的包含展开到此完成 \\
前端与分析 & 语法树、符号表 & 通常不落盘；部分实现有转储产物 & 同一实现的前端；转储名称与格式以其文档为准 \\
中端 & 中间表示 & 通常不落盘 & 与目标机器无关的优化序列 \\
后端 & 汇编文本 & \Sect{s} & 汇编器；也是你手工检查指令选择的第一现场 \\
汇编 & 目标文件 & \Sect{o}；\Win{} 下 \Sect{obj} & 链接器；符号引用仍未解决，产物同时带 \Sect{s} 与 \Sect{o} \\
归档 & 静态库 & \Sect{a}；\Win{} 下 \Sect{lib} & 链接器或归档工具（\Tool{ar} 一类），按成员抽取 \\
链接（共享） & 动态库 & \Sect{so}、\Sect{dll}、\Sect{dylib}；\Win{} 另有导入库 \Sect{lib} & 加载器与链接器 \\
链接 & 可执行镜像 & \Win{} 下 \Sect{exe}；\Nix{} 与 \Mac{} 常无扩展名 & 内核与加载器 \\
调试信息 & 符号与行号映射 & \Sect{pdb}、\Sect{dSYM}、\Sect{debug\_info} 一类节 & 调试器与符号化工具 \\
构建辅助 & 依赖片段 & \Sect{d} & 构建系统，决定谁需要重编 \\
预编译前缀 & 已解析的前缀 & 实现专有产物 & 同一实现的前端（跳过对前缀的重复解析） \\
模块接口 & 已分析的接口 & 实现专有产物 & 同一实现的前端（导入不做文本展开） \\
\end{longtable}

有两处歧义必须点破。其一，\Sect{lib} 这个扩展名在 \Win{} 上同时用于两种完全不同的东西：
它既可能是真正的归档库（成员是一堆 \Sect{obj}，链接期按需抽取），也可能只是导入库
（自己不承载任何实现，只记录“去哪个 \Sect{dll} 找这个符号”）。其二，同一阶段的产物
往往同时带着上下游两种扩展名：汇编这一步的输入是 \Sect{s}、输出是 \Sect{o}，
说“\Sect{o} 属于汇编”与说“\Sect{o} 属于链接的输入”都对，取决于你在数阶段还是数文件。

表中最后两行是第 \ref{ch:cl02-preprocessing} 章的主题，这里先只标注它们在流水线上的位置：
它们跳过的是“解析”这一段工作量，而不是“编译”这一段。

\section{自举是语言的事，载体是实现的事}\label{sec:cl01-carrier}

先给两个词各自的意义，因为它们的中文都常被写成“自编译”。

\emph{自举}（\Mx{boot{\Bk}strapping}）指用目标语言自己写出它的编译器。它有一个绕不开的先有鸡还是蛋
问题：要用 \Mx{L} 写的编译器编译 \Mx{L} 源码，就得先有一个能编译 \Mx{L} 的实现；实践上要么从一个
已有实现接力（三轮构建，逐轮用自己产出下一轮），要么先写一个极小子集来引导。
自举成立说明的是这门语言\emph{表达能力足够}，以及它的生态\emph{愿意付这份维护成本}。

\emph{载体}指这份实现的源码用什么语言写、链接谁的运行时、被什么编译器编译、以及它的
代码生成路径最终落在哪种中级语言或机器码上。载体是\emph{工程事实}，不是语言属性。

作者的立论是这一句：语言可以自举，但实现的载体几乎都是 \CL{}/\CPP{} 或其他中级语言。
下表把这件事摊开看。表里只写被文档与实现自身明确说明过的路径；
具体版本的行为会变，尤其是“实验性后端”那一类，以各实现自己的发布说明为准。

\begin{longtable}{@{}>{\raggedright\arraybackslash}p{0.15\textwidth} >{\raggedright\arraybackslash}p{0.12\textwidth} >{\raggedright\arraybackslash}p{0.11\textwidth} >{\raggedright\arraybackslash}p{0.23\textwidth} >{\raggedright\arraybackslash}p{0.19\textwidth}@{}}
\caption{若干主流实现的“自举”与“载体”对照}\label{tab:cl01-carrier}\\
\toprule
实现 & 源码语言 & 是否自举 & 代码生成路径 & 依赖的非本语言运行时 \\
\midrule
\endfirsthead
\toprule
实现 & 源码语言 & 是否自举 & 代码生成路径 & 依赖的非本语言运行时 \\
\midrule
\endhead
\bottomrule
\endfoot
\Tool{GCC}（C 前端） & \CL & 可自举，首轮仍需外部 \CL{} 编译器 & 自有中端与低位表示，再出机器码 & libc \\
\Tool{GCC}（C++ 前端） & \CPP & 与上同一条流水线 & 同上，与 C 前端共享中端与后端 & libc、\CPP{} 标准库实现 \\
\Tool{Clang} / \Tool{LLVM} & \CPP & 否 & 语法树降到 LLVM IR，由目标后端出汇编 & libc、\CPP{} 标准库实现 \\
\Tool{CPython} & \CL & 否 & 编译到字节码，由解释循环执行 & libc \\
\Tool{Ruby MRI} & \CL & 否 & 字节码虚拟机执行 & libc \\
\Tool{V8} & \CPP & 否 & 宿主语言 \Lang{JavaScript}；自有管线直达机器码，不经通用 IR 库 & libc 与系统 API \\
\Tool{HotSpot}（\Lang{JVM}） & \CPP & 否 & 即时编译分层出机器码，另有字节码解释器兜底 & libc、线程库 \\
\Tool{GNU Guile} & \CL & 否 & 字节码虚拟机执行 & libc \\
\Tool{Lua} & \CL & 否 & 寄存器式字节码虚拟机执行 & libc \\
\Tool{LuaJIT} & \CL & 否 & trace 编译，自有后端直接出机器码 & libc \\
\Tool{BEAM}（\Mx{Erlang}） & \CL & 否 & 字节码虚拟机，运行期另带本机码编译路径 & libc \\
\Tool{CoreCLR}（.NET 运行时） & \CPP & 否 & 中间语言由即时编译分层出机器码 & libc、系统线程库 \\
\Tool{OCaml} & \Mx{OCaml} 与 \CL & 可自举，需外部宿主起步 & 字节码与自有本机码后端两条路 & 以 \CL{} 写的运行时（内存管理、系统调用） \\
\Tool{GHC}（\Mx{Haskell}） & \Mx{Haskell} & 可自举，需外部宿主起步 & 自有中间语言；历史上以 \CL{} 中转，后来以自有后端出机器码 & 以 \CL{} 写的运行时系统 \\
\Tool{Rust} & \Mx{Rust} & 可自举，首轮需已有实现 & 自有中间表示降到 LLVM IR，再由 LLVM 出机器码 & \CPP{} 写的 LLVM、libc \\
\Tool{Go} 的 \Tool{gc} & \Mx{Go} & 可自举，首轮需已有实现 & 自有静态单赋值后端直接出机器码 & 运行时含 \CL{} 与汇编（内存管理、调度器） \\
\Tool{Swift} & \Mx{Swift} & 可自举，首轮需已有实现 & 自有中间表示降到 LLVM IR & \CPP{} 写的 LLVM；\Mac{} 上还依赖 \Mx{Objective-C} 运行时 \\
\Tool{SBCL}（\Mx{Common Lisp}） & \Mx{Common Lisp} & 可自举，需外部宿主起步 & 自有汇编器出机器码 & libc；链接仍需外部 \CL{} 编译器 \\
\end{longtable}

读这张表的正确方式不是数“谁更纯粹”，而是注意三个事实同时成立：
其一，源码语言是 \Mx{Haskell}、\Mx{Rust}、\Mx{Go}、\Mx{Swift} 的实现，仍然要链接 \CL{} 写的运行时
或 \CPP{} 写的代码生成库；其二，多数自举实现的\emph{第一轮}由别的编译器产出，
“自举”描述的是链条能闭合，不是链条不需要起点；其三，代码生成路径上出现 \CL{} 或 LLVM IR
这两个中级表示的次数，远高于直觉。

自举的实际收益不在美学，而在两处工程约束：语言规范必须可被一个实现严格满足（不能靠
“运行时帮忙”糊过去），以及实现不再受外部工具链的发布节奏绑架。代价同样明确——
每次改语言特性都要重跑整条构建链，三轮甚至更多，跨版本兼容矩阵会迅速变大。
本书后文遇到这类成本时只给定性说法（“数量级”“取决于实现”），不给数值。

\subsection{自举的门槛只有一条：图灵完备}\label{sec:cl01-carrier-tc}

上表“是否自举”那一列有一处容易读错的地方：写“否”的那些行记录的是\emph{生态选择}，
不是\emph{语言能力}。这里把两条判据分开。

可行性只有一条门槛：目标语言图灵完备。理由很直接——编译本身就是一次可计算变换：
把源码文本读进内存，按语法规则构造出中间结构，再按规则逐条产出汇编、字节码或
中级表示。既然这是一次可计算变换，一门能够表达分支、循环、递归、数据结构与
函数调用的语言就足以写出它自己的编译器。反过来，一门图灵不完备的语言写不出
自己的编译器，因为总有一段编译过程它表达不出来；这份缺口无法靠工程投入补齐，
只能靠换一门宿主语言。

按这条判据重读表里那些“否”：\Mx{Python}、\Lang{JavaScript}、\Mx{Java}、\Mx{C\#}、
\Mx{Erlang} 与 \Mx{Lua} 全都远过得了这条线。它们的实现不换载体，只是因为另一笔账
更划算——已有的 \CL{}/\CPP{} 实现足够快、足够稳，换载体要重建引导流程、剖析器、
消毒器与部署形态，收益抵不上成本。这是取舍，不是能力不足。

需要提醒的是，门槛管可行性，不管起点。图灵完备只说明链条\emph{能够}闭合，
闭合仍然需要一个起点：从一个已有实现接力，或者先手写一个极小子集来引导。
本节开头已经把这两件事分开了，这里给出的是它们的判据来源。

真正落在这条门槛之外的常用语言极少，最值得点名的是 \Mx{SQL}。它的查询核心是
声明式的：表达的是关系代数加上不动点这一类可定义的集合运算，没有一条通用控制流
可以承载“读一段文本、按语法构造结构、逐条规则输出字节码”这个过程。这个判断只
针对查询核心——标准后来加入了递归查询与持久存储模块，各家实现还嵌入了过程式语言、
外部例程与脚本扩展，是否越过门槛要按具体方言里递归与控制流的表达力判断，
不能用一句“\Mx{SQL} 图灵完备”概括。常见文献对它的定位仍然是声明式查询语言，
而不是通用程序设计语言。

和 \Mx{SQL} 一起常被列进来的还有两类：标记与样式语言（\Mx{HTML}、\Mx{CSS}）
与模式语言（正则表达式）。它们一般不作为程序设计语言讨论，谈自举时不必为它们
单独设一格。数控加工里的 G 代码偶尔会被算作一门语言，它同样不作为计算机语言讨论；
标准化文档只规定基本代码的含义，程序化扩展留给各控制器厂商，各家宏语义互不相同，
因此也没有一份可据以判定表达能力的共同规范。

\begin{note}[title={两条判据不要混用}]
问“能不能用 \Mx{L} 写 \Mx{L} 自己的编译器”，只需要问 \Mx{L} 是否图灵完备，
答案是二值的。问“该不该这么做”，要问工具链成熟度、引导维护成本与人力，
答案是一个随生态变化的排序。本书后文说某个实现“没有自举”时，指的是第二类判断，
不是第一类；而它即使选了第二类，载体这一层仍然由下一节那四条硬约束决定。
两条判据分开记，就不会把“某个实现至今用 \CL{} 写”读成“那门语言写不出自己的编译器”。
附录 \ref{ap:clC-selfhost} 按实现逐条给出起始语言、当前语言与是否自举三列，
可当作这两条判据的取值表。
\end{note}

\section{为什么现实如此}\label{sec:cl01-why}

载体几乎总是 \CL{}/\CPP{} 或其他中级语言，不是因为它们“够快”，而是因为四件事同时压在
每一个实现者身上。

\begin{guideline}[title={选择实现载体的四条硬约束}]
\begin{itemize}
  \item ABI 边界与外部接口成本。操作系统 API、线程库、分配器、图形与网络栈的公开接口
    几乎都是 \CL{} 调用约定。以 \CL{} 为实现语言，跨边界是“直接调”；以别的语言为实现语言，
    要么自带一层 \CL{} 胶水，要么把整个运行时嵌进去，胶水与运行时都是长期维护负债。
    平台 ABI 的细节在第 \ref{ch:cl13-cxx-abi} 章与附录 \ref{ap:clB-abi}。
  \item 与加载器、线程库、异常机制的直接接触。启动代码、线程局部存储、展开表、
    信号处理这些设施在多数平台上只以 \CL{} 形态提供，实现要么用 \CL{} 写这一层，
    要么复用一个已经写好这一层的运行时。
  \item 成熟工具的绑定关系。消毒器、剖析器、静态分析、覆盖率的可用性与这些工具的实现
    语言绑定：能把“被分析程序”与“分析器”放在同一个中间表示上，是极其省事的设计。
    走 \CL{}/\CPP{} 的实现对这类工具是白送的；走自定义运行时时要自己重做一遍。
  \item 可移植性的来源是 \CL{}，不是语言本身。一门语言宣称“可移植到所有 POSIX 平台”，
    实际上是在说“它的实现所依赖的那层 \CL{} 编译器在这些平台上存在”。
    换平台时的第一份工作量，永远是移植运行时与构建脚本，而不是移植语言特性。
\end{itemize}
\end{guideline}

还有第三条理由：\CPP{} 作为载体的收益，恰好落在编译器的数据结构密集区。
编译器与链接器是“大量小对象、大量字符串键、大量树与图”的程序，容器、抽象基类、
移动语义与 \Tp{constexpr} 常量求值这类设施的价值在这类负载上很高；而运行期开销的敏感度
相对低（一次性成本，不是稳态循环）。这也解释了为什么 \Tool{LLVM}、\Tool{Clang}、\Tool{V8}、
\Tool{HotSpot} 与 \Tool{CoreCLR} 都选了 \CPP{} 而不是 \CL{}。

还有一层与此正交的考虑：即使一个实现不用 \CL{}/\CPP{} 写源码，它也常常把\emph{代码生成}
汇合到某个中级表示上——LLVM IR、\Lang{JVM} 字节码、\Mx{Wasm}、或某个平台的汇编文本。
这个中级表示就是不同实现之间的“公共汇合点”，选择它的判据是客观的：
（1）该表示上是否已有足够的优化与验证工具；（2）从该表示到镜像的最后一段路是否由别人维护；
（3）该表示是否会把你锁在某个平台或某个实现上。第（3）条常常在多年后才显现代价。

\section{一次冷启动的依赖链}\label{sec:cl01-coldstart}

流水线跑到链接产出镜像并不是终点。镜像里那些“还没有地址的引用”，以及“地址有了但内容还没准备好”
的部分，构成一条冷启动依赖链。理解这条链的关键是把两个时刻分开：\emph{链接期看到的东西}
与\emph{运行期看到的东西}。

\begin{abibox}[title={链接期可见与运行期可见，两回事}]
\begin{itemize}
  \item 链接期可见：目标文件里的符号表与重定位项、静态库的成员、导入库或链接脚本、
    启动代码（一组由运行库提供的目标文件）、运行时库自身、以及入口符号的解析结果。
    链接器只解决它这一侧的问题，见第 \ref{ch:cl03-resolution} 章。
  \item 运行期可见：镜像的依赖清单、动态加载器、符号插入与延迟绑定、
    重定位的实际执行、线程局部存储的初始化模板、全局构造与析构的注册表、
    以及最终的控制转移。见第 \ref{ch:cl09-dynamic} 章。
  \item 两个时刻的常见混淆：链接期“找到了符号”不等于运行期“符号可用”；
    运行期“加载成功”也不等于“入口已经跑起来”。
\end{itemize}
\end{abibox}

把一次启动按事件顺序排开，形状大致是这样：

\begin{lstlisting}[style=outstyle,caption={一次冷启动的事件顺序（示意描述，非真实工具输出，不含实测数据）}]
# 每段带一个标签：标签属于阶段，不属于某个工具
[link]  符号解析与重定位 -> 镜像成型；未定义引用此时必须全部有个着落
[load]  内核建立地址空间与初始线程 -> 加载器定位镜像 -> 逐个解析依赖清单
[reloc] 执行地址重定位 -> 为每个线程按初始化模板铺开线程局部存储
[init]  执行镜像注册的初始化序列（共享库自身的初始化排在其前）
[entry] 跳到入口符号 -> 运行库启动例程建立 C 运行环境 -> 全局对象构造
[main]  调用语言入口
[exit]  主程序返回 -> 逆序执行析构与注册回调 -> 交由退出路径结束进程
\end{lstlisting}

这条链上有几个角色值得点名。启动代码（\CL{} 生态里常说 \Mx{crt}，即 C 运行时启动目标文件）
负责把操作系统交给进程的原始状态整理成语言入口期望的样子：参数与环境变量的来路、
栈的起点、标准输入输出的初始化。入口符号的确切名字随平台与运行库而变——\Nix{} 与
\Mac{} 的 \Elf{}/\Macho{} 世界里常见 \Sym{\_start}（它属于运行库而不属于语言），
\Win{} 的工具链则由运行库提供形如 \Fn{mainCRTStartup} 的启动例程；本书只在“角色”的意义上
使用这些名字，确切符号以你的链接器映射报告为准。

线程局部存储这一段最容易被忽略：带初始值的 \Tp{thread\_local} 对象在镜像里不是一个现成的实例，
而是一份初始化模板，加载器为每个线程复制一份。这也意味着“某个全局对象的构造函数里
读取线程局部对象”这种写法，其正确性取决于运行期初始化是否已经铺到那一层，而不是
链接期能否解析符号。语言层对这类时序没有承诺，平台 ABI 与实现才是事实来源。

\begin{note}[title={本册范围声明}]
本册讲从源码到镜像的机制、目标文件格式、链接与加载，以及平台 ABI 与工具链用法。
三件事不属于本册：语言层面的未定义行为判定（另一分册专门讨论求值顺序与未定义行为的边界，
本册只在预处理让“同一段源码在不同档位成为不同程序”这一点上回指它）、并行算法的使用方式、
以及性能测量的统计判据（另一分册讲基准测试的口径与置信区间）。
本册因此不给编译时间、产物体积与百分比数据，讲代价一律用定性表述。
\end{note}

\section{每个阶段用什么工具去看}\label{sec:cl01-viewing}

最后给一张“看”的索引。工具名的确切拼写与开关以各自的手册为准；
这里给的是形状，不是命令手册。

\begin{lstlisting}[style=shellstyle,language=ShellCL,caption={各阶段产物用只读工具查看的形状}]
# 预处理产物：肉眼确认包含闭包展开成了什么
less widget.ii

# 汇编文本：指令选择与调用约定的第一现场
less widget.s

# 目标文件：先看节，再看符号，最后看重定位
readelf -hS widget.o
nm -C widget.o
objdump -d -r widget.o

# 静态库：把成员列表与每个成员的符号一起看
ar t libwidget.a
nm -C libwidget.a

# 镜像：三个平台三种视角（选项含义以各自手册为准）
readelf -d app                # 依赖清单与运行期搜索路径
objdump -T app                # 动态符号表
dumpbin /headers app.exe      # Windows 下的同类视角
otool -h -l app               # Mach-O 下的同类视角

# 目标文件加工：裁剪节名、剥离符号、在格式间搬运内容
objcopy --only-keep-debug app app.dbg
strip --strip-debug app

# 汇编器视角：让实现自己吐出它认得的输入形状
as --version
\end{lstlisting}

\begin{bestpractice}[title={按阶段定位问题的最短路径}]
\begin{itemize}
  \item 怀疑预处理：把源文件的包含闭包展开成文本再看一遍，多数“诡异错误”在这一步现形。
  \item 怀疑前端与中端：落盘汇编文本，比较两个优化档位的差异，比读日志有效。
  \item 怀疑链接：拿目标文件的符号表与链接器的映射报告对表，不要凭命名习惯猜。
  \item 怀疑加载：用只读工具列镜像依赖与节布局，再决定是否需要延迟绑定或符号可见性调整。
  \item 怀疑自己的理解：把每一步产物存档成一个可复放的脚本，这是编译器问题报告里
    最有价值的附件。
\end{itemize}
\end{bestpractice}

\section{目标环境的覆盖面：桌面、移动与裸机}\label{sec:cl01-targetenv}

全卷的证据面按“三个平台 × 两套架构”组织（\Nix、\Win、\Mac{} 与 \xSixtyFour、\ArmSixtyFour）。
区分的依据不是市场大小，而是流水线的\emph{形状}会不会变——工具名换了不算变，阶段里没有
加载器、启动不由操作系统接管才算变。

\subsection{上层应用：移动平台走的是同一条流水线}\label{sec:cl01-mobile}

移动平台的上层应用与桌面走同一条流水线；差别只有三处。其一，交叉编译：产出目标 ABI 的
产物，一切由三元组选择。其二，ABI 由平台规定，没有挑选余地（第 \ref{ch:cl10-cc} 章的表在
\ArmSixtyFour{} 一列即是）。其三，分发形态受平台约束：镜像格式、签名与可加载的动态库范围
由平台决定。

\begin{lstlisting}[style=outstyle, caption={目标三元组的四段各决定什么（示意形状，不是某个工具的实际字符串）}]
<arch> - <vendor> - <os> - <environment>
  arch         指令集、数据模型与调用约定
  vendor       多数场合为空，个别发行用它区分形态
  os           系统调用与线程设施的接口形状
  environment  链接哪一份 C 运行时、能否动态加载
\end{lstlisting}

三元组的读法在移动平台上立刻有用：\Mx{Android} native 侧的操作系统段仍然是 \Nix{}，
变化发生在环境段——它决定链接的是 \Mx{Android} 自己那份 C 运行时，
而不是桌面发行里的那一份。也就是说“移动”并不等于“另一套底层机制”，
它多半只是同一套机制换了一份实现。

本卷只在 \Mx{Android} 一侧把系统分层写出来，因为它是“上层应用里混着两份语言与两份 ABI”
的典型形状。按从应用到底层的顺序，与 \CL/\CPP{} 相关的层次是这几层：

\begin{enumerate}
  \item 应用层：界面与业务代码跑在 \Tool{ART} 上，与 \CL/\CPP{} 无关，
        只有需要 native 能力时才向下降一层。
  \item 跨语言边界：\Mx{JNI} 一类的调用接口。这一层的参数与返回值约定、
        引用与对象生命周期由另一份规范管，不属于第 \ref{ch:cl13-cxx-abi} 章的语言层 ABI；
        把两侧混为一谈是这一类集成最常见的错。
  \item native 层：\CL/\CPP{} 源码由该平台的交叉工具链（\Tool{Clang} 一族）编译，
        按各自的 ABI 各出一份共享库，链接产物再打进分发格式。
  \item C 运行时：\Tool{Bionic}——\Mx{Android} 自己的实现，既不是桌面 \Nix{} 上的那一族，
        也不是另一支常见的嵌入式实现。它决定了第 \ref{ch:cl12-endian} 章数据模型
        与实现定义取值的实际答案，同一份源码在这里可能得到不同结论。
  \item 加载层：动态链接器与共享库的搜索路径、可见性策略。现代版本按应用划分
        命名空间，跨应用共用符号需要显式声明，第 \ref{ch:cl09-dynamic} 章讲的
        那一层机制在这里换了主人。
  \item 内核层：与桌面 \Nix{} 同源，所以容器格式、系统调用约定、线程与线程局部存储
        的机制照用，本卷正文可以不加修改地套用。
\end{enumerate}

结论是：移动平台的“移动”只作用在第 3 与第 5 两层；第 1、2 层是另一门语言与另一份规范，
第 4、6 层与本卷对 \Nix{} 的讨论部分重合。另一支移动平台在分层上不同
（镜像格式、系统框架与分发策略各自成一家），但对\CL/\CPP{} 的结论相同，
本卷只在需要对照时提一句，不单独展开。

\subsection{裸机：整段链接与加载换了主人}\label{sec:cl01-baremetal}

标准把一致实现分成两档：\Mx{hosted}（提供完整标准库）与
\Mx{freestanding}（只承诺一个很小的头文件集合）。裸机属于后者，
而两档的区别正好落在这条流水线的后半段上。

前半段完全不变：预处理、语法与语义分析、中间表示、优化、指令选择、汇编器
都照原样工作，只是目标三元组里没有操作系统段。变的是阶段六与阶段七：

\begin{itemize}
  \item 没有动态链接，也没有加载器。符号解析与地址确定全在链接期完成，
        第 \ref{ch:cl09-dynamic} 章的机制在这里整体退场。
  \item 镜像布局由链接脚本描述：代码与只读数据放哪段地址、可写数据放哪里、
        向量表落在哪个固定入口。加载地址与运行地址可以不同，
        于是需要启动代码在初始化阶段把数据从前者搬到后者。
  \item 入口不是操作系统给的。第 \ref{ch:cl09-dynamic} 章末那条“从镜像到 \Fn{main}”
        的启动序列由启动文件接管：设置栈与数据段、搬运初始化数据、清未初始化区，
        然后才跳到入口；标准库那套全局对象登记在很多目标上根本不存在。
  \item 外部性来自寄存器映射与中断，不来自系统调用。地址空间靠映射的 peripherals
        表达，这类访问的语义由目标手册而非语言标准给出。
  \item 标准库换成面向目标的子集，异常与运行时类型信息常被关闭；
        并发与线程局部存储由实时系统或裸实现替代。
\end{itemize}

所以本卷的调用约定、数据布局、字节序、指令选择与寄存器分配这些内容在裸机上照用；
需要替换的是目标文件、静态链接与动态链接那三章的具体形态。这一类的启动细节、
链接脚本写法与目标手册的读法属于另一分册，本卷把边界停在这里。

\subsection{不讨论的那一类：旧世代主机与掌机}\label{sec:cl01-excluded}

旧世代游戏主机与掌机不进本卷，理由不是它们工程价值低，而是它们不满足本卷的取证前提。
那一代设备的系统是厂商自建的：不满足 \Nix{} 与 POSIX 那一族的接口要求，
也没有公开可引用的 ABI 文档；编译器、链接器与运行时只在开发套件里分发，
行为细节要查厂商手册，而手册通常带保密条款。本卷的纪律是每个结论都要能在
公开文档或产物里核对（第 \ref{ch:cl10-cc} 章开头把这条写成了可信边界），
在这一类目标上这条纪律执行不了；把“某代主机上是这么做的”写进正文，
等于让读者依赖一份他拿不到的文档。

同一条判据也排除别的定制目标：需要授权才能使用的工具链、厂商私有汇编器与链接器、
只在内部文档里定义的符号约定。判据只有一句——\textbf{这个环境的 ABI 与工具链行为
能不能公开核对}。能，就在讨论范围内（三个桌面平台、移动平台的 native 层、
裸机与实时系统都满足）；不能，就排除，并且在这里写明被排除。

\begin{note}[title={覆盖范围的用法}]
这一节不是免责条款，而是一张边界图：读第 \ref{ch:cl07-object} 到
第 \ref{ch:cl09-dynamic} 章时，先确认自己的目标环境落在哪一格。
桌面与移动 native 可以直接套用；裸机要换成上一小节那一组等价物；
自建系统那一类则本卷的任何具体名字都不要照抄，只借机制。
\end{note}

本章的形状会在全书反复使用：第 \ref{ch:cl02-preprocessing} 章把阶段二讲透；
第 \ref{ch:cl04-ir} 与第 \ref{ch:cl05-passes} 章进入阶段三与阶段四；第 \ref{ch:cl06-codegen} 章
进入阶段五；第 \ref{ch:cl07-object} 到第 \ref{ch:cl09-dynamic} 章处理阶段六与阶段七；
第 \ref{ch:cl14-ide} 章把“用什么工具去看”扩到编辑器与索引层，
附录 \ref{ap:clA-cli} 收录命令行速查，附录 \ref{ap:clC-selfhost} 回到自举链的构造细节。

\chapter{预处理、头文件与模块}\label{ch:cl02-preprocessing}

第 \ref{ch:cl01-pipeline} 章的阶段图里，第二阶段是唯一一个不产生语义结构、只产生文本的阶段。
它决定了后面所有阶段拿到什么：一份已经把包含闭包展开完、条件分支已经裁决完、宏已经替换完
的文本。本章讲清楚这一层的机制，然后讲它带来的两笔账——重复解析的成本，以及它如何改变语义本身；
最后讲 \cpp{20} 模块想从根上解决的是同一笔账，但解决的是哪一段。

\section{预处理是一门独立的语言层}\label{sec:cl02-layer}

预处理不是“编译的第一步优化”，而是一门运行在记号层上的独立语言。它与主体语言有三点根本差别。

第一，它不认识作用域与类型。宏替换发生在名字查找之前，替换结果只是记号序列；
如果替换后的序列不合法，错误会报在你根本没写过的行上。第二，它不求值。
条件编译只看标识符是否被定义、以及常量表达式求值的规则，不看程序的实际行为，
所以“这个分支永远不会被执行”与“这个分支不会被编译”是两件不同的事，代价完全不同。
第三，它是复制而非引用。\#include 把被包含文件的\emph{全文}插入当前位置，
同一份头在十个翻译单元里就被读十遍，每次都是从头解析。

\begin{table}[htbp]
\centering
\begin{tabularx}{0.76\textwidth}{@{}>{\raggedright\arraybackslash}p{0.20\textwidth} X@{}}
\toprule
指令族 & 它在流水线里做什么 \\
\midrule
\#define / \#undef & 建立或撤销一个宏；带参数的宏做记号层替换，替换结果会再次扫描 \\
条件编译族 & 决定哪些记号序列进入后端；未选中的分支连语法都不必合法 \\
\#include & 展开包含闭包：把另一份文本原样插入当前位置，形成翻译单元的实际大小 \\
\#pragma & 一个开放机制：语义由实现定义，不属于语言承诺。用来控制警告、对齐、
优化提示、运行期特性等；同名 \#pragma 在不同实现上的效果可以完全不同 \\
其余 & 行控制、错误报告、记号拼接与字符串化、\Tp{\_Pragma} 一类变体 \\
\bottomrule
\end{tabularx}
\caption{预处理指令族按“做什么”分类，而不是按关键字罗列}\label{tab:cl02-directives}
\end{table}

\#pragma 值得单独说：它是语言留给实现的扩展点，标准只规定“未识别的 \#pragma 必须被忽略”。
因此同一个 \#pragma 在 \Tool{gcc}、\Tool{clang} 与 \Tool{cl} 上可能一个是关键优化开关、
一个是空操作、一个是警告。任何依赖具体 \#pragma 语义的代码，可移植性都要按“实现专有特性”评估；
需要跨实现时，用该语言标准设施或条件编译分派，而不是指望 \#pragma 名字通用。

\begin{guideline}[title={把预处理当“文本层”而不是“语言层”}]
\begin{itemize}
  \item 判断一个设施属于哪一层，问一句：它有没有类型？没有，多半在预处理层。
  \item 宏不做类型检查，所以它“什么都能接受”，也因此“什么都能接受错的”。
  \item 条件编译改变的是编译进去的程序，不是执行的程序；换档位等于换源码。
  \item 头文件属于文本层，所以它的可见性靠包含顺序，而不是靠声明位置。
\end{itemize}
\end{guideline}

把这一层放回阶段图：阶段二产出的是翻译单元，而从这一刻起，后面所有阶段的工作量都由
这一份文本的大小决定。这就是“翻译单元有多大”这个问题的来源。

\section{头文件的成本模型：翻译单元等于包含闭包}\label{sec:cl02-cost}

一个翻译单元的实际内容不是那几百行源文件，而是它的包含闭包：源文件、它直接包含的每个头、
那些头各自包含的头，加上被条件编译选中的所有分支，全部展平后的文本。包含关系构成一张图，
而闭包展开会把这张图变成一棵树——同一个头出现在两条包含路径上，就会在这棵树里出现两次
（头卫哨或 \#pragma once 让第二次展开变成空操作，但“打开文件、读到卫哨、判定已包含”
这一步仍然要执行，仍是一次文件系统与解析开销）。

标准库头在这里特别贵，原因全是定性的，与它们的写法有关：

\begin{itemize}
  \item 模板密集。标准库的绝大部分实现必须出现在头里，因为模板定义必须对使用点可见。
    读一个容器头，实际是在读它的实现，而不是一份声明。
  \item 相互牵连。一个标准库头常常拉进若干其它标准库头与若干底层设施头，
    闭包大小不由你决定。
  \item 语义要在解析时完成。名字查找、重载决议、模板实参推导、局部特化匹配，
    这些工作都要在读到定义之后才能做；文本没读完，语义就不存在。
  \item 同一实现的前端默认没有跨翻译单元的解析缓存。同一份闭包在十个源文件里就是十次解析，
    这也是编译期加速手段全部围绕“少解析几遍”设计的原因。
\end{itemize}

\begin{note}[title={成本落在哪一段，决定用哪种办法省}]
包含闭包影响的几乎是“解析与分析”这一段的工作量，而不是代码生成或链接的工作量。
因此任何只优化后端的开关（例如更高的优化档位）都会让它更贵而非更便宜：同一份语法树上
要跑更多优化序列。反过来说，减少解析量的一切手段——预编译头、模块、unity 构建——
都不改变生成的机器码语义（这是它们能被接受的前提），改变的只是前端读进去多少文本。
\end{note}

\section{复用已解析的信息：三种做法的机制}\label{sec:cl02-reuse}

三种主流做法本质上只做一件事：把“已经解析过的信息”复用一次。它们的差别在于复用\emph{什么}、
在\emph{什么时候}失效、以及把这条边界推给了谁。

预编译头（各类实现里也常写作 \Mx{PCH}）把一个前缀——通常是若干稳定头——解析成一份实现专有的
产物，后续翻译单元从这份产物起步，跳过对前缀的重复解析。这份产物不是目标文件，
不能跨实现甚至常常不能跨版本使用，它对同一实现的前端而言等于“已经读完的那半本词典”。
前缀里任何一处改动都会让它整体失效，因为产物描述的是解析结果而不是文件时间戳。
它在不同实现上的可用性与产物命名都属实现细节，以该实现文档为准。

unity 构建（也写作 \Mx{jumbo} 或 \Mx{amalgamation}）思路完全相反：不做解析缓存，
而是把多个源文件拼成一个巨大的翻译单元。省下来的是“每个源文件都要付一次完整包含闭包”这笔
固定开销——拼成一个之后，那笔开销只付一次。它的代价在下面逐条列。

模块接口是第三种：把\emph{语义层}的接口信息固化成产物，导入方直接消费已分析的声明，
不再做文本展开。机制细节放在本章下一节。

\begin{longtable}{@{}>{\raggedright\arraybackslash}p{0.14\textwidth} >{\raggedright\arraybackslash}p{0.20\textwidth} >{\raggedright\arraybackslash}p{0.22\textwidth} >{\raggedright\arraybackslash}p{0.24\textwidth}@{}}
\caption{三种“复用已编译信息”的做法对照}\label{tab:cl02-reuse}\\
\toprule
做法 & 复用什么 & 失效条件 & 对增量与链接的影响 \\
\midrule
\endfirsthead
\toprule
做法 & 复用什么 & 失效条件 & 对增量与链接的影响 \\
\midrule
\endhead
\bottomrule
\endfoot
预编译头 & 前缀的解析结果（符号表、模板、重载集） & 前缀内任一头被改动；实现或编译器版本更换 & 增量粒度不变；产物与实现绑定，不进版本控制；不改变目标文件内容与链接语义 \\
模块接口（BMI/PCM 一类接口产物） & 已分析的接口声明本身，而非文本 & 接口单元被改动并重编；导入方需重新消费新产物 & 增量沿“谁导入谁”的依赖图传播；符号身份仍由修饰规则决定，链接期看到的东西不变 \\
unity 构建 & 靠合并翻译单元摊薄包含闭包 & 合并组内任一源文件被改动，整组重编 & 增量粒度显著变粗；静态符号在组内互相可见，可能改变名字查找结果；目标文件数量变少，链接输入的形状随之改变 \\
\end{longtable}

\begin{compare}[title={缓存解析结果与改变输入形状}]
\begin{itemize}
  \item 预编译头与模块接口都是“让同一实现的前端少读文本”，源文件边界不变，链接期不变。
  \item unity 构建是“改变翻译单元的边界”，代价直接落在语义上（静态符号可见性、宏泄漏），
    必须靠纪律控制。
  \item 三者的收益都集中在前端；对代码生成与链接没有任何承诺。
  \item 三者都不能替代“少包含、把实现从头里搬走”这类结构性优化——那才是唯一的根因治理。
\end{itemize}
\end{compare}

unity 构建的副作用值得逐条写清楚，因为它们都会以“代码没问题但构建失败”的形式出现：

\begin{itemize}
  \item 静态符号可见性变化。文件级 \Tp{static} 函数与匿名命名空间本来是每个源文件私有的，
    拼成一个翻译单元后它们同处一个作用域，重名即冲突。反过来，某文件原本能解析到的
    私有名字，现在可能被另一个文件的同名声明抢走，名字查找结果发生实质改变。
  \item 宏泄漏。前一个文件的宏还在，后一个文件的代码就在那份宏的影响下被解析。
    一个常见的“拼接后突然编译失败”的原因，就是某个头里的宏改写了另一个文件里的标识符。
  \item 增量粒度变粗。改动任意一个成员，整个合并组重编；持续集成上的表现是“改一行等一份大文件”。
  \item 错误定位与诊断噪声。行号指向拼接后的文件；同一个头在组内被包含多次时报重复定义。
  \item 峰值内存。一个巨大翻译单元的解析与分析需要同时驻留，实现侧内存占用与
    单个巨型语法树的规模相关，这是它最主要的失败模式。
  \item 调试信息映射。行号表要覆盖拼接后的来源，源映射不对时断点会跳。
\end{itemize}

\begin{lstlisting}[style=shellstyle,language=ShellCL,caption={unity 与分步的命令形状（示意，只用来观察代价）}]
# 分步：三个源文件、三次完整包含闭包、三个目标文件
g++ -std=c++20 -c src/a.cpp -o build/a.o
g++ -std=c++20 -c src/b.cpp -o build/b.o
g++ -std=c++20 -c src/c.cpp -o build/c.o

# unity：build/unity.cpp 由构建系统生成，内容是把这三个源文件包含进来
# 一个翻译单元、一次闭包、一个目标文件；-I 让被拼接的文件能被找到
g++ -std=c++20 -Isrc -c build/unity.cpp -o build/unity.o

# 合并之后先看符号表：原本私有的静态符号是否撞名
nm -C build/unity.o
\end{lstlisting}

CMake 侧的目标级用法只列确证过的命令签名与属性名；预编译头的具体实现兼容性、
以及模块接口在构建系统里的声明形式，都随版本演进，以官方文档为准。

\begin{lstlisting}[style=cmakestyle,language=CMakeCL,caption={目标级用法：预编译头与全局优化提示}]
cmake_minimum_required(VERSION 3.30)
project(cl_pch_demo LANGUAGES CXX)

add_library(widget STATIC widget.cpp util.cpp)
target_compile_features(widget PUBLIC cxx_std_20)

# 前缀：稳定且体积大的头
target_precompile_headers(widget PRIVATE
  <vector>
  <string>
  widget_pch.h)

add_executable(app main.cpp)
target_link_libraries(app PRIVATE widget)

# 跨翻译单元优化开关：按策略与实现支持情况启用
set_target_properties(app PROPERTIES INTERPROCEDURAL_OPTIMIZATION TRUE)

# 变量级同名开关存在（CMAKE_INTERPROCEDURAL_OPTIMIZATION），
# 模块接口在构建系统里的写法随版本变化，两者细节均以官方文档为准
\end{lstlisting}

\begin{lstlisting}[style=outstyle,caption={接口被复用前后，“翻译单元”构成的变化（示意描述，不含实测数据）}]
# 示意：只描述构成，不给时间、体积或百分比
复用前（普通包含）
  主源文件
  + 它直接包含的每个头 + 每个头再拉进去的头（闭包，出现在多条路径上）
  + 标准库头的完整文本与被选中的所有条件分支
  => 同一实现的前端从零解析全部内容；换一个源文件再来一遍

复用后（预编译头 / 模块接口）
  主源文件
  + 一份“已解析前缀”或“已分析接口”（由实现产出，导入方直接消费）
  => 前端只为增量部分重建符号表；机器码与复用前应当等价

失效时的样子
  前缀或接口一旦改动，依赖它的翻译单元全部要重新消费
  增量粒度从“改一个源文件”变成“改一个接口，重编一片”
\end{lstlisting}

\section{\cpp{20} 模块：把接口从文本层搬开}\label{sec:cl02-modules}

模块的目标是把“接口”从文本层移到语义层。它不是一门新语言，它改变了阶段二的输入形状。

\emph{命名模块}是一份显式命名的接口加实现。接口单元产出的不再只是文本约定，
而是一组被分析过的声明；导入方通过名字拿到这些声明，而不是把别人的头抄进自己的翻译单元。

\emph{头单元}是过渡设施：把一个已有的头包装成可导入的实体，让存量代码在不改头的前提下
进入导入路径。它是妥协，不是替代方案——用它意味着那个头仍然要存在，仍然会被实现内部处理。

\emph{接口产物}（\Mx{BMI} / \Mx{PCM} 一类叫法）是编译接口单元时产出的实现专有文件：
已分析的声明序列，加上模块身份与依赖信息。它的名字、扩展名、存放位置随实现而变，
\Tool{Clang} 与 \Tool{GCC} 各行其是，标准另有约定；本书一律称“接口产物”，具体名称以该实现文档为准。

关键的机制差别只有一句：\emph{导入不再做文本展开，编译器可以复用已分析过的接口}。
由此推出三条实际后果：包含顺序不再决定可见性；实现细节可以从接口里搬走，不再被迫放在头中；
重复解析被模块身份挡住。它没有推出的结论也值得写下来：模块不改变名字修饰，不改变布局，
不改变链接期看到的符号集合，也不替你做 ABI 兼容承诺。

\emph{导出}是显式动作：没写导出标记的声明属于模块内部，导入方看不见。这正是模块能把
“实现细节从头里搬走”的原因——可见性由声明位置决定，而不是由文件边界决定。

宏不可导出，原因是层次不匹配：导出工作在语义层，而宏工作在语义层之前的文本层，
等语义建立时宏已经不存在了。因此模块单元里能出现的预处理指令被限制得很窄：
\#include 只能出现在模块声明之前的那一段，实现诊断里常写作
\Tp{global{\Bk}module{\Bk}fragment}。这不是实现偷懒，而是分工：一份用宏做配置的头，
不能指望通过导入获得传递性。

再把两个容易混的概念分开。头单元（\Tp{header{\Bk}unit}）是\emph{把一个既有的头}包装成
可导入的实体，对象是别人的文件。分区（partition）则相反，它把一个模块自己的实现
拆成若干可命名的片段：分区之间可以互相导入，外部只看到主接口聚合的结果，
对象是你自己的模块。头单元解决“存量头怎么进导入路径”，分区解决“一个模块太大怎么组织”，
两者都不解决循环依赖。循环依赖在模块图里是被禁止的：接口单元的依赖必须无环，
因为接口产物的生成顺序本身就构成一次拓扑排序——甲的接口要导入乙的接口，
就得先有乙的产物；环上的产物无法约定谁先产出。想把两个模块互相导入时，
正确解法是把公共部分抽成第三个模块，或者退回普通包含。

\begin{warning}[title={把宏藏进模块声明之前那一段：可移植性地雷}]
模块声明之前那一段（实现诊断里写作 \Tp{global{\Bk}module{\Bk}fragment}）里的声明
仍属于全局模块，它存在的目的是容纳那些必须“在模块之前”看到的东西，例如某些第三方头。
把宏定义塞在这里、再指望导入方受益，是行不通的：宏不会被导出，导入方看不到它。
真正的问题是另一个方向——你在这段里包含的头会带进宏，这些宏在模块外的其它翻译单元
仍然有效，于是行为取决于这段里包的头的顺序。
再加一条：同一个头既被文本包含又被作为头单元导入时，其组合是被标准明确约束的，
实现普遍以诊断或难料的行为对待。把这种混合当技巧，等于把风险留给下一个编译器版本。
\end{warning}

\begin{abibox}[title={模块接口里的模板与 inline 函数，链接期以什么身份存在}]
模块不改变这些东西在链接期的身份。导出的 \Tp{inline} 函数与模板实例仍然要遵守
\CPP{} 的多重定义规则：每个翻译单元都可以产生一份完全等价的定义，链接期需要一种机制
把重复的合并掉。这个合并机制属于目标文件与链接器的实现：\Elf{} 生态常见做法是弱符号
与公共符号，\Win{} 的 \PeFmt{} 生态是 \Mx{COMDAT} 折叠加去重元数据，\Macho{} 生态另有
其约定。换句话说，模块产物让同一实现的前端少读文本，链接期看到的仍然是同一批符号；
具体合并与去重规则见第 \ref{ch:cl08-static-link} 章，名字修饰与可见性属性见
第 \ref{ch:cl13-cxx-abi} 章。
\end{abibox}

\section{预处理对语义的影响}\label{sec:cl02-semantics}

预处理不是“不影响语义的便利设施”。它直接改写代码，因而改写行为。两类影响最容易伤人。

第一类是展开形状。带参数的宏按记号替换，替换体里出现几次，实参就求值几次；
实参里有副作用，副作用就重复。宏体少给一对括号，优先级就变了；
实参里带上逗号或关键字，替换结果可能根本不是你想的那个语法成分。

\begin{lstlisting}[style=cstyle,caption={宏把“一次求值”变成“多次求值”，把类型名当成变量名}]
#define SQUARE(x) ((x) * (x))

int read_and_advance();

int probe(int step) {
    int a = SQUARE(read_and_advance());  /* 调用两次：副作用被复制 */
    int b = SQUARE(step + 1);            /* 无副作用，但算了两遍 */
    return a + b;
}

typedef int *IntPtr;
IntPtr cursor;
#define DIST_BASE(x) ((x) - *cursor)     /* 实参是指针时，减法变成指针算术 */

/* 条件编译：同一份源码在不同档位里是不同的程序 */
#ifdef DEBUG_BUILD
#ifdef NDEBUG
# error DEBUG_BUILD 与 NDEBUG 互斥：这一档根本不存在
#endif
#endif
\end{lstlisting}

第二类是分档。条件编译让“同一份源码”在不同档位、不同平台、不同宏定义集合下对应
\emph{不同的程序}。这条事实有一串后果：

\begin{itemize}
  \item 任何“这段代码是否合法”的判断都要先问“在哪一档里”。同一个表达式可能在调试档合法、
    在发布档因某个分支未被选中而从不接受语义检查（未选中的分支连语法都不必合法）。
  \item 头文件里的条件块如果影响类布局（成员数量、对齐、基类选择），一份头在两个档位下
    产生两种内存布局，链接期不会报错，运行期会以最难查的方式崩。这类判定属于 ABI，
    见第 \ref{ch:cl11-layout} 与第 \ref{ch:cl13-cxx-abi} 章。
  \item 语言层对未定义行为的判定，是对“实际编译出来的那一份程序”做的判定。
    档位一旦改变，被判定的对象就变了：某些在甲档里越界的行为，在乙档里根本不存在。
    未定义行为的边界与求值顺序本身属于另一分册的专章，本章只指出这一点：
    拿源码文本去讨论未定义行为而不指定宏集合，那个讨论没有确定结论。
  \item 优化器读的是展开后的程序。它在被宏复制过的表达式上看到的“重复计算”，
    可能触发更多展开；这也是“加了这个宏之后代码变大了”这类现象的机制解释。
\end{itemize}

\begin{guideline}[title={把预处理从接口里赶出去}]
\begin{itemize}
  \item 能用 \Tp{constexpr} 对象、枚举、\Fn{static\_assert} 表达的，不要用宏。
  \item 必须用宏时，函数式宏一律给每个实参加括号、整个宏体加外层括号，
    并且绝不把有副作用的表达式传进去。
  \item 影响布局的条件编译块，用一条编译期断言把期望布局写进代码，让它无法静默漂移。
  \item 头文件里不做策略选择：策略在实现文件里，头里只留接口。
\end{itemize}
\end{guideline}

\section{本章的取舍}\label{sec:cl02-wrap}

\begin{bestpractice}[title={包含、复用与模块的决策顺序}]
\begin{enumerate}
  \item 先看结构：头里有没有被迫外露的实现？把它搬走，收益比任何编译期加速手段都持久。
  \item 再看闭包：把稳定且体积大的头集中成少数几个入口，让包含图变浅。
  \item 需要复用解析结果时，优先预编译头：它不改语义，不引入新的链接期身份，失效后果明确。
  \item 需要把可见性从文件边界搬进语言时，才考虑模块；同时接受接口产物与实现绑定、
    增量沿导入图传播这两笔新成本。
  \item unity 构建只在“包含闭包主导构建时间”且项目对静态符号与宏泄漏有纪律时使用，
    并把它做成可关闭的开关，而不是唯一构建路径。
\end{enumerate}
\end{bestpractice}

预处理是阶段二，模块改变了这一阶段的输入形状，但都不改变阶段三以后对产物的承诺：
接口产物也好、预编译头也好，最后交给链接器或归档工具（\Tool{ar} 一类）的仍然是同一批目标文件。
这一点在第 \ref{ch:cl07-object} 章之后会被反复用到；
构建系统如何组织这些复用，见第 \ref{ch:cl10-cc} 章与第 \ref{ch:cl14-ide} 章。

\chapter{名字解析、重载决议与模板实例化}\label{ch:cl03-resolution}

前两章把源码变成了一个个预处理之后的翻译单元：宏已经展开，头文件已经拼接，剩下的文本仍然只是一串带类型的名字。本章处理前端最抽象的一段工作：把源码里的名字变成有唯一含义的引用，把同名函数区分开，把模板变成真实的函数与类，并把这一整套判定结果交给下游。对编译与链接这条流水线来说，本章的产物不是目标代码，而是\emph{带类型的、已命名的符号集合}：每一个被引用的名字，要么在这里被解释成一个确定的实体，要么在这里就报出错误。

\section{前端做的三件事}\label{sec:cl03-frontend-duty}

忽略词法与语法细节，一个 \cpp{} 前端在本章关心的层面上只做三件事。

第一件是\textbf{名字查找}：给定源码中的一个名字，确定它指的是哪一批候选实体。查找的结果通常不是一个实体，而是一个\emph{候选集合}——同名函数在这个作用域里可能同时有五个声明，它们都要进入下一步。

第二件是\textbf{类型检查与重载决议}：在候选集合里，按实参类型挑出唯一可行的那个，或者诚实地承认挑不出唯一的一个并报二义性。这一步把“名字加实参”变成“对某个具体实体的引用”。

第三件是\textbf{把语义交给下游}：决议完成之后，前端要能回答三个问题——这次调用指向哪个实体、这个实体在链接期应当叫什么名字、它的定义是不是已经可见。这些答案决定了链接器最后看到什么。

\begin{flow}[title={从源码名字到链接期身份：前端的三段准备}]
源码名字 $\to$ 名字查找得到候选集合 $\to$ 重载决议选定唯一实体 $\to$ 需要定义时触发实例化 $\to$ 生成符号身份并登记到目标文件 $\to$ 链接期由符号表证实引用
\end{flow}

回看第~\ref{ch:cl01-pipeline}~章的阶段图可以更清楚地看到分工：前端负责“这个名字是什么意思”，链接器负责“这个意思有没有实现”。前端不会去别的翻译单元里寻找定义，它只是把找不到的东西标成未定义引用交出去；反过来，链接器也绝不重新做名字查找，它只按字节比对符号名。这条边界是理解本章一切内容的钥匙：前端为链接期\emph{准备身份}，链接期只是\emph{核对身份}。

\section{名字查找：作用域链、限定名与依赖名}\label{sec:cl03-lookup}

\subsection{非限定名沿作用域链向外走}\label{sec:cl03-lookup-unqualified}

非限定名（写成 \Tp{count} 而不是 \Tp{ns::count} 的名字）的查找是一条向外返回的路径。理解它最重要的一点是：查找在\emph{第一层命中即停止}，而不是把各层收集到的候选合并起来。

\begin{flow}[title={非限定名的查找次序（函数体内的调用点）}]
块作用域内的局部声明 $\to$ 所在函数的形参与函数级声明 $\to$ 类作用域（成员名，含基类成员） $\to$ 外层命名空间作用域（由内向外逐层） $\to$ 全局命名空间作用域 $\to$ 命中即停，未命中则报“未声明”
\end{flow}

“命中即停”这条规则有两个工程后果值得记住。其一是遮蔽：内层同名声明会让外层实体在该点完全不可见，此时必须用限定名才能取到外层实体。其二是声明顺序敏感：在某一层作用域里，只有出现在使用点\emph{之前}的声明才参与查找；同名实体的后续声明不会加入这个查找点。这解释了一个常见现象——把一个函数声明往后挪几行，就可能让上面原本正常的调用编译不过。

\subsection{限定名把查找起点交给书写者}\label{sec:cl03-lookup-qualified}

限定名（\Tp{a::b}、\Tp{a::b::c}）不做向外返回。它先要求确定最左侧名字的含义（一个命名空间、一个类，还是一个枚举），然后在它的作用域里查找右侧的名字。因为起点是显式给出的，限定名天然绕开遮蔽，也绕开下一小节的候选集合意外扩展。

需要区分的是成员访问的两种语境：通过对象表达式访问成员时，查找先按对象的静态类型进入类作用域；而在类定义内部直接写成员名，走的是上面那条类作用域一层。二者在诊断信息里常让人误判，因为报错文本可能同时提到两个不同的同名实体，而它们的可见性规则并不相同。

\subsection{ADL 把候选集合扩大到实参所属命名空间}\label{sec:cl03-lookup-adl}

参数依赖查找（ADL）是对非限定查找的一次\emph{候选集合扩大}：当一个名字后面紧跟函数调用形式，且实参的类型定义在某个命名空间或某个类里时，那个命名空间（以及类及其相关外层作用域）中的同名函数也被加入候选集合，即使调用点写的是非限定名、也没有用 \Tp{using} 引入。

\begin{lstlisting}[style=cppstyle,caption={ADL 让一个非限定调用选中别的命名空间里的函数}]
namespace graphics {
  struct Pen { int width; };
  void prepare(Pen p);        // 与实参类型同处 graphics
}

void work() {
  graphics::Pen p;
  p.width = 3;
  prepare(p);                 // 非限定调用：ADL 把 graphics::prepare 加入候选
}
\end{lstlisting}

ADL 的价值在于：与某个类型配套的操作（交换、迭代器推进、流式输出这类）不必写限定名就能被找到，这也是让自定义类型接入标准库能力的常规通道。它的代价是查找结果不再只由“调用点写了什么”决定，而由“实参类型定义在哪里”决定。于是下面这个警告必须常驻心头。

\begin{warning}[title={ADL 会选中实现命名空间里的函数}]
实参类型常常定义在库的内部命名空间里。ADL 会把这些命名空间的同名函数一并当成候选，于是看似普通的非限定调用可能突然选中一个内部实现函数，或者引入一个原本不该参与决议的候选，把唯一解变成二义解。这类问题在依赖第三方类型时尤其隐蔽，因为它会随着库内部命名安排的调整而出现或消失，本地代码一行没改。

写法纪律：凡是在全局命名空间存在同名函数风险的调用点，一律用限定名调用真正想要的那一个，例如显式写出 \Tp{std::} 前缀的形式，不要指望非限定写法总能找到正确实体。确实需要利用 ADL 时，只在明确知道实参类型归属命名空间、且该命名空间内没有多余同名实体的前提下使用。
\end{warning}

\subsection{using 引入改变的是候选集合而不是实现}\label{sec:cl03-lookup-using}

\Tp{using} 声明把某个名字注入当前声明所在的作用域层。它的决议影响是结构性的：注入之后，被引入的名字与当前层已有的同名声明\emph{合并成同一个候选集合}，一起进入三段筛选；它不会替换既有声明，也不会把查找终点挪到原命名空间。

这带来两类常见问题。一类是过度引入：把整个命名空间的同名函数都带进候选集合，其中某个重载在某个实参上恰好与本地候选打成平手，二义性凭空出现。另一类是引入位置不当：写在命名空间层与写在函数体层，命中层次不同，可见集合也不同。修复二义性时，“先检查这里引入了什么”通常比“先读错误提示的后半段”更有效。

\subsection{两阶段查找：依赖名与非依赖名}\label{sec:cl03-lookup-twophase}

模板里的名字查找分两次进行，这是 \cpp{} 模板语义里最容易被忽略的一条规则。

\textbf{第一阶段}在模板定义处进行，处理\emph{非依赖名}：含义不随模板参数变化的名字。它们必须在定义点就能查到；查不到就是错误，即使这个模板从未被任何参数实例化过。

\textbf{第二阶段}在实例化点进行，处理\emph{依赖名}：含义依赖模板参数的名字，例如写作 \Tp{T::value} 的成员访问，或者实参类型由 \Tp{T} 决定的调用。这些名字要到参数代入之后才能查找，此时还会额外做一次 ADL，使用\emph{实例化点可见}的声明。

两阶段查找为什么必须存在？如果所有名字都推迟到实例化时才查，定义点里的拼写错误与缺失声明永远不会被报告，模板也就失去了基本可读契约；如果所有名字都在定义点查，为自定义类型扩展辅助函数的整套机制就失效，ADL 也无从谈起。区分依赖与非依赖，是让两者同时成立的唯一办法。

\begin{lstlisting}[style=cppstyle,caption={同一个模板里的两个名字在不同阶段被查}]
namespace base {
  struct Tag {};
  void describe(Tag t);        // 定义在 base 中
}

template <typename T>
void report(T value) {
  describe(value);             // 依赖名：第二阶段查找，ADL 命中 base::describe
  helper(value);               // 非依赖名：第一阶段查找，定义点不可见即报错
}

int main() {
  base::Tag t;
  report(t);                   // 实例化点：describe 可找到，helper 早在定义处报错
}
\end{lstlisting}

这段代码的报错时机很有教育意义：\Tp{helper} 的错误在编译模板\emph{定义}时就出现，而 \Tp{describe} 能否成立要等到 \Tp{report} 被具体参数使用。“删掉一处不相干的调用，另一处报错就消失了”这类现象，根源往往就在这里。

\section{重载决议三段：候选、可行、最佳}\label{sec:cl03-overload}

\subsection{三段的职责}\label{sec:cl03-overload-stages}

重载决议是纯粹的筛选过程：它不产生新的名字含义，只把候选集合削减到零个、一个或“无法判定”的状态。

\textbf{第一段：候选集合。}来源是名字查找的结果——非限定查找命中层的全部同名声明，或限定名指定范围内的全部同名声明——再并上 ADL 带来的额外函数。此段与实参是否匹配无关，只看名字与作用域。

\textbf{第二段：可行函数。}要求形参与实参的个数对得上（考虑默认实参与省略号形式），并且每个实参到对应形参类型之间存在一条\emph{隐式转换序列}。此段的输出是可行函数集合；集合为空时，诊断是“没有匹配的函数”。

\textbf{第三段：最佳匹配。}对可行集合做两两比较：若某候选在所有参数上都不劣于另一候选，且至少在一个参数上严格更优，它胜出。比较判据是隐式转换序列的等级，见下表。

\begin{table}[htbp]
\centering
\caption{隐式转换序列的等级与决议中的相对位置}
\label{tab:cl03-rank}
\begin{tabular}{>{\raggedright\arraybackslash}p{0.18\textwidth} >{\raggedright\arraybackslash}p{0.26\textwidth} >{\raggedright\arraybackslash}p{0.32\textwidth}}
\toprule
等级 & 典型内容 & 决议中的位置 \\
\midrule
恒等与左值转换 & 同一类型；加或去 \Tp{const}；左值到右值 & 最优，几乎不参与比较 \\
提升与标准转换 & 整型提升、浮点到整、指针到布尔、派生到基 & 中间层，等级越低越优 \\
用户定义转换 & 调用转换构造函数或转换运算符 & 次等：标准转换序列优于它 \\
省略号匹配 & 实参落入省略号参数 & 最差，仅当其余都不成立时使用 \\
\bottomrule
\end{tabular}
\end{table}

上表只给相对次序，不给分数：决议从不计算总分，它做逐参数比较。同一等级内部还有更细的次级判据，例如保留原始取值范围的转换优于丢失精度的转换、跨越更少派生层次的优于跨越更多层次的。这些细节在讨论跨编译器兼容性时会再次出现（参见第~\ref{ch:cl10-cc}~章）。

\subsection{为什么“各优一处”必然退化成二义性}\label{sec:cl03-ambiguity}

第三段的比较是\emph{偏序}而不是全序。设两个候选 \Tp{f1} 与 \Tp{f2}：若 \Tp{f1} 在第一个参数上更优、\Tp{f2} 在第二个参数上更优，就不存在“谁整体更好”的判定。决议没有加权求和机制，也没有“多数参数胜出就当选”的规则，此时前端报二义性，并把两个候选以及各自在哪个参数上更优一并写进诊断。

理解这一点，对修复动作有实际意义：二义性不是靠反复阅读错误信息解决的，而是靠\emph{改变比较结构}解决的。三种常见修复分别落在三段的不同位置——用显式转换把某个实参拉回恒等序列（改第二段），删掉 \Tp{using} 引入或改用限定名（缩小第一段的候选集合），调整重载集合本身（改变第一段的命中层内容）。知道动作作用于哪一段，才知道它为什么有效。

\subsection{完全平手时的两条决胜规则}\label{sec:cl03-tiebreak}

偏序比较有时会给出“不分优劣”的结论。此时还有两条决胜规则，它们的性质与前三段完全不同：不再比较转换序列，而是比较候选\emph{自身的种类}。

其一是\textbf{非模板优先}：若两个候选在每个参数上取的转换序列完全同等，其中一个是普通函数、另一个是模板实例，则选普通函数。直觉在于普通函数的签名由作者写死，模板实例的类型由推导得到，前者更值得信任。

其二是\textbf{更专用者优先}：若两个候选都是模板实例且转换序列同等，则对两个模板做偏序比较，能覆盖另一个候选全部参数形式的更专用者胜出。

两条规则都用尽后仍然平手，才报二义性。这一层对维护的实际意义是：往同名集合里新增一个普通函数，可能瞬间压过一批模板实例；反过来新增一个更专用的模板，也可能悄悄改变既有调用点的选择。两种改动都不必修改任何调用点，却会改变决议结果，因而也会改变第~\ref{ch:cl13-cxx-abi}~章所关心的符号身份。

\section{模板实例化：时机、位置与收拢}\label{sec:cl03-instantiation}

\subsection{隐式实例化点与“只在用到时生成”}\label{sec:cl03-implicit}

函数模板与类模板的成员都不在模板定义处生成。隐式实例化发生在\emph{需要其定义的那个点}：需要函数体时实例化函数定义，需要对象大小时实例化类并生成成员声明，取成员地址或调用成员时实例化那个成员的定义。类模板的成员函数本身仍是模板，它们按各自的调用点逐个实例化。

“只在用到时才实例化”带来三个可观察后果。

其一是错误延迟：未被使用的成员即使完全写错也不会被报出；一旦有代码使用它，一大串错误会突然出现在看似无关的调用点上。其二是编译代价的收益：一个被大量包含却只少量使用的类模板，不会在每个翻译单元生成全部机器码。其三是链接形态：同一条实例化规则在不同时机执行，会让同一个实例在链接期以不同身份出现，见本节之后的结论框。

\subsection{显式实例化定义与显式实例化声明}\label{sec:cl03-explicit}

隐式实例化会在“每个用到它的翻译单元”里各发生一次。对小型实例这无所谓，对大型实例代价明显：同样的代码在每个单元各编译一遍，也各自向链接器提交一份可折叠的定义。

\textbf{显式实例化定义}解决“只编译一次”：在某个单元里用具体类型实参请求把实例生成出来，其余单元不再自行生成它。\textbf{显式实例化声明}解决“其余单元知道别人会生成”：以声明形式告诉每个包含者，该实例由别处提供，不要在本地隐式生成，只需留下引用等待链接。

两者配对使用后，实例就从“分散在多个翻译单元的可折叠符号”收拢成“一个单元里的唯一定义，加若干未定义引用”。维护代价也随之而来：显式实例化清单是一份需要人工同步的契约，新增类型实参时若忘了补上，问题会以链接期未定义引用的形式出现，而不是编译期的一句提示。

\begin{abibox}[title={实例化时机决定链接期身份}]
同一个模板实例，若由隐式实例化在多个翻译单元各自生成，链接器看到的是多份内容相同、允许折叠的定义；若由显式实例化定义集中生成，链接器看到的是一份唯一定义加若干未定义引用。前者依赖链接阶段的可重复定义折叠能力（参见第~\ref{ch:cl08-static-link}~章），后者把折叠问题提前消灭在编译阶段的安排里。两种形态在语义上等价，在链接耗时、目标文件体积与诊断出现时刻上并不等价。
\end{abibox}

顺带一提，头文件里的 \Tp{inline} 函数与类内定义的成员也走“多份可折叠定义”这条路：它们不是模板，却因为同样可能在每个单元各生成一份而需要折叠规则。把这条线索延伸到符号层面，正是本章与第~\ref{ch:cl08-static-link}~章的交界。

\section{符号身份的生成时刻}\label{sec:cl03-mangling}

\subsection{修饰名由前端产生，满足与否由链接期证实}\label{sec:cl03-mangling-when}

链接器比较的是符号名字节。在 C 语言里名字与符号基本一对一，链接期不需要额外信息；\cpp{} 里同名函数、成员函数、模板实例必须各自有独立身份，于是只能由\emph{懂得类型的前端}在生成引用时把类型信息编码进符号名，这个编码习惯上称为修饰名。

时机是关键：修饰名在决议完成后立刻确定，而“这个符号能否被满足”要等到链接期由符号表证实。前端无法验证另一个翻译单元里有什么，它只保证一件事——我对这次调用的理解，会以这个名字表达出来。修饰规则本身属于跨编译器与跨平台接口约定的范畴，本章不展开其编码形式（参见第~\ref{ch:cl13-cxx-abi}~章）；这里只需要一个结论：决议结果变化，符号随之变化。

\subsection{最小例子：同名重载在不同翻译单元里的机理}\label{sec:cl03-mangling-example}

下面用 \Tp{cppstyle} 风格的两个翻译单元，配一段示意符号表输出，说明“同一个函数名在链接期是两个符号”的机理。

\begin{lstlisting}[style=cppstyle,caption={翻译单元甲：声明齐备，但只定义 long 版本}]
namespace app {
  int  total(int n);
  long total(long n);
}

long app::total(long n) { return n * 2; }   // 本单元只定义 long 版本

int probe(int n) {
  return app::total(n);      // 实参为 int：决议选中 int 版本，留下它的引用
}
\end{lstlisting}

\begin{lstlisting}[style=cppstyle,caption={翻译单元乙：只定义 int 版本}]
namespace app {
  int total(int n);
}

int app::total(int n) { return n + 1; }     // 本单元只定义 int 版本
\end{lstlisting}

单元甲的调用点做的是一次完整决议：\Tp{n} 是 \Tp{int}，两个候选都可行，\Tp{int} 版本在第一个参数上取恒等序列因而胜出。甲并没有定义 \Tp{int} 版本，于是它在目标文件里留下一条未定义引用，名字是 \Tp{int} 版本对应的修饰符号。

\begin{lstlisting}[style=outstyle,caption={符号表输出（示意；实际修饰形式依平台与接口约定而定）}]
app 模块（示意）
  甲  未定义引用  需要 total(int)  的修饰符号
  甲  已定义符号  提供 total(long) 的修饰符号
  乙  已定义符号  提供 total(int)  的修饰符号
\end{lstlisting}

链接期逐条比对：甲的未定义引用恰好由乙的定义满足，程序成立。反过来，若把调用点的 \Tp{n} 改成 \Tp{long}，决议就选中 \Tp{long} 版本，甲不再留下跨单元引用，链接甚至不需要乙。可见“改一个实参类型”改变的不只是类型检查结论，而是链接期看到的\emph{符号集合}本身。

\begin{warning}[title={本地改动可能悄悄改变跨单元身份}]
修改形参类型、增加默认实参、把函数移进命名空间、改变 \Tp{const} 修饰，都可能改变修饰名，从而让原本配对的引用与定义不再同名。这类问题的报错出现在链接期（未定义引用或重复定义），而不是编译期，因此“某个翻译单元能单独编译通过”毫无安慰作用。涉及跨单元接口的签名调整，应当连同使用这些接口的所有单元一起重新构建；只重编其中一半，得到的往往是一个看起来能链接、实际行为错误的产物。
\end{warning}

\section{诊断为什么需要保留源码信息}\label{sec:cl03-diagnostic}

如果前端的产物只是符号集合，那么把中间结构丢掉、只留符号名就够了。诊断不是这样。一次“引用了未声明的名字”要成为可用信息，必须能回答三个问题。

第一，\emph{拼写}：报错点附近是否存在一个高度相似的名字？这是拼写建议的基础，它要求保留源码里出现过的原始标识符文本，而不是已决议实体的内部编号。第二，\emph{范围}：错误指向表达式中的哪一段？“哪个实参导致没有匹配”这类提示要能定位到具体子表达式，这要求保留词法单元的位置信息。第三，\emph{候选来源}：这一批函数为什么进入候选？作用域链命中层、\Tp{using} 引入、ADL 扩大的命名空间是三种不同来源；把它们区分开，使用者才可能选对修复动作。

这三类信息都必须在决议\emph{过程中}被保留，因为决议完成后候选集合就消失了：未被选中的候选不留痕迹，链接器也从不关心谁曾经参与比较。同一份分析结果还被开发工具消费：补全、悬停提示、跳转定义、重命名重构，用的都是“名字查找加重载决议”这套判定，只是被要求更快、可增量、能容忍不完整的源码。这条线在第~\ref{ch:cl14-ide}~章展开。至于未定义行为的判定、代码风格规范与性能测量，属于别的分册主题，本章不作承诺。

\begin{bestpractice}[title={模板报错的定位三步}]
第一步，先找实例化点，再读错误链。模板诊断通常会沿实例化位置一路带出多段信息，被标注为“在实例化中”的那一段才真正触发了代码生成，改动应从那里考虑，而不是从错误列表的第一行开始。

第二步，判断错误属于哪个阶段。在定义点就报出的，是非依赖名查找失败或语法问题；只在实例化之后报出的，是依赖名查找失败、可行函数集合为空，或者被使用的成员自身有问题。两类修复动作完全不同：前者改模板，后者改调用点或参数类型。

第三步，缩小实例再验证。把模板参数换成最简单的具体类型，或把被实例化的成员削减到只剩出问题的那个，重新编译一次。实例规模下降后，错误链常常从几十行缩到几行，剩下的定位工作回到普通代码的难度。
\end{bestpractice}

\section{案例分层：一次调用的诊断来自哪里}\label{sec:cl03-case}

\begin{lstlisting}[style=cppstyle,caption={案例源码：两个同族重载与一个参数个数不符的模板}]
namespace svc {
  struct Money { Money(int v); };
  struct Count { Count(int v); };

  void log(Money m);
  void log(Count c);

  template <typename T>
  void log(T first, T second);      // 需要两个实参
}

void emit() {
  svc::log(42);                     // 限定名：候选集合就是 svc 中的三个 log
}
\end{lstlisting}

\begin{note}[title={案例：诊断的三段来源与一次都没发生的实例化}]
把这一处调用按三段拆开，可以看清每一段诊断各自来自哪里。

第一段\emph{候选集合}来自前端的限定名查找：三个同名声明全部进入集合，其中包含那个模板。这一步的产物只受作用域影响，与实参类型无关；若调用点写成别的作用域里的名字，或删掉一个重载，候选集合立刻改变。

第二段\emph{可行性检查}把模板筛掉了：它需要两个实参，调用点只给一个，元数不符，于是根本不存在隐式转换序列。这一步不报错，只把集合从三个削减为两个可行的。

第三段\emph{最佳匹配}才是错误本身：\Tp{int} 到 \Tp{Money} 与 \Tp{int} 到 \Tp{Count} 都只含一次用户定义转换，两者同等级，比较不出优劣；决胜规则也帮不上忙，因为剩下的两个候选都是普通函数。于是报二义性，诊断文本同时列出两个候选。

值得强调的是\emph{没有发生}的事：模板从未被选定，所以实例化一次都没有发生，它的函数体即使写错也不会被报出。反过来说，若把模板改成单参数形式，它对 \Tp{int} 取恒等序列，立即压过两个需要用户定义转换的候选；此时模板胜出、实例化开始，体内的依赖名问题才第一次露面。同一次调用于是可能带出两段来源完全不同的诊断：一段属于决议第三段，一段属于实例化。

修复动作的落点同样分层：加显式转换或换实参类型（作用于第二、三段，立刻把两个同级序列拉开）；删掉或移走一个重载（作用于第一段，缩小候选集合）；调整模板签名（看似与二义性无关，实际会引入实例化诊断这一全新类别，往往让报错变多而不是变少）。
\end{note}

这个例子的实际结论还取决于调用点能否看到 \Tp{svc} 里的名字：若那里既没有引入也没有使用限定名，报错就是“未声明的名字”，属于第一段；能看到之后，报错才升级为二义性，属于第三段。同一个书写形式，报错文本、报错阶段、修复层位全都不同——这正是本章主线的注脚。

\section{语言构造与阶段对照}\label{sec:cl03-matrix}

下面的表把本章讨论过的构造放到同一条时间轴上：决议或实例化发生在哪一段，问题最早在什么时候可见。

\begin{table}[htbp]
\centering
\caption{语言构造：决议或实例化阶段与错误可见时刻}
\label{tab:cl03-matrix}
\begin{tabularx}{0.96\textwidth}{>{\raggedright\arraybackslash}p{0.20\textwidth} >{\raggedright\arraybackslash}p{0.24\textwidth} X}
\toprule
构造 & 决议或实例化发生的阶段 & 错误最早可见的时刻 \\
\midrule
重载函数 & 使用点，名字查找后立即三段筛选 & 编译期：二义性与无匹配 \\
函数模板 & 使用点先参与重载比较，选定后在实例化点生成定义 & 编译期；函数体内部错误推迟到实例化 \\
类模板成员 & 类实例化生成声明，成员在各自调用点实例化 & 编译期；未使用成员的错误可能永不出现 \\
\Tp{inline} 函数 & 使用点决议；定义在各单元分别生成 & 语义问题在编译期，多份定义问题到链接期 \\
静态数据成员 & 类实例化时生成声明，定义需单独存在 & 缺少定义在链接期暴露 \\
虚表 & 分派决议推迟到运行时；表随类定义构建 & 签名不匹配在编译期，覆盖不符在运行期 \\
类型信息实体 & 随类定义与使用点生成，跨单元需统一 & 引用无法满足时在链接期 \\
\bottomrule
\end{tabularx}
\end{table}

这张表值得反复看的不是具体某一行，而是右列的整体分布：越靠近语言抽象高层的构造（决议类问题），错误越早可见；越依赖跨单元一致性的构造（多份定义、单一实体、表与类型信息），错误越往链接期甚至运行期滑动。看到编译期就报出的错误，可以推断它不来自链接安排；反过来，链接期的未定义引用几乎总能在左列找到对应构造。

\section{本章要点}\label{sec:cl03-keypoints}

\begin{itemize}
  \item 前端负责名字的含义，链接器负责含义的实现；本章的产物是带类型的、已命名的符号集合。
  \item 非限定名查找在第一层命中即停；限定名把起点交给书写者；ADL 是候选集合的扩大，需要用限定名纪律来约束。
  \item 两阶段查找区分依赖名与非依赖名，决定了一条错误在模板定义点报还是在实例化点报。
  \item 重载决议是偏序比较，不存在加权求和；两个候选各优一处必然退化为二义性。
  \item 显式实例化定义与声明把实例从多个单元收拢到一个单元，实例化时机因此决定链接期符号形态。
  \item 修饰名由前端在决议完成后产生，签名层面的本地改动会把问题推到链接期才暴露。
  \item 拼写建议、错误范围与候选来源三类诊断信息，必须在决议过程中保留，开发工具消费的是同一份结果。
\end{itemize}

\begin{note}[title={本章的产物与下一章的入口}]
到这里，源码里的每个名字都有了确定的实体归属，每个被使用的模板实例都有了生成安排，每个引用都带上了可以在链接期比对的符号身份。这份“带类型的符号集合”仍然紧贴 \cpp{} 的语言结构：它知道重载、知道模板参数、知道访问控制。下一章开始，这些结构将被逐步丢弃，换成与语言无关、只关心数据流与控制流的中间表示（参见第~\ref{ch:cl04-ir}~章）；届时本章保留的源码信息不再影响正确性判定，而只服务于诊断与开发工具。
\end{note}

\part{中端与后端}
\chapter{中间表示与静态单赋值}\label{ch:cl04-ir}

前端把源码翻译成某种内部形态，后端把某种内部形态翻译成汇编。这句话里最容易被忽略的是
那个“某种”：两端真正对话的语言既不是源码也不是汇编，而是一份被精心设计过的
\textbf{中间表示}（Intermediate Representation，下称 IR）。本章讨论这份中间表示为什么必须存在、
它通常包含哪些要素、静态单赋值形式为什么是绝大多数优化的前提，以及\cpp{} 的对象语义如何在这层
被展开。工具细节与开关名不在本章范围内，本章只谈结构。

\section{为什么需要中间表示}\label{sec:cl04-why}

\subsection{把乘法拆成加法}\label{sec:cl04-why-multiply}

设前端有 $M$ 个（不同语言、不同方言、不同标准级别各算一个），后端有 $N$ 个（不同 CPU 架构、
不同操作系统、不同目标格式各算一个）。若每个前端直接对接每个后端，需要维护 $M \times N$
套翻译逻辑；插入一层共同的目标语言之后，工作量变成 $M + N$ 套。

这个算术差别只是表面理由，更硬的是两条：优化只需写一遍——常量折叠、循环变形、公共子表达式
消除与目标机器无关，放在中间层一次做好，所有前端与后端同时受益；移植成本被隔离——新增一个
后端不必理解任何一门源语言的语法细节，只需理解中间层与目标机器的映射关系。

\subsection{三族实现各有几层}\label{sec:cl04-why-families}

真实实现并不把这套分层结构统一公开，因此讨论只能停留在层次数量与可见性这两个安全话题上。
表~\ref{tab:cl04-layers} 给出三种常见形态的对照，措辞刻意保守：凡文档没写清楚的层次，
本表不猜测，判断依据以各实现自己的文档为准。

\begin{table}[htbp]
  \centering
  \begin{tabular}{>{\raggedright\arraybackslash}p{0.18\textwidth} >{\raggedright\arraybackslash}p{0.20\textwidth} >{\raggedright\arraybackslash}p{0.20\textwidth} >{\raggedright\arraybackslash}p{0.20\textwidth}}
    \hline
    \textbf{形态} & \textbf{机器无关层} & \textbf{机器相关层} & \textbf{对外可见性} \\
    \hline
    一族以两三种中间语言分层 & 先做一种偏高级、类型信息较多的表示，再降到偏低级、语法更紧的表示 & 最后是一种接近指令、带虚拟寄存器与分配痕迹的表示 & 偏高的那层能以文本产出，偏低的那层通常只在调试转储里出现 \\
    \hline
    一族把两层明确分开 & 语言无关表示与机器相关表示是两套数据结构，中间有一次显式切换 & 切换之后仍使用无限虚拟寄存器，把冲突留给后续阶段 & 语言无关那层既是内存结构，也可持久化成文本产物 \\
    \hline
    另一族不公开内部表示 & 从产物可以推断存在一个语言无关阶段 & 只能从汇编输出反推其形状 & 不对外提供可读的中间表示产物 \\
    \hline
  \end{tabular}
  \caption{三族实现的表示层次对照（概念性归纳，层次命名与可见性以该实现文档为准）}
  \label{tab:cl04-layers}
\end{table}

\begin{flow}[title={从源码到汇编的层次序列}]
抽象语法树 $\rightarrow$ 高层语言无关表示 $\rightarrow$ 优化管线 $\rightarrow$
机器相关表示 $\rightarrow$ 汇编文本
\end{flow}

抽象语法树仍带语言语法节点，属于语言相关层；高层语言无关表示把语法压平成一小组指令与
显式类型，语言相关信息被折叠进属性；中间的优化管线（第~\ref{ch:cl05-passes} 章）在这一层
反复改写；机器相关表示才引入目标机器约束；最后吐出汇编文本（第~\ref{ch:cl06-codegen} 章）。
\textbf{本章全部讨论都发生在第三格}，也就是高层语言无关表示这一层。

\section{语言无关表示的共性要素}\label{sec:cl04-common}

不同实现的 IR 语法差别很大，组成要素却高度一致。理解下面四件事，就能读懂大多数 IR 文本。

\subsection{三地址风格的指令}\label{sec:cl04-three-address}

表达式树在 AST 层可以任意嵌套，IR 层几乎都把它拉平成“两个源操作数加一个目的操作数”的形式。

\begin{itemize}
  \item \textbf{单定义。}一条指令产生一个值，该值可被后面引用多次。嵌套表达式被拆成一串
        有序小步，每一步命名一个临时值。
  \item \textbf{显式加载与存储。}内存访问不再是隐式副作用，而是独立一条指令。只有把读内存
        与算数值区分开，别名分析才有对象可分析。
  \item \textbf{分支取代控制流语法。}\Tp{if}、循环、\Tp{switch} 统统变成条件分支加跳转，
        控制流图形化之后才谈得上循环识别。
\end{itemize}

\subsection{显式类型}\label{sec:cl04-explicit-types}

IR 层每条指令上的类型标注必须无歧义。源码里 \Tp{1 + 2} 的类型由语言的整型提升规则决定，
到了 IR 层它已经是一次确定宽度的整数加法。这一层的类型至少承担三件事：决定指令的位宽语义
（截断、扩展、饱和）、决定对象的大小与对齐、决定调用参数如何拆分与传递。类型信息一旦在这层
丢失，后端只能靠猜，而猜出来的宽度差异会直接变成第~\ref{ch:cl11-layout} 章里那种
“同一份源码在不同目标上大小不同”的问题。

\subsection{基本块与控制流图}\label{sec:cl04-basic-block}

\textbf{基本块}是一段只有唯一入口（第一条指令）与唯一出口（最后一条终止指令）的线性指令
序列，中间不允许跳入或跳出。终止指令把一个块接到若干后继块上，这些边构成\textbf{控制流图}。
可达性、支配关系、循环识别、活区间，全都建立在块与边之上。读 IR 时最实用的动作是先只读
每个块的最后一行，把控制流图骨架画出来，再逐块填充内部指令。

\subsection{函数级信息：调用约定与属性}\label{sec:cl04-attrs}

函数在 IR 层是一个带完整签名的对象：参数类型与顺序、返回值形式、调用约定、若干标志（是否
可能抛异常、是否叶子函数、参数指针是否不重叠、对齐要求等）。这些标志不是注释，它们参与决策：
标志缺失，对应变换立刻被关掉。第~\ref{ch:cl10-cc} 章说明调用约定如何在参数传递与寄存器保存两个方向上约束签名。

\subsection{为什么无限虚拟寄存器是合理的}\label{sec:cl04-vregs}

机器无关层的每条指令都可以写一个新名字，寄存器数量看起来不受限制。这不是实现偷懒，
而是刻意把两个问题拆开：

\begin{enumerate}
  \item \textbf{值的生命周期问题。}哪个值在哪些路径上仍活着、能否被替换、能否被常量传播，
        只依赖数据流，与目标机器有几个寄存器无关。
  \item \textbf{有限资源的分配问题。}把任意多个活值塞进固定数量物理寄存器，是图着色或
        线性扫描意义上的组合优化问题。
\end{enumerate}

把第二类问题推迟到寄存器分配阶段，好处是第一类问题的解法不会持续被硬件约束干扰，而且分配
代价\textbf{可测量}：分配器只在信息足够时插入搬运，搬运本身又成为后续优化的输入。这一分工在
后端如何展开是第~\ref{ch:cl06-codegen} 章的主题；机器无关层里很多临时值不是浪费，是留给分配器的自由度。

\section{静态单赋值}\label{sec:cl04-ssa}

\subsection{定义}\label{sec:cl04-ssa-def}

\textbf{静态单赋值}（Static Single Assignment，SSA）要求每个标量值在整个函数里
\textbf{只被定义一次}，使用处只能引用一个明确的定义。未采用 SSA 时，一个变量可能被多次赋值，
于是“此处读到的值由哪一次赋值产生”必须靠分析才能回答。进入 SSA 后，每次赋值产生一个新名字，
合并点上写出一条 \Tp{phi} 指令把候选来源列全：

\begin{lstlisting}[style=outstyle]
x0 = a                  // entry block: one definition
if c goto L else G
L: x1 = b              // branch block: the other definition
G: x2 = phi(x0, x1)    // merge point, sources per predecessor
y  = phi(x0, x2)       // later uses read a single value
\end{lstlisting}

这段是刻意回避具体语法的示意文本，不对应任何真实工具的输入输出格式。真实 IR 里的块名与
\Tp{phi} 书写顺序都必须严格遵守该实现的定义规则，不能凭直觉套用。

\Tp{phi} 不是运行时会执行的机器动作，它表达一个静态事实：此处的值来自哪几条前驱边中的
哪一条，由到达路径决定。合并点上必须写 \Tp{phi}，否则“这个值是谁写的”就没有唯一答案。
合并指令的正确插入位置由支配边界理论给出，这是 SSA 构建算法的核心，本章不展开。

\subsection{为什么绝大多数优化要在 SSA 上跑}\label{sec:cl04-ssa-why}

\begin{table}[htbp]
  \centering
  \begin{tabular}{>{\raggedright\arraybackslash}p{0.20\textwidth} >{\raggedright\arraybackslash}p{0.30\textwidth} >{\raggedright\arraybackslash}p{0.30\textwidth}}
    \hline
    \textbf{关注点} & \textbf{普通形式} & \textbf{SSA 形式} \\
    \hline
    这个值是谁写的 & 向前扫所有路径，必要时做到达-定义分析 & 每个使用点直接指向唯一定义，关系是结构里的边 \\
    \hline
    这个定义被谁用 & 需要全函数扫描 & 定义的 use 列表直接可读 \\
    \hline
    常量与复制传播 & 多次迭代才收敛，容易漏 & 沿 use 边走一遍即可，代价接近线性 \\
    \hline
    死代码判定 & 要判断后面是否还有赋值覆盖它 & 一个定义没有 use，直接删 \\
    \hline
    公共子表达式 & 要先确认中间没有被重新赋值 & 两个候选的定义可比较，无需模拟中间赋值 \\
    \hline
  \end{tabular}
  \caption{同一批问题在两种形式下的处理难度对照（概念性归纳）}
  \label{tab:cl04-ssa-vs-plain}
\end{table}

一句话概括这张表：SSA 把“需要算出来的程序性质”变成“可以直接查表的结构”。定义-使用链是
IR 对象自身携带的边，因此传播类变换不必每次重做一遍数据流分析，既降低实现难度，也降低
运行代价。优化管线里几乎所有标量层面的变换都要求输入是 SSA，有些甚至会在 SSA 被破坏之后
先把它重建出来再继续。

\subsection{SSA 的还原发生在哪一步}\label{sec:cl04-ssa-destruct}

\Tp{phi} 不能直接映射到任何一条真实指令：真实机器没有“从一个集合里挑一个来源”这种操作。
因此 SSA 必须在某个时刻被\textbf{析构}：把每条 \Tp{phi} 变成在对应前驱边上插入的一次复制，
并把值落到具体存储位置（栈上的临时槽位，或直接落到寄存器）。

这一步的位置很关键：太早做，后面的优化失去定义-使用链这个抓手；太晚做，尤其推到寄存器分配
之后再做，边上插入的复制会与已完成的分配互相干扰，常常需要返工。常见做法是放在机器无关层
走完最后一批依赖 SSA 的变换之后、进入机器相关层之前或刚进入时。换句话说，SSA 是给优化器
看的形式，不是给机器看的形态；第~\ref{ch:cl06-codegen} 章接手时，\Tp{phi} 已变成复制与槽位。

\begin{bestpractice}[title={读 SSA 的顺序}]
先只看每个块的终止指令，把控制流图与所有合并点标出来；再单独看 \Tp{phi}，把每个合并点的
候选来源与对应前驱边一一对上；最后才读块内的算术与访存指令。按这个顺序读，几乎不会把
跨路径的值来源弄错。
\end{bestpractice}

\section{标量类型、指针与聚合类型}\label{sec:cl04-types}

\subsection{标量与指针}\label{sec:cl04-types-scalar}

整数类型在 IR 层用显式位宽命名，例如 \Tp{i32}、\Tp{i64}、\Tp{i8}；浮点类型按目标支持的精度
分成单精度、双精度与扩展精度几档；指针类型写成“指向某种类型”的形式。这里的类型是
操作数属性，而不是语言类型系统：语言侧的 \Tp{short}、\Tp{int}、\Tp{long} 在同一目标上可能
被折叠成同一个位宽，语言侧重载与隐式转换在这里已是一次次显式指令。聚合类型的\textbf{布局}
（大小、成员偏移、填充、对齐）不在 IR 层重新发明，而由目标 ABI 决定，第~\ref{ch:cl11-layout} 章
给出完整规则；IR 层只把成员访问表达成一次求址运算，而不是靠符号名直接下标。

\subsection{聚合类型如何被拆成标量}\label{sec:cl04-scalar-replace}

源码里一个结构体变量在 SSA 层的理想形态是“若干彼此独立的字段值”，各有自己的定义与使用链。
把结构体拆成字段级操作正是这一类 pass 的职责，其步骤大致是：

\begin{enumerate}
  \item 确认该聚合对象的地址从未泄漏（没有被取址逃逸、没有被当作整块传给外部函数）；
  \item 把对字段的读写改写为对独立标量值的读写；
  \item 把这些新标量并入 SSA，让它们获得与其他局部变量同等的传播与消除待遇；
  \item 只有拆分失败时才保留整块形态，交给后端按聚合访问处理。
\end{enumerate}

收益很实在：字段化之后，一个从未被读的字段立刻表现为一个没有 use 的定义，可以被删除；
结构体填充字节也不再需要在每次整体拷贝时被搬运。反过来，一旦地址逃逸，整块拷贝必须保守保留，
这也是第~\ref{ch:cl06-codegen} 章里“看似等价的赋值为什么搬运了额外字节”问题的根源之一。

\subsection{不透明指针化的动机}\label{sec:cl04-opaque-ptr}

早期风格的指针类型会把指向元素的完整类型一路写下来，形成任意深度的链。这种精细标注在
类型检查阶段有用，代价是同一份事实被每条指令重复携带一遍：访存地址、元素宽度、求址结果
类型、函数签名里的形式类型，说的往往是同一件事。

演进方向是让指针成为不携带元素信息的统一类型，只在\textbf{真正需要位宽的那一条指令}上标注
访问宽度。动机可以归结为一句话：\textbf{为了减少类型系统在指令里的重复}。重复越少，指令形式的
组合爆炸越少，前端与优化器之间需要同步的类型细节也越少。这属于哪一代接口、以什么方式选择
形态，本册不给出版本号与名称，请以该实现文档为准。

\section{\cpp{} 语义如何落到中间表示}\label{sec:cl04-cxx-semantics}

一门带对象生命周期、继承与异常的语言，必须在机器无关层被展开成一组可被优化器处理的显式操作，
同时又不能展得太过火。这一节是全章最有压力的部分。

\subsection{构造与析构成为显式调用}\label{sec:cl04-ctor-dtor}

对象构造与析构在 IR 层表现为一次普通调用：被调函数指向该类型的构造函数或析构函数，第一个
参数是对象地址；带虚基类或需要补正的场景走带变体说明的形式。成员按声明顺序构造、在同一
作用域内逆序析构，这一顺序在 IR 层就是一串已经排好的调用序列。

关键收益是：构造与析构一旦是可见调用，优化器就能对它做内联，能把相邻的“构造后马上拷贝再
析构”合并，能在确认对象生命周期完全内联时把整块内存操作折叠成若干标量赋值。代价是 IR
体积变大、指令数变多、可读性下降。

\subsection{虚调用成为取指针加间接调用}\label{sec:cl04-virtual-call}

虚函数调用的展开形态是三步：从对象头取出表指针、按索引取出函数指针、发起一次间接调用。
这三步在 IR 层彼此独立可查，因此优化器有机会把索引常量化的取表动作折叠成直接引用，或者
在能证明动态类型时把整条链替换成一次具体调用。更进一步的推测性选择会先做一次类型判断，
成功走直接调用、失败退回原来的间接调用；这类改写必须配合第~\ref{ch:cl13-cxx-abi} 章的规则，
才能既改写又保持链接兼容。类型信息（多态对象的动态类型描述、跨共享库所需的类型识别）在
ABI 层定义，IR 层只负责引用它们，这里不复述名字修饰与表的字节布局。

\subsection{异常边界成为带清理动作的着陆区}\label{sec:cl04-landing-pad}

可能抛出异常的调用在 IR 层带有两条出路：正常返回边，以及指向\textbf{着陆区}的异常边。
着陆区拿到异常对象与选择器，据此决定是清理本地对象、继续向外传播，还是进入某个处理块。
清理动作是前面那串析构调用的另一份排列，它与正常退出路径上的析构调用必须\textbf{逐字一致}，
否则会表现为抛出时少析构了一个对象。

这个结构还决定一件事：函数里只要存在任何可能的异常路径，寄存器分配与调度就要保守得多，因为
控制流可能从任意调用边跳到着陆区。是否采用表驱动的零成本异常模型影响外层结构，IR 层看到的仍
是调用带异常边这一形态；这些表落在哪个段、由谁注册取决于目标格式，第~\ref{ch:cl07-object} 章讨论。

\subsection{语言层契约与优化能力之间的取舍}\label{sec:cl04-asif}

把\cpp{} 语义降到 IR，本质上是在“保留语言契约”与“给优化器留出改写空间”之间找平衡。
语言标准用未定义行为把一部分判断权交还给实现：编译器有权假定某些坏情况不会发生，并据此化简；
这类假定的完整边界在别册的\cpp{} 分册里以 as-if 论断详细讨论，本章不交叉引用。结构层面的事实是：
\textbf{若降到 IR 时过早丢掉可观察行为的信息，后端就无从判断哪些改写是被允许的}，因此 IR 层还要携带“这个值可能被外部观察”一类属性，共同定义优化器的合法边界。

\begin{warning}[title={IR 上的两处常见误解}]
其一，虚表指针被优化掉不等于虚函数机制被取消：链接期仍可能有别的翻译单元通过同一多态类型
发起间接调用，类型信息与表的发射决策不完全取决于本函数内的内联结果。
其二，异常展开的 IR 形态与最终落盘的运行时表不是同一份东西：IR 里的清理序列会被整理成表项，
表项形式由目标格式与异常模型决定；看到 IR 里有两条出路，不等于产物里一定有对应的显式跳转。
\end{warning}

\section{内建函数：语言扩展进入 IR 的正门}\label{sec:cl04-intrinsics}

有一类操作同时满足三个条件：语言标准不规定它、目标机器上它有直接对应物、优化器知道它的语义。
例如“假定这个地址按某个字节数对齐”“这个加法不允许溢出”“在这里插入内存屏障，顺序不可跨越”
“把这块内存当整块搬运”。这类操作通常以\textbf{内建函数}一族的身份进入 IR：写法是调用、
名字是函数，但语义不是普通函数调用。它们的价值来自“既像调用、又可被理解”这一双重身份：

\begin{itemize}
  \item \textbf{对前端}：新增语言扩展不必给指令格式扩条，一次内建函数调用即可表达，
        IR 的形式集合保持稳定。
  \item \textbf{对优化器}：语义被登记在案，优化器知道它是纯的、可交换的、可折叠成常量、
        可证明无溢出，因而能安全地参与删除与合并。
  \item \textbf{对后端}：目标机器有对应指令就展开成指令，没有就退化成一次真实外部调用，
        正确性不依赖目标机器。
\end{itemize}

正因如此，内建函数是语言扩展进入 IR 的正门：新语义先进这一族，而不是先改指令格式。
用得过多，源码就与具体优化器的假设绑死，可移植性与可预期性随之变差。

\begin{guideline}[title={内建函数使用原则}]
只在语义确实无法用普通语言构造表达、或明确需要向后端传递一个额外假定时使用；每个使用点都
要写下“去掉它程序是否仍然正确”的答案。假定类内建函数尤其危险：它把正确性的举证责任从
实现转移到使用者身上。
\end{guideline}

\section{把表示层变成可观察的产物}\label{sec:cl04-observe}

IR 只有能被打印出来，才从实现内部数据结构变成工程师手里的工具。三种常见做法：

\begin{enumerate}
  \item \textbf{把 IR 作为文本产物输出。}前端与优化之间的一次文件往返，让 IR 成为可版本化、
        可比较、可测试的工件。
  \item \textbf{保留值名以便阅读。}多数实现允许转储时保留来自源码的名字，使 \Fn{sum} 里的
        累加值写成有语义的名字，而不是一串编号。
  \item \textbf{把 IR 交给独立的优化器与后端分别执行。}这一步把编译器拆成可查工具：可以只
        优化不生成，可以只生成不再优化，可以把某一次产物存下来做前后对比。
\end{enumerate}

下面这段用中性动作名描述这种拆分。它是\textbf{示意占位}，不对应任何真实命令行，具体工具名
与开关请以该实现文档为准。

\begin{lstlisting}[style=shellstyle]
to-ir      source.c      ->  prog.ir
optimize   prog.ir       ->  prog.opt.ir
gen-asm    prog.opt.ir   ->  prog.s
assemble   prog.s        ->  prog.o
link       prog.o libs   ->  program
\end{lstlisting}

\subsection{一段示意 IR 与对应源码}\label{sec:cl04-observe-sample}

先看一段普通的 \Tp{cstyle} 源码：

\begin{lstlisting}[style=cstyle]
int sum(int a, int b) {
    int acc = a;
    int i = 0;
    while (i < b) {
        acc += i;
        i += 1;
    }
    return acc;
}
\end{lstlisting}

再看同一逻辑的 IR 风格文本。\textbf{以下是示意}：结构接近真实文本 IR，但块名、指令顺序与
类型细节都刻意简化，不提供任何真实工具的输出编号，不应当作某个实现的规范文本照抄。

\begin{lstlisting}[style=outstyle]
define i32 @sum(i32 %a, i32 %b) {
entry:
  %guard = icmp slt i32 0, %b
  br i1 %guard, label %loop, label %exit

loop:
  %i   = phi i32 [ 0, %entry ], [ %i.next, %loop ]
  %acc = phi i32 [ %a, %entry ], [ %acc.next, %loop ]
  %acc.next = add i32 %acc, %i
  %i.next   = add i32 %i, 1
  %again    = icmp slt i32 %i.next, %b
  br i1 %again, label %loop, label %exit

exit:
  %result = phi i32 [ %a, %entry ], [ %acc.next, %loop ]
  ret i32 %result
}
\end{lstlisting}

对着这段示意文本可以读出本章前面的全部要素：三地址指令（每条只有一个目的值）、显式类型
（\Tp{i32} 出现在每条指令上）、基本块与终止指令（每块最后一行是分支或返回）、\Tp{phi} 合并
（三处合并点各列出候选来源与前驱块）。拿它核对真实转储时请注意，合法 \Tp{phi} 插入位置
可能与直觉不同，前驱边名称也必须精确匹配。

\begin{note}[title={示意性说明}]
本节所有 IR 与命令文本都是示意，只用于说明结构与阅读方法。任何具体语法、块命名规则、类型名、
属性拼写，都必须回到所用实现的文档核对；本册不替它做断言。
\end{note}

\section{IR 与目标无关，但不等于可移植}\label{sec:cl04-abi-ir}

\begin{abibox}[title={目标无关的 IR 里为什么已经有 ABI 的痕迹}]
IR 层的指令集与寄存器数量无关、与字节序无关、与机器码编码无关；但 ABI 决定的三件事在 IR 层
就已存在，形式是属性或参数：
（一）按值传结构体的拆分方式，要么在 IR 层展开成逐个标量参数，要么带属性留给后端展开，
两种选择都会在函数签名上留下痕迹；
（二）可变参数的登记要求，调用点必须说明这是一次变参调用并给出所需的形式参数信息，否则
目标平台的寄存器窗口无法被正确准备；
（三）对齐要求，参数与返回值的对齐在 IR 层以属性或求址宽度表达，进入目标层之后才成为指令
选择依据。
\end{abibox}

由此得出一个经常被误读的结论：\textbf{IR 与目标无关说的是指令格式与寄存器模型，不是函数签名与布局}。
同一份 IR 文本换到另一个目标上重编译，会得到不同的拆分、对齐与调用序列；把 IR 文件当作跨平台
可移植的产物来存档或分发，等于把一次 ABI 决策冻结在它并不适合的环境里。真正可移植的是源码加构建配置，不是 IR。

第~\ref{ch:cl12-endian} 章讨论字节序时还会看到同类现象：IR 里的整数运算对字节序不敏感，
但一旦涉及整块内存的搬运与比较，字节序就已经渗进来了。层次抽象提供的是\textbf{关注点的分离}，
不是“这一层与硬件完全无关”。

\section{读 IR 找优化原因时的三类误判}\label{sec:cl04-misread}

\begin{warning}[title={三类常见误判}]
（一）\textbf{把优化前的形态当成语义。}看到 IR 里有一次结构体整体拷贝，就断定这里必然搬运
若干字节。实际那只是尚未字段化的中间状态，后续可能整块消失；IR 的某一步不是承诺，
只有可观察行为才是承诺。
（二）\textbf{忽略管线顺序的影响。}同样两个变换，先做 A 再做 B 常常把机会消耗掉：某个化简
把循环变成后端更喜欢的形态，却也消除了另一个更有力变换的匹配条件。只盯“为什么没优化”
而不去看“它在这一步被谁先改写掉了”，多半找不到原因。
（三）\textbf{把内联后的行号信息当作独立函数。}调试信息里的行号与函数名在转储中常以位置
标注出现，一个从未被独立发射的函数会带着自己的名字出现在调用方体内，很容易误读成确实
生成了这个函数；判断是否真存在独立实体要看链接符号表，第~\ref{ch:cl08-static-link} 章给出
核对方法。
\end{warning}

\begin{note}[title={本章范围声明}]
本章只讨论中间表示的结构、SSA 形式及其与\cpp{} 语义、ABI 的关系，\textbf{不涉及任何具体优化开关}。
开关的名称、级别组合与效果比较属于第~\ref{ch:cl05-passes} 章；命令行层面的编译驱动选项见
附录~\ref{ap:clA-cli}；符号与 ABI 细节见第~\ref{ch:cl13-cxx-abi} 章与附录~\ref{ap:clB-abi}。
\end{note}

\section{本章要点}\label{sec:cl04-summary}

\begin{itemize}
  \item 中间表示把 $M \times N$ 的翻译问题拆成 $M + N$，真正的收益是优化与移植各只需写一遍。
  \item 三地址指令、显式类型、基本块与控制流图、挂在函数上的调用约定与属性是共性四要素；
        无限虚拟寄存器是把资源分配推迟到后端的刻意选择。
  \item SSA 把数据流分析结果变成结构本身，传播与消除类变换因此便宜运行；\Tp{phi} 必须在
        进入机器相关层之前还原成边上的复制与槽位。
  \item 聚合类型字段化让结构体获得与标量同等的优化待遇，地址逃逸是这条路径的开关；指针
        不透明化的动机是减少类型系统在指令里的重复。
  \item \cpp{} 的构造析构、虚调用、异常清理展开成显式调用与额外控制流边，边界由可观察行为定义；
        内建函数则让标准不规定而优化器已知语义的操作从正门进入 IR。
  \item IR 可打印、可保留名字、可被独立优化器与后端分别消费，但它携带 ABI 决策，因此不是
        可跨平台分发的可移植产物。
\end{itemize}

\chapter{优化管线、整程序优化与反馈驱动}\label{ch:cl05-passes}

第~\ref{ch:cl01-pipeline}~章把编译画成一条七步流水线，第~\ref{ch:cl04-ir}~章说明中间
表示长什么样。本章接管的是中间那一段最不像“翻译”的部分：优化。它没有语义承诺可言，
全部是策略——先做什么、后做什么、在多大范围内做、根据什么做。同一份源码在同一个
编译器上换一档设置，机器码可以面目全非；而语言层面看到的程序仍然完全一样。理解这一
段的关键不是记住几十个开关，而是记住三件事：优化以 \Mx{pass} 为单位组织、整程序优化
拆掉的是翻译单元边界、反馈驱动优化改的是编译器的先验而非源码。

\section{Pass 管线的组织方式}\label{sec:cl05-passorg}

\subsection{函数级与跨函数级}\label{sec:cl05-scope}

优化按可见范围分两层。函数级 \Mx{pass} 只看到一个函数体：它可以在这里任意重排、
合并、删除，而不必询问任何人的意见，代价是它看到的信息只有这么多。跨函数级 \Mx{pass}
要处理调用边界：内联把被调者的函数体搬进调用者，从而把一次跨函数问题降格为一次函数
内问题；过程间的常量传播、参数值分布推断、返回值推断则要在调用图上传播事实，需要一个
能遍历整个模块（甚至整个程序）的组织结构。

这两层的分界不是实现细节，它决定了优化能力的上限。一个只在函数级工作的 \Mx{pass}
看不到调用点传来的常量，于是同一行代码在 \Tool{main} 里被折成常量、在别处保留为
运行期计算，两种结果都“正确”，但只有一种是想要的。想要后者，就得把信息送到跨函数
这一层——要么靠源码（模板、\Tp{constexpr}、头文件里的函数体），要么靠 \Mx{LTO}。

\subsection{分析结果与改写}\label{sec:cl05-analysis}

现代管线普遍把“看出事实”和“改动程序”拆开。分析 \Mx{pass} 只回答问题：这个指针是否
为空、这个循环的边界是什么、这个值在哪些使用点还活着、这个内存区域是否会被别名引用。
改写 \Mx{pass} 消费这些答案并产生新的 \Mx{IR}。这样切分带来两个直接好处：同一个事实在
一轮里的多个改写之间复用；改写之后一旦程序结构变了，相关分析结果被显式作废
（invalidate），下一个消费者被迫重算，而不会读到过期答案。

读到过期事实是优化器最典型的错误来源，也是为什么管线里必须有明确的失效协议：凡是
改写了控制流或内存访问的 \Mx{pass}，都要声明自己弄脏了哪些分析。理解这一层，才能理解
为什么“只开某一个 \Mx{pass}”经常复现不出完整的优化效果：那个 \Mx{pass} 依赖的前置事实
是别人算出来的。

\subsection{常见 pass 族的职责}\label{sec:cl05-passlist}

下面列的是通行名，描述的是职责而非某个实现里的标识符；某个具体实现用什么名字报告
自己的 \Mx{pass}，以它的文档与管线打印为准。

\begin{itemize}
  \item 标量替换与内存提升：把聚合对象（结构体、数组、临时量）的逐成员访问拆成
        独立的标量值，再把只在一处定义、多处读取的临时内存提升为 \Mx{SSA} 上的值，
        使后续 \Mx{pass} 能把它当常量或表达式看。
  \item 全局价值编号：跨基本块识别语义等价的表达式，让同一表达式只求值一次，其余
        使用点改写为同一个值；它顺带完成一部分规范化。
  \item 死代码消除：删除结果无人使用的计算，以及由已知条件推出的不可达分支；
        只在没有可观察副作用时才允许删除。
  \item 循环不变量外提：把每次迭代结果相同的计算移到循环之前，用一次求值换掉
        $n$ 次求值。
  \item 循环展开与向量化：展开减少循环控制开销并拉开指令级并行窗口；向量化把
        多个元素的同种运算合成一条单指令多数据操作，两者都受成本模型约束。
  \item 内联：把调用替换为被调函数体的副本，是所有跨函数优化的入口动作；它同时
        是体积膨胀的主要来源。
  \item 尾调用与尾部合并：把返回位置为尾调用的调用改为跳转以复用栈帧；把若干分支
        尾部相同的块合并成一个共同尾部以省体积。
\end{itemize}

\begin{guideline}[title={顺序本身就是策略}]
\Mx{pass} 之间存在解锁与屏蔽关系：内联把常量送进被调体，之后的价值编号与死代码消除
才有东西可删；先做规范化，改写才认得模式；反过来，过早的展开会让循环结构变得不再
匹配向量化前提。因此“同一批 \Mx{pass}，换一个顺序”不是等价的实现选择，而是不同的
优化器。调档位时你实际调的是一整条顺序化策略。
\end{guideline}

\section{优化档位之间到底差了什么}\label{sec:cl05-levels}

档位不是强度刻度盘，而是策略包。约定俗成的 \Sw{-O0} 到 \Sw{-O3} 加上两个尺寸档，
差别通常是“这一族 \Mx{pass} 的开与关”，而不是某个 \Mx{pass} 参数被调大调小。
下表按取向概括，具体某档在某个实现里打开了哪些 \Mx{pass}，属于实现选择，会随版本
变动，该实现提供对应说明文档，名称与清单以其为准。

\begin{table}[H]
\centering
\begin{tabularx}{0.86\textwidth}{@{}>{\raggedright\arraybackslash}p{0.11\textwidth} >{\raggedright\arraybackslash}p{0.21\textwidth} X@{}}
\toprule
档位 & 取向 & 结构与可调试性的代价 \\
\midrule
\Sw{-O0} & 忠实映射源码结构，尽量不改变语句与变量的对应关系 &
  断点逐行命中、局部变量随时可查；几乎不做任何跨语句推理，运行代价最高。 \\
\Sw{-O1} & 保守的标量优化：代价小、编译时间增加不多的那几族 &
  语句开始合并，单步执行出现跳跃；变量可能已进寄存器。 \\
\Sw{-O2} & 打开多数常规优化，含循环变换与函数内并行性挖掘 &
  常规发布配置的常见选择；调试信息依赖位点归并，回溯仍可工作。 \\
\Sw{-O3} & 在 \Sw{-O2} 之上追加激进变换与更激进的内联预算 &
  体积膨胀与编译时间上升，热路径可能因展开而失去局部性。 \\
\Sw{-Os} & 尺寸优先：偏向减少代码体积的改写与更保守的展开 &
  面向嵌入式与缓存受限场景；不是“更慢”，而是换了目标函数。 \\
\Sw{-Oz} & 极限尺寸：只要不增大体积就不做扩张性变换 &
  体积最小，通常牺牲最多运行期代价。 \\
\bottomrule
\end{tabularx}
\caption{六档优化取向的机制差别（定性描述，不含任何实测收益数字）}\label{tab:cl05-levels}
\end{table}

两点容易搞错。第一，尺寸档不等于慢档：当瓶颈是指令缓存缺失时，缩小代码本身可能就是
最快的优化，这类判断属于性能测量分册。第二，档位不改变语言承诺：它不引入新的可观察
行为，也不允许程序依赖档位存在——依赖某个档位才有正确性的代码，本来就依赖未定义行为。

\subsection{可观察性：谁改了哪一行}\label{sec:cl05-observable}

优化不可见时会带来两类问题：期望的优化没发生却无人知道，不期望的改写发生了却找不到
出处。因此各家实现普遍提供两类能力，术语各异但形态相同。

\begin{itemize}
  \item 决策报告（remark 一类）：把“我内联了这个调用”“这个循环没被向量化，原因
        是成本模型认为标量更快”作为诊断输出打印，并带上源码位置与对应的 \Mx{pass}
        标识。报告通常可分三级：做到的、没做成的、做出判断所依据的分析。
  \item 管线与 \Mx{IR} 转储：打印 \Mx{pass} 顺序表，或在每个 \Mx{pass} 前后把某个函数
        的 \Mx{IR} 落到文件里，供 diff。这类能力用于定位，不用于日常构建。
\end{itemize}

\begin{lstlisting}[style=shellstyle, caption={可观察性用法：决策报告与汇编读回（示意命令）}]
# 一个家族：把优化决策当诊断打印出来，分做到的/没做成的/依据的分析三级
clang++ -O2 -Rpass=inline -Rpass-analysis=loop-vectorize -c widget.cpp
# 另一个家族提供等价的转储与报告能力，名称以该实现文档为准，不跨实现套用
# 最稳的兜底路径：直接读汇编本身
g++ -O2 -S widget.cpp -o widget.s
\end{lstlisting}

报告类开关与转储类开关都分属具体实现，别指望互相通用；能力描述是共同的：都能回答
“哪个 \Mx{pass} 动了哪个位置”。若某个实现没有对应开关，就用汇编读回与 \Mx{IR} 转储
这两条兜底路径。

位点信息（debug location）是可观察性的另一半。优化后的一条机器指令可能对应源码里
根本不存在的表达式，管线靠元数据把它映射回某个位置；当多个来源被折叠成同一个值时，
位置只能取其一。这就是高优化档位下“断点不落在你以为的那行”“变量显示为已优化掉”的
机制原因，而不是调试器坏了。

\section{整程序优化}\label{sec:cl05-lto}

默认模式下，每个翻译单元独立编译：编译器只看得见本单元，跨单元调用对它而言是一个
不透明的函数指针。链接器能合并代码，但它看到的是机器码，无法重新推理语义。\Mx{LTO}
做的事是：把编译推迟到链接期，或者更准确地说，把各单元的中间表示随目标文件一起
交出去，在链接阶段合并成一份足够大的 \Mx{IR}，然后继续跑优化管线。

\subsection{边界拆掉之后发生什么}\label{sec:cl05-ltobenefit}

四件事在跨单元之后才成为可能：

\begin{enumerate}
  \item 跨单元内联：调用点与被调体分处两个 \CPP{} 源文件，合并后它们进了同一个模块，
        内联 \Mx{pass} 不再受文件边界限制。
  \item 跨单元常量传播：某个在 A 单元里是运行期参数的值，可能在 B 单元的调用点上
        始终是常量；合并后才看得见这条边。
  \item 去虚化：虚调用要在运行期查表，前提是编译器不知道实际类型；一旦程序里某个
        基类的全部派生类都可见且可穷举，调用就能折成直接调用。这是典型的全程序判断。
  \item 死代码消除与折叠：没人调用的模板实例、没被使用的虚函数体、常量条件里的
        整个分支，都要在全程序可见性建立之后才能安全删除。
\end{enumerate}

\subsection{薄 LTO 与厚 LTO}\label{sec:cl05-thickthin}

两条路线的区别是：链接期到底还留多少活要并行做。

\begin{compare}[title={厚 LTO 与薄 LTO}]
厚（full）LTO 在链接期把全部单元的中间表示合成一个大模块，一次跑完整管线，优化视野
最完整，代价是链接期串行度高、内存峰值大、增量构建几乎失效。薄（thin）LTO 不合并中间
表示，只在链接期建立一张全局索引，据此做跨单元的轻量决策（主要是内联与可见性调整），
然后把每个单元各自的原代码生成任务并行分派到多个线程或子进程。薄 LTO 构建快得多，
但跨单元的标量优化视野比厚 LTO 浅。二者是取舍关系，不是新旧关系。
\end{compare}

增量构建与 \Mx{LTO} 的冲突值得单独提醒：一旦优化决策依赖“全程序可见”，任何一处改动
都可能改变这些决策，于是链接期要重做的工作不再是常数。实践中常见的折中是开发期关
\Mx{LTO}，发布与性能基线构建开；这与构建目录分离、缓存策略属于工程流程话题。

\subsection{为什么 LTO 会改变符号可见性}\label{sec:cl05-symbols}

合并 \Mx{IR} 之后，编译器第一次能对“这个符号在整个程序里还有谁用”给出确定答案，
于是它会把原本外部可见的函数降格为内部可见，进而做更激进的改写：删除未被使用的形参、
把返回结构改成寄存器约定、把只在一处调用的函数直接并入调用者。这类改写让符号的
\Mx{ABI} 形状发生变化，而符号名不变——这正是第~\ref{ch:cl08-static-link}~章强调过的
风险面：静态库里的目标文件一旦经过 \Mx{LTO} 重编，它就不再与同名的非 \Mx{LTO} 目标
文件语义等价，混链时要么全部一致，要么按实现的规则明确拒绝。

另一件由 \Mx{LTO} 改变的事是链接期折叠。模板实例、\Tp{inline} 函数、\Tp{constexpr}
函数在每个使用它的单元里都留下一份相同内容的定义，传统链接靠重复定义规则与段合并
把它们收敛成一份（机制见第~\ref{ch:cl08-static-link}~章）；\Mx{LTO} 则在 \Mx{IR} 层就
认出它们是同一个函数，折叠发生在优化之前而不是之后，因而被折进来的实例还能享受
跨单元内联。

\begin{lstlisting}[style=shellstyle, caption={LTO 的最小链路：编译期留 IR，链接期做优化}]
# 编译：产物里的机器码很少，主体是随对象文件走的可重定位 IR
g++ -O2 -flto=auto -c a.cpp -o a.o
g++ -O2 -flto=auto -c b.cpp -o b.o
# 链接：合并 IR、跨单元内联、去虚化、死代码消除，最后才出机器码
g++ -O2 -flto=auto a.o b.o -o app
# 检查单元里是否真的留了 IR（也可以只读节名列表）
llvm-readelf --section-headers a.o | grep -i bitcode
\end{lstlisting}

\begin{lstlisting}[style=cmakestyle, caption={构建配置里同时给编译与链接开 LTO}]
target_compile_options(app PRIVATE $<$<CXX_COMPILER_ID:GNU,Clang>:-flto=auto>)
target_link_options(app PRIVATE $<$<CXX_COMPILER_ID:GNU,Clang>:-flto=auto>)
# 只给编译加而漏掉链接，是最常见的一次无效尝试
\end{lstlisting}

\begin{longtable}{@{}>{\raggedright\arraybackslash}p{0.15\textwidth} >{\raggedright\arraybackslash}p{0.20\textwidth} >{\raggedright\arraybackslash}p{0.20\textwidth} >{\raggedright\arraybackslash}p{0.22\textwidth}@{}}
\caption{三种跨翻译单元信息传递方式的机制对照}\label{tab:cl05-crossTU}\\
\toprule
方式 & 编译器看到什么 & 链接期看到什么 & 何时失效 \\
\midrule
\endfirsthead
\toprule
方式 & 编译器看到什么 & 链接期看到什么 & 何时失效 \\
\midrule
\endhead
\bottomrule
\endfoot
头文件与模板源码 &
  函数体与类型定义在编译期就完整可见，可当场内联与折叠 &
  每单元一份同内容定义，靠重复定义规则与段合并收敛 &
  定义未随声明暴露；头文件与实现分离后跨单元视野消失 \\
\Mx{LTO} &
  编译期只保留 \Mx{IR}，不做最终代码生成 &
  一份合并的程序级 \Mx{IR} 或一张全局索引，重跑优化 &
  单元混用非 \Mx{LTO} 产物；链接器与编译器不匹配；增量构建退化 \\
模块接口 &
  编译期直接读取被依赖模块的语义信息，无需文本包含 &
  基本只看到已编译的模块与目标文件，无需再合并 \Mx{IR} &
  接口变更迫使依赖方重编；跨模块深度优化仍需实现提供额外机制 \\
\end{longtable}

\section{反馈驱动优化}\label{sec:cl05-pgo}

优化器所有的判断都建立在猜测上：这个分支多半走哪边、这个循环大概跑几次、这个调用值
不值得内联。猜测来自静态结构与代价模型。反馈驱动优化（\Mx{PGO}）改变的就是这一点：
先跑一次程序，把真实执行过的路径统计出来，重编时把这些统计作为先验喂给管线。

\subsection{两条数据来源：插桩与采样}\label{sec:cl05-pgosource}

\begin{itemize}
  \item 插桩：先编一个带计数器的版本，运行时记录边（基本块之间的一次转移）的执行次数，
        退出时落盘。计数精确，代价是这一次构建的体积与速度都明显劣化，且程序里多了
        写文件的副作用。
  \item 采样：用不修改程序的手段（硬件性能计数器或定时中断的调用栈）收集样本，再把它
        映射回函数与分支。开销低、可以来自生产环境，代价是统计量稀疏且带偏差，
        需要按权重还原。
\end{itemize}

两类的产物形态不同，因此在实现里通常对应两套互不兼容的 profile。粗略地归类：
一类实现把自己的原始计数文件转换成一份转换后的合并数据供重编读取；另一类实现直接
消费采样数据文件。两类实现各用自己的 \Mx{profile} 格式，扩展名、存放位置与转换步骤
都以各自文档为准，跨实现通用是不成立的假设。

\begin{lstlisting}[style=shellstyle, caption={PGO 的三步形态：生成、运行、消费（示意）}]
# 1) 生成一次带插桩的构建
g++ -O2 -fprofile-generate -c hot.cpp -o hot-instr.o
g++ -O2 -fprofile-generate hot-instr.o -o app-instr
# 2) 用有代表性的负载跑它，让计数落盘到 profile 目录
./app-instr --workload typical-input.bin
# 3) 重编：把这些计数作为分支与内联决策的先验
g++ -O2 -fprofile-use -c hot.cpp -o hot.o
g++ -O2 -fprofile-use hot.o -o app
\end{lstlisting}

\subsection{profile 与代码、编译器版本绑定}\label{sec:cl05-profilebind}

\Mx{profile} 不是一张通用地图。它按函数标识与函数内部结构来挂计数：常见做法是同时
记录函数体的结构摘要，重编时校验摘要，不一致就丢弃这份数据并报警。于是三种改动都会
让 profile 失效：源码结构性变化（函数拆了、循环重写）、编译选项改变函数划分（尤其是
内联与 \Mx{LTO} 决策）、编译器版本升级改变了结构摘要的算法。工程上应把 profile 视为
与提交历史绑定的构建输入，随源码一起版本管理，并让生成步骤进入构建脚本而不是靠人记得跑。

\subsection{代表性不足时优化为什么反而变差}\label{sec:cl05-pgorisk}

\Mx{PGO} 把预算集中到它认为的热路径上：热路径被展开、被内联、被铺平，冷路径被移出
或干脆按“不会发生”处理。如果收集 \Mx{profile} 时的负载不能代表真实负载，这种集中
就成了押注：

\begin{itemize}
  \item 内联预算错配：在测试里每轮都调一次的小函数被激进内联，真实调用者众多，
        代码体积膨胀、指令缓存局部性变差。
  \item 分支布局押错方向：罕见分支被当成主路径重排，常走的分支反而要跳两次。
  \item 冷路径消失：只在真实故障条件下走的路径被压缩得几乎不可读，甚至因“假设不会
        到达”而失去原本的形状。
  \item 循环决策反向：短循环被判定不值得向量化，长循环被激进展开，与真实分布相反。
\end{itemize}

\begin{warning}[title={profile 失真比没有 profile 更糟}]
没有反馈时优化器按保守的先验工作，退化幅度有限；有反馈时它会自信地重分配资源，
把错误先验放大。因此上线 \Mx{PGO} 的判据不该是“跑分高了”，而该是“收集 \Mx{profile}
的负载与生产负载在函数与路径分布上确实接近”，并且每次 \Mx{profile} 更新都要重新核对
这件事。这条判断本身属于性能测量分册的方法，本册只解释机制。
\end{warning}

\section{链接之后的重排}\label{sec:cl05-postlink}

\Mx{PGO} 与 \Mx{LTO} 都在代码生成之前起作用，改的是指令如何产生。还有一类做法接受
已链接好的可执行文件，把它连同运行期观测到的执行路径一起再处理一遍：重排函数的存放
顺序、调整函数内部基本块的排列、把首次访问就需要的数据放到相邻位置，目标是让热代码
在地址空间里连续，减少取指缺失与地址翻译开销，顺带把冷代码推到文件末尾便于丢弃。
这类工具常以 \Mx{BOLT} 之名被提到，它们读入的是链接产物与采样数据，输出的仍是
可执行文件，不重新编译源码。

差别只在阶段位置，作用面也因此不同。

\begin{flow}[title={三种反馈回路的阶段位置}]
源码级反馈：开发者或静态分析在写代码时给出提示（分支提示、布局属性、手工内联取舍），
影响进入前端的语义，改动可见于源码；编译期反馈：\Mx{PGO} 与 \Mx{LTO} 在代码生成前
重排优化决策与模块边界，产物结构由编译器负责；链接后反馈：二进制优化工具在完整
可执行文件上重排代码与数据，不参与语义，只搬动已有的块。三者可以叠加，但每一层都要
单独验证：链接后重排会改变函数地址与调试信息的对应关系，剖析与回溯工具需要重新对齐。
\end{flow}

\Mx{BOLT} 一类做法的前提是二进制里的符号、重定位与展开信息足够完整：一旦被 \Tool{strip}
处理或表被丢弃，重排与回溯都会受限。这也解释了为什么这类流程要求在发布产物之外保留
带信息的副本（相关目标文件结构见第~\ref{ch:cl07-object}~章）。

\section{与语言承诺的交互}\label{sec:cl05-promises}

\begin{warning}[title={许可不是保证}]
三件事在优化器与语言之间是不对称的。第一，\Tp{inline} 只是许可合并定义并暗示可以考虑
内联，编译器完全可以不内联；\Tp{constexpr} 只是保证在常量上下文中求值得到常量表达式，
运行期调用照样可能在机器码里算一遍。两者都不是性能承诺，验证要看汇编或 \Mx{remark}。
第二，优化可以合法地假设未定义行为不会发生，因而删掉它认为“永远为真”的检查：一旦
程序里有未定义行为，被删掉的正是那层本来用来报错的保护。这类推理属于 as-if 框架，
机制细节见未定义行为分册的 as-if 相关章节，本册不重复展开。
第三，\Tp{inline} 的另一种含义——它在链接期造成多份同定义与折叠——是发布配置必须
关心的事，见第~\ref{ch:cl08-static-link}~章；至于“什么时候该写 \Tp{inline}、多长的函数
不该放进头文件”这类风格判断，属于代码风格分册，本章不涉及。
\end{warning}

\begin{note}[title={本册的边界：只讲机制}]
本章讲反馈如何改变编译器的决策、优化在哪些范围内做。计时方法、微基准设计、统计
显著性判据、缓存与分支预测的实测流程，全部属于性能测量分册。本册刻意不给任何收益
数字：具体百分比与倍数依赖代码、负载、机器与版本，任何写进手册的数字都会在下一个
版本变成误导。判断“该不该开 \Mx{PGO}”要自己在目标负载上测，机制只告诉你为什么结果
可能和你的直觉相反。
\end{note}

\subsection{把三件事分开记}\label{sec:cl05-recap}

\Mx{LTO} 扩大视野，\Mx{PGO} 改进先验，链接后优化搬动成品。它们互不替代：只有
\Mx{LTO} 而无反馈，跨单元优化仍按猜测分配预算；只有反馈而无 \Mx{LTO}，热路径上的
跨单元调用仍然留在函数边界之外；两者都有，链接后重排仍然能改善指令缓存的布局质量，
因为它面对的是最终地址分配结果。

\begin{bestpractice}[title={反馈类优化的落地顺序}]
先固定基线：同一提交、同一编译器版本、同一构建目录，记录产物结构与符号表。再逐层加：
先开高优化档，确认汇编里预期发生的改写真的发生了；再引入 \Mx{LTO}，检查链接期日志里
合并与代码生成的工作量，确认没有与非 \Mx{LTO} 目标混链；最后引入 \Mx{profile}，并把
profile 的生成负载写进构建脚本，让每次刷新都有明确出处。一次只加一层，出问题才能定位
到层。开关的准确名称以所用实现的文档为准，不确定的项就先查文档再动手。
\end{bestpractice}

\Mx{pass} 顺序与档位是编译器的内部策略，\Mx{LTO} 与 \Mx{PGO} 是把程序作为一个整体来
优化的两种手段。下一章进入最后一个阶段：当所有决策做完，抽象的 \Mx{IR} 操作要落到具体
指令和具体寄存器上，那里的取舍不再由“语义是否等价”决定，而由目标 \Mx{ABI} 与指令集合
决定。

\chapter{指令选择与寄存器分配}\label{ch:cl06-codegen}

前两章讲的是编译器在哪一层做决策；本章讲的是决策落地的地方：抽象的 \Mx{IR} 操作要
变成具体的指令、具体的寄存器、具体的栈帧位置。这一段是整条流水线里唯一同时受三个
因素约束的阶段——目标指令集合的表达能力、\Mx{ABI} 规定的寄存器职责、以及程序员对
可调试性与可剖析性的期待。第~\ref{ch:cl01-pipeline}~章给了阶段划分，本章把它填实。

\section{后端在管线的哪一段}\label{sec:cl06-stage}

\begin{flow}[title={从 IR 到可提交汇编的六步}]
合法性检查（这个操作目标能不能直接表达）→ 指令选择（用哪些指令覆盖这张图）→
前置机器调度（在虚拟寄存器上拉开并行度）→ 寄存器分配（虚拟换物理，决定谁溢出）→
序言与尾声插入、栈槽物化（帧大小此刻才确定）→ 后置机器调度（在真实寄存器与调用
约定约束下微调）。每往后一步自由度就少一分：选择阶段还能换指令形态，分配阶段之后
只剩下搬动顺序，而序言一旦插入，帧布局就写进了二进制。
\end{flow}

\section{指令选择的三条途径}\label{sec:cl06-isel}

\Mx{IR} 里的一条加法与目标上的一条加法之间，往往不是一一对应。选择器要解决的真正
问题是：给定一串语义等价的候选指令序列，选哪一串。三条途径就是这个“怎么选”的答案。

\subsection{直接模式匹配}\label{sec:cl06-match}

逐条 \Mx{IR} 操作查表，命中一条模板就吐一条指令。优点是实现简单、结果可预测、编译
快。缺点是它看不见邻接关系：“加载、位移、相加”在 \Mx{IR} 里是三个节点，在
\xSixtyFour{} 上其实可以折进一条带变址寻址的指令；比较与分支是两条独立指令，在
\ArmSixtyFour{} 上则可以合成一条带比较的分支。逐条匹配放弃了这类合并机会。

\subsection{DAG 覆盖}\label{sec:cl06-dagcover}

把一个基本块（或更大范围）内的相关操作表示成有向无环图，用一套模板库对这张图做
“树覆盖”：每个模板声明自己能覆盖什么样的子图，选择器遍历图中每个节点，把可用模板
的代价累加起来，用动态规划求出一张覆盖整张图、总代价最小的模板集合。
这条途径的价值在于它会主动合并节点——寻址模式折进访存、常量装载与运算合并、
条件与分支融合，同时能避开那些代价高的组合（比如需要额外寄存器拷贝的形态）。
代价是模板表达能力的上限：模板描述模式，但很难描述“整段一起考虑”的约束。

\subsection{带成本模型的搜索}\label{sec:cl06-search}

第三条途径把选择当成一个全局优化问题：不再逐节点局部求最优，而是枚举候选覆盖并
结合寄存器压力、后续调度空间、溢出概率这些非局部指标来打分。它最可能选出“对分配器
友好”的形态——局部指令数略多、但少占用跨调用活跃的寄存器，总代价反而更低。工程
代价是要维护可信的成本模型，并且搜索规模必须被限定，否则编译时间不可控。

\begin{compare}[title={三条途径的取舍}]
逐条匹配：可预测、编译快，放弃指令合并，代码质量下限最低。图覆盖：主流做法，靠模式
与代价合并相邻操作，能表达绝大多数硬件特性，但视野以局部结构为限。全局搜索：视野
最完整，能把寄存器压力纳入选择决策，代价是成本模型的可信度与搜索开销。三者并不互斥：
实际实现常常是“覆盖为主、成本模型定夺、局部保留逐条兜底”。
\end{compare}

\section{目标描述：同一套后端如何认识新机器}\label{sec:cl06-targetdesc}

后端不该为每个目标重写一遍。可复用的部分（图构造、覆盖算法、分配器、调度器）与
目标相关的部分之间，靠一层“目标描述”隔开。这层描述至少要回答四类问题。

\begin{itemize}
  \item 合法性：目标没有原生形式时，这个操作怎么拆。64 位立即数装载拆成两条；
        没有原生除法的机器把除法拆成一条调用；宽向量洗牌拆成几条窄的重排与拼接。
        拆法本身也是一种选择，拆完还要重新过一遍合法性。
  \item 成本：每条指令形态消耗什么资源——占几个发射槽、几个执行端口、需要多少个
        临时寄存器、属于哪一类寄存器（通用、向量、标志）。成本模型的粒度直接决定
        上一节那条搜索途径有没有意义。
  \item 调度描述：延迟、吞吐、功能单元与依赖类型，前置与后置调度都读它。
  \item 约束表达：寄存器类别与别名、调用约定的寄存器职责、汇编输出的方言
        （同一段选择结果要能按 AT\&T 或 Intel 写法吐出文本）。
\end{itemize}

\subsection{LLVM 的两条选择路径}\label{sec:cl06-llvm-two}

\Tool{LLVM} 在同一个后端框架里提供两条并存的路径，不是新旧替换而是分工。

\begin{itemize}
  \item 数据依赖图方式（\Tool{SelectionDAG} 一类）：先做类型合法化，把 \Mx{IR} 下降到
        与目标类型系统对齐的节点图，再在这张图上做 DAG 覆盖，最后落到机器级表示。
        它的优点是模式表达能力成熟、覆盖面广；限制在于下降过程会过早丢掉 \Mx{IR}
        里的信息，而覆盖决策被限定在图上能看到的局部结构。
  \item 机器级 \Mx{IR} 上的全局选择（\Tool{GlobalISel} 一类）：跳过那张专门的数据依赖
        图，一开始就进入机器级表示，选择、 legalization 与合并都在同一层做，因而能对
        更大范围（整个函数或大块区域）统一求解。它的收益在于决策视野与信息保真，
        代价是要为目标逐条补齐机器级规则。
\end{itemize}

\subsection{GCC 的机器描述文件}\label{sec:cl06-gccmd}

\Tool{GCC} 走的是另一条组织路线：目标相关的指令以机器描述文件表达，文件里每条指令
是一个带模式的模板，模式用寄存器类、内存形式、立即数范围这类约束字母刻画，模板体
写成 RTL 形式；选择阶段在函数内部的 RTL 表示上匹配这些模板，并用模板里声明的代价
与约束决定采纳与否。寄存器类、条件码与调用约定同样在这层描述里定义。

层级上的差别值得记住：一条途径把目标知识主要放在“模式与代价”里，另一条把目标知识
主要放在“合法化规则与指令定义”里；两者都要回答合法性、成本、调度三件事，只是分工
粒度不同。本章刻意不写具体文件与模板语法细节：那属于某个实现的内部约定，版本之间会
变，动手前以对应文档为准。

\section{伪指令、约束与延迟落槽}\label{sec:cl06-pseudo}

选择做完，屏幕上出现的仍然不是最终汇编。三个原因。

\begin{enumerate}
  \item 虚拟寄存器。选择结果里的操作数属于“寄存器类”，不是具体的 \Sym{rax} 或
        \Sym{x7}。它们在分配阶段才被换成物理寄存器或被搬到内存。
  \item 伪指令。有些形式语义正确但机器不能直接执行：一次跨模块调用可能被写成一条
        伪指令，展开后是取地址、设置寄存器、间接转移若干条；大常量可能被写成一条
        “装载任意立即数”，实际展开依赖值的大小；带条件的跳转与比较的融合形态在
        帧指针被引入后还要重新展开成实际序列。
  \item 延迟落槽。栈槽、参数保存区、对齐填充的位置与大小，都要等寄存器分配知道
        “有几个值必须搬到内存”之后才能定；帧指针是否保留、返回值放哪里，也在这一步
        才物化。这就是为什么序言与尾声通常在分配之后插入。
\end{enumerate}

约束的作用在分配阶段才显出分量。约束字母规定的东西包括：这个操作数必须是哪一类
寄存器、必须是内存形式还是立即数、必须与另一个操作数同处一个寄存器、必须落在某个
指定寄存器上。要求过多则分配器没有自由度，只能靠拷贝或溢出绕开；要求过少则指令
模板无法成立。选择阶段留下的约束与分配阶段的可行性是同一枚硬币的两面。

\begin{guideline}[title={调度的两次，各有各的理由}]
分配前的调度只看虚拟寄存器，不受“这条指令要占用哪个物理寄存器”与“函数调用会破坏
哪些寄存器”的干扰，适合把独立操作拉开、为向量化与流水线腾出窗口；缺点是一些依赖
真实寄存器编号的优化（例如相邻指令共享同一寄存器的约束）无法判断。分配后的调度看
得到全部真相，代价是它的移动范围受限于已经定下的帧布局与拷贝指令。一次改写只在
其中一次生效很常见，看不出这点就会误判“调度没起作用”。
\end{guideline}

\section{寄存器分配的两大流派与默认选择}\label{sec:cl06-regalloc}

分配问题本身：给定若干虚拟寄存器与它们的活跃区间，把它们映射到数量有限的物理
寄存器，映射不完的部分落到内存，总代价（溢出与搬运的开销，加上调用约定要求的保存
恢复）最低。

\subsection{图着色}\label{sec:cl06-coloring}

构造冲突图：节点是虚拟寄存器的活跃区间，两个区间在任一程序点上同时活跃就连一条边；
边还额外标记“是否必须不同寄存器”之类的约束。可用物理寄存器数为 $k$ 时，反复删除度数
小于 $k$ 的节点（这些节点在回填时总能找到颜色），剩下的高度数节点集合要么整体可着色，
要么必须挑出牺牲者。牺牲者的选择按代价启发式：跨调用活跃、循环内的区间更贵，叶区域
的更便宜。落进内存之后，分配器可以在区间交界处分拆（split）：只在一段用寄存器，另一
段用内存，让一个长区间变成两个短区间，把冲突解开。

这条途径的优点是有清晰的图论语言可以讨论冲突与代价；缺点是构图与求解的成本随函数
规模迅速上升，且冲突图上的度数统计对“一个寄存器被反复使用”这种情形不敏感。

\subsection{线性扫描}\label{sec:cl06-linscan}

不构图，把所有区间按起始位置排序，然后沿程序点单向扫描：新区间开始活动时给它一个
空闲物理寄存器；没有空闲时，在正在活动的区间里挑一个牺牲者把它搬到内存；区间结束时
归还寄存器。为“活动顺序”保留的判定让它能识别一段区间里最多有多少个值同时活跃，
避免不必要的溢出。

它的性能是可预期的线性，因此出现在需要即时编译的场合。质量上它不如全局求解细致：
区间一多、跨度一长，牺牲者选择就更容易押错，代价表现为额外的搬运指令；对区间内多次
使用的情形需要额外机制（间隔、拆分或重加载）弥补。

\subsection{贪心与优先级搜索}\label{sec:cl06-greedy}

现代编译器里常见的默认形态是介于两者之间的一种：以活跃区间为单位，按优先级逐个
决定插入点。优先级来自一组代价估计——区间是密集还是稀疏、是否自干涉（同一区间在
某段还要两个寄存器）、所处循环的嵌套深度、是否跨越调用、候选物理寄存器是
caller-saved 还是 callee-saved（后者附带序言与尾声的保存恢复代价）。每插入一个区间，
冲突信息被增量维护，代价高的区间可以被拆散到别处。

这条路线既保留了图着色里“代价驱动”的语言，又继承了线性扫描的增量结构，因此实践中
通常取代了经典图着色；拆分与溢出的策略仍然按图着色那一套代价模型计算。

\begin{longtable}{@{}>{\raggedright\arraybackslash}p{0.14\textwidth} >{\raggedright\arraybackslash}p{0.18\textwidth} >{\raggedright\arraybackslash}p{0.22\textwidth} >{\raggedright\arraybackslash}p{0.22\textwidth}@{}}
\caption{三种寄存器分配路线的机制对照}\label{tab:cl06-regalloc}\\
\toprule
流派 & 决策单位 & 溢出与拆分 & 适用场景 \\
\midrule
\endfirsthead
\toprule
流派 & 决策单位 & 溢出与拆分 & 适用场景 \\
\midrule
\endhead
\bottomrule
\endfoot
图着色 & 冲突图上的节点（活跃区间） & 度数超限时按代价挑牺牲者；在区间交界处拆分以解除冲突 &
  函数规模适中、希望显式讨论冲突结构时 \\
线性扫描 & 按起始排序的区间，单向扫描 & 无空闲时选活动区间做牺牲；靠间隔与重加载补救 &
  编译时间必须可预期，或需要即时编译的场合 \\
贪心与优先级 & 单个区间，插入时增量维护冲突信息 & 按优先级与保存寄存器类别决定落点，可反复拆分 &
  当前多数后端编译期的默认取向 \\
\end{longtable}

多数实现允许在同一目标上切换分配策略，但切换手段属于实现细节：有的按目标选定，
有的按优化档位改变默认，有的提供实验性选项。具体名称以所用实现的文档为准；本章
只保证一件事——不同分配器的差别体现在“谁被搬到内存、什么时候拆”，而绝不体现在
程序语义上。

\section{ABI 与分配的耦合}\label{sec:cl06-abi}

分配器的代价模型不是抽象的：它的每一项都由调用约定折算出来（约定的完整内容见
第~\ref{ch:cl10-cc}~章）。三条耦合最要紧。

\subsection{caller-saved 与 callee-saved 决定溢出代价}\label{sec:cl06-saved}

caller-saved 寄存器在调用之后不保证还保有原值，callee-saved 寄存器则由被调用者负责
保持。于是一个值的活跃区间如果跨越调用点，它只能占 callee-saved 寄存器或者内存；
如果整个函数的 callee-saved 都被占用，新来者必须挤掉某个已占用的寄存器，代价立刻
变成两条：本函数序言里多一条保存、尾声里多一条恢复，加上搬进搬出内存的指令。
反过来说，活跃区间不跨调用时，用 \Sym{rax} 这一类比用 \Sym{rbx} 这一类便宜得多，
因为前者不产生任何帧内开销。同一个函数在 callee-saved 集合更小的平台上，溢出次数
往往更高，这不是分配器差，而是约定给的池子小。

\subsection{影子空间}\label{sec:cl06-shadow}

\xSixtyFour{} 上的 \Win{} 约定要求调用者在每个调用点之前，为被调用者的前几个参数
寄存器预留一段固定栈空间，即使参数全部通过寄存器传递、这段空间根本不会被读。
这条规则对分配产生两个可见后果：任何有调用的函数都不能是零栈帧，序言要移动栈指针；
被调用者可以放心把参数溢出到这段空间而不必自己检查。尾调用与这段空间的交互还额外
限制了可做的改写。叶函数（不含调用）不受影响，这也是叶函数优化特别强的原因之一。

\subsection{红色区域}\label{sec:cl06-redzone}

\xSixtyFour{} 上的另一族约定允许函数在不改变栈指针的前提下，使用栈指针以下一小段
固定字节做临时存放（这段区域被称为红色区域），前提是没有别的执行体会踩到这里。
收益是省去两条栈指针操作，代价是只能在确定不会被信号或异步处理打断的位置使用：
被调函数会移动栈指针并覆盖这段区域，因此它只对叶函数的最后一段临时量有意义。
这就是同一个分配决策在两个平台上给出不同指令序列的机制原因。

\begin{abibox}[title={同一份代码在三族约定下的分配差异（定性）}]
\Nix{} 与 \Mac{} 的 \xSixtyFour{} 约定用一组较丰富的参数寄存器，caller-saved 集合
相对宽松，分配器在跨调用区间上更愿意直接用寄存器；\Win{} 的 \xSixtyFour{} 约定自带
影子空间，帧布局与栈对齐因此更保守；\ArmSixtyFour{} 上参数与结果的寄存器分类和条件
标志的处理都不同，条件码不能作为可分配值参与区间求解，比较与分支的融合形态也因此
改变。三族的寄存器命名与职责在第~\ref{ch:cl10-cc}~章逐项给出，这里只强调一句：
溢出次数与序言长度是约定的函数，不是代码质量的函数。
\end{abibox}

\section{栈帧、帧指针与可调试性}\label{sec:cl06-frame}

帧指针的作用是把调用栈变成一条可走通的链表：每个帧里保存上一个帧的位置，回溯工具
只要顺着这条链就能列出调用者。它同时占用一个通用寄存器，因此在寄存器紧张的平台上
常常被省掉（省略能力与档位相关，实现各异）。

省掉之后，栈布局本身不再能告诉回溯工具“这个帧从哪来、返回地址在哪、保存的寄存器
在什么偏移”，唯一的信息来源是编译期写进二进制的展开表：\Elf{} 里的
\Sect{eh\_frame} 与 \Sect{text} 段中的调用框架信息、\PeFmt{} 里的
\Sect{pdata} 与 \Sect{xdata}、\Macho{} 里的紧凑展开记录（三者的格式细节见
第~\ref{ch:cl07-object}~章与第~\ref{ch:cl13-cxx-abi}~章）。异常抛出时寻找捕获点也
走同一批表：栈上每一层是谁、返回地址在哪里、哪些寄存器需要恢复，全靠表项回答。

于是可以准确说出“省略帧指针会让展开失败”的边界条件：不是省略本身导致失败，而是
省略之后表缺失或不可用。三种典型情况：产物被处理掉展开节（第~\ref{ch:cl08-static-link}~章
讨论过 \Sym{strip} 类工具的取舍）；代码不在表的覆盖范围内（运行期动态生成的机器码、
被改写过的帧、越界或已失效的返回地址）；表与实际代码不匹配（不同版本的目标文件混在
一个产物里）。这时剖析器拿不到完整调用链，异常穿越该帧时无法定位捕获点。

\begin{warning}[title={帧指针、栈对齐与剖析的三角关系}]
两件事常被混在一起，其实各自独立。第一，栈对齐是硬约束：多数 \xSixtyFour{} 约定要求
调用点处栈指针按固定粒度对齐，序言里保存了几个寄存器、又额外移动了多少栈指针，共同
决定这个不变量是否成立；少一条或少两条保存就可能让后续访存跨缓存行，甚至让依赖对齐
的向量化访存直接出错。第二，可剖析性是运维约束：省略帧指针后，采样剖析器只能靠展开
表或帧记录重建调用链，表不全时热路径会被记到错误的父帧上。这就是为什么不少发行版
会在发布产物里保留帧指针：用一点寄存器压力换可观测性。这个取舍属于性能测量分册的
议题，本章只给出机制。
\end{warning}

\section{验证路径：怎么知道后端做了什么}\label{sec:cl06-verify}

指令选择与寄存器分配没有中间态可查，最直接的证据是文本形式的机器码。两条读法：
从源码直接编出汇编文本，或从已链接产物反汇编回读。前者看得到分配器留下的伪指令
与调用约定骨架，后者看到的是最终布局、符号与真实地址。

\begin{lstlisting}[style=shellstyle, caption={两条读法：源到汇编、产物到反汇编（示意命令）}]
# 源码直接编成汇编文本：读选择与分配的结果
g++ -O2 -S -masm=intel widget.cpp -o widget-intel.s
clang++ -O2 -S widget.cpp -o widget.s
# 已链接产物反汇编：读最终地址、符号与真实字节
objdump -d --no-show-raw-insn ./app | grep -A40 "<heavy>:"
llvm-objdump -d --symbolize-operands ./app
dumpbin /DISASM app.exe                  # Windows 下的等价手段
\end{lstlisting}

同一份选择结果在不同方言下写法完全不同。三个坐标要记住：\xSixtyFour{} 的文本形式有
AT\&T 与 Intel 两种，两者都由同一后端产出，只是输出方言不同；\ArmSixtyFour{} 的
\Nix{}/\Mac{} 汇编器使用自己的固定写法，寄存器以 \Sym{x} 与 \Sym{w} 两组宽度命名，
源与目的的排列顺序与 ARM 传统文档习惯不一致；\Win{} 的 \Tool{MASM} 风格在符号前缀、
段声明与寄存器写法上又自成一套。

\begin{lstlisting}[style=asmstyle, caption={同一函数两种 x86-64 方言的对照（示意，非编译器实测输出）}]
# 写法一：AT&T（GAS 默认，Linux 与 macOS 常见）
# 源在前、目的在后；寄存器带 % 前缀，立即数带 $ 前缀
add_offset:
    leaq    8(%rdi), %rax
    addl    %esi, %eax
    ret
# 写法二：Intel / MASM 风格
# 目的在前、源在后；寄存器无前缀，内存操作数写成方括号形式
add_offset:
    lea     rax, [rdi+8]
    add     eax, esi
    ret
\end{lstlisting}

\begin{table}[H]
\centering
\begin{tabularx}{0.84\textwidth}{@{}>{\raggedright\arraybackslash}p{0.15\textwidth} >{\raggedright\arraybackslash}p{0.20\textwidth} X@{}}
\toprule
维度 & 常见差异 & 读汇编时的判据 \\
\midrule
操作数顺序 &
  AT\&T 是源、目的；Intel 与 MASM 是目的、源 &
  看第一条 \Fn{mov} 哪一侧是目标，先确认方向再判断改写是否等价 \\
寄存器与立即数记号 &
  一族给寄存器名与立即数各加一个前缀记号，另一族两种都不加；\ArmSixtyFour{} 用 \Sym{x0} 与 \Sym{w0} 区分宽度 &
  宽度后缀（\Sym{b}、\Sym{w}、\Sym{l}、\Sym{q}）与前缀记号混用时最容易读错 \\
伪指令与段声明 &
  汇编器家族各自规定段的声明形式、符号可见性与对齐指示的写法 &
  跨方言搬运文本前先确认伪指令是否被目标汇编器接受，不然错误在汇编期才暴露 \\
符号名 &
  是否在名字前加下划线、是否区分大小写由目标格式与约定决定 &
  反汇编里核对符号要配合第~\ref{ch:cl13-cxx-abi}~章的装饰规则，不能直接比对文本 \\
\bottomrule
\end{tabularx}
\caption{三平台汇编方言的阅读维度（差异类别，不逐条列举语法）}\label{tab:cl06-dialect}
\end{table}

序言与尾声是分配与约定共同作用的结果，读反汇编时最先看它们：被保存的
callee-saved 寄存器列表、栈指针的移动量、返回值落在哪里。

\begin{lstlisting}[style=outstyle, caption={反汇编里的序言与尾声（示意结构，非实测字节）}]
# 一个使用 callee-saved 寄存器并有调用的函数
push   rbx
push   r12
push   r13
sub    rsp, 0x20            # 帧大小在此刻才确定，含对齐填充
...
call   helper
...
add    rsp, 0x20
pop    r13
pop    r12
pop    rbx
ret
# 三条 push 加一次 sub：callee-saved 的保存数量会影响 rsp 对齐，两者要一起算
\end{lstlisting}

\section{外部片段：绕过分配器的那道口子}\label{sec:cl06-embedded}

编译器的一切代价模型都建立在“它知道每条指令碰了哪些寄存器”之上。手工嵌入的汇编
片段或外部机器码文件不提供这个信息，于是选择与分配阶段把它当成一个不透明的值：
它读写什么由程序员在约束表里声明。声明不全的后果不是立刻崩溃，而是让分配器做出
错误的重用。

\begin{abibox}[title={破坏 callee-saved 的片段：症状为什么出现在无关位置}]
设想一段外部片段在执行时改用了某个 callee-saved 寄存器，却没有在约束表里声明这个
破坏。分配器相信这个寄存器的内容仍然属于原来的虚拟寄存器，于是后续代码直接拿它做
别的用途——或者更常见的是，认为不必在序言里保存它，也不必在尾声里恢复它。程序在
这段函数内部看起来完全正常，出错的是它的调用者：调用者的循环变量在返回后被换成了
别人的值，或者调用者自己的栈帧内容被破坏，症状因此出现在距离这段片段几条指令甚至
好几层帧之外，有时只在特定迭代次数或特定优化档位下才现形。

正确做法是把事实全部登记在约束表里：哪些寄存器被写坏就列入破坏项，读写内存就在
约束里表示出来，需要固定落在某个寄存器就用对应的约束字母要求，不希望片段被重排或
删掉就使用带 volatile 语义的形式，多次执行的片段还要保证输入输出与临时值的关系
写清楚。各家内嵌汇编的具体语法互不兼容，有的实现甚至不提供，因此这里只给类别：
输出、输入、破坏三张表必须完整。约束字母与寄存器类别属于实现文档，动手前查它，
不要凭别家的写法套。校验方法回到本章的验证路径：编出汇编看这个函数的序言里有没有
应有的保存，反汇编链接产物看调用者有没有被破坏的痕迹。
\end{abibox}

\begin{lstlisting}[style=cstyle, caption={约束表要如实登记副作用（GNU 家族风格示意，其他实现语法不同）}]
/* 目的：在写入之后给编译器一个不可穿透的重排屏障 */
void publish(int *slot, int value) {
    *slot = value;
    __asm__ __volatile__ ("" : : : "memory");  /* 声明这段代码可能读写内存 */
}
/* 只做屏障时不写任何操作数，副作用全部落在破坏项里，
   这样分配器不会把活跃值放在可能被片段改坏的位置 */
\end{lstlisting}

\begin{bestpractice}[title={改档位或加内联汇编后的最小验证三步}]
第一步，把源码编成汇编文本读一遍：只看这个函数的序言与尾声、参数与返回值的落点、
以及预期发生的那处改写，一眼就能发现“保存的寄存器列表变了”或“改写根本没发生”。
第二步，反汇编已链接产物，按符号核对最终代码：分配与折叠决策在链接期还可能被
\Mx{LTO} 或段合并改动，产物里的形态才是跑起来的形态。第三步，跑一遍带展开信息的
回归：用带符号与展开表的构建触发一次异常穿越多层帧，或让剖析工具采一次调用栈，
确认回溯链完整。三步各查一层，缺一步就有一个失败模式没人守着。
\end{bestpractice}

寄存器是有限的物理资源，\Mx{ABI} 是分配它时必须遵守的规则，展开表是事后找回上下文的
唯一凭据。本章的三条结论是：指令选择由目标描述里的合法性与成本决定，寄存器分配的
代价由调用约定折算，可调试性由是否把信息编进二进制决定。下一章进入这些决策的落盘
形式：目标文件如何记录段、符号与重定位。

\part{产物、链接与加载}
\chapter{目标文件与符号表：\Elf、\PeFmt 与 \Macho}\label{ch:cl07-object}

一个编译单元离开编译器之后，并不会立刻成为可执行程序，而是先成为一份“半成品”。
这份半成品同时欠着两个答案：程序里的地址该落在哪里，程序里的名字该指向谁。
第 \ref{ch:cl06-codegen} 章把指令选好、把寄存器分配完，交出来的仍然是一个带空洞的文件。
本章讲这些空洞长什么样，以及 \Elf、\PeFmt、\Macho 三种容器分别用什么结构去描述它们。

\section{目标文件的两件未决之事}\label{sec:cl07-pending}

第一件未决之事是地址。代码里一条取地址的指令、一张表里的一个指针槽位，
它们在编译时只能写成占位值：编译器不知道这段代码最终会落在内存的哪一页，
甚至不知道它引用的那个变量会被放进哪个文件。于是目标文件必须携带一份
“修补清单”：在哪个偏移、按什么规则、把哪个符号的值填进去、再加多少。
这份清单就是重定位记录，第 \ref{sec:cl07-reloc} 节展开。

第二件未决之事是名字。跨翻译单元的引用在编译时只知道一个字符串，
不知道它最终对应哪段机器码。目标文件因此必须携带一份符号表：
哪些名字由本文件定义、哪些名字只是被引用而本文件没有定义、
每个名字绑定为全局还是局部、强度是强是弱、可见性是否对外隐藏。
链接器把一堆目标文件拼成镜像的过程，就是把第一份清单逐项兑现、
把第二份清单里的未决名字逐一绑定的过程，详见第 \ref{ch:cl08-static-link} 章。

值得单独说清的一件事是：可执行文件与目标文件不是两种格式，而是同一族格式的
两种角色。三种容器各自都是这样，但“同一族”的含义并不相同。
在 \Elf{} 里，可重定位文件与可执行文件、共享对象共用同一个文件头结构，
文件头里的类型字段区分身份（可重定位文件的取值通常写作 \Tp{ET\_REL}，
完整取值以 \Elf{} 规范文档为准）：目标文件带节头表、符号表和重定位节，
一般没有程序头表；可执行文件反过来，加载器只认程序头表。
\PeFmt{} 这边，目标文件是纯粹的 COFF 记录：节表加符号表加重定位记录；
镜像在同样的 COFF 骨架外面补上 MS-DOS 头与可选头，再由数据目录把导出表、
导入表、基址重定位表等大块元数据串起来。头一份在对象里缺席、在镜像里必备，
这正是“同一族的两种角色”。\Macho{} 则把两种角色压进同一个结构：
头部加一串加载命令，头部的文件类型字段声明这份文件是对象、可执行文件还是
动态库（具体常量名以系统头文件与 \Macho{} 格式文档为准），
加载命令的种类与数量随角色伸缩，骨架本身不变。

\begin{guideline}[title={读目标文件的三条主线}]
拿到任何一份目标文件，只看三条主线：符号表回答“谁是谁”，
重定位记录回答“哪里要填数”，节表与属性回答“内容最终归哪块内存”。
三种容器换了三层皮，这三条主线一根都没少。
\end{guideline}

\section{三种容器的骨架对照}\label{sec:cl07-containers}

\subsection{\Elf：链接视图与加载视图并排摆}\label{sec:cl07-elf}

\Elf{} 文件以固定长度的文件头开头，随后是内容本体，表通常排在文件尾部。
它有两张彼此独立的表。节头表把文件切成一组节：\Sect{text}、\Sect{data}、
\Sect{rodata}、各条重定位节、符号表 \Sect{symtab} 与字符串表 \Sect{strtab}，
这是链接器工作的视图。程序头表把同一批字节切成若干按内存布局的段：
每条目声明“把这批字节映射到虚拟地址区间某处、给什么权限”，
这是加载器工作的视图，只在可执行文件与共享对象里出现。
两张表描述的是同一堆字节，划分方式却可以不同。
节名、段名之外的细节（例如段的类型取值 \Tp{PT\_LOAD}）以 \Elf{} 规范为准。

\subsection{\PeFmt：头叠头，再用目录索引}\label{sec:cl07-pe}

\PeFmt{} 镜像从一段兼容用的 MS-DOS 头开始，其尾部字段（通常写作 \Tp{e\_lfanew}）
指出真正的 \PeFmt{} 签名所在。签名之后依次是标准 COFF 文件头（机器类型、
节数量、时间戳、可选头大小、文件特征），可选头（\Mx{PE32} 与 \Mx{PE32plus}
两个魔数分别对应 32 位与 64 位，内含镜像基址、入口相对虚拟地址、子系统类型），
然后是节表。每个节条目带名字与一组特征位（可读、可写、可执行，以及
“内容不参与镜像”这类标志）。可选头末尾是一串数据目录：
导出表、导入表、基址重定位表、调试目录、异常数据目录等各占一个目录项，
目录项只是“相对虚拟地址加大小”的指针，真正的结构散落在各节里。
对象文件（\Tp{.obj}）没有 MS-DOS 头也没有可选头，只有 COFF 文件头、节表、
符号表与按节挂出的重定位记录；文件尾部还挂一张字符串表存放长符号名。

\subsection{\Macho：一条头部接一串加载命令}\label{sec:cl07-macho}

\Macho{} 的骨架是“固定头部 + 变长加载命令序列”。头部声明魔数（区分 32/64 位
与字节序）、文件类型与加载命令的条数和总长度。每条加载命令自带类型与长度，
链接器与加载器按顺序读过去，认识的处理、不认识的按规则跳过，
这让格式本身可以演化。段级命令（\Mx{64} 位下通常写作 \Tp{LC\_SEGMENT\_64}）
一次声明一段虚拟地址区间及其在文件中的范围，段里再嵌若干节：
代码段里的 \Tp{\_\_text}、只读数据里的 \Tp{\_\_const}、未初始化数据节等。
\Macho{} 是段套节的两层结构，而 \Elf{} 的段与节是同一批字节的两种切法。
符号表、字符串表与重定位信息集中存放在 \Tp{\_\_LINKEDIT} 段：
命令序列里只放它们的文件偏移与长度，内容本体在段尾排成一串。
调试信息则留在各对象文件里，由专门工具事后归拢，见 第~\ref{sec:cl07-debug} 节。

下面两段输出是同一个两文件小程序在 \Nix{} 与 \Win{} 上的产物骨架。
为免误导，这里只保留字段名与层次，真实数字（地址、尺寸、条数）一律从略，
内容均属示意，不是本仓库实测。

\begin{lstlisting}[style=outstyle, caption={示意：\Elf{} 可重定位文件的节头骨架（字段名层次，数值从略）}]
# 示意输出：readelf -SW 的裁剪版，只保留名字、类型、标志三个字段
# Section Headers（目标文件角色：有节头表，无程序头表）
  .text      PROGBITS   AX     # 代码：分配、可执行、不可写
  .rela.text RELA       ---    # 针对 .text 的重定位记录，每条带加数
  .data      PROGBITS   WA     # 已初始化可写数据
  .bss       NOBITS     WA     # 未初始化数据：只记尺寸，不带文件内容
  .rodata    PROGBITS   A      # 只读数据与常量池
  .eh_frame  PROGBITS   A      # 展开与回溯表（角色见第 13 章口径）
  .comment   PROGBITS   ---    # 产出本文件的工具链签名串
  .symtab    SYMTAB     ---    # 符号表：名字走 strtab
  .strtab    STRTAB     ---    # 符号名与节名字符串池
  .shstrtab  STRTAB     ---    # 节名字符串池
# 符号表里的三种典型条目（对应下文的符号视图一节）
# 本文件定义的函数  段内偏移 + 大小 + GLOBAL
# 本文件引用的函数  UND 标记，等链接器兑现
# 未定义全局变量    COMMON 身份，尺寸先记着
\end{lstlisting}

\begin{lstlisting}[style=outstyle, caption={示意：同一程序在 \Win{} 上的对象与镜像骨架对照（层次视图）}]
# 示意输出：COFF 对象（probe.obj，dumpbin /headers /symbols 的裁剪版）
  COFF FILE HEADER        machine / number of sections / size of optional header = 0
  SECTION TABLE           .text  .data  .bss  .rdata   各带 characteristics
  SYMBOL TABLE            名字表 + 每个符号的节号与节内偏移
  LINE / RELOC RECORDS    按节挂出的重定位记录：符号 + 偏移 + 类型
# 链接后的镜像（probe.exe，dumpbin /headers 的裁剪版）
  DOS HEADER              e_lfanew 指向真正的头
  PE SIGNATURE + FILE HEADER
  OPTIONAL HEADER         PE32 或 PE32 加宽位形态 / image base / entry point
  SECTION TABLE           .text CODE 可执行不可写
  DATA DIRECTORIES        export / import / base reloc / debug 各一指针
# 目录之外：调试信息不在镜像里，落在同名的 .pdb 文件（见 7.6 节）
\end{lstlisting}

三种容器对同一批概念的叫法，汇总在下表。“链接前后是否存在”一列说的是
从可重定位文件变成可执行镜像时，这个结构还在不在。

\begin{longtable}{@{}>{\raggedright\arraybackslash}p{0.13\textwidth} >{\raggedright\arraybackslash}p{0.17\textwidth} >{\raggedright\arraybackslash}p{0.17\textwidth} >{\raggedright\arraybackslash}p{0.17\textwidth} >{\raggedright\arraybackslash}p{0.13\textwidth}@{}}
\caption{三种容器的结构对照：概念、叫法与存续}\label{tab:cl07-containers}\\
\toprule
概念 & \Elf{} 里的叫法 & \PeFmt{} 里的叫法 & \Macho{} 里的叫法 & 链接前后 \\
\midrule
\endfirsthead
\toprule
概念 & \Elf{} 里的叫法 & \PeFmt{} 里的叫法 & \Macho{} 里的叫法 & 链接前后 \\
\midrule
\endhead
\bottomrule
\endfoot
文件骨架 & 文件头 + 两张表 & 头叠头 + 节表 + 数据目录 & 头部 + 加载命令序列 & 两种角色同在 \\
加载视图 & 程序头表（段） & 无独立表；可选头加基址重定位 & 段级加载命令 & 前者仅镜像有 \\
内容切法 & 节头表分节 & 节表分节 & 段下再嵌节，两层 & 前后都有 \\
符号表 & \Sect{symtab}（另有动态符号表） & 对象有 COFF 表；镜像剩导出导入表 & 链接编辑区的符号表 & 发布前多被剥 \\
字符串池 & \Sect{strtab} 与节名表 & 文件尾字符串表 & 链接编辑区字符串区 & 对象里必在 \\
重定位 & 每节配一条重定位节 & 对象按节记录；镜像只留基址重定位 & 链接编辑区重定位信息 & 基址类才保留 \\
调试信息 & \Sect{debug\_info} 一族内嵌节 & CodeView 记录另存 \Tp{.pdb} & 调试留在对象，镜像带归拢线索 & 发布前剥离 \\
展开与回溯 & \Sect{eh\_frame} 等表 & 异常数据目录指向的表 & 展开信息节 & 前后都可能带 \\
工具链签名 & \Sect{comment} 存编译器串 & 信息并入符号与节 & 并入加载命令与链接编辑区 & 常与调试同剥 \\
\end{longtable}

\begin{compare}[title={一张表还是两张表}]
\Elf{} 把“链接时怎么看”与“加载时怎么看”做成两张表，切换角色时旧表可以整张丢掉；
\PeFmt{} 与 \Macho{} 各只有一套节或段的清单，靠头部字段和链接器改写来兼任两种角色。
这一区别直接决定了排障时的取材点：\Elf{} 上看“丢了哪张表”，
\PeFmt{} 上看“哪个目录项空了”，\Macho{} 上看“哪条加载命令没被理解”。
\end{compare}

\section{节与段的语义}\label{sec:cl07-sections}

内容分类的骨架在三家是一致的：代码、只读数据、已初始化可写数据、未初始化数据，
再加上元数据一类。对应到惯用名：\Sect{text} 装指令；\Sect{rodata} 装字符串常量、
跳转表与虚表这类只读对象；\Sect{data} 装带初始值的可写全局；
\Sect{bss} 装不带初值或初值为零的可写对象；\Sect{comment} 装“我用哪个编译器
产出你了”的签名串。元数据一族还包括构造与析构函数指针表、
动态注册表等，各家节名互异，这里不逐一列举，以各格式规范文档为准。

\Sect{bss} 值得多说一句：它在文件里通常不占任何字节，节头里只登记尺寸与地址需求，
加载时由系统给一片清零内存。“程序体积比数据总量小”这类直觉，
在未初始化大数组上最容易被反直觉打穿：几十兆的缓冲数组在文件里几乎免费，
代价只在运行时。调试信息节一族（\Sect{debug\_info}、\Sect{debug\_line} 等）
不参与运行时布局，链接器把它们从镜像里摘出去或整体剥离，见 第~\ref{sec:cl07-debug} 节。

展开表一族管的是“异常抛出或崩溃回溯时，怎么从当前指令反推调用栈”。
\Nix{} 与 \Mac{} 把它做成带表的内嵌节（\Nix{} 下常见 \Sect{eh\_frame}），
\Win{} 的 64 位约定则把展开信息做成镜像内被数据目录指点的表；
这些节与表的确切命名规则与各家的口径细节不在本章展开，
第 \ref{ch:cl13-cxx-abi} 章按平台讲展开与 ABI 的约定。
这里只要记住一条：展开表是被链接期合并、被运行期查询的只读数据，
它跟着代码段走，剥离它等于放弃跨模块回溯的能力。

\begin{note}[title={权限与分段是一件事的两面}]
编译器把“初值为零或不带初值的全局”划入 \Sect{bss}，把“带初值的全局”划入
\Sect{data}，把字符串常量划入 \Sect{rodata}，这不是文件整理癖，
而是决定了加载后三块内存各自的读写权限。写字符串常量之所以是未定义行为，
物理上就源于这个划分；语言层面的口径见未定义行为讨论的第 \ref{ch:cl06-codegen} 章后话。
\end{note}

\section{符号视图：名字的世界}\label{sec:cl07-symbols}

符号表条目承载的信息可以归纳为四问：名字是什么；由本文件定义还是只被引用
（未定义条目通常以 \Tp{UND} 或空值标记）；绑定是全局、局部还是弱
（\Tool{nm} 视图里 \Tp{LOCAL} 与 \Tp{GLOBAL}、\Tp{WEAK} 是三个字母档位）；
可见性是默认还是隐藏。隐藏意味着“这个全局名字只在本最终模块内部有效”，
动态链接期对外不可见，第 \ref{ch:cl09-dynamic} 章讲它的后果。

COMMON 符号要单独讲。C 语言允许 tentative definition：
\Sym{int counter;} 放在文件作用域、没有初始值，同一个名字在多个翻译单元里
各出现一次是合法程序。若每个目标文件都把 \Sym{counter} 记成普通全局定义，
链接器必然报多重定义。于是这类符号被登记为 COMMON 身份：
“本文件对这个名字主张一块多大的可写空间”，链接器把所有主张合并，
取最大尺寸分配一块归入 \Sect{bss}。这就是“多份 tentative 定义可以合并”的机制；
\CPP{} 里对应的合并机制是第 \ref{ch:cl08-static-link} 章的可折叠符号，
两者语义来源不同，别混为一谈（语言口径回指第 \ref{ch:cl03-resolution} 章）。

五个查看工具各管一片，工具名都是确证的，参数以各自的帮助文档为准：
\Tool{readelf}（含 \Tool{llvm-readelf}）读 \Elf{} 的文件头、节头表、
程序头表与符号表；\Tool{objdump} 兼顾节布局与反汇编，重定位清单可用它的
按节转储选项看；\Tool{nm} 只给符号清单，是最快的“这个库里到底有没有它”手段；
\Tool{dumpbin} 是 \Win{} 工具链的对应物，按头、符号、重定位分几个开关输出；
\Tool{otool} 是 \Mac{} 侧的对应物，加载命令、动态库依赖与反汇编各一个开关。
跨平台排查时，\Tool{llvm-objdump} 与 \Tool{llvm-nm} 对三种容器都能读，
适合在同一台机器上对照三种格式。

\begin{lstlisting}[style=shellstyle, caption={查看命令序列：同一个程序在三个平台上的取材点（示意流程）}]
# 第 0 步：在任何平台上先固定事实——未定义符号清单
gcc -c probe.c -o probe.o
nm probe.o                    # U 未定义 / T 定义在代码段 / C 为 COMMON
readelf -SW probe.o           # 节头表：名字、类型、标志三列
readelf -sW probe.o           # 符号表：值、大小、绑定、可见性、所指节
objdump -h probe.o            # 节的尺寸与标志的紧凑视图
objdump -r probe.o            # 重定位清单：偏移 + 类型 + 目标符号

# 第 1 步：同一骨架走一遍（示意流程，具体参数以版本为准）
clang -c probe.c -o probe.o
otool -h probe.o              # Mach-O 头部：文件类型与加载命令计数
otool -l probe.o              # 逐条加载命令：段、节、符号表位置
nm -m probe.o                 # 符号的类型与所归节
# Windows：先编译出 .obj，再让 dumpbin 分三个视角看
cl /c probe.c
dumpbin /headers probe.obj    # COFF 文件头与节表
dumpbin /symbols probe.obj    # 符号表
dumpbin /relocations probe.obj # 按节挂出的重定位记录
# 链接成镜像后再看一次，体会“两种角色”的差别：
link probe.obj /OUT:probe.exe /MAP
dumpbin /headers probe.exe    # 这次会多出可选头与数据目录
\end{lstlisting}

\begin{abibox}[title={“未定义符号”与“符号明明在”}]
链接器报 \Sym{undefined symbol: draw\_widget}，而你用 \Tool{nm} 在库里
明明看到了一个叫 \Sym{draw\_widget} 的定义。三条最常见的原因，各对应一种核对法。
其一，名字修饰不同：引用方按 \CPP{} 装饰名找 \Sym{\_ZN6Widget4drawEv} 这样的条目，
而库里那份是 \CL{} 编译的裸名——定义方或引用方少了一层 \Sw{extern “C”} 声明
（回指第 \ref{ch:cl03-resolution} 章），核对办法是给 \Tool{nm} 加反装饰选项
（\Tool{nm} 与 \Tool{objdump} 均支持，参数名以帮助文档为准）把两边都还原成
源码名再对表。其二，可见性被隐藏：库构建时把该符号标成隐藏，
它存在于符号表却不进导出集合，动态链接找不到它——用带导出视图的转储
（\Tool{dumpbin} 的导出开关或 \Tool{nm} 加动态表选项）确认“在表里”与“在表外”。
其三，类型不匹配：引用要函数、库里定义成了同名变量，或反过来，
符号表里的类型字段能区分，工具输出里那列字母（\Tp{FUNCTION} 与 \Tp{OBJECT}
一档）就是干这个的。三条都排不掉，再怀疑库版本与搜索路径。
\end{abibox}

\section{重定位项：一条记录里装了什么}\label{sec:cl07-reloc}

一条重定位记录固定承载四样东西：修补位置（哪个节的哪个偏移，或相对节首的偏移）、
目标符号（要填谁的值）、类型（怎么算：填绝对地址、填相对跳转距离、
还是填表槽位编号）、加数（在符号值上再加一个常数）。四种信息缺一，
链接器就只能猜。为什么目标文件里那些地址字段是占位：编译期一切布局未定，
指令只能先按“差值等于多少”的假设写成全零或全一，链接器兑现重定位时
把真值补进去——占位不是缺陷，是“先编译、后组装”这个流水线前提的必然。
顺带一提，链接器还会把一部分决策就地改写掉：能直接定下来的引用在输出里
不再保留记录，只有加载时才可能变的部分（镜像基址、符号的实际落位）
才会被写进最终产物的修补清单。

三种平台的重定位类型族，这里只用描述性说法，逐条编号不确证就不写。
\Elf{} 的重定位节分两型：带显式加数字段的（节名以 \Tp{rela} 为词干）与
把加数就地藏在占位字里的；类型名形如 \Sym{R\_X86\_64\_PC32} 一族，
按目标架构划分，完整清单以各架构的过程接口规范（psABI）文档为准。
\PeFmt{} 的对象文件按节携带重定位记录，字段是符号加节内偏移加类型；
链接成镜像后，绝大多数重定位被就地兑现，只有基址重定位表保留下来，
记录“某页内哪些位置要按实际加载基址与首选基址之差做加法”——
这正是地址随机化能在 \Win{} 上工作的底层机制。
\Macho{} 把重定位信息放进链接编辑区，分“按虚拟内存地址修补”与
“按页内偏移修补”两类，各自的类型编码同样以格式文档为准；
共享库与可执行文件在这一点上角色互换得比 \Elf{} 复杂（回指第
\ref{ch:cl09-dynamic} 章）。

\begin{warning}[title={绝对寻址让镜像失去可重定位性}]
链接器若在生成阶段就把绝对地址烧进代码（比如非位置无关的目标码被直接
挂进需要随机基址的产物），后果要么是一整张巨大的基址修补表拖慢加载，
要么是干脆拒绝链接并抱怨地址越界。看到“重定位数量异常巨大”或
“必须按首选基址加载”这类报错，先回符号表与节视图确认：
是哪条引用逼出了绝对地址，而不是先去翻链接参数。
\end{warning}

\section{调试信息：内嵌与分离两种走法}\label{sec:cl07-debug}

调试信息存在形式只有两大族。一族是 \Mx{DWARF}：类型、变量位置、行号映射
编码成多段结构化数据，\Elf{} 与 \Macho{} 把它们作为 \Sect{debug\_info}、
\Sect{debug\_line} 等节内嵌在目标文件里；进入最终镜像前，链接器通常
把这些节整体丢弃或剥离，运行时调试器再从别处取回。
另一族是 \Mx{CodeView} 记录加 \Tp{.pdb} 文件：\Win{} 工具链把类型信息与
源文件映射写进与镜像同名、扩展名为 \Tp{.pdb} 的独立文件，
镜像里只留一个指路的调试目录项。“调试信息与镜像分离”不是一时兴起：
\Mx{PDB} 体量常与镜像同数量级，内嵌会拖慢分发，而调试器本来就按路径查找它。
\Mac{} 走的是折中：调试信息留在各对象文件里，链接器不搬运，
由 \Tool{dsymutil} 事后按符号表引用把它们归拢成 \Tp{.dSYM} 包，
镜像只带一枚唯一标识供配对。

\Tool{strip}（及其各家等价物）删除的是符号表与调试节，不删机器码。
由此产生一条张力：发布产物越干净越小，事后剖析与崩溃回溯就越难还原符号。
性能剖析与回溯符号的取舍属于第 \ref{ch:cl15-endtoend} 章端到端流程的一环，
测量方法本身在性能测量分册里展开，这里只给结论式的操作口径：
剥离可以照做，但每个发布产物必须能按唯一标识或版本号找回到未剥离副本，
否则线上的火焰图只会是一串十六进制地址。

\begin{bestpractice}[title={发布产物的符号保全}]
按“镜像剥离、副本归档”的成对做法发布：随包产物用 \Tool{strip} 或各家等价工具
去掉符号与调试节；同版本的未剥离镜像、\Tp{.pdb} 或 \Tp{.dSYM}、
链接器导出的映射文件（如 \Tool{link} 的映射开关产物）一起归档，
归档键用镜像内的唯一标识或“版本号加提交号”的复合键。
排障时的第一步从“用 \Tool{addr2line} 把地址翻译回行号”开始，
前提就是这份归档还在。
\end{bestpractice}

\chapter{静态链接：符号解析、可折叠符号与一致性}\label{ch:cl08-static-link}

静态链接是整条流水线里最像“装配”的一步：上游交来一堆各自独立编译的目标文件
与若干库，链接器要把名字对名字、地址对地址地拼成一个自足的整体，
交出去的时候不再欠任何外部引用（动态链接欠的那部分除外，见第
\ref{ch:cl09-dynamic} 章）。本章讲装配的两件正事、库的两种形态带来的顺序问题、
可折叠符号为什么必须存在，以及一致性规则在链接期的两种表现。

\section{链接器的两件正事}\label{sec:cl08-two-jobs}

第一件是符号解析：把每个未定义引用绑定到全体输入里唯一的一份定义上。
输入文件之间没有主次，绑定关系只由“谁有定义、谁只有引用”决定，
第 \ref{ch:cl07-object} 章的符号视图正是这件事的账本。
第二件是地址分配与布局：把各输入文件的同类节合并（所有目标文件的
\Sect{text} 并成一块代码区，所有 \Sect{data} 并成一块数据区），
给合并后的每一段安排最终地址，然后逐项兑现重定位记录。
两件正事之间有严格的先后：布局未定时符号值还是相对量，
所以解析先行、布局居中、兑现收尾。
还要认清链接器身份的宽度：同一个程序常被拿来干“没干完”的活——
半链接（partial linking）把一组目标合成一个更大的目标文件，产物仍是可重定位
身份，重定位记录只是被部分兑现；创建归档更是只动符号索引不动地址。
“链接器”这个名字描述的是一套阶段能力，不是“必然产出可执行文件”的承诺。

\begin{flow}[title={静态链接内部的四个阶段}]
读入与解析：把全部目标文件与候选库的符号表并入一张全局符号字典；
库的处理：按引用情况决定哪些库成员进入输入集合（见下一节）；
合并与布局：按段类型与对齐要求给每块内容定地址，入口符号在此定位；
兑现：按定下的地址逐条改写重定位，写出最终镜像或归档。
其中第三、四阶段对 \Elf、\PeFmt、\Macho 三家只是叫法不同，动作一致。
\end{flow}

\section{库的两种形态与链接顺序}\label{sec:cl08-archives}

交给链接器的输入有两种身份。目标文件是全量参与的：只要列在命令行里，
它的每个符号都进入解析集合，列与列之间的先后毫无意义。
归档（archive，\Nix{} 与 \Mac{} 下通常叫 \Tp{.a}，\Win{} 下叫 \Tp{.lib}）
则完全不同：它是打包在一张目录表后面的若干目标成员，链接器按“按需抽取”
的规则处理——只有当某个成员能提供当前悬而未决的未定义符号时，
它才被拉进解析集合，否则整个成员原封不动留在包里。

按需抽取带来一个经典故障：链接顺序。多数链接器按命令行从左到右扫描，
扫到归档时，此刻“未决”的引用集合由它左边的输入决定。
\Tp{a.o} 依赖 \Tp{b.o}、\Tp{b.o} 又回调 \Tp{a.o} 的情形，把这两个
目标打包成 \Tp{ab.a} 放在 \Tp{a.o} 与 \Tp{b.o} 之间，就可能谁也没能触发对方抽取，
报出“未定义符号”；而把两个目标文件直接列出来再链接却成功。
同理，“引用在后面、库在前面”也会扫空：库被扫到时引用方还没入场，
未决集合里不含那个名字，成员不被抽取，后面的引用无处可绑。
把相关目标强制全部纳入的能力各家都有提供：\Nix{} 与 \Mac{} 的链接器提供
“包裹整个归档、全量抽取”的开关（\Elf{} 世界常写作 \Opt{whole-archive}
一族，\Win{} 的 \Tool{link} 有对等能力），用于注册表模式这类
“成员看似没人引用、实则被构造函数挂载”的库。具体开关拼写以版本为准。

\subsection{归档的另一个身份：导入库}\label{sec:cl08-importlib}

\Win{} 上同一种扩展名背负了两种身份：静态库归档，以及动态库的导入库。
导入库不是归档，它的“成员”只为链接期提供指路信息，让调用方在链接时
把对 \Tp{.dll} 的引用登记进镜像的导入表。本章的规则只适用于前者：
按需抽取、顺序敏感、全量开关；导入库的解析发生在动态链接那一侧，
口径见第 \ref{ch:cl09-dynamic} 章。混用两种 \Tp{.lib} 报出的错误长得几乎一样，
最快的区分办法是看文件开头有没有归档目录。

\begin{lstlisting}[style=shellstyle, caption={归档的创建与两种链接顺序（示意流程）}]
# 打包：把三个目标做成归档，再补索引（ar s 与 ranlib 等价于建索引）
ar rcs libcodec.a encode.o decode.o tables.o
ar t libcodec.a               # 列成员，确认打包内容
nm -A libcodec.a              # 带成员名前缀看符号，U 与 T 一目了然

# 顺序一：库在前、引用在后——decode.o 用到 encode.o 里的符号时，
# 归档已被扫过，可能报未定义
g++ decode.o -L. -lcodec
# 顺序二：引用在前、库在后——扫到归档时未决集合已含全部需求
g++ -lcodec decode.o          # 仍然不对：方向反了
g++ decode.o -L. -lcodec      # 这才是要的形状
# 循环依赖或挂载型库：强制全量抽取（GNU ld 与 clang 前端写法）
g++ main.o -Wl,--whole-archive libcodec.a -Wl,--no-whole-archive
\end{lstlisting}

\begin{compare}[title={目标文件与归档成员，谁进解析集合}]
目标文件列出的瞬间就“全员入场”，符号解析、重复定义检查、折叠机会
都按全量发生。归档成员则始终处于“候补”状态：没被抽中的成员连符号都不入字典，
也就不参与重复定义检查——同一个符号在包外目标里有一份定义、
在包内成员里还有一份更强的定义时，很可能两者都不报错，
最终链接进来的那一份取决于扫描顺序。这类“静默二选一”是排版本问题的重灾区。
\end{compare}

\section{弱符号与可折叠符号}\label{sec:cl08-weak}

有一类定义天生会出现多份一模一样的拷贝：类模板的成员函数实例（每个使用它的
翻译单元都可以合法地再生成一遍）、声明为 \Sw{inline} 的函数（多单元内联展开后，
取地址的那个单元仍要留一份实体）、虚表与虚函数元数据、\Sw{typeinfo} 与
\Sw{typeinfo\_name}（异常与 \Fn{dynamic\_cast} 的运行时凭据，见第
\ref{ch:cl13-cxx-abi} 章）。如果这些符号按普通全局强定义登记，
任何跨单元的正常使用都会撞上多重定义错误。
于是它们被登记为另一身份：弱符号或可折叠符号。链接器遇到多份弱定义时
任选其一（通常取第一个或尺寸最小的那个），其余丢弃，不再报错。
归档的“按需抽取”也因此有了配合：同一定义被抽进两次也不会互相冲突。
在 \PeFmt{}/COFF 的世界里，承担同样职责的机制是折叠分组的节记录
（常以 \Mx{COMDAT} 称呼）：每个可折叠实体打包成“一个节加一个折叠头”，
折叠头声明选择规则，链接器按规则留一份、扔其余；
\Mx{ELF} 与 \Macho{} 则直接用符号的弱绑定来登记同一意图。
两套登记法在源码层不可见，差别只在链接器读到的账本格式。

\begin{lstlisting}[style=cppstyle, caption={必须以“多份可任选”身份存在的定义长什么样}]
template <typename Tag>
struct Counter {
    static long total;              // 声明：每个实例化各有一份
    inline static long bump() {     // inline 静态成员：实体可出现多次
        return ++total;
    }
};
template <typename Tag> long Counter<Tag>::total = 0;   // 旧式的盘外定义

struct Base { virtual void run(); virtual ~Base(); };
void use(Base* b) { b->run(); }     // 虚表与 run 的实体同样要能折叠
\end{lstlisting}

折叠的另一层是相同代码合并（ICF）：把内容完全相同的两个函数体合并成一个，
调用方两路都指向同一地址。收益实在——体积与指令缓存占用都下降；
语义代价却隐蔽：函数指针相等性。若程序里同时存在两个不同的函数，
只因机器码碰巧相同就被折叠成一个，那么“取这两个函数的地址比较是否相等”
就会从假变成真。标准允许实现做这种合并，但允许不等于随便：
对函数指针身份有语义依赖的代码（第 \ref{sec:cl08-odr} 节末尾会再遇到它）
必须显式关掉折叠。各家开关名不一，能力描述为主，拼写以链接器文档为准。

还有一层安全阀常被忽略：折叠的判定单位是实现选择的，多数实现按“整节内容
加引用结构”比对，而编译器把不同实体塞进同名折叠分组的概率并不为零——
调试构建通常顺手禁用折叠，正是为了避免断点停进“别人的函数体”这种
指鹿为马的现场。发布构建开启折叠前，先确认代码里没有依赖
函数地址互不相等来区分身份的逻辑，这比记开关名重要得多。

符号种类与规则汇总如下表。“折叠可行”列说的是多份同名定义出现时，
链接器能不能自行选一份了事。

\begin{longtable}{@{}>{\raggedright\arraybackslash}p{0.14\textwidth} >{\raggedright\arraybackslash}p{0.18\textwidth} >{\raggedright\arraybackslash}p{0.21\textwidth} >{\raggedright\arraybackslash}p{0.12\textwidth} >{\raggedright\arraybackslash}p{0.15\textwidth}@{}}
\caption{符号身份与链接期冲突规则}\label{tab:cl08-symbols}\\
\toprule
种类 & 谁能引用到 & 同名多份时的规则 & 折叠可行 & 备注 \\
\midrule
\endfirsthead
\toprule
种类 & 谁能引用到 & 同名多份时的规则 & 折叠可行 & 备注 \\
\midrule
\endhead
\bottomrule
\endfoot
全局强定义 & 全体输入与最终导出 & 两份强定义即多重定义错误 & 否 & 普通函数与非模板全局属于此类 \\
弱符号 & 全体输入 & 强与弱并存取强；全弱则任选一份 & 是 & 模板实例与 \Sw{inline} 实体的常用登记法 \\
COMMON & 全体输入 & 主张合并取最大尺寸落未初始化区 & 是 & \CL{} 暂定义合并的专用通道 \\
局部符号 & 仅本目标文件 & 各文件互不相干，天然多份并存 & 不适用 & 文件作用域 \Sw{static} 与内部链接实体 \\
隐藏可见 & 仅本最终模块 & 模块间同名隐藏符号互不可见 & 视实现 & 不进导出集合，动态期查找绕开它 \\
\end{longtable}

\begin{note}[title={弱与折叠是身份，不是质量标记}]
把符号登记为弱或可折叠，是“多份等价定义在语义上就是同一实体”的实现机制，
不表示链接器替你做了择优判断。等价的前提由编译期保证：
同一个模板在两个单元里实例化，要求两边的声明一致（下一节），
链接器只做名字与字节，不做语义比对。
\end{note}

\begin{note}[title={顺序故障的两种解法别混用}]
解循环依赖有两条路：把同一批目标与归档在命令行上重复列出（多给扫描器几次机会），
或全量抽取归档（绕开机会主义）。前者只改变扫描次序，成员依旧按需抽取，
产物里可能仍缺“没人引用但必须在场”的注册成员；后者改变入场规则，
一次到位，代价是全量体积与全量初始化。选哪条看库的写作风格，而不是报错的字面形状。
\end{note}

\section{一致性规则在链接期的表现}\label{sec:cl08-odr}

\CPP{} 的一致性规则要求：同一实体在不同翻译单元里的定义必须逐字等价
（语言口径在第 \ref{ch:cl03-resolution} 章）。链接器看不到语言层面的定义，
它只比对符号名与节内容。于是违规有两种下场，都不在违规处报错。
第一种，同名不同实体：两个单元各自看到不同版本的类布局
（差一个成员、差一重继承、宏开关不同），各自生成自己的成员函数，
符号名却一模一样——弱定义折叠或“强与弱并存取强”会让其中一份胜出，
而两份对字段偏移的假设不同，症状出现在读写这些字段的任何无关代码里，
崩溃点离病因十万八千里。第二种，被折叠的两段代码指针相同而身份预期不同，
见上一节的函数指针相等性讨论。链接器无法诊断这类问题，是因为它的比较
只到符号名与字节内容层面：名字相同、哈希不同，它顶多给你多重定义错误
或任选一份；名字相同、哈希也相同（比如布局差异常落在填充里），它连疑心都不起。
这也是为什么一致性检查只能靠构建纪律与编译期断言，见第
\ref{ch:cl15-endtoend} 章的流程收口。

\begin{warning}[title={跨单元宏配置不一致：最常踩的那条}]
最普遍的病根不是手改了两份头文件，而是同一头文件在不同单元里展开成不同形状：
A 单元带着某个功能宏编译，B 单元没带，共享的头里那个条件编译块让同一个类
少了一个成员；再或者不同单元链接了同一库的不同宏配置的预编译产物。
链接期一声不响，运行期在随机位置越界。
对策是把配置集中成构建层可见的单一来源：共享的编译宏由构建系统统一注入
（\Tool{CMake} 层用接口目标传播，见下一节），
禁止在源码里就地加开关改结构体形状；发布归档保留各单元的宏清单以便对质。
\end{warning}

事发之后的取证也全在链接期产物里：映射文件给出每个符号来自哪个输入成员，
把崩溃偏移落到符号上就能回推出事主属于哪个单元；再对照各单元的编译参数清单，
两边宏形状的差异常常一眼可见。这类工具备受冷落，是因为大家习惯把映射文件
当调试符号的替代品——它其实是“解析结果的黑匣子”，多重定义疑云、
静默二选一、布局突变，都先在它的成员归属列里露出马脚。

\section{布局控制：链接器为什么提供这层可编程性}\label{sec:cl08-layout}

链接器除了两件正事，还管一批布局参数：节合并规则（哪些输入节并进哪个输出段、
按什么次序）、对齐要求（段与段边界对齐到多少字节、超对齐数据的显式要求）、
入口符号（从哪个符号开始执行，以及运行时入口前的初始化钩子挂哪几个）、
栈大小（\Win{} 镜像把栈预留量写进可选头，是可确证的能力）、
镜像基址与它的随机化策略。这些参数多数有默认值，默认值的口径属于实现选择，
不属于语言承诺。
对齐参数里最容易吃亏的是“过对齐”：类型自己的对齐要求由编译期写进目标文件
（第 \ref{ch:cl06-codegen} 章的产物里带着它），链接器合并节时必须保住最严的那个
要求，否则 SIMD 载荷指令会在运行期以难以复现的方式报错。
这也是“把同类节并成一块”并非无脑拼接的原因之一：
合并次序、间隙填充、边界对齐都要回头影响指令能否合法执行。

把段精确放到某个物理地址、为不同存储区各写一块规则的做法，
靠的是链接脚本或各家等价的描述语言——那属于嵌入式与硬实时分册的主题，
本章不展开语法。这里只回答“链接器为什么要提供这层可编程性”：
通用平台上它服务极少数特例（代码预热、大页映射、性能实验）；
受限平台上它是刚需（向量表必须在固定地址）。可编程性的共同前提是
布局阶段本来就独立于解析与兑现，四个阶段的拆分（见本章第一节的框）
让“重排段”只是重跑后半程。端到端各阶段如何串起来，第
\ref{ch:cl15-endtoend} 章有一张总装图。

\section{可见性与导出：静态链接里的“对外”}\label{sec:cl08-visibility}

静态链接完成时，每个符号随结果镜像只有两种去向：进入导出集合（对外部可见、
可被动态链接器解析），或不进入。隐藏可见性走的就是第二条路：符号留在
内部符号表供调试与重定位使用，但不进导出与动态符号表，动态链接期
对其它模块彻底不可见。收益在 \ref{ch:cl09-dynamic} 章展开——查找集合变小、
抢占面变小、代码还能更紧凑——但决定是在静态链接这一步做出的：
库里没导出的符号，后面任何链接都变不出来。
隐藏还与折叠相互作用：内部实现符号如果被登记成可折叠的多份弱定义，
隐藏并不能替你把“任选一份”的选择逻辑讲清楚——折叠发生在链接期，
可见性发生在导出面，两本账要分开记，排障时也分开查。

语言与工具链入口方面，可确证的名字有这些：\Tool{gcc} 与 \Tool{clang}
经 \Sw{\_\_attribute\_\_} 提供 \Sw{visibility} 属性一族（取值形态为
\Sw{default}、\Sw{hidden}、\Sw{protected}，并可用编译开关
\Sw{-fvisibility=hidden} 改默认值）；\Win{} 工具链经 \Sw{\_\_declspec}
的 \Sw{dllexport} 与 \Sw{dllimport} 声明导出方向；
\Tool{CMake} 用 \Sw{generate\_export\_header} 把这套标记生成到头文件里。
各家完整清单以各自文档为准，不要跨平台抄属性。

\section{\Tool{CMake} 层的写法}\label{sec:cl08-cmake}

构建系统层能替链接期消解大部分隐患，抓手是接口与实现的分离：
链接需求拆成三档归属，自用而已、随接口传导给使用方、纯接口要求，
私有依赖不进入使用方的传导清单，“我自己链接什么”与“用我的人
要跟着链接什么”被显式分开，顺手也限制了导出集合的规模；
区分三档归属的语法载体是 \Sw{target\_link\_libraries}。

跨单元优化把链接期的折叠与解析推迟到编译期合流（见第
\ref{ch:cl04-ir} 章）；开启途径同样只有两条，都在构建层：目标的同名属性，
或者全项目范围的全局构建变量 \Sw{CMAKE\_INTERPROCEDURAL\_OPTIMIZATION}，
两者的取舍是风格问题不是语义问题，工具链是否支持则由构建系统自行探测。

\begin{lstlisting}[style=cmakestyle, caption={目标级写法：私有依赖、可见性默认值与跨单元优化}]
add_library(codec STATIC encode.cpp decode.cpp tables.cpp)
target_compile_options(codec PRIVATE -fvisibility=hidden)
target_compile_definitions(codec PUBLIC CODEC_FORMAT_V2)
target_link_libraries(codec PUBLIC util)
target_link_libraries(codec PRIVATE zlib)

set(CMAKE_INTERPROCEDURAL_OPTIMIZATION ON)
add_executable(app main.cpp)
target_link_libraries(app PRIVATE codec)
\end{lstlisting}

上例里 \Sw{zlib} 是 \Sw{PRIVATE} 依赖：它进入 \Sw{codec} 的链接命令，
却不传播给 \Sw{app}；\Sw{CODEC\_FORMAT\_V2} 是 \Sw{PUBLIC} 定义，
使用方必须同带，这正是防住上一节“宏配置不一致”的机制化写法。
\Sw{-fvisibility=hidden} 改了默认可见性后，导出要用显式标记白名单化，
否则整个库链接进可执行文件时一切照常，单独打成动态库那天才会发现“全空了”。
再补一层构建形态的选择：对象库（编译出目标文件、但不归档不链库）
把“谁进解析集合”的决定推迟到最终可执行目标那一行，多个共享同一组实现的
目标因此少踩“静态初始化重复”的坑；归档库则把决定提前固化。
两种形态没有优劣，只有“解析集合在哪一层成型”的差异，
与本节接口三档的选择是同一枚硬币的两面。

\begin{lstlisting}[style=outstyle, caption={示意：nm -A 视角下归档成员与引用方的符号账目（数值从略）}]
# 示意：archive 成员的符号身份；字母含义见第 7 章符号视图一节口径，地址位已抹去
encode.o:  T encode_frame          # 强定义：函数
decode.o:  T decode_frame
           U encode_frame          # 引用：等归档抽取兑现
           W Counter<1>::bump()    # 弱/可折叠：inline 实体
tables.o:  B crc_table             # 数据定义（未初始化区归属）
main.o:    C shared_scratch        # COMMON 主张：与别处合并取最大
# 链接后的镜像：U 应清零，W 只剩一份，T 与导出表里的名字对得上
\end{lstlisting}

\begin{abibox}[title={模板静态数据成员：两种相反的报错}]
同一个模板静态数据成员，工程 A 报“多重定义”，工程 B 报“未定义符号”，
看起来互相打脸，其实是同一条链的两端。报多重定义：成员的定义以非模板方式
（或借助不受折叠保护的旧链接路线）同时出现在多个单元的输入里，
两份都是强定义，折叠规则救不了强对强。报未定义：单元里只有声明，
定义在另一个从未被链接进来的翻译单元里，又恰好没有任何单元触发隐式实例化；
若写了 \Sw{extern} 模板声明抑制实例化（回指第 \ref{ch:cl03-resolution} 章
的声明与定义分界），而显式实例化定义所在的单元没进链接命令行，
未决引用就无处兑现。修法沿同一条链走：定义收拢到显式实例化定义那一个单元，
其余单元靠 \Sw{extern} 模板声明拿声明、模板本体保证可实例化，
或者升级到把成员声明为内联变量、让“定义可多份”成为语言承诺
（语言口径见第 \ref{ch:cl02-preprocessing} 章的翻译单元讨论收尾）。
\end{abibox}

\begin{bestpractice}[title={把“谁可见”写成构建层默认}]
库代码默认隐藏、按需导出：\Nix{} 与 \Mac{} 下用编译开关把默认可见性设为隐藏，
导出面用显式标记白名单化；\Win{} 下用导出清单或生成头划定边界。
每个对外库固定一张“导出符号清单”做回归比对——导出集合一夜翻倍，
多半是内部符号泄漏，下一版动态链接期的查找成本与抢占事故都要记在它头上。
折叠则反过来：默认不开，只有确认没有函数指针身份依赖的代码才开启。
\end{bestpractice}

\chapter{动态链接与加载执行}\label{ch:cl09-dynamic}

前面两章把链接归结为一次性的闭合操作：目标文件之间的引用在静态链接期全部兑现，
产物拿到的是一个已经确定的地址空间布局。动态链接打破了这个假设——它在镜像被放进
内存之后仍然保留一部分未完成的绑定工作，交给加载器在运行期执行。
本章讨论的核心只有一句话：\emph{编译期不知道最终地址}。
这句话在数据访问与函数调用两条路径上引出不同的补偿机制，也引出本章其余的全部话题：
惰性解析、平台加载器差异、线程局部存储、以及安全属性在链接产物里的落点。

本章假定读者已掌握第~\ref{ch:cl07-object}~章的段与重定位表结构，
以及第~\ref{ch:cl08-static-link}~章的符号解析与初始化顺序结论；
平台 ABI 的细节留给第~\ref{ch:cl10-cc}~章及其后的章节。

\section{位置无关的四个层次}\label{sec:cl09-pic-layers}

\subsection{从固定基址到共享对象}

按编译产物对最终加载地址的依赖程度，常见的链接产物可排成四个层次。

第一层是\textbf{需要固定基址的代码}。镜像内记录的地址被视为唯一的装载地址，
加载器只能把镜像映射到那个地址，否则拒绝执行。它实现简单，
但地址不可移动意味着随机化无从谈起，也决定了它无法与其他镜像共享物理页。

第二层是\textbf{位置无关目标文件}。编译器等产物时即假定最终地址未知，
生成的重定位项全部可以经由数据段中的指针表中转来完成，
文本段本身不含任何需要改写指令流的绝对地址。
这类目标文件既可参与共享对象的构建，也可参与位置无关可执行镜像的构建。

第三层是\textbf{位置无关可执行镜像}。它由位置无关目标文件链接而成，
整体可被映射到任意地址；由于它不依赖共享对象对全局变量的抢占式绑定，
其符号查找语义反而比普通的动态可执行镜像更接近静态链接的直觉。

第四层是\textbf{共享对象或 DLL}。它必须位置无关，因为同一份文件可能在不同的进程中、
同一进程的不同地址上被映射；它还必须容忍自己依赖的符号在别处、甚至比自己更晚确定。

\subsection{两条路径：数据与函数}

最终地址未知这件事，对数据访问和函数调用的影响方式并不对称。

访问\emph{本镜像内的数据}时，问题在于绝对地址无法写死在指令里。
通行的补偿方式是让访问代码基于寄存器中的镜像基址加上段内偏移来定位数据，
前提是每个函数序言都执行同样的基址建立步骤，这就是基址寄存器方案。

访问\emph{本镜像外的函数}时，问题在于连目标地址本身都可能要到运行期才确定。
补偿方式是把跳转目标放进数据段的一个表项：指令流只引用表项地址，
表项内容由加载器在解析符号后填写。两张表分工不同：一张面向全局变量的指针中转，
一张面向函数入口的跳转存根。两者都是数据，因此都能被改写；指令流则应尽量保持只读。

表~\ref{tab:cl09-trampoline} 按访问对象归类，标出每种补偿机制的角色、
地址最终在何时确定、以及能否推迟。

\begin{table}
  \centering
  \caption{位置无关产物中各类访问的补偿机制}
  \label{tab:cl09-trampoline}
  \begin{tabular}{>{\raggedright\arraybackslash}p{0.14\textwidth} >{\raggedright\arraybackslash}p{0.24\textwidth} >{\raggedright\arraybackslash}p{0.16\textwidth} >{\raggedright\arraybackslash}p{0.16\textwidth}}
    \toprule
    要访问什么 & 补偿机制的角色 & 何时确定 & 能否延迟 \\
    \midrule
    本镜像数据 & 基址寄存器加段内偏移，指令不含绝对地址 & 映射完成即确定 & 否，偏移编译期已知 \\
    外部数据 & 指针表存条目地址，改写表项而不改写指令 & 加载器解析符号时 & 是，可首次使用时再填 \\
    外部函数 & 跳转存根表加指针表，先跳表再跳转 & 加载器解析符号时 & 是，可首次调用时再填 \\
    本镜像函数 & 直接按镜像内偏移调用 & 链接期 & 无需延迟 \\
    线程局部变量 & 间接层数取决于访问模型，见第~\ref{sec:cl09-tls}~节 & 依模型而定 & 部分模型可以 \\
    \bottomrule
  \end{tabular}
\end{table}

不对称的根源在于：数据访问改写的是\emph{指针}，函数访问改写的是\emph{跳转目标}，
后者天然更适合延迟，因为一条调用指令在被执行之前无人关心它的目标。

\section{惰性解析与立即解析}\label{sec:cl09-lazy}

\subsection{把解析推迟到首次调用的收益}

延迟函数地址的解析，收益是启动期的。符号解析涉及跨镜像的名称查找，
名称查找的成本随依赖图规模与符号表长度增长，而进程生命周期内并非所有导出函数
都会被调用。把一个从未执行的调用点的解析彻底省掉，
对依赖面宽的镜像是可观的启动时间节省，同时也省掉解析期间的内存写操作。

代价是错误暴露的位置后移。符号缺失在立即解析模式下于加载时报错、进程尚未开始执行；
在惰性模式下则要跑到那一行调用才暴露，且往往表现为跳进存根后由加载器现场补查。

\begin{warning}[title={惰性解析给调试带来的困难}]
  缺符号的错误发生在第一条走到该调用点的线程上，而不是发生在启动时：
  崩溃位置离真正的成因很远，栈回溯里只能看到跳转存根与加载器的补查入口，
  看不到调用方原本想表达的业务语义。多线程下还可能只由某一条执行路径触发，
  测试环境难以复现。排查时先确认镜像是否允许延迟解析，再在受控环境里把它切成
  立即解析：若切换后启动即报错，成因即锁定。具体切换手段属平台与工具链能力，
  名称以该平台文档为准。
\end{warning}

\subsection{把重定位提前完成}

与惰性方向相反的做法，是让加载器在把控制权交给应用之前完成全部重定位，
即立即解析。它把符号缺失、依赖不齐一类问题全部压回启动边界，
代价是启动时间增加，收益是可调试性与可预测性。

在立即解析的基础上，进一步的安全取向是把重定位分两批处理：先完成的表项集中放在
只读、或只写后转只读的区域内，让承载函数指针的数据页在运行期不可写；
文本段则自始至终可执行且不可写。这类布局的意义在于缩小运行期可被改写的地址范围。
这里是能力描述而非操作说明：哪些重定位类型能进入只读段、由什么开关触发，
各平台与链接器实现差异较大，不列名称，以平台文档为准。

\begin{note}[title={只读重定位段的角色}]
  只读重定位段不改变解析语义，只改变解析完成后的内存属性。
  它的价值在安全属性（第~\ref{sec:cl09-security}~节）与异常展开表两处会反复出现：
  凡是运行期需要保持干净的数据结构，越早在加载期固化越好。
\end{note}

\section{\Nix：加载器与符号抢占}\label{sec:cl09-linux}

\subsection{加载器的职责与搜索顺序}

在 \Nix 类系统上，加载器负责按段表把各段以正确权限映射进地址空间、递归加载依赖、
完成重定位与符号解析，最后调用入口。可执行文件通过解释器路径字段指定由哪个加载器
完成这些工作；动态链接的可执行镜像没有加载器配合就无法启动。

共享对象的搜索顺序由两类信息共同决定：环境变量提供运行期的临时干预通道；
镜像内记录的路径提供随产物分发的、更稳定的信息。需要严格区分的是，
\textbf{记录在镜像里的路径有两种不同的语义}。

第一种是构建期路径：它是链接时用于定位依赖副本的路径记录，属于历史与校验信息。
运行期搜索可能仍会参考它，但把它当作部署承诺并不安全——构建目录不存在于目标机器上。

第二种是运行期相对路径：它以基于镜像自身位置的占位符表达依赖所在目录的相对关系，
让镜像可以随整个目录一起搬迁。这类记录才是部署时真正应该写入的路径。

两者存放在不同的字段里，工具在写入时也会做去重与取舍；混淆它们会造成
开发机正常、部署环境找不到库的经典故障。字段名与写入策略以目标平台的加载器文档为准。

\begin{compare}[title={环境变量与镜像内记录}]
  环境变量适合开发机上临时验证某一份库副本，它的作用域是用户会话，
  不随产物分发，也不进版本控制；镜像内记录的路径适合部署，
  它随文件走，可在打包流水线上校验。把部署所需的路径信息寄托在环境变量上，
  等于把产物正确性寄托在运维记忆上。
\end{compare}

\subsection{抢占式链接与可见性}

默认情形下，全局符号名在所有已载入镜像中做一次查找：名字若已在先载入的镜像中出现，
后续镜像对该名字的引用会绑定到先载入者的定义上。这个行为称为抢占式链接。
它保证同一进程内一个全局名字只有一个实例——这是 \cpp 全局对象、单例、
运行时状态不被复制一份的前提，也与第~\ref{ch:cl08-static-link}~章静态链接
一次性闭合的语义形成对照：静态链接下重复定义在链接期报错，
动态链接下重复定义在运行期被合并。

\begin{warning}[title={抢占是有代价的}]
  抢占意味着镜像内自己的函数定义可能在运行期被别的镜像的同名定义顶替，
  因此连对本镜像内部函数的调用也无法在编译期固化为直接调用，
  必须经表中转，这既增加一次间接，也让内联与局部优化失效。
  反过来，若一个镜像把本应内部使用的符号留在可抢占集合里，
  它就可能被应用层无意中的同名符号劫持——这类故障只表现为行为诡异，
  不表现为崩溃，最难排查。
\end{warning}

可见性属性用来关闭抢占：把符号标记为非抢占后，链接器就能按第~\ref{sec:cl09-pic-layers}~节
的机制直接生成表内固定条目甚至直接调用；标记为隐藏或更严格则同时限制其可导出性。
属性写法属于编译器扩展语法，各编译器支持度不一，这里只给概念不展开语法。
与之正交的另一手段是导出清单：不进接口集合的符号根本不进入动态符号表，
既不被抢占也不占用导出名。

\section{\Win：DLL 搜索、加载器锁与延迟加载}\label{sec:cl09-windows}

\subsection{可控与不可控的搜索次序}

\PeFmt 可执行镜像在运行期按固定次序搜索 DLL。次序中有一段不可配置：
应用目录、系统目录等位置的先后由操作系统决定，属于安全敏感的默认策略，
不该指望修改它来获得可用性。可配置的部分集中在进程级的搜索路径管理接口上：
可以显式关闭不需要的搜索位置，也可以追加若干目录。
把可控部分用足、把不可控部分当常量对待，是 Windows 部署的第一原则。

\subsection{依赖的递归解析与调试信息分离}

加载一个 DLL 时会递归加载它自己记录的依赖项，任何一层缺失都会在调用加载接口的
边界上统一表现为找不到模块，错误信息并不区分缺口发生在哪一层。
因此定位缺库必须先读依赖列表，而不是逐个尝试路径。

调试信息与镜像分离存放，是因为镜像需要保持精简、只含可执行所需的运行时信息，
而符号记录体积大、更新频繁、且不应进入分发件。这与第~\ref{ch:cl07-object}~章的结论
一脉相承：调试用的符号信息与运行期的符号、重定位信息服务于两个读者，
前者是调试器，后者是加载器。

\subsection{加载器锁}

进程内维护 DLL 状态的核心锁通常被称为加载器锁，系统在持有该锁的过程中执行每个
DLL 的入口点代码。由此产生两条限制。其一，入口点里不应执行可能反向依赖加载器状态的
操作，例如再动态加载别的 DLL 并等待其完成，否则有死锁风险。其二，入口点内执行的
用户代码越少，系统启动与模块加载的可预测性越高。
工程结论是：把需要复杂初始化的逻辑推迟到应用主线程显式调用的初始化函数里，
让 DLL 的入口点只做必要的、无回环的登记。

\subsection{延迟加载}\label{sec:cl09-delayload}

\begin{compare}[title={普通隐式加载与延迟加载}]
  普通隐式加载在进程启动时就让全部依赖就位，接口调用直达导入表，
  代价是启动开销与对部署完整性的强约束；
  延迟加载把某个 DLL 的首次解析推迟到程序真正调用它的第一个函数时，
  由链接器生成的辅助存根负责现场触发加载。
  它适合可选功能与重量级组件，不适合核心路径——它把找不到库的错误从启动时
  移到运行时，与第~\ref{sec:cl09-lazy}~节讨论的惰性解析在调试代价上同构。
\end{compare}

\section{\Mac：载入与地址修正两阶段}\label{sec:cl09-macos}

\Macho 镜像的启动在概念上分两个阶段：先\textbf{载入}，把各段按映射信息落成内存图像；
再\textbf{修正}，把镜像内需要按实际地址调整的指针与符号引用就地改写。
把修正作为独立步骤讲清，是因为它同时决定启动开销与安全属性。

对启动开销：修正是一次遍历修正信息表、逐条写内存的过程。工作集越大、修正条目越多，
加载阶段越长；共享缓存与固定映射一类手段的目的，正是把大量镜像的修正工作一次性摊销。

对安全属性：修正期间必然存在既承载指针又允许写的窗口，页面必须临时可写，
这恰好是攻击面。让修正阶段结构化、可校验（修正信息以紧凑表项而非自由散布的
重定位形式存在），才谈得上在修正完成后把相应页转为只读，
也才谈得上让随机化在不破坏语义的前提下选择基址。
换言之，修正独立成阶段，是把不可信窗口做成可控窗口的前提。

共享对象的符号解析次序与同名镜像的隔离策略，在不同系统与不同加载器版本上有差异，
写为以加载器文档为准。与 \Nix 一致的一条原则仍然成立：
部署路径信息应当随镜像分发，而不是散落在环境变量里。

\section{线程局部存储}\label{sec:cl09-tls}

\subsection{三种访问模型：差别是几层间接}

语言层提供一个每线程独享副本的变量存储期；实现层要在不知道变量最终地址、
也不知道访问发生在哪个线程的前提下定位它。常见的实现被冠以三种模型之名，
它们的差别可以归结为一个问题：访问一个线程局部变量需要几层间接。

\textbf{局部模型}：地址在单个镜像内即可完全确定，访问只需一次基址计算，
没有表项可被抢占，适用于不被共享的可执行镜像。

\textbf{初始执行模型}：地址由静态分配的线程块加固定偏移确定，一次表加一次偏移，
只对启动时即就位的镜像成立。

\textbf{通用执行模型}：地址要通过运行期句柄查找才能确定，间接层数最多，
但这是唯一能覆盖运行期动态载入模块的模型。

需要强调：上述名字是\emph{实现层的叫法}，源自平台处理器 ABI 文档，不是语言概念。
语言只承诺存储期与每线程副本，不承诺间接层数。同一平台的不同二进制类型要求不同模型，
编译单元所选模型若与最终产物类型不匹配，在共享对象构建中会直接报错而非降级运行。

\subsection{与静态链接期望的冲突}

\Tp{thread\_local} 的副本靠运行时为每个线程构造线程块来兑现，构造时机与镜像的加载
方式耦合。静态链接下线程块布局在链接期就是常量，运行期按固定尺寸分配即可；
动态链接下运行期载入的模块会带来新的线程局部变量，线程块尺寸必须能扩展。
某些平台对延迟载入模块可占用的线程块容量设有上限，超限的镜像在加载时报错；
具体上限与查询方式以平台文档为准，这里不写任何数字。
工程纪律与可见性一致：共享对象里的线程局部对象数量应当是受审查的接口面，
而不是随手使用的临时存储。

\begin{warning}[title={线程局部对象与加载期执行}]
  带非平凡构造的 \Tp{thread\_local} 对象会把代码推入首次访问路径，
  该路径可能在任意线程上触发，也可能与加载器状态纠缠。
  把它当普通全局对象的初始化来写，是跨线程偶发故障的常见来源。
\end{warning}

\section{安全属性在链接产物里的落点}\label{sec:cl09-security}

安全属性不是运行期的魔法，它从链接期就把决定写进了镜像。
可以按链接器在镜像里写了什么、加载器据此做什么来逐项核对。

地址随机化。链接器写的是镜像允许非固定基址这一事实：文件头中的基址语义，
加上段重定位信息确实存在。加载器据此做的是选择一个随机基址映射，
并按镜像内的修正信息把地址相关项调整到位。产物若退回固定基址，
链接器写不出可重定位性，加载器就无从随机。

数据执行保护。链接器写的是各段的可执行与可写属性，即段表里的权限位。
加载器据此按写权限映射页面，使数据页不可执行。链接期要警惕的是将代码与
承载函数指针的数据混进同一个既可写又可执行的段，那等于主动放弃这条属性。

控制流完整性一族。这一族的共同思路是让运行期的间接跳转只能落在被许可的地址集合内。
镜像侧的贡献是在链接期产出合法目标集合的记录与可供校验的表项。
不同平台采用的记录形式不同，机制描述到此为止，特性名称、开关名称与指令级细节
以平台文档为准；加载器或运行库据此在装载后校验，或在跳转存根处检查。

\begin{note}[title={异常展开表与回溯信息为何必须在镜像里可查}]
  异常传播与栈回溯要回答三个问题：当前指令地址落在哪个函数内、它的活动记录多大、
  callee-saved 寄存器被保存在何处。这些问答发生在运行期，而调试符号在发布产物里
  通常已被剥离，因此答案必须以压缩表格的形式留在镜像的一个专门段里供按需查询。
  静态链接可以把展开信息合并进最终镜像；动态链接则要求每个镜像各自携带、
  并能被加载器登记，否则展开穿越镜像边界时查不到表。
  展开规则的 ABI 约定见第~\ref{ch:cl13-cxx-abi}~章，
  调用约定对活动记录形状的影响见第~\ref{ch:cl10-cc}~章，
  段的容器格式见第~\ref{ch:cl07-object}~章。
\end{note}

\section{从镜像到 \Fn{main} 的启动序列}\label{sec:cl09-startup}

把前面各节串起来，一个动态链接镜像从被拉起至交出语言控制权经过下面七步。

\begin{flow}[title={执行镜像到 \Fn{main} 的启动序列}]
  \begin{enumerate}
    \item 映射各段：按段表把代码段与数据段以各自权限落地。
    \item 重定位与修正：改写指针表中转表与跳转存根表的条目。
    \item 线程局部存储准备：分配并清零线程块，登记延迟初始化队列。
    \item 依赖解析：递归载入依赖镜像，按可见性规则完成符号绑定。
    \item 动态初始化：依次执行各镜像的预初始化与初始化入口。
    \item 入口：跳转到语言启动例程，准备参数与环境。
    \item 交回语言运行时：完成语言层初始化后调用 \Fn{main}。
  \end{enumerate}
\end{flow}

三点说明。第一，映射与修正的先后是刚性的：指针表条目在被改写之前页面必须已可写，
而文本段自始不可写。第二，依赖解析与动态初始化互相穿插——被依赖镜像的初始化
先于依赖它的镜像。第三，交回语言运行时之后，才轮到第~\ref{ch:cl01-pipeline}~章所述的
语言层初始化，之后 \Fn{main} 才被调用。

\begin{warning}[title={静态初始化顺序只在单个镜像内有保证}]
  语言只规定单个翻译单元内、以及单个镜像内的静态对象初始化次序
  （单元内顺序由谁决定，第~\ref{ch:cl08-static-link}~章已讨论）。
  跨镜像的顺序——哪个 DLL 或共享对象的全局对象先构造——\textbf{不由语言规定}，
  它由依赖图拓扑与加载器实现决定，而依赖图会随链接顺序、
  是否启用延迟解析而改变。因此：镜像 A 的全局对象不得读写镜像 B 的全局对象状态；
  跨镜像的共享状态必须使用显式初始化函数，或首次访问时才构造的函数内静态量。
  这不是风格偏好，而是可移植性的边界条件。
\end{warning}

\section{跨平台共享库的构建配置}\label{sec:cl09-cmake}

下面这份配置把本章各议题落到构建系统上，只用确证的命令与变量签名，
注释说明每一段解决本章的哪个问题。

\begin{lstlisting}[style=cmakestyle,caption={共享库的最小跨平台目标配置（示意）},label={lst:cl09-cmake}]
# 第四层产物：SHARED 让库名与前后缀按平台落成共享对象形态
add_library(mylib SHARED src/mylib.cpp)

# 位置无关层次：共享对象必须位置无关，目标级属性覆盖全局默认
set_target_properties(mylib PROPERTIES
    POSITION_INDEPENDENT_CODE ON)

# 外部符号从哪来：链接期把依赖登记进镜像本身，
# 而不是等运行期靠环境变量去撞
target_link_libraries(mylib PRIVATE some::other)

# 抢占与可见性：只让走导出宏的符号进入导出集合
target_compile_definitions(mylib PRIVATE MYLIB_BUILDING)

# 头文件目录随目标传递，避免部署后手工拼装包含路径
target_include_directories(mylib PUBLIC
    ${CMAKE_CURRENT_SOURCE_DIR}/include)

# Windows 默认不导出任何符号；打开总导出会削弱显式导出纪律
set(CMAKE_WINDOWS_EXPORT_ALL_SYMBOLS OFF)
\end{lstlisting}

两个变量的作用域要分清。\Sw{CMAKE\_POSITION\_INDEPENDENT\_CODE} 设定的是各目标
位置无关属性的全局默认值，目标级属性优先于它；
\Sw{CMAKE\_WINDOWS\_EXPORT\_ALL\_SYMBOLS} 只影响 \Win 平台的导出行为，
相当于把本章第一条纪律交给工具链代劳，原型阶段之外都不该打开。
两者都不改变本章前面讨论的任何运行期语义。

\subsection{运行期定位缺库的三种原因}\label{sec:cl09-missinglib}

三种失败原因在现象上完全相同，加载器只报\textbf{找不到库或找不到入口点}；
在概念上它们分属本章与第~\ref{ch:cl07-object}~章的不同视图。

其一是搜索路径不含它，属于搜索顺序议题：镜像的依赖记录里名字是对的，
只是当前次序命不中那份副本。其二是它自己的依赖缺失，属于递归解析议题：
报错的是最上层镜像，缺口在下一层，视图落在第~\ref{ch:cl07-object}~章的依赖字段上。
其三是导出集合里根本没有那个符号，属于可见性与导出议题：库本体加载成功，
是符号绑定失败，视图同样落在第~\ref{ch:cl07-object}~章的动态符号表上。

\begin{abibox}[title={运行期找不到库：一条可执行的排查路径}]
  按下列顺序收集事实，每一步都在排除上面三类原因之一。

  \begin{enumerate}
    \item 离线读取待加载镜像记录的依赖名称列表。若列表里没有报错的那个名字，
      说明失败的镜像不是它，而是更深层的依赖，回到本步骤处理那个更深的镜像。
      这一步把原因二与原因一分开。
    \item 若列表里有该名字，确认磁盘上被找到的那份文件确实是预期的构建产物：
      它的导出集合里应包含报错中缺失的符号名。不含则属原因三，
      处理办法回到编译配置（导出宏或可见性属性），而不是改环境变量。
    \item 前两步都通过，才检查原因一：当前环境的搜索路径是否包含该文件所在目录，
      以及镜像内记录的那条路径属于构建期路径还是运行期相对路径
      （第~\ref{sec:cl09-linux}~节的两种语义）。
      部署环境的修复应落在镜像内记录的路径上，临时验证才允许用环境变量。
  \end{enumerate}

  这个顺序的理由与第~\ref{sec:cl09-delayload}~节相同：最上层的报错不区分成因层次，
  而依赖列表与符号表是可以离线读到的客观事实，搜索路径却是环境属性。
  读取这两个字段的工具名称各平台不同（\Nix 类系统看动态段与符号表，
  \Win 用镜像转储工具，\Mac 用面向 \Macho 的检查工具），此处不列命令，以各自手册为准。
\end{abibox}

\begin{guideline}[title={跨平台库的三条链接期纪律}]
  \textbf{一、显式导出。}导出集合必须是一份审查过的接口清单，用导出宏或可见性属性
  声明，而不是把全部符号开放出去。它同时关闭了抢占带来的行为诡异，
  也让上面排查路径里的原因三变成可通过配置静态检查的错误。

  \textbf{二、不依赖搜索顺序。}任何两个目录不得存放同一依赖名的不同副本；
  版本差异用文件名表达而不是用目录先后表达。搜索顺序是环境属性，
  随部署机器变化，不可版本控制。

  \textbf{三、把运行期路径写进镜像。}若安装布局允许整体搬迁，
  就把基于镜像自身位置的路径计算写进镜像记录；环境变量只作为开发机上的
  临时手段存在，并明确记录在流水线配置里，以免它成为无人认领的隐式依赖。
\end{guideline}

\begin{note}[title={本章小结}]
  本章把一个成因贯穿到底：编译期不知道最终地址。数据经指针表中转、
  函数经跳转存根表中转；推迟解析换来启动速度，也换来错误在后延的位置暴露；
  抢占与可见性决定谁绑定谁；三个平台的差异集中在搜索次序、修正阶段、
  线程块容量三处；安全属性是链接器写进镜像的事实与加载器据此动作的配对；
  启动序列的最后一段之前，跨镜像的初始化顺序不由语言负责。
  下一章把视角从系统边界移向语言边界：调用约定如何规定参数的去留与名字的变形。
\end{note}

\part{平台 ABI}
\chapter{调用约定：系统 V、MS \xSixtyFour、\ArmSixtyFour 与 \Lang{RISC-V}}\label{ch:cl10-cc}

\begin{abibox}[title={调用约定定义}]
调用约定（calling convention）是处理器接口（ABI）的核心子集，它规定：参数
放在哪些寄存器或栈槽；返回值放在哪里；哪些寄存器由被调用方保存；栈如何
对齐与清理；以及名字如何修饰（修饰另见第~\ref{ch:cl13-cxx-abi}章）。
同一个 C 或 \cpp 签名在不同约定下的机器级含义完全不同——这正是跨平台
链接失败与静默数据损坏的根源。
\end{abibox}

\section{可信边界声明}\label{sec:cl10-trust}

本章表格与正文只列出经过确证的项。凡未确证或存在版本差异的格子，统一写
``以对应处理器接口文档为准''。读者在将任何结论用于生产环境之前，应核对
对应平台的官方 ABI 规范。汇编与输出片段均标注``示意''，不声称是实测结果。

\section{跨平台对照主表}\label{sec:cl10-main-table}

\Nix{} \xSixtyFour 系统 V、\Win{} \xSixtyFour、\ArmSixtyFour 一族（\Nix、\Mac、\Win
共用同一套整型与浮点参数寄存器）以及 \Mac{} \xSixtyFour（沿用系统 V 的参数传递）
和 \Lang{RISC-V} 的寄存器组对照如下。

\begin{longtable}{>{\raggedright\arraybackslash}p{0.14\textwidth}>{\raggedright\arraybackslash}p{0.10\textwidth}>{\raggedright\arraybackslash}p{0.09\textwidth}>{\raggedright\arraybackslash}p{0.08\textwidth}>{\raggedright\arraybackslash}p{0.08\textwidth}>{\raggedright\arraybackslash}p{0.09\textwidth}>{\raggedright\arraybackslash}p{0.07\textwidth}>{\raggedright\arraybackslash}p{0.07\textwidth}}
\caption{跨平台调用约定对照}\label{tab:cl10-abibig}\\
\toprule
约定 & 整数指针参数 & 浮点参数 & 整型返回 & 浮点返回 & 被调方保存 & 栈对齐 & 特殊区 \\
\midrule
\endfirsthead
\toprule
约定 & 整数指针参数 & 浮点参数 & 整型返回 & 浮点返回 & 被调方保存 & 栈对齐 & 特殊区 \\
\midrule
\endhead
\bottomrule
\endfoot
\Nix{} \xSixtyFour 系统 V
& \Tp{rdi} \Tp{rsi} \Tp{rdx} \Tp{rcx} \Tp{r8} \Tp{r9}
& \Tp{xmm0} 至 \Tp{xmm7}
& \Tp{rax} 附带 \Tp{rdx}
& \Tp{xmm0} 至 \Tp{xmm1}
& \Tp{rbx} \Tp{rbp} \Tp{r12} 至 \Tp{r15}
& 16 字节
& 128 字节红色区域 \\[4pt]
\Win{} \xSixtyFour
& \Tp{rcx} \Tp{rdx} \Tp{r8} \Tp{r9}
& 同整数前四个
& \Tp{rax}
& \Tp{xmm0}
& 以文档为准
& 16 字节
& 32 字节影子空间 \\[4pt]
\ArmSixtyFour 一族
& \Tp{x0} 至 \Tp{x7}
& \Tp{v0} 至 \Tp{v7}
& \Tp{x0} 必要时 \Tp{x1}
& \Tp{v0} 必要时 \Tp{v1}
& \Tp{x19} 至 \Tp{x28} \Tp{fp} \Tp{lr}
& 16 字节
& 无红色区域 \\[4pt]
\Mac{} \xSixtyFour
& 同系统 V
& \Tp{xmm0} 至 \Tp{xmm7}
& \Tp{rax}
& \Tp{xmm0}
& 以文档为准
& 16 字节
& 以文档为准 \\[4pt]
\Lang{RISC-V} 组
& 以文档为准
& 以文档为准
& 以文档为准
& 以文档为准
& 以文档为准
& 以文档为准
& 以文档为准 \\
\end{longtable}

表中未展开的 \Lang{RISC-V} 细节与 \Win{} \xSixtyFour 被调用方保存寄存器完整列表，
一律以对应处理器接口文档为准。表中 ``以文档为准'' 是 ``以对应处理器接口
文档为准'' 的缩写。

\subsection{寄存器分配的直觉}\label{sec:cl10-reg-intuition}

系统 V 用六个整数/指针寄存器加八个浮点寄存器，给大多数函数提供了充足的
寄存器传参空间。MS\xSixtyFour 只给四个参数位置（\Tp{rcx}、\Tp{rdx}、\Tp{r8}、
\Tp{r9}），且整型与浮型共用同一组——第五个参数及以后的所有参数一律走栈。
\ArmSixtyFour 把八个整数/指针寄存器 \Tp{x0} 至 \Tp{x7} 和八个浮点寄存器
\Tp{v0} 至 \Tp{v7} 分开，同时被调用方保存寄存器占 \Tp{x19} 至 \Tp{x28} 十个
加上帧指针与链接寄存器，留给调用方自由使用的寄存器比例较小。

\begin{flow}[title={参数去向}]
签名确定后查约定，按类别分配寄存器，
寄存器用完则推栈，栈对齐 16 字节，
系统 V 调用方清栈，MS 不要求被调方清。
\end{flow}

\section{三个特殊区域}\label{sec:cl10-special-zones}

\subsection{红色区域（Red Zone）}\label{sec:cl10-redzone}

\begin{abibox}[title={红色区域}]
系统 V 规定被调用方可在 \Tp{rsp} 以下 128 字节范围内暂存数据，无需先将
\Tp{rsp} 下移。此区域只有叶子函数（不再调用其他函数）可安全使用。
原因：任何后续 \Fn{call} 会把返回地址压到 \Tp{rsp} 处，信号处理程序
也可能向 \Tp{rsp} 以下写入上下文，二者都会覆盖红色区域。
\end{abibox}

\Win{} \xSixtyFour 没有红色区域。信号或中断到达时栈帧保护机制以对应处理器
接口文档为准。

\subsection{影子空间（Shadow Space）}\label{sec:cl10-shadow}

\begin{abibox}[title={影子空间}]
\Win{} \xSixtyFour 约定要求调用方在每个 \Fn{call} 指令之前，在栈上预留
32 字节（四个 8 字节槽位），用于被调用方将 \Tp{rcx}、\Tp{rdx}、\Tp{r8}、
\Tp{r9} 中保存的参数溢出到已知位置。被调用方可自由读写这 32 字节而不
修改 \Tp{rsp}。影响：每个函数调用至少需要 32 字节额外栈空间，即使参数
全部走寄存器；嵌套调用时栈帧膨胀成线性比例。
\end{abibox}

\subsection{间接返回（Hidden Pointer）}\label{sec:cl10-indirect}

当返回类型超过两个通用寄存器宽度时（例如 \cpp 中超过 16 字节的结构体），
调用方在栈上分配一块返回缓冲区，将其地址作为隐藏的第一参数传入：
\begin{itemize}
\item 系统 V：隐藏指针占 \Tp{rdi}（第一个整数参数位），
      其余参数依次后移一个寄存器。
\item MS\xSixtyFour：隐藏指针占 \Tp{rcx}，其余参数后移。
\item \ArmSixtyFour：隐藏指针占 \Tp{x0}，其余参数后移。
\end{itemize}
被调用方写入缓冲区后将同一地址复制到 \Tp{rax}（或 \Tp{x0}）作为返回值。
此机制与第~\ref{ch:cl11-layout}章中结构体大小决定能否寄存器返回直接相关。

\section{结构体与类的传参}\label{sec:cl10-struct-pass}

\subsection{按字段类别拆分（系统 V）}\label{sec:cl10-sysv-split}

系统 V 对结构体传参执行分类算法：每个字段被归类为整数或 SSE（浮点），
然后分别计入整型参数序列和浮点参数序列。若两类均不超过可用寄存器数且
结构体不超过 16 字节，可用一对寄存器传递（如 \Tp{rdi} 与 \Tp{xmm0}）；
任何一类超出可用数量则整体走内存。

\subsection{MS 与 \ArmSixtyFour 的做法}\label{sec:cl10-ms-arm-struct}

MS\xSixtyFour 对结构体传参采用简化策略：大小 1/2/4/8 字节的结构体放入
对应的通用寄存器（按整型处理）；其余一律在调用方栈上复制一份，传递副本
地址。不拆分到浮点寄存器。

\ArmSixtyFour 的规则介于两者之间：允许整数与浮点分别占不同寄存器组，
具体分类细节以对应处理器接口文档为准。

\subsection{\cpp 的 \Fn{this} 指针}\label{sec:cl10-this}

在三种约定下，\cpp 非静态成员函数都把 \Fn{this} 作为第一个不可见参数传入，
占据第一个整数/指针寄存器（系统 V 的 \Tp{rdi}、MS 的 \Tp{rcx}、
\ArmSixtyFour 的 \Tp{x0}）。

\begin{warning}[title={传值布局依赖}]
同一个 C 签名 ``\Tp{f(struct S s)}''，在系统 V 下按类别拆分至整型
与浮点寄存器组，在 MS\xSixtyFour 下若大于 8 字节则走内存。若头文件只
声明了结构体大小而未暴露字段排列，两种编译器的内存布局要求就不同（回指
第~\ref{ch:cl11-layout}章）。跨平台共享的头文件必须保证字段顺序与对齐
在两端一致。
\end{warning}

\section{可变参数}\label{sec:cl10-variadic}

\subsection{约定如何决定变参形状}\label{sec:cl10-vararg-shape}

系统 V 的 \Fn{va\_list} 内部记录了已使用的 \Tp{rdi} 至 \Tp{r9} 中哪些寄存器
被用到，以及后续参数从栈的哪个偏移开始读取——因为 \Tp{al} 被规定为
传参时使用的 \Tp{xmm} 寄存器数（用于被调方保存浮点寄存器到溢出区）。
MS\xSixtyFour 的 \Fn{va\_list} 形状不同，它记录影子空间与栈参数的混合
位置。

\subsection{跨平台转发为何只能由编译器生成}\label{sec:cl10-forward}

可变参数的物理布局完全由调用约定决定，不同平台的 \Fn{va\_list} 内部
表示互不兼容。手写转发代码必须在每种约定下知道参数是散落在寄存器还是
栈上。只有编译器能根据目标约定生成正确的保存与恢复逻辑。

\subsection{\Tp{va\_list} 的角色}\label{sec:cl10-va-list-role}

\Tp{va\_list} 是标准库提供的不透明类型，本质上是当前参数游标的
平台编码。它属于约定层的产物而非语言语义——不同编译器对同一 \Tp{va\_list}
的内部字段排列可能不同，因此绝不能将 \Tp{va\_list} 跨 DLL/DSO 传递后
在另一端继续消费。

\section{语言层入口：属性与关键字}\label{sec:cl10-lang-entry}

\subsection{属性语法与类型修饰符}\label{sec:cl10-attr-vs-modifier}

\Win 编译器族用类型修饰符声明约定：
\begin{itemize}
\item \Tp{\_\_cdecl}（默认）：调用方清理栈。
\item \Tp{\_\_stdcall}：被调用方清理栈——主要活在 32 位 \Win 生态。
\item \Tp{\_\_fastcall}：历史上把前两个参数放寄存器，
      现已被其他修饰符取代。
\end{itemize}
\Lang{GCC} 与 \Lang{Clang} 一族用属性语法，声明在函数或函数指针类型上。
例如 \Fn{sysv\_cc}（系统 V）与 \Fn{ms\_abi}（MS 约定）作为
\Lang{GCC} 的 target 属性参数出现。

\subsection{外部语言链接}\label{sec:cl10-extern-lang}

C 语言链接声明（\Tp{extern}~``C''）影响的是修饰名（name mangling），使
链接器能找到纯 C 名。它并不改变调用约定本身——在 \Win{} \xSixtyFour 上
C 链接的函数仍然按 MS\xSixtyFour 约定调用。名字修饰详见
第~\ref{ch:cl13-cxx-abi}章。

\subsection{声明划界}\label{sec:cl10-scope}

本章讲调用约定的机器级机制。互操作接口怎么写（C 链接封装、COM 接口、
FFI 框架）属于别的分册。

\section{三个可观测后果}\label{sec:cl10-consequences}

\subsection{栈回溯与异常展开}\label{sec:cl10-unwind}

\begin{abibox}[title={回溯与展开}]
\Win{} \xSixtyFour 的异常展开使用 unwind info 表（编码在 \PeFmt 中）；
\Nix{} \xSixtyFour 的 \Elf 使用 CFI（.eh\_frame 节）。两者都记录每个函数
对哪些寄存器做了保存、保存位置相对 \Tp{rsp}（或 \ArmSixtyFour 的 \Tp{sp}）
的偏移。若手写片段保存寄存器却未生成对应展开信息，异常处理器无法恢复
调用者栈——症状出现在无关的 try/catch 块（回指第~\ref{ch:cl13-cxx-abi}章
与第~\ref{ch:cl09-dynamic}章）。
\end{abibox}

\subsection{破坏被调用方保存寄存器}\label{sec:cl10-clobber}

如果手写汇编或外部二进制（如第三方库中的优化子程序）修改了 \Tp{rbx}、
\Tp{r12} 至 \Tp{r15} 而未恢复，损坏不会立即崩溃——函数返回值可能仍然正确。
症状出现在调用者后续使用这些寄存器的代码处（回指第~\ref{ch:cl06-codegen}
章中寄存器分配）。排查极其困难，因为崩溃位置与破坏位置之间可能有数十
层调用。

\subsection{跨镜像回调}\label{sec:cl10-callback}

DLL/DSO/动态库传递的函数指针若约定不一致，参数寄存器会被解释为栈参数，
返回值可能被当作隐藏指针地址。后果是内存损坏或段错误。跨镜像回调必须
把约定冻结在接口里：使用显式约定属性声明函数指针类型（回指
第~\ref{ch:cl08-static-link}章与第~\ref{ch:cl15-endtoend}章）。

\section{对照示例：函数序言}\label{sec:cl10-prologue}

以下两段汇编均为示意，展示同一签名 ``\Tp{int sample(int a, double b)}''
在不同约定下的序言。

\subsection{系统 V 序言（示意）}\label{sec:cl10-prologue-sysv}

\begin{lstlisting}[style=asmstyle, caption={系统 V 序言（示意）}]
; int sample(int a, double b)
; a in edi, b in xmm0
sample:
    push    rbp
    mov     rbp, rsp
    sub     rsp, 16
    mov     [rbp-8], rdi
    movsd   [rbp-16], xmm0
    ; ... body ...
    mov     rsp, rbp
    pop     rbp
    ret
\end{lstlisting}

\subsection{MS 序言（示意）}\label{sec:cl10-prologue-ms}

\begin{lstlisting}[style=asmstyle, caption={MS 序言（示意）}]
; int sample(int a, double b)
; a in ecx, b in xmm1
sample:
    push    rbp
    mov     rbp, rsp
    sub     rsp, 48
    mov     [rbp-8], ecx
    movsd   [rbp-16], xmm1
    ; ... body ...
    mov     rsp, rbp
    pop     rbp
    ret
\end{lstlisting}

\begin{abibox}[title={序言差异}]
系统 V：被调用方可用红色区域避免修改 \Tp{rsp}（仅叶子函数）。
MS\xSixtyFour：调用方必须为每个调用预留 32 字节影子空间，被调用方可在
其中存取参数。栈帧大小因此至少多 32 字节。两种约定的对齐要求都是
16 字节。
\end{abibox}

\section{回调接口的冻结写法}\label{sec:cl10-callback-freeze}

\begin{lstlisting}[style=cstyle, caption={回调接口冻结模板（示意）}]
#include <stdint.h>
#include <stddef.h>

typedef int32_t (*cb_t)(uint32_t ver,
                        uint64_t id,
                        void *ud);

typedef struct {
    uint32_t api_ver;
    uint32_t struct_sz;
    cb_t     on_event;
    void    *user_data;
} frozen_iface_t;

_Static_assert(sizeof(frozen_iface_t) == 32,
    "size frozen");
_Static_assert(offsetof(frozen_iface_t,
    on_event) == 8, "off frozen");
_Static_assert(offsetof(frozen_iface_t,
    user_data) == 16, "off frozen");

int32_t register_handler(
    const frozen_iface_t *iface);
\end{lstlisting}

\begin{bestpractice}[title={跨镜像纪律}]
\begin{enumerate}
\item 只用确证宽度的 C 类型（\Tp{uint32\_t}、\Tp{int64\_t} 等
      \Tp{stdint.h} 定宽类型），不用 \Tp{int}、\Tp{long}、\Tp{size\_t}。
\item 不跨越传 \cpp 对象：构造函数、虚表、异常类型在不同编译器版本下
      布局不同。只传 \Tp{void*} 加显式 \Fn{create} 与 \Fn{destroy} 函数。
\item 显式传用户数据指针：不依赖 TLS 或全局变量做回调上下文传递。
\item 接口带版本号字段与结构体大小字段：接收方可检测 ABI 不匹配并
      返回明确错误码，而非静默读越界。
\end{enumerate}
\end{bestpractice}

\begin{warning}[title={本机测试不足}]
在一台机器、一个编译器版本上测试通过不能证明接口 ABI 正确。常见陷阱：
\begin{itemize}
\item 目标平台 \Tp{long} 宽度为 4（\Win{} \xSixtyFour 的 LLP64 模型），
      而开发机为 8（\Nix 的 LP64 模型），结构体偏移全变。
\item 调试优化关闭时寄存器保存规则被偶然遵守，开启 \Fn{-O2} 后
      分配器重新使用被破坏的寄存器。
\item 链接顺序或符号解析改变时，同一调用可能绑定到不同约定的实现。
\end{itemize}
``能跑''只是充分证据的缺失，而非正确性的证明。
\end{warning}

\begin{note}[title={三章关系}]
本章讲调用约定——参数与返回值在寄存器/栈之间如何移动。
第~\ref{ch:cl11-layout}章讲数据布局——类型在内存中如何排列。
第~\ref{ch:cl12-endian}章讲字节序——多字节数如何填入连续地址。
三章合起来才能回答：同一份源码在不同平台上到底是什么。任何一项
不对齐都会导致链接成功但运行时静默损坏。
\end{note}

\section{小结}\label{sec:cl10-summary}

\begin{abibox}[title={速查}]
\xSixtyFour 系统 V：六整数寄存器 + 八浮点寄存器，128 字节红色区域，
被调方保存 \Tp{rbx}、\Tp{rbp}、\Tp{r12} 至 \Tp{r15}。
\Win{} \xSixtyFour：四参数寄存器，32 字节影子空间，第五参数起走栈。
\ArmSixtyFour：八整数 \Tp{x0} 至 \Tp{x7} + 八浮点 \Tp{v0} 至 \Tp{v7}，
被调方保存 \Tp{x19} 至 \Tp{x28} 与 \Tp{fp}、\Tp{lr}，无红色区域。
三种约定栈对齐均为 16 字节。间接返回使用隐藏指针占第一参数位。
\end{abibox}

跨平台开发的 ABI 正确性来自对约定的精确遵守，而非对``通常情况''的依赖。
任何未列入本章确证项的格子，请以对应处理器接口文档为准。

\chapter{数据布局：对齐、打包与位域}\label{ch:cl11-layout}

同一组成员，换个声明顺序，结构体就换了大小；同一个位段，换个编译器，就换了位置。
布局是源码与机器之间最容易“看起来一样、实际上不同”的一层：
它既由语言标准写下几条硬规则，又留了成片的实现定义地带，
还直接决定 \Fn{memcpy} 的字节流、共享内存的两侧视图与原子操作能否落地。
本章把这三股力量分开：哪些是标准承诺的，哪些是实现自由，哪些是你可以主动改写的。
三条主线——成员重排、打包指示、过对齐——分别对应第 \ref{sec:cl11-layout-algo}、
第~\ref{sec:cl11-packing} 节、\ref{sec:cl11-align-concurrency} 节；
位域作为“实现定义地带”的典型样本单独成节（第 \ref{sec:cl11-bitfields} 节）。
字节序与数据模型是布局的另一半，放在第 \ref{ch:cl12-endian} 章；
布局假设如何跨编译单元冻结，回指第 \ref{ch:cl10-cc} 章的接口纪律。

\section{布局算法：编译器如何摆放成员}\label{sec:cl11-layout-algo}

标准对记录布局写死了四条规则，它们是本章一切讨论的地基。
第一，非静态数据成员按其声明顺序获得递增的内存地址；
第二，每个成员被放到满足其对齐要求的地址上，不满足就插填充字节；
第三，整个结构体的对齐要求等于所有成员对齐要求的最大值；
第四，结构体的总大小必须向上取整到自身对齐的整数倍，
这样当它作为数组元素连续排布时，第二个元素的第一个成员才依然对齐。
尾部填充正是第四条规则的产物，它不服务于任何成员，只服务于“数组”这个用法。

这四条承诺之外，就是实现的领地：填充到底插几字节、
匿名联合与位段如何借用填充空间、空基类占不占位置，都随实现与目标平台走。
因此下面所有讨论只比较相对关系，不写绝对的字节数；
需要绝对值时，用第 \ref{sec:cl11-cmake-probe} 节的方法在配置期探测。

\begin{lstlisting}[style=cppstyle, caption={同一组成员、两种声明顺序（对照示意；只断言相对关系）}]
struct Header {
    char tag;    // 先放对齐要求小的
    int  kind;   // 对齐要求大的成员夹在中间
    char mark;   // 后放对齐要求小的
};

struct HeaderSorted {
    int  kind;   // 对齐要求最大的成员排第一
    char tag;
    char mark;
};

// 四个成员集合相同,所以结构体对齐相同:整体对齐取成员对齐的最大值,
// 与声明顺序无关。这条断言只依赖本章开头列出的规则,可以写进代码。
static_assert(alignof(Header) == alignof(HeaderSorted),
              "structure alignment is the max member alignment");

// sizeof 的绝对值取决于各类型的对齐要求,属于实现领地,不写死数字。
// 一般关系:Header 中 kind 之前要为它垫出对齐空间,两个 char 分处两端,
// 各自制造填充;HeaderSorted 把 kind 顶到地址 0,两个 char 可以贴着
// 它连续排布,尾部再统一补齐。成员越多、对齐跨度越大,这种重排的
// 收益越明显——语义完全等价,布局却可能更小。
\end{lstlisting}

重排成员为什么能改变大小而不改变语义？
因为“声明顺序递增”只约束相对次序，不约束绝对位置：
把对齐要求最大的成员提到最前面，它天然落在满足自己的地址上，
其后的小对齐成员可以见缝插针地共享它留下的尾部空间；
而把大对齐成员夹在两个小成员中间，就会在它前后各逼出一段填充。
两种写法生成的对象在语言层完全等价——成员名、类型、初值语义都不变，
变的只有各成员的偏移与总大小。这正是布局的隐蔽之处：
语义等价的两个版本，字节视图并不等价，跨镜像传对象时就会露馅。

\begin{guideline}[title={布局算法的四条硬规则}]
按声明顺序放置、按成员对齐垫到位、整体对齐取最大、总大小补到对齐的整数倍。
四条都是标准的承诺；填充的具体数量与位域的落点不是。
任何写进代码的布局假设，只允许引用这四条规则能推出的相对关系。
\end{guideline}

\section{语言层的布局控制}\label{sec:cl11-lang-control}

\subsection{alignas、alignof 与 std::max\_align\_t}\label{sec:cl11-alignas}

\cpp{11} 起，语言给了显式干预对齐的一对工具：
说明符 \Tp{alignas} 提高（只提高、不降低）某个类型或变量的对齐要求，
运算符 \Fn{alignof} 查询一个类型的对齐要求。
两者都是标准设施，可以安全地出现在 \Fn{static\_assert} 里。
基线类型能对齐到多远，由 \Tp{std::max\_align\_t}（\Tp{<cstddef>}）划定：
它是所有“标量类型”都能对齐到的最大值，
要求比它更宽的对齐属于扩展能力，实现可以拒绝（\Tp{alignas} 超限是病态程序）。
平台 psABI 通常还会写明各类型的自然对齐，具体数值以对应平台 psABI 为准。

\subsection{空基类优化、访问控制与 offsetof}\label{sec:cl11-stdlayout}

第二条暗线是“哪些类允许你问偏移”。\Fn{offsetof} 只对 standard-layout 类
有标准定义；标准布局大致要求：非静态数据成员具有相同的访问控制、
基类与数据成员不交错、没有虚函数与虚基类。
这个前提不是官僚主义：非标准布局的类里，不同访问控制段的相对顺序是实现定义的，
编译器甚至可以重排这些段——此时“成员偏移”本身就不是一个确定答案，
\Fn{offsetof} 自然无从保证。
空基类优化（EBO）则展示另一面：一个不含数据的基类可以不占存储空间，
让派生类看起来像“没有继承过一样紧凑”。
但做不做 EBO 是实现与 ABI 的选择，不属于标准的强制承诺；
\cpp{17} 起的 \Tp{[[no\_unique\_address]]} 属性
把“这个成员可以压扁”的决定权交给了写代码的人，同样不保证生效顺序以外的布局。
一句话：能问偏移的类是被标准划了圈的，圈外的类连布局都是实现自由。

\begin{lstlisting}[style=cppstyle, caption={过对齐与缓存行粒度的标准写法（C++17 设施）}]
#include <new>
#include <atomic>

struct Counter {
    std::atomic<unsigned> value;
};

// 把对齐抬到“破坏性干扰粒度”:相邻对象大概率落在不同缓存行,
// 降低伪共享。该常量由 <new> 提供,数值属于实现选择,
// 只能当提示用,不能当硬件事实写死。
struct alignas(std::hardware_destructive_interference_size)
    PaddedCounter {
        Counter slot;
};

static_assert(alignof(PaddedCounter) >= alignof(Counter),
              "alignas only ever raises alignment");
\end{lstlisting}

\begin{note}[title={标准常量的可用性}]
\Tp{std::hardware{\Bk}\_destructive{\Bk}\_interference{\Bk}\_size} 与
\Tp{std::hardware{\Bk}\_constructive{\Bk}\_interference{\Bk}\_size} 由 \cpp{17}
的 \Tp{<new>} 提供，但数值是实现定义的，个别标准库实现或构建配置可能不提供。
把它们当“库给的提示值”，两侧都要用同一实现编译；跨工具链共享内存时，
宁可像样例那样由配置期探测统一注入（第 \ref{sec:cl11-cmake-probe} 节）。
\end{note}

\section{打包：把对齐要求压下去}\label{sec:cl11-packing}

网络协议头、磁盘格式、硬件寄存器这些“外部合同”常要求成员紧挨着排，
不允许填充存在。语言标准没有“压缩对齐”这个设施，于是出现了两族非标准写法。
属性式：在类型声明上挂 \Tp{\_\_attribute\_\_((packed))}，
GCC 与 \Tool{clang} 一族支持这一形态，但对各属性的支持程度随版本演进，
移植时以对应编译器文档为准；\Win{} 的 MSVC 工具链则以指示式为主。
编译指示式：\Tp{\_\_pragma(pack(N)} 与 \Tp{pack(push, N)}、
\Tp{pack(pop)}（MSVC 与 GCC 一族都识别 \Tp{\_\_pragma(pack(...))}），
从指示出现处起改变其后成员的默认对齐上限；
再往上一档是命令行全局打包，如 MSVC 的 \Sw{/Zp2} 一类开关，
把整个编译单元所有结构体的默认对齐上限统一压低。

\begin{compare}[title={属性式与指示式：两种打包写法}]
属性式（\Tp{\_\_attribute\_\_((packed))}）逐类型生效，意图写在类型旁边，
漏标一个类型不会被误伤；指示式（\Tp{\_\_pragma(pack(...))}、\Sw{/Zp}）
按范围生效，一行覆盖一批类型，方便顺手，也最容易忘关：
\Tp{pack(push} 不配对 \Tp{pack(pop)}，打包状态就泄进了后面的头文件。
两族都不是标准设施，跨厂商移植时要逐目标确认支持形态。
\end{compare}

\begin{warning}[title={取 packed 成员地址在语言层就有问题}]
打包压低了类型对齐，但成员自身的类型仍声称较高的对齐要求。
一旦你取它的地址——绑一个引用、传给一个按对齐优化过的函数——
就得到了一只“声称 4 字节对齐、实际可能在奇数地址”的指针，
后续解引用踩进未对齐访问：有的平台硬件陷入，有的平台性能断崖，
而按标准口径这已落入未定义行为的地带（对齐承诺与后果的完整链条
见未定义行为分册的对齐主题章）。编译器通常直接报错拦截取地址，
但那只是实现加的前门栓，不是语言承诺。
纪律：packed 结构体只当“字节容器”用，进出字段一律走逐字节读写或
\Fn{memcpy}，不要让它的成员地址逃逸进普通指针。
\end{warning}

全局打包档位比逐类型打包更危险，危险的根源不在“变小”，而在“变 ABI”：
它改变该编译单元内所有结构体的对齐与大小，而结构体布局正是跨编译单元
接口的一部分。一个单元带 \Sw{/Zp}、另一个单元不带，两边对同一个头文件
算出的 \Fn{sizeof} 与 \Fn{offsetof} 就不同，函数签名里的“传结构体”会
按两套不同的字节数搬运数据——链接器不会报错，它会安静地成功。
这正是第 \ref{ch:cl10-cc} 章要求接口两侧共享同一组编译开关的原因之一。

\section{位域：四件事几乎都是实现定义}\label{sec:cl11-bitfields}

位域是 C 与 C++ 里承诺最少、误会最多的一块。拿不准的口径先摆出来：
第一，底层类型——除整数类型与 \Tp{bool}（\cpp{20} 起）外还能用什么、
零宽度位域怎么起对齐分隔作用，实现定义成分不少；
第二，可符号性——\Tp{int} 作底层类型的位域是否带符号，标准允许实现选；
第三，字段在存储单元内的分配顺序——从高位往低位还是从低位往高位，
实现定义；第四，字段放不进当前存储单元剩余空间时，是丢弃剩余另起单元，
还是跨单元切开，实现定义。四条连起来意味着：
位域声明刻画的是“按位分组访问”的语法便利，不是内存里的地图。

\begin{warning}[title={位域布局不是协议格式}]
用位域结构体去“映射”硬件寄存器或网络报文的字段顺序，
是在把三串实现定义的选择当成合同条款：字段分配方向、跨单元策略、
符号性都随编译器与目标平台走，换了工具链、甚至换了底层类型，
同一份声明的比特落点就可能整体镜像。
标准从未承诺位域与任何外部布局一致；\cpp{20} 的 \Tp{<bit>} 字节序设施
（第 \ref{ch:cl12-endian} 章）也不为位域兜底。
跨平台合同请用下一节的显式整数加移位掩码表达。
\end{warning}

可移植的替代方案把“位”显式化：一个确证宽度的整数当存储单元，
每个字段用移位与掩码读写，宽度关系用 \Fn{static\_assert} 钉死。

\begin{lstlisting}[style=cstyle, caption={显式整数与移位掩码：位布局自己写死（C11）}]
#include <stdint.h>

/* 存储单元:一个确证宽度的整数,而不是位域声明 */
#define DESC_READY_SHIFT  0u
#define DESC_READY_MASK   (0x1u << DESC_READY_SHIFT)
#define DESC_LEN_SHIFT    8u
#define DESC_LEN_MASK     (0xFFu << DESC_LEN_SHIFT)

_Static_assert(sizeof(uint32_t) == 4u,
               "mask constants assume a 32-bit storage unit");

uint32_t desc_encode(uint32_t ready, uint32_t len) {
    uint32_t word = 0u;
    word |= (ready & 0x1u)  << DESC_READY_SHIFT;
    word |= (len   & 0xFFu) << DESC_LEN_SHIFT;
    return word;
}

uint32_t desc_decode_len(uint32_t word) {
    return (word & DESC_LEN_MASK) >> DESC_LEN_SHIFT;
}
\end{lstlisting}

样例里的 \Fn{desc\_encode} 对任意目标平台产生同一个 32 位整数，
字段的比特位置由常量写死，不再有任何实现定义环节；
落盘前若要固定字节顺序，再叠加第 \ref{ch:cl12-endian} 章的显式字节序转换。
这正是“位域不能保证、显式整数能保证”的分界线。

\section{对齐与并发的接口}\label{sec:cl11-align-concurrency}

对齐的第一位用户其实是并发。两个原子变量落在同一条缓存行时，
对其中一个的写会让另一个所在的行失效，逻辑上互不相干的线程互相拖慢——
伪共享。前文的 \Tp{alignas} 过对齐把对象抬到破坏性干扰粒度，
是让“每线程一个计数器”真变成“每线程一行”的标准做法。
对齐与原子还有更硬的一层关系：多数平台的原子指令只对
“自然对齐且尺寸在指令集合覆盖范围内”的对象提供无锁保证；
过对齐的 16 字节原子对、带 packed 标记的原子字段，
在部分平台根本编不出单条原子指令，要么退化为加锁实现，
要么直接陷入——具体哪些尺寸无锁、哪些会陷入，随第 \ref{ch:cl10-cc} 章
讨论过的目标三元组而变，代码里的探测手段是
\Tp{is\_always\_lock\_free} 一类常量（\Tp{<atomic>}）与配置期断言，
实测口径也以该章为准。

\section{构建层：配置期探测并写入头文件}\label{sec:cl11-cmake-probe}

布局假设不该靠口头传承。\Tool{CMake} 把探测做成配置期的一次性仪式：
\Tool{CheckTypeSize} 模块的 \Tool{check\_type\_size} 命令问编译器
“某个类型有多大”，把答案存成缓存变量并生成 \Tp{HAVE\_\_SYMBOL} 形态的宏；
\Tool{CheckCXXSourceCompiles} 模块的 \Tool{check\_cxx\_source\_compiles}
命令把任意源码（包括用 \Fn{static\_assert} 写的布局断言）交给编译器试编译，
成败写入一个缓存布尔量；指针宽度则由现成的
\Tp{CMAKE{\Bk}\_SIZEOF{\Bk}\_VOID{\Bk}\_P} 变量给出。
最后用 \Tool{configure\_file} 把结果展开进一份生成的头文件，
源码侧只读这份头，不再各自猜。

\begin{lstlisting}[style=cmakestyle, caption={配置期探测布局事实并注入编译定义}]
include(CheckTypeSize)
include(CheckCXXSourceCompiles)

check_type_size("long" SIZEOF_LONG LANGUAGE C)
check_cxx_source_compiles([[
    #include <cstddef>
    int main() {
        static_assert(alignof(std::max_align_t) >= 4,
                      "baseline scalar alignment");
        return 0;
    }
]] BASELINE_ALIGN_OK)

message(STATUS "pointer size (bytes): ${CMAKE_SIZEOF_VOID_P}")
message(STATUS "long size: ${SIZEOF_LONG}, baseline align probe: ${BASELINE_ALIGN_OK}")
target_compile_definitions(app PRIVATE
    SIZEOF_LONG=${SIZEOF_LONG})
\end{lstlisting}

样例中 \Fn{check\_type\_size} 与 \Fn{check\_cxx\_source\_compiles}
的模块名与命令签名是 \Tool{CMake} 的稳定接口，
变量按各自命令的文档落进缓存；头文件模板（\Tp{layout{\Bk}\_config.h.in}）
里用 \Tp{@SIZEOF{\Bk}\_LONG@} 这类占位接收，命令细节以所用
\Tool{CMake} 版本文档为准。断言探针的好处是失败发生在配置日志里，
而不是某个平台的深夜生产事故里。

\section{三种改变布局的手段对照}\label{sec:cl11-means-table}

\begin{longtable}{@{}>{\raggedright\arraybackslash}p{0.14\textwidth} >{\raggedright\arraybackslash}p{0.24\textwidth} >{\raggedright\arraybackslash}p{0.34\textwidth}@{}}
\caption{三种“改变布局”的手段：何时生效、会破坏什么}\label{tab:cl11-layout-means}\\
\toprule
手段 & 生效时机与范围 & 主要破坏对象 \\
\midrule
\endfirsthead
\toprule
手段 & 生效时机与范围 & 主要破坏对象 \\
\midrule
\endhead
\bottomrule
\endfoot
重排成员 & 编译期，逐结构体，语义等价 & 本身不破坏规则；破坏的是“对端声明了同一顺序”的假设与旧二进制布局兼容 \\
打包指示 & 编译期，属性式逐类型，指示式与 \Sw{/Zp} 按范围覆盖 & 类型对外的 ABI（大小与偏移）、成员地址的对齐承诺、原子无锁前提 \\
过对齐（\Tp{alignas}） & 编译期抬类型对齐，动态分配时也要求分配器满足 & 对未知齐要求的对端不透明；两侧探测口径不一即产生布局分歧 \\
\end{longtable}

表里三行的共同点是：它们改变的都是“字节视图”，
而语言层的等价性保证从不延伸到位图上。
凡是要把字节视图当合同用的场合——共享内存、镜像序列化、文件格式——
三行里的“破坏对象”就全列变成风险清单，逐条要有对应检查。

\section{填充字节的对外泄漏}\label{sec:cl11-padding-leak}

\begin{abibox}[title={把带填充的结构体整体写出去}]
一个用 \Fn{memcpy} 或文件接口按 \Fn{sizeof} 整块写出的结构体，
被写进文件的除了命名字段，还有对齐垫出来的填充字节。
标准没有承诺填充必须清零——它往往是栈上或堆上残留的旧字节。
于是：同一平台同一编译器的两次导出可能字节不同，哈希与签名不稳定；
跨平台比较更直接失败，因为两侧的填充位置本就不同。
把这类对象当“值”做逐位比较、进缓存键、做差分同步，都会中招。
\end{abibox}

修复路径二选一。其一，逐字段序列化：写出时一个字段一个字段地编码，
读入时逐字段解码，填充根本不进入字节流；
其二，保留整体写入，但先把填充变成确定值——
用一个清零后的对象做落盘缓冲（样例思路：\Fn{memset} 整个对象再赋字段值，
或以零初始化聚合后再逐字段赋值），并且两侧接受“比较的是规范化字节”这一约定。
前者健壮、慢一点；后者快，但对打包方式敏感，
和第 \ref{ch:cl10-cc} 章的接口冻结纪律配套使用才安全。
同理适用于跨镜像传递：共享内存段的两侧若用不同工具链构建，
先对齐“填充与打包口径”，再谈字段偏移；哈希输入里掺进填充，
等于把内存残留当协议字段。

\begin{bestpractice}[title={把布局假设写进代码}]
凡依赖布局的接口，用三件套把假设显式化：\Fn{static\_assert} 钉住
\Fn{sizeof} 与 \Fn{alignof} 的可证关系，配置期探测（第
\ref{sec:cl11-cmake-probe} 节）钉住跨平台差异，序列化层逐字段读写
钉住对外合同。断言里只写本章四条规则能推出的相对关系；
绝对数值一旦要写，来源必须是探测，而不是“我这台机器上是 12”。
\end{bestpractice}

本章的地基是“标准只承诺相对关系，绝对布局是实现与平台联合决定的”。
下一章处理这个地基上最容易出事的两类跨平台差异：字节序与数据模型。

\chapter{字节序与数据模型}\label{ch:cl12-endian}

布局的另一半，是两个各自独立、却常被混为一谈的问题：
字节序问“一个多字节对象的各个字节在内存里按什么顺序摆放”，
数据模型问“\Tp{int}、\Tp{long}、指针这些基本类型在目标平台上有多宽”。
两者都是跨平台差异的高发地：前者决定你读写的每一个外部格式，
后者决定你在接口边界写下的每一个类型名。
上一章（第 \ref{ch:cl11-layout} 章）把“字节视图何时成为合同”讲清楚，
本章在这个合同上补全字节序与宽度两个条款；
接口两侧如何共享同一口径，回指第 \ref{ch:cl10-cc} 章的编译口径纪律。

\section{先把两个问题分开}\label{sec:cl12-two-questions}

第一问是内存内序：一个 \Tp{uint32\_t} 变量存放在内存里的四个字节，
高位字节落在低地址还是高地址。这是机器与 ABI 的属性，语言标准对它不作承诺，
只规定运算语义——你在小端机器上写 \Tp{x = 0x11223344u}，
\Fn{printf} 出来永远是同一个数值，字节顺序只是它背后的落地方式。
第二问是外部序：文件、网络报文、共享内存段这些“离开本进程的字节流”里，
约定第几个字节承载哪个位置的值。这是格式的条款，属于设计文档，不属于 CPU。

\begin{guideline}[title={“我的机器是小端”回答不了“这个文件怎么读”}]
内存序回答“本机对象如何落地”，外部序回答“这份字节流如何约定”。
读一个外部格式，唯一正确的依据是格式规范本身；机器字节序只决定
读出后要不要做交换，永远不决定字段边界。把两者混在一起的典型后果，
是“在我机器上能读”的解析器换台机器就整体镜像了字段。
\end{guideline}

\begin{lstlisting}[style=outstyle, caption={同一份字节的两种解释（示意，非实测输出）}]
# 文件头四字节在盘上的顺序永远是唯一的: 44 33 22 11
od -An -t x1 data.bin            # 原始字节: 44 33 22 11
od -An -t x4 --endian=big data.bin    # 按 32 位单元、大端解释: 44332211
od -An -t x4 --endian=little data.bin # 按 32 位单元、小端解释: 11223344
# 变的是解释方式,不变的是盘上字节;--endian 一类开关
# 以所用 coreutils 版本文档为准
\end{lstlisting}

样例想说明的顺序只有一个：字节是事实，整数是解释。
调试外部格式时先看 \Tool{od}、\Tool{hexdump} 一族的原始字节，
再谈“这个数是几”；解释开关只用于比对假设，不用于下结论。

\section{探测与转换：语言与库给了什么}\label{sec:cl12-probe}

\subsection{std::endian 与编译期探测}\label{sec:cl12-std-endian}

\cpp{20} 把“探测本机字节序”标准化为 \Tp{std::endian}（\Tp{<bit>}）：
枚举值 \Tp{little}、\Tp{big} 加一个混合字节序的 \Tp{mixed}，
常量 \Tp{native} 表示当前编译目标的字节序。
它可以出现在 \Fn{static\_assert} 与 \Tp{if constexpr} 里，
这是标准第一次让“字节序假设”能写进编译期检查。

\begin{lstlisting}[style=cppstyle, caption={编译期字节序探测与交换（std::byteswap 为 C++23 设施）}]
#include <bit>
#include <cstdint>

constexpr bool target_little =
    (std::endian::native == std::endian::little);

// 把一个按大端字节流拼进来的值换算为本机数值。
// std::byteswap 来自 <bit>(C++23);较旧的标准用编译器内建函数,
// 再旧的用移位掩码自己写(见下一小节的 C 样例)。
constexpr std::uint32_t from_be32(std::uint32_t raw) {
    if constexpr (target_little) {
        return std::byteswap(raw);
    } else {
        return raw;
    }
}

static_assert(sizeof(std::uint32_t) == 4u,
              "fixed-width type is the precondition");
\end{lstlisting}

\subsection{交换设施的三层生态}\label{sec:cl12-swap-layers}

标准库之外，字节交换有三层设施，角色各不相同。
编译器内建函数族——GCC 一族的 \Tp{\_\_builtin\_bswap16}、
\Tp{\_\_builtin\_bswap32}、\Tp{\_\_builtin\_bswap64}，MSVC 的
\Tp{\_byteswap\_ushort}、\Tp{\_byteswap\_ulong}、\Tp{\_byteswap\_uint64}
（后者声明于运行库头）——它们只保证“值反转”，生成的通常是单条
\Tp{bswap} 一类指令，是手写可移植层最常落脚的一层；
较新的标准库提供了 \Fn{std::byteswap} 等类型化设施（\cpp{23}，\Tp{<bit>}）；
而网络字节序函数族（\Fn{htonl} 一类）服务于套接字地址这一具体场景，
把“协议约定”硬编码进了函数名。三层的正确性相同，可移植面不同。

\begin{lstlisting}[style=cstyle, caption={逐字节拼装与分宽度交换：两种都正确的写法}]
#include <stdint.h>
#include <stdio.h>

/* 写法一:逐字节拼装。与本机字节序无关,永远正确。 */
uint32_t load_be32(const unsigned char *p) {
    return ((uint32_t)p[0] << 24) | ((uint32_t)p[1] << 16) |
           ((uint32_t)p[2] << 8)  | (uint32_t)p[3];
}

/* 写法二:值反转。用于字节已经在内存里拼成整字、只差镜像的场合。 */
uint32_t swap32(uint32_t v) {
#if defined(__GNUC__) || defined(__clang__)
    return __builtin_bswap32(v);
#elif defined(_MSC_VER)
    return _byteswap_ulong(v);
#else
    return ((v & 0x000000FFu) << 24) | ((v & 0x0000FF00u) << 8) |
           ((v & 0x00FF0000u) >> 8)  | ((v & 0xFF000000u) >> 24);
#endif
}

void report(uint64_t n) {
    /* 格式宏 PRIu64 来自 <inttypes.h>,见下一小节 */
    printf("%" PRIu64 "\n", n);
}
\end{lstlisting}

\subsection{定宽类型、格式宏与格式化出口}\label{sec:cl12-cstdint}

字节序讨论的前提是“先有一个宽度确证的整数”。\Tp{<cstdint>}（\CL{} 侧
\Tp{<stdint.h>}）提供的 \Tp{uint32\_t}、\Tp{int64\_t} 一族
只在实现能提供该精确宽度时才定义——没有 \Tp{uint64\_t} 的平台是合法的，
所以跨镜像接口用它们之前，第 \ref{sec:cl11-cmake-probe} 节的探测环节不可省。
打印侧的配套是 \Tp{<inttypes.h>} 的格式宏族：
\Tp{PRIu64}、\Tp{PRId64}、\Tp{PRIx32} 这类宏展开为该平台上
对应类型在 \Fn{printf} 家族的正确转换说明尾部，与格式串主体
在预处理期拼接——先写百分号引导符、紧跟宏名，不要把宏名塞进字符串——
这比按 \Tp{ld}、\Tp{lld} 一类长度手工猜测转换说明要稳。
\cpp{20} 的 \Fn{std::format}（\Tp{<format>}）是格式化侧的新出口：
类型化参数不再需要手工拼接转换说明，定宽整数的输出随类型自动正确；
它的运行时库覆盖度随版本演进，能否使用以工具链文档为准。

\begin{compare}[title={交换设施的三层分工}]
逐字节拼装（移位掩码）最笨也最稳，跨任何编译器语义一致；
编译器内建函数次之，依赖一族宏名，用第 \ref{sec:cl12-swap-layers} 节
的预处理探测包一层即可；标准库设施（\Fn{std::byteswap}）与语言特性
（\Tp{std::endian}）最干净，但受语言版本门槛。
三层的选择本质是“你的可移植面画在哪”。
\end{compare}

\section{数据模型：ILP32、LP64、LLP64 与更罕见的名字}\label{sec:cl12-datamodel}

模型名的构词法一行讲完：末尾数字是指针宽度，
名字里的字母标出与指针同宽的类型集合——\Tp{S}、\Tp{I}、\Tp{L}、
\Tp{LL}、\Tp{P} 分别对应 \Tp{short}、\Tp{int}、\Tp{long}、
\Tp{long long} 与指针；没进名字的整型退回较窄档位，
而 \Tp{long long} 按 \CL{} 标准（C99 起）至少 64 位。
ILP32 里 \Tp{int}、\Tp{long}、指针都是 32 位；
LP64 只把 \Tp{long} 和指针抬到 64；LLP64 不动 \Tp{long}，
而是让 \Tp{long long} 与指针到 64。另有两个罕见名字：ILP64
（\Tp{int} 也 64 位）与 SILP64（连 \Tp{short} 都 64 位），
历史上见于超级计算机与专有平台，细节以对应平台 ABI 文档为准，
现代主流桌面与服务器目标实际只有前三种。

\begin{longtable}{@{}>{\raggedright\arraybackslash}p{0.12\textwidth} >{\raggedright\arraybackslash}p{0.07\textwidth} >{\raggedright\arraybackslash}p{0.07\textwidth} >{\raggedright\arraybackslash}p{0.09\textwidth} >{\raggedright\arraybackslash}p{0.08\textwidth} >{\raggedright\arraybackslash}p{0.26\textwidth}@{}}
\caption{四种数据模型规定的类型宽度（单位：位）}\label{tab:cl12-data-models}\\
\toprule
模型 & \Tp{int} & \Tp{long} & \Tp{long long} & 指针 & 常见落点 \\
\midrule
\endfirsthead
\toprule
模型 & \Tp{int} & \Tp{long} & \Tp{long long} & 指针 & 常见落点 \\
\midrule
\endhead
\bottomrule
\endfoot
ILP32 & 32 & 32 & 64 & 32 & 32 位目标：\Nix{} 的 arm、i386，\Win{} 的 32 位进程 \\
LP64 & 32 & 64 & 64 & 64 & \Nix{} 与 \Mac{} 的 64 位用户态（psABI 规定） \\
LLP64 & 32 & 32 & 64 & 64 & \Win{} 的 64 位用户态 \\
ILP64 与 SILP64 & 64 & 64 & 64 & 64 & 罕见；历史超级计算机与专有平台 \\
\end{longtable}

为什么 \Win{} 的 64 位没有随大流改 \Tp{long}？
因为 \Tp{long} 在 32 位 \Win{} 时代就是 32 位，海量 Win32 API 签名、
结构体字段与位标志把它当 32 位宽度用。进 64 位时改它等于重写整套 API，
微软选择只把指针（和句柄这类指针宽度类型）翻倍，\Tp{long} 原地不动——
这就是 LLP64。同一段“\Tp{long} 装指针”“\Tp{long} 当 64 位容器”的代码，
在 LP64 与 LLP64 上是两种命运。

\begin{warning}[title={long 与 size\_t 混用的三类故障}]
其一，截断：LP64 到 LLP64（或反向）之间用 \Tp{long} 中转 \Tp{size\_t}
或指针宽度，值在 32 位一侧被削顶，且不报错。
其二，格式化：\Tp{ld} 长度的转换说明在 LP64 打 64 位、在 LLP64 打 32 位，同一格式串
两侧含义不同；\Tp{z} 修饰符一族（\Tp{zu}、\Tp{td} 等）的支持度还随运行时库版本
与编译器年代变化，历史包袱主要在老 \Win{} 运行时。三类里最难查。
其三，位移：\Tp{1L << 40} 里的 1 在 \Tp{long} 为 32 位的模型下
移位量超出宽度，直接落入未定义行为；位图与掩码代码换模型即碎。
\end{warning}

\section{可移植宽度的写法纪律}\label{sec:cl12-width-discipline}

纪律可以压缩成四句。容量用 \Tp{size\_t}：它是可移植代码里唯一的
“对象大小”类型，不借 \Tp{long} 代打。差值用 \Tp{ptrdiff\_t}：
指针算术的结果类型，与 \Tp{size\_t} 同宽但带符号，两者混用另生一坑。
要精确宽度用定宽类型：\Tp{uint32\_t} 一类，前提是探测确认其存在。
跨镜像接口只用确证宽度的类型加显式填充：结构体字段的宽度合同
由第 \ref{ch:cl11-layout} 章的布局纪律与这里的宽度纪律共同构成；
\CPP{} 还要叠加名字修饰的稳定策略，见第 \ref{ch:cl13-cxx-abi} 章，
接口冻结的一般原则回指第 \ref{ch:cl10-cc} 章。
指针宽度差异（32 位与 64 位之间）是数据模型的另一半：它改变
可寻址空间、改变 \Tp{uintptr\_t} 的意义、改变第 \ref{ch:cl07-object} 章
讨论过的重定位类型集合，跨位宽移植是另一场手术，不在本章展开。

\section{字节序与 ABI 的交叉}\label{sec:cl12-endian-abi}

\begin{table}[htbp]
\centering
\caption{跨字节序时：什么跟着变，什么不跟着变}\label{tab:cl12-endian-scope}
\begin{tabularx}{0.80\textwidth}{>{\raggedright\arraybackslash}p{0.20\textwidth} X X}
\toprule
对象 & 随本机字节序变吗 & 处理纪律 \\
\midrule
文件与网络字节流 & 不变——格式规范规定外部序 & 显式转换（第 \ref{sec:cl12-probe} 节三层设施任选） \\
内存对象的 \Fn{memcpy} 原始导出 & 变——字节视图是本机落地物 & 别导出；逐字段序列化（第 \ref{ch:cl11-layout} 章） \\
位域与整数的位视图 & 变——字段落点还叠加实现定义 & 显式整数加移位掩码，位域不出镜像 \\
名字修饰与符号表字符串 & 不变——ASCII 序列无字节序 & 跨端调试符号可安全共享 \\
目标文件格式头与节表 & 不由本机决定——\Elf{} 与 \Macho{} 在文件里自带字节序声明，\PeFmt{} 按规范固定为小端 & 用理解目标格式的工具，随文件解释而非随本机解释 \\
\bottomrule
\end{tabularx}
\end{table}

表的下半行解释一个常见误解：“大端机器上的 \Elf{} 文件拿到小端机器就读不了”。
读不了的是把文件头当本机整数直接解释的代码；
规范化的工具链与调试器按文件自身的字节序声明解释字段，两侧看同一份文件
得到同一张符号表。反过来说，混合大小端的调试现场真正麻烦的
是数据侧：核心 dump 里的整数、串口日志抓下来的报文、对端固件写出的
遥测文件，都需要逐流确认外部序后手工重解释，
这就是这类场合离不开 \Tool{od} 一族与专用格式脚本的原因。

\begin{note}[title={字节序常与架构耦合出现}]
实践中大端目标（部分网络处理器、历史 \Mx{PowerPC} 与 \Mx{SPARC} 年代）
很少与 \xSixtyFour{}、\ArmSixtyFour{} 的主流小端组合共处，
但 ARM 一族允许按配置切大小端模式，双端模式的目标确实存在。
所以探测不能停在“这台开发机是小端”：交叉编译时
\Fn{std::endian::native} 描述的是编译目标的字节序，与开发机无关。
\end{note}

\section{原子、内存模型与字节序}\label{sec:cl12-atomics}

多字节原子操作的第一个前提是地址对齐——这条前提在大小端两侧同样成立：
无论文档把高字节排在低地址还是高地址，原子对都要落在整对象起始地址，
第 \ref{ch:cl11-layout} 章的对齐纪律不因字节序而放松。
真正随字节序变的是“拆分后的世界”：把一个宽字段拆到若干存储单元
（无锁数据结构里的字长计数器、seqlock 里的版本与载荷）时，
哪个子字段落在低单元、哪个落在高单元，随字节序镜像；
依赖“先写低字后写高字”这类顺序推理的无锁协议，
换端不重推导就是错的。字节序在语言与库层可见的部分到
\Fn{std::endian} 为止，更深的“内存模型承诺什么”属于
未定义行为分册的内存模型主题，本章只标注边界：
字节序是表示问题，原子是同步问题，两者在“拆字”处交叉但不混同。

\section{案例：一个二进制容器格式的头}\label{sec:cl12-container}

\begin{abibox}[title={容器格式设计：把字节序与宽度写进合同}]
设计一个跨平台容器格式的最小头：4 字节魔数、4 字节版本、4 字节头长，
共 12 字节，外部序固定为大端，全部字段使用确证宽度类型，
读入方先验魔数、再拒未知版本、再核对头长，任何一步失败都不碰载荷。
四条纪律各挡一类事故：显式大端挡掉“本机序直读”的镜像错误；
显式版本挡掉格式演化后的静默错读；确证宽度挡掉数据模型漂移；
拒绝未知版本挡掉把未来字节当现在字段。反例是把 \Fn{sizeof} 或
本机 \Fn{uint64\_t} 直接写进头——那是把第 \ref{ch:cl11-layout} 章
与本章的所有坑一次性埋进格式里。
\end{abibox}

\begin{lstlisting}[style=cstyle, caption={容器头：逐字节读入与三重校验（示意格式）}]
#include <stdint.h>

/* 盘上布局(大端外部序):magic[4] | version(4) | hdr_size(4) */
/* load_be32：即本章前文“逐字节拼装”的写法一 */

int fmt_header_check(const unsigned char *buf, uint32_t avail) {
    uint32_t version;
    uint32_t hdr_size;
    if (avail < 12u) {
        return -1;                    /* 短读:连定长头都不够 */
    }
    if (!(buf[0] == 0x46u && buf[1] == 0x4Fu &&
          buf[2] == 0x52u && buf[3] == 0x4Du)) {
        return -2;                    /* 魔数不符:根本不是这个格式 */
    }
    version  = load_be32(buf + 4);
    hdr_size = load_be32(buf + 8);
    if (version != 1u && version != 2u) {
        return -3;                    /* 未知版本:拒绝而非猜测 */
    }
    if (hdr_size != 12u) {
        return -4;                    /* 头长自检:防版本混叠 */
    }
    return 0;
}
\end{lstlisting}

注意样例从头到尾没碰任何多字节加载：四个 32 位字段里，
魔数逐字节比，其余两个用 \Fn{load\_be32} 拼装，
缓冲区里的原始字节序与本机字节序全程无关——
这就是“外部序由格式规定，本机序只影响内存落地”的落地写法。

\section{收口：跨平台读写的三条纪律}\label{sec:cl12-closing}

\begin{bestpractice}[title={外部格式读写三纪律}]
永远显式转换：读入与写出都过一道确证方向的手（拼装、内建函数或
\Fn{std::byteswap}），让“依赖本机序”的代码在编译期就无处藏身；
永远显式宽度：边界上的每个字段用 \Tp{uint32\_t} 一族命名，
容量与差值用 \Tp{size\_t}、\Tp{ptrdiff\_t}，打印用 \Tp{PRIu64} 族或
\Fn{std::format}；永远校验版本：魔数、版本、头长三关前置，
未知版本宁可拒绝。三条纪律对应的外向观察点——哪些字段要查、
查完记在哪——并入第 \ref{ch:cl15-endtoend} 章的观察点清单。
\end{bestpractice}

至此，布局、字节序与数据模型三块拼图合拢：对象的字节视图有了合同，
整数的宽度有了口径。下一章把镜头拉回代码层——\CPP{} 在布局之上
还要叠一层名字与语义的合同，即名字修饰与 ABI。

\chapter{\CPP{} ABI：名字修饰、虚表与展开表}\label{ch:cl13-cxx-abi}

\CPP{} 程序在链接期与加载期看到的从来不是源码，而是源码的编码结果：一个名字如何变成符号、
一个对象里挂着哪几张表、一次抛出如何找到它对应的清理动作，都由一套跨模块共享的约定决定。
这套约定就是 ABI（二进制接口约定）。本章只讲其中属于 \CPP{} 的那一半，把调用约定、参数与
返回值传递、目标文件格式与加载语义留给第 \ref{ch:cl10-cc} 章。

写作口径先钉住三句。第一，凡是本章写出的修饰前缀、节名与符号名，都属于两族主流实现长期
稳定的公共部分，但它们是“实现事实”而非“语言承诺”，换实现、换版本都可能改动。第二，落在
具体字符上的规则一旦不能确证，本章只描述这个角色要编码什么，不复述编码表。第三，所有工具
输出都是示意，不是本仓库的实测值。

\section{两层 ABI，以及本章与调用约定章的分工}\label{sec:cl13-scope}

把 ABI 想成一层是不准确的，至少要拆成两层来谈，否则排障时找不到归属。

\subsection{语言层 ABI}\label{sec:cl13-scope-lang}

语言层 ABI 规定的是“同一个类型在不同翻译单元里必须长成同一个样子”所需的全部信息：类内部的
成员排布与虚表指针放在哪里、多继承与虚继承时子对象如何被寻址、运行时类型信息以什么结构存在、
异常抛出与栈展开靠哪张表推进、带非平凡构造与析构的全局对象由谁在何时登记。这些内容与操作
系统无关，却决定了两个分别编译的目标文件能否共享同一个类型。

\subsection{系统层 ABI}\label{sec:cl13-scope-os}

系统层 ABI 规定的是“函数之间怎么说话，程序怎么被装进内存”：寄存器与栈的调用约定、结构体
按值传递时的落位规则、目标文件格式、符号解析时机、加载与重定位。第 \ref{ch:cl10-cc} 章那张
对照表讲的就是这一层；第 \ref{ch:cl07-object} 章讲容器结构，第 \ref{ch:cl11-layout} 章讲
布局探测手段，第 \ref{ch:cl09-dynamic} 章讲加载期发生了什么。

\begin{guideline}[title={本章与调用约定章的分界线}]
判断问题归属只需一句口诀：换一个 \CPP{} 编译器就会变的东西属于语言层 ABI，换一个操作系统
或换一个架构就会变的东西属于系统层 ABI。虚表槽位序号、修饰名形态、展开表位置属于前者；
关键字约定的细节、参数落在哪个寄存器属于后者。两侧同时满足才有二进制兼容，只满足一侧
得到的“看起来能跑”，是下一次崩溃的伏笔。
\end{guideline}

\section{名字修饰：为什么必须有它}\label{sec:cl13-mangling}

\CL{} 时代一个函数名就是一个符号，链接器只需按字符串比对。\CPP{} 打破了这个前提：同一作用域
里可以有多个同名函数，同名符号可以分布在多个命名空间中，同一个模板对不同实参会生成不同的
实体。链接器不认识类型系统，它只认字符串，所以编译器必须把“这个名字的全部区分信息”压进一个
唯一字符串，这一步就叫名字修饰。

\subsection{修饰必须编码的四件事}\label{sec:cl13-mangling-input}

其一，重载：形参类型的序列必须进名字，否则同名不同参的两个函数会在符号层撞车。
其二，作用域：命名空间与嵌套类的路径必须进名字。其三，模板：实参列表必须进名字，而且函数
模板与变量模板的实例还要能表达“这是同一个模板的第几份实体”。其四，语言链接：以 \CL{} 语言
链接声明的函数必须退回到裸名字，这样目标文件里的符号才能被 \CL{} 侧代码与既有库按原名解析。
语言链接的写法细节与调用约定的组合属于另一分册，本章只保留机制理由：语言链接是编译器决定
“要不要修饰”的那一个开关。

另有两类信息通常也被编进名字：引用与 const 组合在签名上的区别、成员函数的 cv 限定与引用
限定。反过来说，返回值类型一般不进名字，于是“仅返回值不同”的重载在源码层就不被允许，在符号
层也就必然冲突；这是一条从语言规则直连到链接期行为的因果链。第 \ref{ch:cl03-resolution} 章
讲的符号解析时机，正是这条链的下游。

\subsection{两族方案的形态对照}\label{sec:cl13-mangling-families}

\Nix{} 与 \Mac{} 上主流编译器共享一族方案（下文称表式前缀族），\Win{} 上的另一族实现
（下文称自定前缀族）形态完全不同。本章只描述形态，不逐字符复述规则，因为字符级的编码表
既长又随版本微调。

\begin{compare}[title={两族修饰形态的差异}]
表式前缀族：所有 \CPP{} 符号带一个统一短前缀，之后是作用域与实体类别的标记字符，标识符用
“长度加原文”的方式串起来，因此不需要为奇怪字符做转义。例如 \Sym{\_ZN6Widget4drawEv}，
\Tp{6Widget} 与 \Tp{4draw} 的长度前缀让解码器能直接切出名字，尾部的 \Tp{v} 表示无参。
自定前缀族：以一个可辨识的字符起头，把类名、函数名、调用约定、参数类型、返回类型按自己的
类型编码表串起来，末尾再挂限定标记；符号里会出现问号与 at 号一类可读性极低的连接字符。
共同点：两族都把“能否唯一区分重载”放在第一位。不同点：两族互不认账，同一份源码在两侧
得到的符号串毫无对应关系，也没有任何转换表可以在事后把一族改成另一族。
\end{compare}

\subsection{把修饰名还原成源码形式}\label{sec:cl13-demangle}

工程上几乎不会手读修饰名，都要靠解码入口。可确证的入口有三类：命令行的 \Tool{c++filt}
（\Tool{binutils} 家族）与 \Tool{llvm-cxxfilt}（\Tool{LLVM} 家族），把标准输入里的修饰名
换成源码形式；\Tool{nm} 的 \Sw{-C} 开关在列出符号时顺手解码；表式前缀族还把解码能力做进了
运行时，\Fn{\_\_cxa\_demangle} 是其 ABI 文档规定的入口，异常信息里的可读类型名就来自它。
\Win{} 侧的对应物是 \Tool{undname} 这一类反装饰工具，只能解自定前缀族自己的名字。

\begin{warning}[title={修饰名不是跨实现的合同}]
一旦把某个具体的修饰名写进脚本、白名单或测试断言，就等于把约定交给了某个编译器的某个版本。
可行的用法只有两种：同一实现内做同版本自比较，或解码后比较源码形式的签名。需要跨实现稳定
的对外符号，请退回到 \CL{} 语言链接那一层，不要用 \CPP{} 名字对外承诺。
\end{warning}

\section{虚表、子对象与运行时类型信息}\label{sec:cl13-vtable}

先划一句界线，因为它决定这一节里哪些内容可以当事实引用、哪些必须回查文档。
虚表与虚表指针是\emph{实现方案}，不是语言要求。标准只规定虚函数的语义：按对象的动态类型
选出最终覆盖者、构造与析构期看到的是“只到当前层”的那个视图、\Tp{dynamic\_cast} 与
运行时类型信息要能回答哪些问题。它没有规定这些信息放在哪里、用几张表、按什么序号索引，
也没有规定对象里必须多出一个指针。原则上实现可以完全不用表——把派发折叠成静态解析或分支，
同样满足语义；这类改写在优化里确实发生，第 \ref{sec:cl04-cxx-semantics} 节的警告框里
已经为它划过一条线：某个虚表指针被优化掉，不等于虚函数机制被取消。

那为什么本卷仍以表为主线？因为主流实现都选了这一个方案，彼此只有细节不同。
桌面与移动的三个实现族、\xSixtyFour{} 与 \ArmSixtyFour{} 两套架构上的主流编译器，
给出的都是“对象里放指针、类型旁放一张按序号排列的表”；差别集中在表项里放不放偏移修正、
运行时类型信息挂在表头还是表尾、构造期另备辅助表还是在写指针时换表、符号名长什么样。
于是这一层的口径是：\textbf{机制可以依赖，布局与名字必须按实现族核对}。

同一句口径适用于本章其余部分。名字修饰的形态、展开表放在哪个节、弱符号与重复实体的合并
策略、线程局部存储选哪种访问模型、\Sw{\_\_cxa\_} 一族符号由谁提供——这些都不是语言标准
规定的，而是主流实现固定下来的。本卷只说明主流实现给出的那一套，凡写到具体形态，
出处一律是“某实现族的文档与它的产物”，不写成语言规则。读到别处把这些形态当成
\CPP{} 语法来讲的，按这条口径打折。

\subsection{对象里放什么，表里放什么}\label{sec:cl13-vtable-layout}

带虚函数的类，其对象里会多出至少一个指针，通常称为虚表指针。它指向一张按索引排列的函数
地址表：编译期把每个虚函数映射成一个固定的槽位序号，运行期用“虚表指针加槽位偏移”完成派发。
这张表不是纯函数地址数组，至少还要能承载下面这些东西：

\begin{itemize}
\item 类型信息的反向指针：从虚表能找回描述该类型的记录，\Tp{dynamic\_cast} 与 \Tp{catch}
的类型匹配都要用它；
\item 偏移修正项：多继承与虚继承时，某个槽位真正需要的 \Tp{this} 可能不是当前子对象的地址，
表项附近要能表达“先加多少字节再跳转”；
\item 继承图记录：\Tp{dynamic\_cast} 沿继承图找目标类型时用的基类数组与层级描述，运行时
类型信息的可读名字串也在这里；
\item 构造期视图：基类构造期间虚表指针必须指向“只到当前层”的那张表，否则在构造函数里调用
虚函数会派发到尚未构造完成的派生实现。为此有的实现另备一张供构造与析构期使用的辅助表
（表式前缀族的文档把它叫 VTT），具体形态以该实现的文档为准。
\end{itemize}

\subsection{多继承：同一对象里的多个视图}\label{sec:cl13-multi-inherit}

多继承时一个对象里排着多个基类子对象，也就可能有多个虚表指针；把派生类指针向上转成某个
非首个基类的指针，得到的地址不是对象起始处。机制层面最疼的后果是回调错位：

\begin{lstlisting}[style=cppstyle, caption={多继承下的子对象地址调整：机制示意}]
struct IObserver { virtual void onEvent(int) = 0; };
struct IRenderer { virtual void draw() = 0; };

struct Widget : IObserver, IRenderer {
    void onEvent(int) override;   // 它落在第二个子对象里
    void draw() override;
};

// post 只按 IObserver 这一个视图给出指针
void post(IObserver* o) { o->onEvent(7); }

void Widget::onEvent(int) {
    // 到这里 this 已经是调整过的 IObserver 子对象视图。
    // 若把这个指针原样存成 void* 再按 Widget* 读回来，
    // 后续访问的成员就整体偏移了一个子对象的宽度。
}
\end{lstlisting}

编译期插入的调整入口（业界常称 thunk）让“同一个函数从不同子对象入口进来都能拿到正确的
\Tp{this}”成立。手工把指针转成 \Tp{void*} 存起来、再按另一个类型读回来，就绕过了这套调整，
错位由此产生，而且往往在很久之后才表现为一次奇怪的崩溃。

\subsection{虚继承：多一层间接}\label{sec:cl13-vbase}

虚继承要保证共享基类在最终派生对象里只有一份，但这一份的位置取决于整张继承图的形状，
而继承图可能在别的翻译单元里继续被派生下去。于是编译期不能再把共享基类的地址写成对象内的
固定偏移，只能记一个“如何找到它”的间接量，访问时多走一跳。这层间接的代价是三重的：每次
访问共享基类多一次取址，虚表槽位要带上额外的偏移项，对象的构造顺序信息也必须集中到最终
派生类一处（表式前缀族为此区分完整对象构造与子对象构造两套入口）。

\subsection{这些机制如何以符号出现}\label{sec:cl13-symbols}

上面每张表都是一个数据符号，会出现在目标文件里参与链接与去重。表式前缀族的形态大致如下
（示意，实际字符以所用版本输出为准）：

\begin{lstlisting}[style=outstyle, caption={虚表与运行时类型信息对应的符号（示意）}]
# nm -C widget.o 的输出骨架：三张表加一个调整入口
_ZTV6Widget      -> vtable for Widget
_ZTI6Widget      -> typeinfo for Widget
_ZTS6Widget      -> typeinfo name for Widget
_ZThn8_N6Widget4drawEv  -> thunk to Widget::draw()
# 同一份源码在自定前缀族下叫法完全不同（示意）：
# 虚表是一个 ??_7 开头的符号，类型信息是若干 ??_R 开头的记录
\end{lstlisting}

这些符号的链接语义与普通函数并无不同：虚表与类型信息记录是否唯一实例、是否被合并、是否
带弱属性，取决于实现与可见性设置；模板实例化产生的虚表在链接期收束，靠的正是第
\ref{ch:cl08-static-link} 章讲的多副本合并机制。把四组对应关系列进表：

\begin{longtable}{@{}>{\raggedright\arraybackslash}p{0.16\textwidth} >{\raggedright\arraybackslash}p{0.20\textwidth} >{\raggedright\arraybackslash}p{0.20\textwidth} >{\raggedright\arraybackslash}p{0.24\textwidth}@{}}
\caption{\CPP{} 机制在链接期留下的四类可见物（角色对照，非穷举）}\label{tab:cl13-linkface}\\
\toprule
机制 & 编译期写下的数据 & 链接期的符号形态 & 典型失效方式 \\
\midrule
\endfirsthead
\toprule
机制 & 编译期写下的数据 & 链接期的符号形态 & 典型失效方式 \\
\midrule
\endhead
\bottomrule
\endfoot
虚函数派发 & 槽位序号到地址的表 & 虚表数据符号与调整入口 & 中间插入虚函数导致序号错位 \\
运行时类型信息 & 类型记录与继承图 & 类型信息与名字串符号 & 跨镜像时同一类型存在两份地址 \\
子对象调整 & 偏移修正项 & 带偏移的调整入口 & 手工转指针绕过调整 \\
静态初始化 & 初始化与销毁登记 & 登记表与守卫变量 & 加载时机与顺序假设被打破 \\
\end{longtable}

\subsection{关掉运行时类型信息与异常时的连带后果}\label{sec:cl13-no-rtti}

工具链上有两类开关能把上面的数据整块省掉：关闭运行时类型信息与关闭异常支持（常见写法是
\Sw{-fno-rtti} 与 \Sw{-fno-exceptions} 一族）。省掉的直接后果是符号集合变小、镜像变小，
间接后果更值得注意：一旦某个头文件里的类被两侧以不同开关编译，一侧发出了虚表与类型信息、
另一侧引用不到，就会出现只在链接期暴露的未定义符号；这属于配置不一致，不属于代码缺陷。
库对外发布时把这类开关状态一并说明，比事后解释未定义符号便宜得多。

\section{表驱动展开与异常}\label{sec:cl13-unwind}

\subsection{把 \Tp{catch} 编译成数据}\label{sec:cl13-lsda}

主流的展开方案是表驱动的：编译器在编译期就把“哪些指令地址范围可能抛出、对应的清理动作
在哪、能匹配哪些异常类型”写成只读数据，随函数一起放进镜像。运行期抛出时，展开器从当前
栈帧的返回地址出发，用地址区间查到清理入口，依次执行析构、判断类型是否匹配，匹配则把
控制权交给处理块。整个过程不需要在栈上留下可执行的手写记录，也不要求每个函数自带一条
登记链。

与之对照的另一族思路是“在栈上走一遍”：每帧在建立时把恢复地址登记进一条链，抛出时沿链
回溯并执行登记的清理逻辑。它的优点是无需镜像附带表，缺点是每个函数都要多几条建立与拆除
指令、干净返回路径上也带着这份开销，而且栈一旦被破坏回溯就跟着失控。现代桌面工具链的
主流选择是表驱动，两条路径要回答的问题却相同：清理谁、匹配什么类型、控制权交给谁。

三个角色的名字值得记住，排障时会遇到：描述本帧该做什么的语言相关数据区、每帧共用的展开
描述（如何恢复栈指针与被调用者保存寄存器）、以及由展开器回调的每函数人格例程（负责判断
类型并决定是否停下）。前两者决定“能不能爬栈”，第三者决定“能不能落到正确的处理块”，
缺任何一个的症状完全不同：前者缺失通常是终止，后者错配是异常穿过帧却无人认领。

抛出一侧的运行时入口也是 ABI 的一部分，排障时会直接在符号表里看到它们：

\begin{itemize}
\item \Fn{\_\_cxa{\Bk}\_throw}：把异常对象交出去，并启动展开；
\item \Fn{\_\_cxa{\Bk}\_begin{\Bk}\_catch}：在处理块入口把指针调整到正确的子对象；
\item \Fn{\_\_cxa{\Bk}\_rethrow{\Bk}\_exception} 与 \Fn{\_\_cxa{\Bk}\_call{\Bk}\_terminate}：
  分别对应重抛与终止；
\item 底层沿栈帧回溯的入口在表式前缀族里另有其名，职责只是逐帧询问该帧要不要清理。
\end{itemize}

这些入口的存在说明异常不是纯语法糖，而是一组有符号、可被链接、也可被裁掉的运行时。

\subsection{三种平台把展开信息放在哪}\label{sec:cl13-unwind-where}

\begin{table}
\centering
\begin{tabular}{@{}>{\raggedright\arraybackslash}p{0.12\textwidth} >{\raggedright\arraybackslash}p{0.34\textwidth} >{\raggedright\arraybackslash}p{0.32\textwidth}@{}}
\toprule
平台 & 展开信息的落点 & 排查时的第一手证据 \\
\midrule
\Nix{} & \Sect{eh\_frame} 承担通用帧描述，\Sect{gcc\_except\_table} 承担语言相关数据 &
\Tool{readelf} 看相关段是否存在、有没有被 \Tool{strip} 一并带走 \\
\Win{} & 表式方案把帧描述与语言数据放进 \Sect{pdata} 与 \Sect{xdata}；不同实现的默认策略
并不一致 & \Tool{dumpbin} 的段表与头信息 \\
\Mac{} & \Tp{\_\_unwind\_info} 用紧凑编码描述帧，语言相关数据可落在
\Tp{\_\_gcc\_except\_tab} & \Tool{otool} 的段与节列表 \\
\bottomrule
\end{tabular}
\caption{展开信息在三种目标文件格式中的位置（角色对照；节名与策略以所用实现文档为准）}
\label{tab:cl13-unwind}
\end{table}

补一句归属关系：这些表与代码段的对应关系必须在链接期保持有效。跨对象合并代码、把函数搬进
另一个段、开启代码剪裁或把模块改成静态库，展开表都要被相应改写。改写在“半边生效”时的症状
是抛出直接终止而不是被捕获，而不是普通的访问违例；诊断面属于第 \ref{ch:cl07-object} 与第
\ref{ch:cl08-static-link} 章，本册的端到端走查见第 \ref{ch:cl15-endtoend} 章。

\subsection{跨镜像抛出为什么要求两侧互相识别}\label{sec:cl13-cross-image}

一次抛出的完整链条跨越四个所有者：抛出方的代码、异常对象的分配与销毁约定、展开器、
捕获方的类型判断。链条上任一环版本不一致都会出问题：

\begin{itemize}
\item 异常对象由抛出方的运行时分配、往往由捕获方一侧销毁，两侧若连着不同的 \CL{} 运行时
副本，堆就不属于同一个分配器，销毁时的断言或崩溃几乎必然出现；
\item 类型匹配靠比较类型信息记录。两族做法都有：比地址（快，但要求同一类型在所有镜像里
合并成一份）与比名字串（能容忍多副本，代价是解码与比对）。若一侧用地址比较而镜像里存在
两份同名的类型记录，捕获就会失败；
\item 抛出与重抛的运行时入口必须落到同一个实现，否则登记与查询的语义错位；
\item \Win{} 上常见故障可以归结为一条机制：多个动态运行库副本同时驻留，抛出方的展开表
查不到、或查到后语言相关数据里的类型串与捕获方不同族，于是异常穿过若干帧却无人认领，
最后由终止钩子收场。表面现象是“明明有 catch 却直接退出”，与代码逻辑无关。
\end{itemize}

\begin{note}[title={展开信息也是一种被链接决定的产物}]
异常能否工作，取决于表的地址范围与最终镜像里的代码地址对得上。因此把“异常行为随构建开关
变化”当成编译器 bug 之前，先看产物里展开相关的段还在不在、覆盖范围有没有被裁窄。
这一步和第 \ref{ch:cl07-object} 章的段审计是同一次操作。
\end{note}

\section{把 ABI 冻住的四条工程手段}\label{sec:cl13-stable}

\subsection{把接口退回 \CL{} 层}\label{sec:cl13-cboundary}

只用确证宽度的整数类型、显式的指针与函数指针、外加一个用户数据指针，这套组合在两族实现下
都能得到稳定的符号与布局。第 \ref{ch:cl11-layout} 章讲过如何探测布局，这里的要点是边界上
不要出现依赖实现的成员：宽度和符号都不确定的基本类型、默认底层类型的枚举、以及随实现变化
的宽字符类型。

\begin{lstlisting}[style=cstyle, caption={插件边界的最小 ABI 形状：句柄、版本与函数指针表}]
#include <stdint.h>

/* 不透明句柄：它的宽度就是指针宽度 */
typedef struct PluginHost PluginHost;

typedef struct PluginVTable {
    uint32_t abi_major;        /* 不兼容改动 */
    uint32_t abi_minor;        /* 只往末尾追加 */
    int32_t  (*init)(PluginHost* h, void* user);
    void     (*shutdown)(PluginHost* h, void* user);
    /* 新成员只能追加在这里，永不插中间 */
} PluginVTable;

/* 以 C 语言链接：符号是不修饰的裸名，两族实现都能解析 */
PluginVTable* plugin_entry(void);
\end{lstlisting}

这张表能成立的前提，是宿主只按已知的偏移读取前若干项，且不假设结构体总宽度。把虚函数放进
同一张表的那一刻，本章前面讲的槽位序号规则就会接管，末尾追加不再是无害操作。

\subsection{只暴露不透明句柄}\label{sec:cl13-opaque}

把实现类藏进创建函数返回的句柄之后，宿主侧既看不到成员也看不到虚表，类的布局变化就不再
影响边界。这条思路的价值不在于封装漂亮，而在于它把布局问题从二进制层搬回了源码层：改私有
成员只需重新编译实现文件。风格层面的命名与写法讨论属于代码风格分册，本章只给机制理由：
只要类型布局跨过边界，任何成员增删都是 ABI 破坏。

\subsection{显式版本化接口}\label{sec:cl13-versioning}

比“加字段”可靠的是“换一个接口”。给接口一个标识符（编号或字符串），宿主按自己认识的最大
版本向插件索要对应的那张表；索要不到可降级的版本就明确拒绝加载。版本必须出现在符号或查询
入口里，而不是只写在注释与文档里，因为只有前者能在构建产物上被核对。

\subsection{符号版本与可见性标签}\label{sec:cl13-symversion}

\Elf{} 支持给动态符号打版本标签：链接脚本里以版本节点分组，导出符号在 \Tool{objdump}
\Sw{-T} 下带上版本串，同一库可以让多套签名共存。\Macho{} 侧的相近机制是两级命名空间，
符号解析时记录它来自哪个镜像，天然抑制了跨镜像抢占；\PeFmt{} 侧没有同一形状的标签机制，
工程上靠导出清单、序号与并排运行的运行库目录达到近似目的。语言层还有一类可见性属性可以把
符号从导出集合里剔出去，写法与默认行为随实现不同，以该实现的文档为准；机制目的只有一句：
不该被看见的符号就不要出现在动态符号表里。

\subsection{不要跨边界传标准库对象与分配器}\label{sec:cl13-no-std-across}

标准库容器与字符串的布局、以及它们的模板实例化所依赖的内部分配策略，都是实现私有的；
两侧版本不同就可能一个在扩容、另一个按旧布局读。同理，一侧分配的内存被另一侧释放，要求
两侧共用同一个分配器。这两条与第 \ref{ch:cl09-dynamic} 章的加载期符号抢占、第
\ref{ch:cl11-layout} 章的探测手段是同一个问题的三个侧面，边界上传裸指针加长度、
或上传一个由被分配方自己提供的销毁函数，才是可核对的做法。

\begin{bestpractice}[title={把对外符号集合当作契约来测}]
把“应当导出的符号列表”写成版本库里的一个文件，构建末尾拿它与实际导出集合做双向差异：
多出来的算泄漏，缺失的算破坏。这条门禁比任何评审都可靠，因为它比较的是产物而不是意图。
列表本身用解码后的源码形式保存，避免把修饰名固化成第二份契约。
\end{bestpractice}

\begin{abibox}[title={插件末尾加了一个虚函数，旧宿主崩在派发上}]
现象：插件重新编译后载入旧宿主，调用无关函数一切正常，第一次经过新虚函数的派发就崩溃。
根因不在新增的那个函数本身：虚函数的地址是按槽位序号取的，而序号由类里虚函数出现的顺序
决定。在类中间（或在被继承的基类里）插入一个虚函数，会把其后所有槽位整体后移；宿主仍按
旧序号发指令，插件表里同一序号已经是另一个地址，跳过去就是执行错误目标。
为什么“末尾追加”对 \CL{} 层的函数指针表安全、对含虚函数的类不安全：前者的宿主只按已知
偏移读取，尾部追加的项没人读，序号不存在于二进制里；后者的序号写在宿主已经编译好的调用点
中，新增一个参与派发的成员会改变整个序号空间的排布。即便新虚函数恰好排在末尾、宿主恰好不
调用它，是否安全也依赖具体实现的排序与继承细节，不是可靠承诺，不要拿它当设计依据。
可验证的修法有四：一，把新能力放进一张全新的函数指针表，由版本化入口索要，永不动已发布的
类；二，接口类一旦发布就视为不可扩展，任何虚函数改动都换一个新接口标识；三，构建时把两侧
虚表符号解码后对比序号序列，差异纳入门禁；四，宿主在加载插件时先校验 ABI 主次版本，
版本不符直接拒绝，而不是继续调用。
\end{abibox}

\section{验证路线：用符号视图核对契约}\label{sec:cl13-verify}

四步走完，一步一类命令：核对导出集合、检查应当私有的符号是否泄漏、确认展开信息仍在、
把两侧修饰名解码后比签名。工具名确证可用，输出数值与格式随版本变动，下面只是形状示意。

\begin{lstlisting}[style=shellstyle, caption={符号视图核对（示意）}]
# 1) 对外承诺看动态符号表，静态符号表只是编译残留
nm -D --defined-only libplugin.so | sort > out/exports.actual
# 2) 展开信息还在不在
readelf -SW libplugin.so | grep -E 'eh_frame|except'
# 3) 应当私有的符号有没有泄漏
nm -D --defined-only libplugin.so | c++filt | grep 'detail::'
# 4) 两侧签名比较：先解码再比，别比修饰名原文
nm -C -D --undefined-only host.bin | sort > out/host.need
dumpbin /EXPORTS plugin.dll > out/plugin.exports
llvm-cxxfilt < out/host.need > out/host.demangled
\end{lstlisting}

最后一步值得多说一句：把两侧符号都还原成源码形式再逐字段对齐，能提前发现三类分歧——
参数多了 const 限定、返回类型被改动、模板实参不同导致两个不同实例。直接比修饰名原文只能
得到“不一样”，说不出哪里不一样，也无法在两族方案之间做任何比较。反过来，若某条链路的失败
只在一侧出现，请先确认它是不是压根没导出、被裁掉、还是名字被两侧实现修饰成了两个完全不同的
实体，这三种情况在符号视图里长得就不一样。

\begin{guideline}[title={本章的四条可执行纪律}]
一，对外边界只用不修饰的名字与确证宽度的类型。二，已发布的类不再新增参与派发的成员。
三，任何跨镜像的对象、异常与内存所有权都要有明确的分配者与销毁者。四，把导出集合与解码后的
签名纳入构建门禁，让 ABI 破坏在链接期就被发现，而不是在客户机上。
\end{guideline}

\part{工具链与端到端}
\chapter{编辑器与 IDE 如何消费编译器的产物}\label{ch:cl14-ide}

编辑器与集成开发环境并不是另一套语言实现，它们是编译流水线的一组读者。本章讨论的对象是
“读者要读什么”：诊断、真实编译参数、前端分析结果、符号与调试信息、构建系统的调度决定。
语言服务协议本身的报文与实现细节属于另一分册，这里只谈被消费的产物；命名、格式化与提示
风格属于代码风格分册。

口径沿用全书：本章不给出任何实测数值，工具输出一律标示意；凡涉及具体开关，都以两族主流
工具长期公开的那几个为准，其余写成角色描述。

\section{立论：只允许存在一份前端}\label{sec:cl14-one-frontend}

集成开发环境要显示错误波浪线，就必须解析源码、展开宏、按包含路径找头文件、推断类型；
这正是编译器前端已经完成的工作。如果环境自己再实现一遍，就出现两个前端，而两个前端不可能
长期保持同结论：宏配置不同步、标准档不同步、包含路径不同步、平台三元组不同步，任何一项都
足以让两边对同一段代码给出不同判断。于是“编辑器说不报错、编译器报错”这类分歧不是偶发故障，
而是架构选择的必然结果。

可行的方向只有两条，本章后面逐一展开：其一，复用同一个前端的分析结果，让编辑器与环境拿到
的诊断就是编译器那一份；其二，让环境读到与构建完全一致的编译参数，从而把自己的猜测空间
压到零。两条都做不到时，至少要指定其中一条为唯一真源，否则误报无法根治，只能逐条关规则、
加豁免，最后环境退化成一个什么都不报的文本窗口。

\begin{flow}[title={编辑期消费的产物链}]
构建配置决定编译参数；编译参数决定预处理结果与类型布局；前端在同一批参数上产出诊断与符号
索引；构建产物留下符号表与调试信息。环境的每一项功能都必须能倒着挂回这条链的某一环，
挂不回去的功能就是在自己重演一遍语言。
\end{flow}

\begin{guideline}[title={判定归属的一句话}]
同一段代码在编辑器与编译器下结论不同，先当作构建配置缺陷来查，不要当作编辑器故障来查。
两边读的是同一套源码，差异只可能来自参数与前端版本这两处。
\end{guideline}

\section{机器可读的诊断}\label{sec:cl14-diagnostics}

\subsection{诊断行的稳定形状}\label{sec:cl14-diag-shape}

绝大多数环境解析的第一件产物，仍是编译器打出的那几行文本。它的形状长期稳定：位置、
严重级、消息，某些实现还在末尾附上诊断编号或触发它的开关名。

\begin{lstlisting}[style=outstyle, caption={诊断行的两种形状（示意，非本仓库实测）}]
# 人读形状：位置 + 严重级 + 消息 + 归因提示
src/parser.cpp:117:24: error: no member named 'value'
  in 'std::optional<int>'
src/parser.cpp:118:9: warning: unused variable 'rc'
  [-Wunused-variable]
include/tree.h:42:1: note: previously declared here
# 结构化形状：同一批信息，字段名固定，可携带条目编号与范围
# 常见字段：文件 / 行 / 列 / 严重级 / 消息 / 诊断码 / 范围 / 子诊断链
\end{lstlisting}

两者的差别不在好不好看，而在能不能被程序稳定消费。仅人读的格式换版本就可能改分隔符、
改换行折法、插入颜色控制字符；结构化输出把字段名固定下来，才允许环境做聚合、去重、跨文件
归并与门禁判断。较新版本的工具可以在人读与机器可读两种输出之间切换，具体取值以所用版本的
帮助为准。

\subsection{为什么必须有列号与范围}\label{sec:cl14-column-range}

只有行号时，波浪线只能画在行首，环境无法高亮真正的表达式，也无法把“参数不匹配”落到具体
那个实参上。列号让位置可锚定，起始列与结束列（或显式的范围标注）让锚定有宽度；缺少范围时，
环境只能自己重新解析一次来决定高亮到哪里，而这一次重新解析恰恰破坏了只允许一份前端的前提。

带范围的诊断还能把多段提示串成一条链：主诊断指出错误，附注指出“此处定义”“此前声明”
“实参在此匹配”，环境据此提供跳转到相关位置。没有这条链，同一个错误的十几种表现就只能靠
用户自己找。行列信息是否进入产物，也受调试信息开关影响；较新的实现里列信息可以单独决定是否
记录，以对应版本文档为准。

\subsection{编号与归因为什么是自动化的前提}\label{sec:cl14-diag-code}

一条诊断若只说“这里错了”，环境最多把它显示出来；若它带着来源（哪条规则、哪个开关把它
打开的），环境才能做三件事：按规则聚合、按文件抑制、把可自动修复的那部分挑出来。人读格式
末尾那个方括号里的开关名，作用就是把“该改代码”与“该改构建配置”分开。缺少归因时，
用户只能去猜是哪条开关触发的，最后常见结果是把整项检查关掉。

\begin{compare}[title={人读格式与结构化输出}]
人读：面向终端，允许颜色、换行折叠、上下文引用行与装饰性提示；优点是当场可读，缺点是字段
边界没有承诺。结构化：面向程序，字段名与层级属于接口的一部分，可携带条目编号、范围、修复
建议与来源链；缺点是要显式打开，且不同实现提供的字段集合不完全一致。
稳妥做法：交互界面读结构化，日志与归档同时保留一份人读文本，两份来自同一次运行而不是两次
编译。
\end{compare}

\section{真实编译参数是前提：编译数据库}\label{sec:cl14-compdb}

\subsection{它是什么}\label{sec:cl14-compdb-what}

编译数据库是一份把构建系统里每个目标的编译命令原样抄录下来的清单，业界通行的形态是
\Sym{compile{\Bk}\_commands{\Bk}.json}：一条记录对应一个翻译单元，通常含工作目录、源文件、
完整命令行（或参数数组），有的还记录目标产物。清单里真正要紧的信息是定义、标准档、包含路径
与目标三元组，四者共同决定了预处理结果与类型布局。有了它，环境就能以与构建完全相同的参数
做一次解析，不必猜。

\begin{lstlisting}[style=cmakestyle, caption={导出编译数据库：变量与目标属性两条入口}]
# 只对支持该特性的生成器有效
set(CMAKE_EXPORT_COMPILE_COMMANDS ON)
add_library(parser STATIC parser.cpp scan.cpp)
target_compile_definitions(parser PRIVATE PARSER_STRICT=1)
target_compile_features(parser PRIVATE cxx_std_20)
target_include_directories(parser PUBLIC include)
set_target_properties(parser PROPERTIES
  EXPORT_COMPILE_COMMANDS ON)
\end{lstlisting}

入口有两条：全局变量把数据库写到构建目录，目标属性把单个目标的命令补进去。两条都要求
构建配置是真实生成的，而不是手工拼出来的；具体取值与适用生成器以 \Tool{CMake} 的文档为准。

\begin{lstlisting}[style=shellstyle, caption={取数据库并据此做一次与构建等价的解析（示意）}]
cmake -G Ninja -S . -B build \
  -DCMAKE_EXPORT_COMPILE_COMMANDS=ON
ninja -C build -t compdb > build/compile_commands.json
ln -sf build/compile_commands.json .
jq -r '.[].file' build/compile_commands.json | wc -l
# 覆盖面：数据库里有几个翻译单元，与实际参与构建的单元数是否相符
\end{lstlisting}

\Tool{Ninja} 的 \Sw{-t compdb} 从构建图里导出条目；对不导出该文件的生成器，也可以用拦截
编译命令的第三方工具生成同一份清单。无论走哪条路，判据只有一个：数据库里的命令行与实际
构建用到的命令行等价。等价性要靠比对建立，而不是靠看起来差不多。

\subsection{数据库里没有的字段一样致命}\label{sec:cl14-compdb-fields}

编译命令只是参数的载体。环境要还原出可解析的一次翻译，还得知道命令的工作目录（相对包含
路径才能被正确解释）、编译器可执行文件本身（不同实现自带不同的标准库头路径），以及响应文件
的展开结果。数据库忠实地抄下了命令行，但抄的是短形式还是展开后的长形式，取决于构建工具的
生成方式；环境若不做同一套展开，参数集合就会差一批。

\subsection{缺了它会发生什么}\label{sec:cl14-compdb-missing}

没有数据库时环境只能猜：猜包含路径、猜标准档、猜定义。猜错的直接后果是三类噪声——找不到
头文件的连锁错误、把新标准的语法当错误、宏关掉之后本不该参与编译的分支被解析出来。更隐蔽
的是目标三元组与布局：环境按本机默认目标解析，交叉编译项目的类型宽度就对不上，于是所有
依赖宽度的诊断集体失真（机理见第 \ref{ch:cl11-layout} 章）；宏配置不一致则让环境解析到与
构建不同的代码分支（见第 \ref{ch:cl02-preprocessing} 章）。

\begin{warning}[title={环境自带的默认参数不是构建参数}]
把包含路径与定义手填进环境的设置项，等于维护第二份构建配置。它会随时间失真，而且失真的
时候没有任何报错：环境只是安静地按旧参数解析，用户看到的现象是突然多出来或突然消失的一批
提示。要么从数据库生成，要么把这份设置纳入同一条门禁，与数据库逐字段比对。
\end{warning}

\section{复用前端的两条路径}\label{sec:cl14-two-roads}

\subsection{把前端做成库}\label{sec:cl14-frontend-lib}

第一条路径是让语言服务进程链接编译器前端本身，用它解析、用它出诊断，环境只负责展示与编辑
辅助。收益是同一份解析器、同一份诊断；代价是环境的生命周期与前端版本绑死——前端升级即环境
升级，并且需要为只做到语义分析为止的路径提供入口（多数实现都留了这个入口，因为它同时服务
于静态检查工具）。这条路径能回答的是所有“这段代码在当前参数下合不合法、这个表达式是什么
类型、这个名字声明在哪里”的问题。

\subsection{从构建产物反查}\label{sec:cl14-from-artifacts}

第二条路径是不再解析源码，而是读构建留下的产物：符号表告诉环境哪些实体真实存在（第
\ref{ch:cl07-object} 章）；调试信息把类型、成员布局与源码位置对应起来；\Win{} 一侧的符号
文件与 \Mac{} 一侧需要先归并再使用的调试信息，各有一套收集流程。这条路径能回答的是“最终
镜像里到底有什么”：模板是否真的被实例化、内联之后符号是否消失、这段源码对应哪个地址。

它的边界同样清楚：产物只反映编译完的东西。未参与构建的代码、被条件编译关掉的分支、尚未
保存的编辑内容都不在其中；被裁掉的符号也在其中缺席，于是跳转指向空处。

\begin{table}
\centering
\begin{tabular}{@{}>{\raggedright\arraybackslash}p{0.26\textwidth} >{\raggedright\arraybackslash}p{0.26\textwidth} >{\raggedright\arraybackslash}p{0.26\textwidth}@{}}
\toprule
环境要回答的问题 & 由哪个产物回答 & 何时失效 \\
\midrule
这行代码在当前参数下是否合法 & 前端解析结果，即同一份诊断 & 参数与数据库不同步、前端版本落后 \\
这个标识符声明在哪 & 前端的符号索引 & 索引未覆盖该目录，或新文件未入库 \\
最终镜像里有没有这个实体 & 符号表与导出清单 & 符号被裁掉、被局部化或只存在于内联中 \\
成员在对象里的偏移 & 调试信息中的类型记录 & 未开启调试信息，或优化后描述不完整 \\
跨实现调用是否兼容 & 解码后的修饰名，见第 \ref{ch:cl13-cxx-abi} 章 & 只解码了一侧，或两侧不同族 \\
构建实际用了哪些定义 & 数据库条目与响应文件展开结果 & 数据库过期，或来自另一套配置 \\
\bottomrule
\end{tabular}
\caption{环境提出的问题与回答它的产物（本册视角：问题到产物的映射）}
\label{tab:cl14-products}
\end{table}

两条路径不是二选一。前端给编辑期的即时反馈，产物给构建后的事实核对；两者的差异本身就是
排障信号——前端说存在而产物里没有，多半是没参与构建、被裁或只在内联中存在；产物里有而前端
找不到，多半是索引没覆盖或数据库缺了条目。

\section{增量、索引与背景任务}\label{sec:cl14-indexing}

\subsection{可重放的编译前缀}\label{sec:cl14-pch}

编辑期的响应预算是几十毫秒，而一次完整解析要把项目可见的头文件全部读一遍。解决办法是把
这部分工作冻结成可重放的前缀：先用同一组参数解析公共头，把解析结果写成中间文件，后续每次
解析直接从中间文件恢复，只有增量部分重新走一遍。这就是预编译头一类机制在语言服务里的角色。
它在编辑期的价值不是编译更快，而是让每次击键都能触发一次完整的语义解析，而不是退化成只看
括号匹配的文本检查。

代价是前缀与参数强绑定：定义、标准档、包含路径任何一项变化都必须重做前缀，否则恢复出来的
符号与当前源码不一致，表现为跳到一个已经不存在的声明。因此前缀的键必须是参数指纹，而不能
只是文件名；用文件名做键的实现在切换构建配置后必然给出陈旧结果。

\subsection{背景索引与前台响应的取舍}\label{sec:cl14-index}

跳转、补全与引用查找需要全工程符号库，这只能由后台索引提供。取舍很清楚：后台索引占用资源，
结果可能滞后；前台只做局部解析，即时但视野只到当前翻译单元。合理的设计让前台用局部结果先
给出波浪线，引用查找与跨文件跳转用后台结果，并把索引可能过时这件事显式标出来。反过来，
让每次保存都触发全工程重解析，等于放弃了增量，通常会在几分钟后被用户关掉。

索引的正确性判据不在索引自身，而在它是否与被消费的产物一致：索引里某个类有四个虚函数，
而最终镜像里对应的槽位数与之不符，就要怀疑索引跨越了一次接口变更（回到第
\ref{ch:cl13-cxx-abi} 章讲的槽位序号规则）。

\subsection{与调试器的接缝}\label{sec:cl14-debug}

调试器是这条流水线的另一个读者：它按地址反查源码位置，靠的是产物里的调试信息，而不是环境
的索引。断点打到模板实例化后被内联吞掉的位置、断点命中却显示不出成员值（布局描述与运行时
不一致）这两类现象，属于同一个消费关系出了问题。环境若要显示此处的值，应当以调试信息为准，
而不是以自己的解析结果为准。

\section{与构建系统的集成}\label{sec:cl14-buildsys}

\subsection{调度交给构建工具}\label{sec:cl14-scheduling}

环境不应该自己维护一套编译循环。批量构建的依赖图、并行度、失败传播与增量判定，构建工具
已经算过一遍；环境重复实现一遍，就会得到与真实构建不同的结果。合理分工是：环境负责单文件的
即时解析，构建工具负责整个工程的一次构建，两者共用同一份编译数据库。

\begin{lstlisting}[style=shellstyle, caption={编辑期与构建期各走各的入口（示意）}]
# 编辑期：只到语义分析为止，参数从数据库原样取
clang++ -fsyntax-only -std=c++20 \
  @build/parser.rsp src/parser.cpp
# 构建期：依赖清单交给构建工具，增量判定不靠时间戳猜
g++ -c -MMD -MP -MF build/parser.d \
  src/parser.cpp -o build/parser.o
ninja -C build -j 8
\end{lstlisting}

缓存工具（\Tool{ccache} 一类）复用的也是编译命令的等价性，判据同样是参数集合；参数指纹不同步
时，缓存命中率下降与索引结果失真会同时出现，那是同一个根因的两种症状。

\subsection{响应文件与命令行长度上限}\label{sec:cl14-response}

大项目的编译命令动辄几百条包含路径与定义，会顶到操作系统的命令行长度上限；\Win{} 一侧
单条命令行的字符数限制是常见瓶颈。两族工具都支持把参数写进文件再交给编译器，常见写法是把
文件名以 at 号前缀传给命令行。这对环境有两个含义：一是从数据库读到的命令行可能只是指向
响应文件的短命令，展开之后才是真实参数；二是环境若自己拼接命令行，必须实现同一套展开规则，
否则参数集合会不一致。

\subsection{并行时诊断的归并}\label{sec:cl14-parallel}

并行编译会让不同翻译单元的诊断交错出现。归并的正确做法是按翻译单元分组，而不是按到达顺序
直出：结构化输出天然带文件字段，按字段分桶即可；仅人读的文本则要靠上下文行前缀重建归属。
同一份源码被多个目标以不同参数编译（不同定义、不同标准档）时，会出现同一行两条结论不同的
诊断——这不是重复噪声，而是参数分裂的信号，应当反过来指向数据库的覆盖问题：数据库里是否给
同一源文件登记了多条命令，环境是否只取了第一条。

\subsection{增量失效的常见原因}\label{sec:cl14-invalidate}

四类最常见：宏配置变化（定义集合一改，条件编译分支整片换掉，回指第
\ref{ch:cl02-preprocessing} 章）；头文件改动面过大（被大量翻译单元包含的头触发全量重做，
前缀与索引同时作废）；参数与文件集变化（数据库重建后旧命令的指纹不再匹配）；以及环境自己的
索引没跟上构建（产物已更新而索引仍是旧的）。前三类由构建系统判定，第四类才是环境自身的
问题。把四类的失效时间记进同一张表，是排查为什么突然全红的最短路径。

\begin{abibox}[title={同一段代码，IDE 与编译器结论不同}]
现象：环境里满屏提示，命令行构建却正常通过；或者反过来，环境一片安静，构建报错。
原因集中在四类不同步：一，标准档——环境按较旧的语言档解析，新语法被判成错误；二，包含
路径——构建靠生成器与响应文件展开，环境只用自己设置里那几条；三，平台三元组——交叉构建的
目标宽度与本机默认目标不同，依赖宽度的诊断集体失真；四，宏定义——项目自定义的开关决定了
一批分支是否参与解析，环境读不到时按默认值猜。
修法只有一条：以编译数据库为唯一真源，让环境、语言服务、前缀与索引全部从它取参数。
验证步骤（描述性）：一，重新生成数据库并确认报错的源文件确实在其中；二，取该文件那一条
命令原样手工执行一次，确认复现的是构建侧的结论；三，逐项比对环境与构建实际使用的标准档、
定义集合、包含路径与目标三元组，差异即病因；四，参数变化后让前缀与索引重建，确认跳向不存在
声明的现象消失。四类都排除后仍有分歧，才轮到怀疑前端版本差异。
\end{abibox}

\begin{bestpractice}[title={把环境报错当构建配置缺陷排查的三步}]
第一步，确认参数来源：环境是否真的读到了编译数据库，条目是否覆盖报错的那个文件，数据库的
时间戳是否晚于最近一次构建配置改动。第二步，复现构建侧命令：从数据库里取出该文件那一条
命令，原样执行一次；若结果与构建一致而与环境不一致，问题在参数传递而不在代码。第三步，
对齐四件套并重建缓存：标准档、定义集合、包含路径与目标三元组逐项比对，任一改动后重建前缀
与索引，再观察诊断是否收敛。三步走完仍有分歧，才去看前端版本与实现差异。
\end{bestpractice}

\begin{note}[title={本册与语言服务分册的边界}]
协议里有哪些请求、报文如何编码、增量文本同步怎么做、服务端如何管理会话，这些属于另一分册。
本册只回答一个问题：编辑器与环境消费的是编译流水线里的哪些产物，以及这些产物在什么情况下
会失真。代码风格层的命名与格式化属于代码风格分册。反过来，凡是环境自己重新实现一遍语言的
功能，都不在本册认定的正当范围之内。
\end{note}

\section{把这条链路走通}\label{sec:cl14-endtoend}

把本章与前几章串起来，就得到一份端到端检查表。每一行都问“谁是真源”，因为分歧从来不发生在
代码里，而发生在两份真相之间。

\begin{longtable}{@{}>{\raggedright\arraybackslash}p{0.12\textwidth} >{\raggedright\arraybackslash}p{0.22\textwidth} >{\raggedright\arraybackslash}p{0.22\textwidth} >{\raggedright\arraybackslash}p{0.22\textwidth}@{}}
\caption{环境侧的端到端检查表：产物、真源与失真信号}\label{tab:cl14-checklist}\\
\toprule
环节 & 被消费的产物 & 唯一真源 & 失真时的典型信号 \\
\midrule
\endfirsthead
\toprule
环节 & 被消费的产物 & 唯一真源 & 失真时的典型信号 \\
\midrule
\endhead
\bottomrule
\endfoot
参数 & 编译数据库条目 & 构建系统的生成结果 & 环境多出一批找不到头文件的错误 \\
诊断 & 结构化输出通道 & 编译器前端 & 同一行两侧结论不同 \\
前缀 & 可重放的解析结果 & 参数指纹 & 跳向已删除的声明 \\
索引 & 全工程符号库 & 与构建相同的文件集合 & 跳转指向空处或出现重复项 \\
事实核对 & 符号表与调试信息 & 最终镜像 & 实体有无与环境判断不符 \\
\end{longtable}

第 \ref{ch:cl15-endtoend} 章把这条链路做成一次完整走查，附录 \ref{ap:clA-cli} 的命令清单
可以作为它的执行手册，附录 \ref{ap:clB-abi} 提供符号与修饰层面的对照。

\chapter{端到端演练：从源码到镜像}\label{ch:cl15-endtoend}

前十四章把流水线拆成零件：预处理、名字查找、中间表示、优化、代码生成、目标文件、
链接、加载、约定、布局、装饰、工作台。本章把它们装回去。我们只搭一个
小得不能再小的工程，但它必须有本章开头承诺的四样东西：两个源文件组成的可执行
目标、一个静态库、一个共享库、一条冻结成 \CL{} 层的跨镜像接口。然后从
构建描述出发，一路走到镜像被装载执行，每一步都把产物拿出来看一次。

演练的目的不是得到一个好的设计，而是建立一条排查路径。走完一遍之后，你应当能
回答三个问题：某条符号是在哪一步产生的、它在哪一步被解析、它的地址在哪一步被
写死。也要能回答第四个：当这两个问题在不同平台上有不同答案时，差异是从哪一行
构建描述或哪一次工具调用开始的。

\section{演练设定：一个够小但五脏俱全的工程}\label{sec:cl15-setup}

文件布局按最小工程摆放，接口层刻意薄到只有一层：

\begin{itemize}
\item \Tp{src/main.c} 与 \Tp{src/transform.c}：可执行目标的两个翻译单元。
      它们的分工不是装饰：\Tp{transform.c} 里放对 \CL{} 入口的调用，
      \Tp{main.c} 只负责调用 \Fn{transform\_run} 一次并返回。
\item \Tp{lib/codec.c}：编解码实现，编成静态库 \Tp{libcodec.a}。
\item \Tp{lib/filter.c}：过滤层实现，编成共享库 \Tp{libfilter}。
\item \Tp{include/codec.h} 与 \Tp{include/filter.h}：对外接口。头文件里的每个
      入口都在 \CL{} 语法内声明，因此 \cpp{} 单元包含它们时不会重新装饰。
\end{itemize}

构建描述只写必要信息。下面这段是本章反复引用的那份，风格 \Tp{cmakestyle}，
只用签名稳定的命令与变量：

\begin{lstlisting}[style=cmakestyle, caption={演练工程唯一的构建描述：两个源文件的可执行目标、一个静态库、一个共享库}]
cmake_minimum_required(VERSION 3.16)
project(cl15_e2e LANGUAGES C)

add_library(codec STATIC lib/codec.c)
add_library(filter SHARED lib/filter.c)
add_executable(app src/main.c src/transform.c)

target_include_directories(codec PUBLIC include)
target_include_directories(filter PUBLIC include)

set_target_properties(codec PROPERTIES POSITION_INDEPENDENT_CODE ON)
set_target_properties(filter PROPERTIES POSITION_INDEPENDENT_CODE ON)

# 共享库吃掉静态库：这是本章观察点的主要来源
target_link_libraries(filter PUBLIC codec)
target_link_libraries(app PRIVATE filter)
\end{lstlisting}

这段描述里已经埋了三个伏笔。\Tp{STATIC} 与 \Tp{SHARED} 决定第 7、8、9 阶段
怎么走；\Tp{POSITION\_INDEPENDENT{\Bk}\_CODE} 决定同一批目标文件能不能同时喂给
可执行文件和共享库；\Tp{PUBLIC} 与 \Tp{PRIVATE} 决定 \Sym{codec\_compute\_checksum}
最终出现在哪几个产物里。三者都是构建层面的选择，不是语言承诺。

\begin{abibox}[title={为什么接口层要冻结成 \CL{} 而非 \cpp{}}]
镜像之间只交换 \CL{} 层的东西：整数、指针、定长结构体、显式的长度与容量参数。
\cpp{} 的类布局、虚表指针位置、异常传播机制、模板实例的归属都随版本与
优化等级改变，把任一种写进跨镜像接口就等于把决定权交给了对侧的编译器。
冻结的做法是：头文件里的入口在 \CL{} 语法内声明，\cpp{} 侧通过一层极薄的包装调用
它们；两侧再各自用 \Fn{static\_assert} 把 \Fn{sizeof} 与 \Fn{alignof} 钉住。
这样 \Win{} 侧的 \Mx{MSVC C++} 与 \Nix{} 侧的 \Mx{GCC C++} 才可能在同一个
\Sym{filter\_open} 上达成共识，理由见第~\ref{ch:cl13-cxx-abi}章。
\end{abibox}

\begin{warning}[title={位置无关代码不是可选项，是链接期的入场券}]
共享库的目标文件必须能在任意加载基址下工作，因此编译期就得关掉绝对地址。
若把同一份 \Tp{codec.o} 既喂给共享库又喂给可执行目标，而编译时没开位置无关，
失败会在最后一个阶段才出现，报错内容却是链接器的地址溢出或重定位类型不支持。
全局写法 \Tp{CMAKE\_POSITION\_INDEPENDENT{\Bk}\_CODE} 与上面那段里的逐目标写法作用范围
不同：前者覆盖整个工程，后者只覆盖点名的那一个目标，两者语义相同。
这类症状的根源在第一个阶段，回指第~\ref{ch:cl06-codegen}章的地址模式选择。
\end{warning}

\section{七个阶段逐个走}\label{sec:cl15-stages}

先把整条链摆成一张表。表里的工具按平台通用写法给出，具体平台的等价物在
第~\ref{sec:cl15-platform}节之外不再重复；阶段是概念阶段，不等于工具调用次数。

\begin{flow}[title={本章要走完的七步}]
预处理与头闭包 $\rightarrow$ 前端分析 $\rightarrow$ 中间表示与优化 $\rightarrow$
指令选择与寄存器分配 $\rightarrow$ 目标文件与符号 $\rightarrow$
静态链接与折叠 $\rightarrow$ 动态库与加载配置。
\end{flow}

\begin{longtable}{@{}>{\raggedright\arraybackslash}p{0.16\textwidth} >{\raggedright\arraybackslash}p{0.20\textwidth} >{\raggedright\arraybackslash}p{0.20\textwidth} >{\raggedright\arraybackslash}p{0.22\textwidth}@{}}
\caption{七步演练：产物、查看手段与对应章节}\label{tab:cl15-stages}\\
\toprule
阶段 & 产物 & 看它的工具 & 本册解释 \\
\midrule
\endfirsthead
\toprule
阶段 & 产物 & 看它的工具 & 本册解释 \\
\midrule
\endhead
预处理与头闭包 & 单个展开后的源文件 & \Tool{cpp} 或 \Sw{-E} 模式 &
第~\ref{ch:cl02-preprocessing}章 \\
前端分析 & 抽象语法树与符号表 & 只打印树或语法树的模式 &
第~\ref{ch:cl03-resolution}章 \\
中间表示与优化 & 一份函数级中间表示 & 只输出中间表示的开关 &
第~\ref{ch:cl04-ir}章、第~\ref{ch:cl05-passes}章 \\
指令选择与分配 & 汇编文本 & \Sw{-S} 模式加 \Tool{objdump} 反汇编 &
第~\ref{ch:cl06-codegen}章 \\
目标文件与符号 & 可重定位目标文件 & \Tool{readelf}、\Tool{nm}、\Tool{objdump} &
第~\ref{ch:cl07-object}章 \\
静态链接与折叠 & 归档文件、合并后的代码 & \Tool{ar}、\Tool{nm}、\Tool{size} &
第~\ref{ch:cl08-static-link}章 \\
动态库与加载配置 & 共享库与镜像 & \Tool{ldd}、加载命令查看工具 &
第~\ref{ch:cl09-dynamic}章 \\
\end{longtable}

逐阶段补充三件表里放不下的事。

预处理阶段的产物是头闭包的度量。把每个翻译单元展开一次，行数与文件数立刻告诉
你隐式依赖有多重；这一步看到的 \Sym{codec\_compute\_checksum} 只是文本，
还不存在于任何符号表里，细节见第~\ref{ch:cl02-preprocessing}章。

前端分析结束时的符号表是本轮演练第一个观察点。名字查找此时已经完成：能被
解析的引用变成内部符号，不能解析的留下一个未定义引用，附带它期望的签名与来源。
第~\ref{ch:cl03-resolution}章讲的就是这张表怎么建，本章只把它当作可读产物。

静态链接阶段的折叠是第二个观察点。把 \Tp{codec.a} 装进 \Tp{libfilter} 时，
链接器按引用从归档里抽出需要的成员，抽中的成员其代码进入共享库的
\Sect{text}，未被引用的成员根本不进产物。因此符号表里的条目数不等于
归档成员数；第~\ref{ch:cl08-static-link}章把这条规则拆得更细。

\begin{lstlisting}[style=outstyle, caption={同一份构建在三个阶段留下的符号线索（示意输出，非本仓库实测值）}]
# 可重定位目标文件：引用是欠账，地址是占位
main.o:  UND  filter_open
main.o:  GLOBAL filter_open  (在 filter.o 中定义为 text)
# 静态归档：成员表 + 索引，索引决定谁被抽出
libcodec.a (codec.o):  GLOBAL codec_compute_checksum
# 共享库：定义成为导出项，引用改走延迟绑定机制
libfilter: GLOBAL filter_open        UND codec_compute_checksum
\end{lstlisting}

\begin{lstlisting}[style=asmstyle, caption={链接前的调用点仍保留外部名字，链接后同一位置是相对偏移或表项（示意，非本仓库实测值）}]
    .text
filter_open:
    endbr64
    call  codec_compute_checksum    # 目标地址未知，重定位项随指令一起留下
    ret
\end{lstlisting}

这段汇编写在汇编文本层面最能说明一件事：未解析的引用不是一个错误状态，而是
目标文件的标准形态。它靠一条重定位记录把绑定责任交给下一阶段，第~\ref{ch:cl07-object}
章定义记录结构，第~\ref{ch:cl08-static-link}章定义谁来消费它。

\section{三个平台各走一遍}\label{sec:cl15-platform}

同一份构建描述在三个平台会展开成三串工具调用。下面三段流程块都是示意，
工具名是各平台的通用形态，参数刻意少写；能确定的只有阶段顺序与产物类型。
这一节的覆盖面按第 \ref{sec:cl01-targetenv} 节的口径限定：桌面三个平台各走一遍；
移动平台的 native 层走的是同一套阶段，变的是三元组与分发形态；裸机要换掉后半段。

\subsection{\Nix{}：\Tool{gcc} 族一路到底}\label{sec:cl15-nix}

归档工具用 \Tool{ar}，索引刷新用 \Tool{ranlib}；\Elf{} 侧的查看手段齐全，
\Tool{readelf} 看节与段，\Tool{nm} 看符号，\Tool{objdump} 同时看两者与反汇编，
\Tool{ldd} 看运行期的依赖闭包。

\begin{lstlisting}[style=shellstyle, caption={\Nix{} 走法：把每个中间产物都留在磁盘上（示意流程）}]
# 阶段一：预处理与头闭包
gcc -E src/main.c -o main.i
# 阶段二到四：前端、优化、指令选择，停在汇编文本
gcc -S src/main.c -o main.s
# 阶段五：汇编，得到可重定位目标文件
gcc -c lib/codec.c -o codec.o
gcc -c lib/filter.c -o filter.o
gcc -c src/main.c -o main.o
gcc -c src/transform.c -o transform.o
# 阶段六：静态库归档
ar rcs libcodec.a codec.o
ranlib libcodec.a
readelf -SW libcodec.a
nm -C libcodec.a
# 阶段七：共享库，把静态库吃进去
gcc -shared -o libfilter.so filter.o codec.o
nm -D libfilter.so
# 可执行目标只依赖共享库
gcc -o app main.o transform.o -L. -lfilter
ldd app
\end{lstlisting}

\subsection{\Win{}：\Tool{cl} 与 \Tool{link} 的分工}\label{sec:cl15-win}

\Win{} 的分工更硬：\Tool{cl} 只管到目标文件，产出 \Tp{.obj}；归档交给 \Tool{lib}，
得到 \Tp{.lib}；镜像生成交给 \Tool{link}。查看工具是 \Tool{dumpbin}，
它的开关是文档化的，能看头、符号、导出表、重定位表与依赖列表；
\Tool{editbin} 用来读改镜像头里的基址相关字段。

\begin{lstlisting}[style=shellstyle, caption={\Win{} 走法：编译、归档、链接三步各自独立（示意流程）}]
cl /c lib/codec.c lib/filter.c src/main.c src/transform.c
lib /OUT:codec.lib codec.obj
link /DLL /OUT:filter.dll filter.obj codec.lib /IMPLIB:filter.lib
cl /c src/main.c
link /OUT:app.exe main.obj transform.obj filter.lib
dumpbin /HEADERS app.exe
dumpbin /SYMBOLS codec.lib
dumpbin /EXPORTS filter.dll
dumpbin /DEPENDENTS app.exe
\end{lstlisting}

\Win{} 的调试信息与镜像天生分离：编译选项 \Sw{/Zi} 产生的信息进独立的
\Tp{.pdb}，链接选项 \Sw{/DEBUG} 只在镜像里留下指向它的记录。于是
\Tool{dumpbin} \Sw{/SYMBOLS} 看到的内容与调试器看到的类型信息不是同一份数据，
这个区别在第~\ref{ch:cl07-object}章里被拆成两小节。

\subsection{\Mac{}：\Tool{clang} 族与 \Macho{} 加载命令}\label{sec:cl15-mac}

\Mac{} 用 \Tool{clang} 与 \Tool{clang++} 编译，静态归档可用 \Tool{ar}
或 \Tool{libtool}，共享库走 \Sw{-dynamiclib} 并得到 \Tp{.dylib} 后缀。
查看工具是 \Tool{otool}：\Sw{-l} 打印加载命令，\Sw{-L} 打印依赖，
\Sw{-tV} 打印带注释的代码段；符号交给 \Tool{nm}，依赖路径的改写交给
\Tool{install\_name\_tool}。

\begin{lstlisting}[style=shellstyle, caption={\Mac{} 走法：加载命令是 \Macho{} 的第一现场（示意流程）}]
clang -E src/main.c -o main.i
clang -c lib/codec.c -o codec.o
libtool -static -o libcodec.a codec.o
clang -dynamiclib -o libfilter.dylib filter.o codec.o
nm -gU libfilter.dylib
otool -l libfilter.dylib | grep LC_
clang -o app main.o transform.o -L. -lfilter
otool -L app
dsymutil app
\end{lstlisting}

\begin{abibox}[title={三个平台把同一件事放在不同位置}]
\begin{tabular}{@{}>{\raggedright\arraybackslash}p{0.16\textwidth} >{\raggedright\arraybackslash}p{0.26\textwidth} >{\raggedright\arraybackslash}p{0.26\textwidth}@{}}
\toprule
事项 & \Nix{} / \Elf{} & \Win{} / \PeFmt{} \\
\midrule
静态归档 & \Tp{.a} 加索引 & \Tp{.lib}，与导入库同名后缀 \\
共享库 & \Tp{.so} 加 soname & \Tp{.dll} 加 \Tp{.lib} 导入库 \\
代码位置 & \Sect{text} 节 & \PeFmt{} 的节与段表 \\
调试信息 & 多数情况留在产物里 & 分离到 \Tp{.pdb} \\
\bottomrule
\end{tabular}
\end{abibox}

\Mac{} 那一列的位置由加载命令承担：段与依赖都写进 \Tool{otool} \Sw{-l} 能看到的
命令项里。三个后缀、三种元数据放法，讲的仍是同一件事，第~\ref{ch:cl07-object}章
与第~\ref{ch:cl11-layout}章分别覆盖结构与布局。

\section{观察点清单}\label{sec:cl15-watch}

\begin{guideline}[title={把问题定位到阶段，而不是定位到报错文本}]
符号有产生、解析、绑定三个时刻；重定位项有生成、消费、消失三个时刻；地址有编译期
占位、链接期固定、加载期改写三种状态。任何一条报错，先落到这九个坐标之一，
再选工具。落不到坐标的报错，才需要回到附录~\ref{ap:clA-cli}去把中间产物全部留下。
\end{guideline}

\begin{longtable}{@{}>{\raggedright\arraybackslash}p{0.24\textwidth} >{\raggedright\arraybackslash}p{0.26\textwidth} >{\raggedright\arraybackslash}p{0.28\textwidth}@{}}
\caption{本章要看清的四组坐标}\label{tab:cl15-watch}\\
\toprule
观察点 & 什么时候成立 & 怎么确认 \\
\midrule
\endfirsthead
\toprule
观察点 & 什么时候成立 & 怎么确认 \\
\midrule
\endhead
符号产生 & 前端分析完成、目标文件写出之后 & 同一阶段用 \Tool{nm} 前后各看一次 \\
符号解析 & 静态链接阶段闭合，跨镜像项推迟到加载期 & \Tool{nm} 区分定义项与未定义项 \\
重定位消失 & 静态链接闭合的同一时刻 & 链接前后分别 dump 重定位表 \\
地址固定 & 链接期写死，或加载期被整体改写 & \Win{} 看镜像基址与重定位表 \\
\bottomrule
\end{longtable}

四组坐标各补一句。

符号在哪一步产生：前端把声明变成符号表条目，汇编器把它写进目标文件的符号表。
在此之前它只是文本；在此之后它带上类型、绑定与所在节。若你在目标文件里找不到
某个入口，原因通常在第一个阶段而不是第二个阶段，因为那时它已经被优化掉了或
根本没被声明为外部可见，见第~\ref{ch:cl07-object}章。

重定位项在哪一步消失：静态链接器把项消费掉，把结果写进指令或数据，然后
\Sect{text} 不再携带它。\Win{} 的基址重定位是例外：它记录的是整段偏移的修正方法，
而不是单个符号的绑定，所以镜像里既有被消费的项也有留给加载器的项，
第~\ref{ch:cl08-static-link}章与第~\ref{ch:cl11-layout}章各讲一半。

地址什么时候固定：静态链接阶段，若产物不要求可重定位，链接器把地址写死。
什么时候又被改写：加载阶段。共享库的加载基址随地址空间布局随机化改变，
因此表项内容在运行期与文件里不一致，这是正常现象而不是损坏。第~\ref{ch:cl09-dynamic}
章讲机制，\Win{} 上的基址重定位讲代价：偏移表与基址重定位表都只在加载期生效，
把它们误当成符号表就会得出产物缺符号的错误结论。

为什么看得到符号不等于看得到调试信息：符号表回答某个地址属于哪个名字，
调试信息回答这个地址上那份对象有哪些字段、行号与变量映射。前者是链接的输入，
后者是调试器的输入；\Tool{strip} 只删前者之外的东西时，调试器仍然打不开类型。
\Win{} 因为两者天然分离，最先暴露这个区别，见第~\ref{ch:cl07-object}章。

\begin{note}[title={优化会让观察点变得不像源码}]
内联之后，源文件里的调用点在汇编里不再存在；循环被向量化之后，
\Sect{text} 里的内容与源码语句顺序无关。符号表仍然可靠，因为它按定义登记，
不按执行路径登记。用汇编文本来推断语义之前，先在一个低优化等级下看一次，
理由见第~\ref{ch:cl05-passes}章。
\end{note}

\section{故障定位路线}\label{sec:cl15-triage}

\begin{compare}[title={同一句报错的两种处理方式}]
被动处理是照着报错文本搜索关键词；主动处理是先问这条报错属于哪个阶段，
再把该阶段的产物单独打印出来。前者依赖别人遇到过同样的问题，
后者依赖本章这张分派表。分派表的每一行都对应一个真实阶段的边界，
所以它可以直接拿去用。
\end{compare}

\begin{longtable}{@{}>{\raggedright\arraybackslash}p{0.26\textwidth} >{\raggedright\arraybackslash}p{0.32\textwidth} >{\raggedright\arraybackslash}p{0.20\textwidth}@{}}
\caption{按症状分派：先怀疑什么，回到哪一章}\label{tab:cl15-triage}\\
\toprule
症状 & 首先怀疑 & 回到 \\
\midrule
\endfirsthead
\toprule
症状 & 首先怀疑 & 回到 \\
\midrule
\endhead
未定义符号 & 库的给出顺序，或者根本没给出这条依赖 &
第~\ref{ch:cl08-static-link}章 \\
符号冲突，或行为随链接顺序变化 & 多份定义构成一致性违规 &
第~\ref{ch:cl08-static-link}章 \\
运行期报找不到库 & 搜索路径与依赖链的深度 &
第~\ref{ch:cl09-dynamic}章 \\
崩在返回地址，或被调用方保存寄存器被踩 & 调用约定被破坏 &
第~\ref{ch:cl10-cc}章、第~\ref{ch:cl06-codegen}章 \\
跨镜像传对象之后崩溃 & 两侧 \Mx{ABI} 未冻结 &
第~\ref{ch:cl13-cxx-abi}章 \\
工作台与编译器的结论不一致 & 两侧编译参数不同步 &
第~\ref{ch:cl14-ide}章 \\
外部格式读写在另一台机器上错乱 & 字节序与整型宽度假设 &
第~\ref{ch:cl12-endian}章、第~\ref{ch:cl11-layout}章 \\
\bottomrule
\end{longtable}

\begin{warning}[title={最贵的两类症状都不在链接期报错}]
最后四行都属于这一类：链接器看不出调用约定被破坏，也看不出跨镜像的对象布局
两侧不一致。它们到运行期才以崩溃的形式出现，因此唯一可靠的预防办法是把假设
写进代码：用 \Fn{static\_assert} 钉住尺寸与对齐、在接口层导出一个版本号函数、
两侧用同一份头文件而不是同一份头文件的两份拷贝。
\end{warning}

\begin{bestpractice}[title={排障时先复制产物，不要复制整个工程}]
把出问题的那一个目标文件、那一个归档、那一条链接命令单独拿出来重放。
能重放到最小就是定位到了阶段；重放不出来，说明问题还在构建描述层。
这个习惯比任何工具都值钱，因为绝大多数链接问题的根因在顺序而不在内容。
\end{bestpractice}

\section{两条收束论断}\label{sec:cl15-close}

\begin{note}[title={自举与载体是两件事}]
语言可以自举：一个 \cpp{} 编译器用 \cpp{} 写，输入是自身源码，输出是自身的新副本，
链条在逻辑上闭合，第~\ref{ch:cl01-pipeline}章按阶段拆开过这条链。但实现总要有
载体，而载体几乎总是 \CL{}、\cpp{} 或其他中级语言：它需要一位已经存在的实现
把第一份副本造出来。自举说的是语言能描述自己的全部语义，载体说的是实现依赖
一个在先的编译者，两个论断并不互相否定。把第一份副本做出来这件事叫
\Sw{bootstrap}，本册附录~\ref{ap:clC-selfhost}专门处理它的路子与坑。
\end{note}

\begin{flow}[title={流水线上没有黑箱}]
本章七个阶段每个都有可查产物：展开后的源文件、语法树或符号表、中间表示、
汇编文本、目标文件、归档与镜像。因此编译器与链接器不是黑箱，而是一串
有明确接口、可以用文本形式核对的组件。排查的正确起点不是加日志，
而是把中间产物拿出来看：附录~\ref{ap:clA-cli}把每个阶段的开关按平台列出，
附录~\ref{ap:clB-abi}把符号与重定位类型做成对照表。有了这两样，
绝大多数问题在第二个阶段就被否证，不必走到第三个阶段。
\end{flow}

\section{全册地图}\label{sec:cl15-map}

最后一张表把十四章正文与三个附录各自回答的问题，配上本章演示过的观察手段。
它按阅读顺序给出，但应当倒过来用：先从症状找到工具，再从工具找到章节。

\begin{longtable}{@{}>{\raggedright\arraybackslash}p{0.24\textwidth} >{\raggedright\arraybackslash}p{0.30\textwidth} >{\raggedright\arraybackslash}p{0.22\textwidth}@{}}
\caption{全册地图：每章回答的问题与对应工具}\label{tab:cl15-map}\\
\toprule
章 & 它回答的问题 & 观察工具 \\
\midrule
\endfirsthead
\toprule
章 & 它回答的问题 & 观察工具 \\
\midrule
\endhead
第~\ref{ch:cl01-pipeline}章：流水线全景 & 一次构建经过哪些阶段 & 分阶段产物 \\
第~\ref{ch:cl02-preprocessing}章：预处理 & 宏与条件编译何时展开 & \Sw{-E} 模式 \\
第~\ref{ch:cl03-resolution}章：名字查找 & 引用与声明如何配对 & 符号表 \\
第~\ref{ch:cl04-ir}章：中间表示 & 语言特性如何被规范化 & \Mx{IR} dump \\
第~\ref{ch:cl05-passes}章：优化通道 & 哪个通道改写了什么 & \Mx{IR} 前后对照 \\
第~\ref{ch:cl06-codegen}章：代码生成 & 指令、寻址、寄存器归属 & \Sw{-S} 模式 \\
第~\ref{ch:cl07-object}章：目标文件 & 节、符号、重定位如何组织 & \Tool{readelf}、\Tool{dumpbin} \\
第~\ref{ch:cl08-static-link}章：静态链接 & 定义如何选出、项如何消费 & \Tool{nm}、\Tool{ar} \\
第~\ref{ch:cl09-dynamic}章：动态链接 & 符号如何延迟绑定 & \Tool{ldd} \\
第~\ref{ch:cl10-cc}章：调用约定 & 参数与返回值如何跨越 & 反汇编 \\
第~\ref{ch:cl11-layout}章：内存布局 & 对象在镜像里放成什么样 & 头 dump \\
第~\ref{ch:cl12-endian}章：字节序 & 多字节整型如何落到字节 & \Tool{od}、\Tool{xxd} \\
第~\ref{ch:cl13-cxx-abi}章：\cpp{} \Mx{ABI} & 名字装饰如何成为承诺 & \Tool{c++filt} \\
第~\ref{ch:cl14-ide}章：工作台 & 谁在真正决定编译参数 & 构建日志 \\
第~\ref{ch:cl15-endtoend}章：端到端演练 & 上述全部如何串成一条链 & 本章七步 \\
附录~\ref{ap:clA-cli}：命令行 & 每阶段的确切开关 & 全部 \\
附录~\ref{ap:clB-abi}：\Mx{ABI} 速查 & 符号与重定位类型对照 & \Tool{nm}、\Tool{objdump} \\
附录~\ref{ap:clC-selfhost}：自举 & 第一份实现从哪来 & 分阶段产物 \\
\bottomrule
\end{longtable}

\begin{guideline}[title={读完这一章应当带走的东西}]
一份构建描述、七个可查产物、九个坐标、一张分派表。四样齐了，编译器与链接器
就从黑箱降级成一组命令行；而命令行是可以被重放的，可重放的东西就可以被排查。
\end{guideline}

% ============================================================
% 附录 A：命令行与产物查看工具速查
% 以表格为主体。工具名、节名、字段名只收确证项；
% 记不准的条目直接不入表，不用相近的名字凑行数。
% 回指：第 1 章流水线、第 3 章符号解析、第 7 至 10 章。
% ============================================================
\appendix
\chapter{命令行与产物查看工具速查}\label{ap:clA-cli}

正文按阶段讲机制，这里按“手边有什么产物、想看清什么”重新排一次。
两张主表都不解释原理，只把对应关系摆出来：第一张把流水线阶段和三个实现里的
命令名对齐，第二张反过来，从问题出发找工具。第三小节解释一件让人反复困惑的事：
同一个镜像里，修饰名、动态符号、调试符号三类计数为什么从来不相等。

本章的结论只有一句：先确定问题落在哪一层，再决定用哪个工具。
跳过这一步的排障，本质上是在猜（见第 \ref{ch:cl01-pipeline} 章）。

\section{阶段、三个实现里的命令名与这一步留下的产物}

同一件事在三个实现里通常有三个名字，但产物是同一类东西：一个可单独查看的中间文件。
下表的“产物”一列写的是这一阶段稳定留下的对象，不写具体开关组合；
开关属于实现选择，随版本变，正文第 \ref{ch:cl02-preprocessing} 章之后各章都强调过这点。

\begin{longtable}{@{}>{\raggedright\arraybackslash}p{0.13\textwidth} >{\raggedright\arraybackslash}p{0.26\textwidth} >{\raggedright\arraybackslash}p{0.24\textwidth} >{\raggedright\arraybackslash}p{0.12\textwidth}@{}}
\caption{流水线阶段与三平台工具名对照}\label{tab:apA-stages}\\
\toprule
阶段 & 三个实现里的名字 & 这一步留下的东西 & 回指 \\
\midrule
\endfirsthead
\toprule
阶段 & 三个实现里的名字 & 这一步留下的东西 & 回指 \\
\midrule
\endhead
\bottomrule
\endfoot
预处理 & \Tool{gcc}、\Tool{g++}、\Tool{clang}、\Tool{clang++} 的只预处理用法；
\Win{} 的 \Tool{cl}；\Mac{} 与 \Nix{} 同为前者 &
展开后的翻译单元文本：宏已替换、头文件已并入、行指示保留 &
第 \ref{ch:cl02-preprocessing} 章 \\
编译到汇编 & 同上四个驱动；\Win{} 仍走 \Tool{cl}；
单独的生成器 \Tool{llc} 与 IR 级工具 \Tool{opt} &
一个汇编文件；换 \Tool{opt} 时先落一份中间表示 &
第 \ref{ch:cl04-ir}、\ref{ch:cl06-codegen} 章 \\
汇编 & \Tool{as}、\Tool{gas}、\Nix{} 侧 \Tool{llvm-mc}；
\Win{} 的 \Tool{ml}；\Mac{} 走 \Tool{clang} 内部同一套 &
目标文件：\Elf{} 或 \PeFmt{} 或 \Macho{} 的对象文件 &
第 \ref{ch:cl07-object} 章 \\
归档 & \Tool{ar}（配 \Tool{ranlib}）；\Win{} 的 \Tool{lib}；
\Mac{} 的 \Tool{libtool} &
静态库：成员表加索引，符号尚未解决 &
第 \ref{ch:cl08-static-link} 章 \\
链接 & \Tool{ld}、\Tool{ld.lld}；\Win{} 的 \Tool{link}；
驱动仍是编译器命令 &
可执行镜像或共享对象；同时决定符号可见性集合 &
第 \ref{ch:cl03-resolution}、\ref{ch:cl09-dynamic} 章 \\
镜像查看 & \Tool{readelf}、\Tool{objdump}；\Win{} 的 \Tool{dumpbin}；
\Mac{} 的 \Tool{otool}、\Tool{vtool} &
段与节的布局、入口、格式版本、加载项 &
第 \ref{ch:cl07-object}、\ref{ch:cl11-layout} 章 \\
符号查看 & \Tool{nm}、\Tool{readelf}、\Tool{objdump}；\Win{} 的 \Tool{dumpbin}；
\Mac{} 的 \Tool{nm}、\Tool{otool} &
定义、引用、弱符号、局部符号的清单 &
第 \ref{ch:cl03-resolution} 章 \\
修饰名还原 & \Tool{cxxfilt}；\Tool{nm} 与 \Tool{readelf} 的还原开关；
\Tool{llvm-cxxfilt} &
把装饰名读回成声明；不改变任何ABI承诺 &
第 \ref{ch:cl13-cxx-abi} 章 \\
加载期诊断 & \Nix{} 的 \Tool{ldd}；\Win{} 以 \Tool{dumpbin}{} 的导入视图为主；
\Mac{} 的 \Tool{otool}、\Tool{vmmap} &
实际被装载的依赖闭包与地址空间布局 &
第 \ref{ch:cl09-dynamic} 章 \\
体积与裁剪 & \Tool{size}、\Tool{strip}、\Tool{objcopy}；\Win{} 的 \Tool{editbin} &
按段统计的体积；调试信息被移除后的镜像 &
第 \ref{ch:cl15-endtoend} 章 \\
\end{longtable}

\Tool{strings} 不在上表里，因为它不属于任何阶段：它按字节找可打印串，
既不认节也不认符号，只在“这段数据到底从哪来”时当最后一招。
同理 \Tool{file} 只回答“这是什么格式”，一次调用就够了。

\section{我想看什么，就用哪个工具，关心哪几个字段}

第二张表把问题当输入。字段名一栏只写各格式里确实存在的容器名，
具体的开关组合留给你自己的手册页：同一工具在不同实现里的短选项并不完全一致，
把这一层写死反而是这一册不该做的事。

\begin{longtable}{@{}>{\raggedright\arraybackslash}p{0.22\textwidth} >{\raggedright\arraybackslash}p{0.24\textwidth} >{\raggedright\arraybackslash}p{0.28\textwidth}@{}}
\caption{按问题查工具与关心字段}\label{tab:apA-goal}\\
\toprule
想确认什么 & 用哪个工具 & 关心哪几个字段 \\
\midrule
\endfirsthead
\toprule
想确认什么 & 用哪个工具 & 关心哪几个字段 \\
\midrule
\endhead
\bottomrule
\endfoot
谁定义了这个符号 & \Tool{nm}、\Tool{readelf}、\Tool{objdump} &
绑定为全局、类型不是未定义、节号指向真实节而不是特殊索引 \\
谁引用了它却没定义 & \Tool{nm}、\Tool{readelf} &
未定义一类的符号项；\Tool{nm}{} 里的小写与大写区分 \\
导出集合有没有泄漏内部符号 & \Tool{readelf}、\Tool{nm}；\Win{} 的 \Tool{dumpbin}；
\Mac{} 的 \Tool{nm} &
动态符号表里可见的全局项；\Win{} 看导出表；\Mac{} 看两层级命名空间的引用 \\
镜像的段结构是否变了 & \Tool{readelf}、\Tool{objdump}；\Tool{dumpbin}；\Tool{otool} &
节名、标志位（可执行、可读、可写）、对齐、按虚拟地址排出的顺序 \\
加载期的承诺 & \Tool{readelf}；\Tool{dumpbin}；\Tool{otool} &
\Elf{} 的程序头视图、\PeFmt{} 的数据目录、\Macho{} 的加载命令序列 \\
展开信息是否真在 & \Tool{readelf}、\Tool{objdump}；\Tool{dumpbin}；\Tool{otool} &
\Sect{eh\_frame} 与 \Sect{gcc\_exception\_table}；\Tp{.pdata} 与 \Tp{.xdata}；
\Tp{\_\_unwind\_info} \\
调试信息是否分离 & \Tool{readelf}、\Tool{objdump}；\Tool{dumpbin}；\Mac{} 的 \Tool{dsymutil} &
\Sect{debug\_info}{} 一类的节是否存在；\Win{} 的独立符号文件；
\Mac{} 的 \Tp{dSYM}{} 包 \\
只读数据是否被搬进可写段 & 同段结构那一行 &
段与节的归属、重定位记录是否出现在只读段里 \\
\end{longtable}

\section{为什么三类输出的数量从来不相等}

一次排障里最常见的误判，是把三个数字互相印证：\Tool{nm}{} 的行数、动态符号表项数、
调试信息里的条目数。它们统计的是三件不同的事，容器也不同。

\begin{table}[htbp]
\centering
\begin{tabularx}{\textwidth}{@{}>{\raggedright\arraybackslash}p{0.16\textwidth} >{\raggedright\arraybackslash}p{0.16\textwidth} X@{}}
\toprule
输出 & 所在容器 & 为什么和另外两个数不一样 \\
\midrule
修饰名 & 符号表与字符串表里的名字本身 &
它是一次编译的产物：局部符号、编译器生成的辅助名、被内联掉后只剩名字的函数都在里面；
还原成声明也不改变链接期的键 \\
动态符号 & 动态符号表与导出/导入结构 &
链接器按可见性与导出集合筛过一轮；同一份定义可能只留一个动态项，
弱符号与版本信息还会让计数进一步偏离 \\
调试符号 & 一组单独的节或单独的文件 &
类型、词法作用域、行表、调用位置各自成套，条目数量比可执行代码里的函数名多一个量级；
一旦剥离或分离，这部分可以从镜像里完全消失 \\
\bottomrule
\end{tabularx}
\caption{三类输出的归属容器与计数差异来源}\label{tab:apA-three}
\end{table}

\begin{compare}[title=三个计数各自能说明什么]
符号表项数只能说明这一次编译留下了哪些名字；动态符号项数只能说明加载期愿意承认哪些名字；
调试条目数只能说明调试器能看到多少结构。三者相等反而是巧合，
用它们的差值定位问题（比如导出集合突然变大）才是正确用法。
\end{compare}

\begin{warning}[title=数量差不是错误，缺了才是]
剥离调试信息后镜像照常运行；动态符号表被误删一项才会崩。
把 \Tool{strip}{} 之后的输出跟原始镜像逐项比对，只会制造噪声。
第 \ref{ch:cl11-layout} 章说过：布局承诺写在段与节里，不写在名字里。
\end{warning}

\subsection*{一次快照该抓哪几条命令}

下面的命令顺序就是本章第一张表的顺序。它的价值不在于命令本身，
而在于同一份产物、改动前后各抓一次，然后做文本比对。

\begin{lstlisting}[style=shellstyle, caption={改动前后各跑一次：符号、段、加载项三份快照}]
# 快照一：改动前
g++ -std=c++20 -O2 -c widget.cpp -o widget.o
ar rcs libwidget.a widget.o
g++ app.cpp libwidget.a -o app
nm -C widget.o            > sym-before.txt
readelf -SW widget.o      > sect-before.txt
dumpbin /headers widget.obj > pe-before.txt
otool -l widget.o         > macho-before.txt
# 快照二：改动后把上面几条再跑一遍，然后比对
diff -u sym-before.txt sym-after.txt
\end{lstlisting}

\begin{lstlisting}[style=outstyle, caption={三份快照长什么样：字段名示意，非某次实测输出}]
# 一：符号 —— 关注绑定、类型、节号
ELF: symbol name ... GLOBAL WEAK  text
PE:  ...  External Section  ...
# 二：段 —— 关注名称、大小、对齐、标志
Mach-O: segname __TEXT  ...  sectname __text
# 三：加载项 —— 关注依赖列表与展开信息所在容器
\end{lstlisting}

\section{产物之间能互相回答什么}

第一、二张表都默认你知道手里拿的是哪一层产物。这一层最容易出错：
同一份代码在目标文件、静态库成员、最终镜像三个阶段能回答的问题并不相同，
而工具往往照常输出、只是答案无效。下面这张表就是把“能问什么”按产物分开列一遍。

\begin{longtable}{@{}>{\raggedright\arraybackslash}p{0.16\textwidth} >{\raggedright\arraybackslash}p{0.30\textwidth} >{\raggedright\arraybackslash}p{0.28\textwidth}@{}}
\caption{三种产物各自能回答与不能回答的问题}\label{tab:apA-artifacts}\\
\toprule
手里的产物 & 在这里能确认 & 在这里确认不了 \\
\midrule
\endfirsthead
\toprule
手里的产物 & 在这里能确认 & 在这里确认不了 \\
\midrule
\endhead
\bottomrule
\endfoot
展开后的翻译单元文本 &
宏展开的结果、头文件是否真的被并入、条件编译走到了哪一支 &
任何与布局、符号绑定有关的事：这一步还没有目标文件 \\
目标文件（单个） &
这个编译单元里定义了哪些名字、引用了哪些未定义的名字、代码与数据各占多少 &
最终镜像的段布局、导出集合、依赖闭包：这些都由链接决定 \\
静态库 &
成员清单与索引是否存在、某个成员里有没有想要的定义 &
谁最终会被拉进镜像、拉进来之后符号变成什么可见性 \\
可执行镜像与共享对象 &
装载要用的全部承诺：段、加载项、动态符号、展开信息与调试信息的去留 &
单个符号原本属于哪个编译单元：这要靠构建记录，不能靠镜像反推 \\
分离出去的调试包与符号文件 &
源码行号、类型布局、调用位置是否还能还原 &
运行时是否真的需要它：多数平台上它缺失只会影响调试体验 \\
\end{longtable}

\begin{compare}[title=为什么不能在目标文件上判断导出集合]
目标文件里的全局名字是“我愿意对外提供这些”，镜像里的导出集合是“加载器承认这些”。
中间隔着可见性策略、版本脚本或导出清单、以及链接器对未使用成员的取舍。
在目标文件上数出来的数字，只能证明链接前那一步做了什么，
第 \ref{ch:cl03-resolution} 章里那条“最终答案写在镜像里”就是针对这件事。
\end{compare}

\subsection*{一次排障该留哪五份}

留档的目的只有一个：让下一次的输出能和这一次比对。所以五份快照都要带产物路径，
不要只留工具名。

\begin{enumerate}
\item 符号快照：定义、引用、弱符号三类都要在里面，别只留已定义的部分。
\item 段与节的快照：名称、大小、对齐、标志四项齐全，才看得出“搬错了段”这类问题。
\item 加载承诺的快照：\Elf{} 的程序头视图、\PeFmt{} 的数据目录、\Macho{} 的加载命令序列，
取各自平台的那一份就够。
\item 依赖闭包的快照：加载期实际解析到的文件列表，与构建时声明的依赖做差。
\item 构建参数本身：所用开关与档位要一起留，否则前四份快照无法归因。
\end{enumerate}

\begin{flow}[title=从报错回到工具的顺序]
报错说找不到符号，回到符号快照；说段不可执行或写入只读页，回到段快照；
说加载不到依赖，回到加载承诺与依赖闭包两份快照；
说调试器停在错误行，回到调试信息去留那一份。
每一条路径都能在这两张表里查到对应工具，不需要凭记忆。
\end{flow}

\section{排查顺序}

\begin{guideline}[title=排查顺序三条]
\begin{enumerate}
\item 先看符号，再看布局，最后看加载。顺序不能反：符号缺失时讨论段对齐毫无意义，
而加载期的差异（依赖闭包、两层级命名空间）往往在符号与布局都正确时仍然存在。
\item 改档位前后各取一次快照并做文本比对，只比对的这两次才算数：
一次调用的输出没有参照物，无法区分“本来如此”与“这次变了”。
\item 比之前先确认比的是同一个产物：目标文件、静态库成员、最终镜像是三份东西，
第 \ref{ch:cl08-static-link} 章里那条“链接后符号集合才是最终答案”在这里同样成立。
\end{enumerate}
\end{guideline}

\begin{bestpractice}[title=把工具名写进构建脚本而不是聊天记录]
把这一节命令固化成一个目标：产物路径、快照文件名、比对方式都写进构建脚本，
让每次改动自动生成一对可比的输出。排障的成本从来不在命令，而在“当时到底跑的是哪一份”。
具体开关与工具的组合见第 \ref{ch:cl14-ide} 章与第 \ref{ch:cl15-endtoend} 章。
\end{bestpractice}

% ============================================================
% 附录 B：三平台两架构的 ABI 速查
% 只把确证过的寄存器与宽度写进表；格子里出现“文档为准”就是
% 本节明确放弃的那部分，不要把它当成缺漏来补。
% 回指：第 10 章调用约定、第 11 章布局、第 12 章字节序、第 13 章 ABI。
% ============================================================
\chapter{三平台两架构的 ABI 速查}\label{ap:clB-abi}

这一册正文讲调用约定为什么这样设计，本章把结论压缩成四张表：寄存器落点、
数据模型、对齐与布局、各类内容在三种文件格式里的容器。四张表都只回答“落在哪里”，
不回答“为什么”，理由请回看第 \ref{ch:cl10-cc} 章与第 \ref{ch:cl11-layout} 章。

\begin{abibox}[title=这张表的可信边界]
表里直接写出的寄存器名与宽度是确证项，可以照着写内联汇编、照着核对反汇编结果；
写成“以对应文档为准”的格子是本节主动放弃的，请回到对应平台的处理器补充 ABI 文档核对。
这张表不是任何一处布局承诺的依据：真正要依赖的东西（结构体布局、对齐、语言链接）
必须用 \Fn{static\_assert} 与 \Fn{alignof} 写进代码里，理由见第 \ref{ch:cl11-layout} 章。
字节序只影响跨边界数据的解释方式，与寄存器落点无关，见第 \ref{ch:cl12-endian} 章。
\end{abibox}

\section{调用约定主表}

第一张表只管参数与返回值的落点。整数与指针走同一组寄存器，这是三套主流约定共同的特点；
浮点是否复用同一批参数槽，是 \Nix{} 侧与 \Win{} 侧最实质的差异之一。

\begin{longtable}{@{}>{\raggedright\arraybackslash}p{0.12\textwidth} >{\raggedright\arraybackslash}p{0.22\textwidth} >{\raggedright\arraybackslash}p{0.16\textwidth} >{\raggedright\arraybackslash}p{0.14\textwidth} >{\raggedright\arraybackslash}p{0.12\textwidth}@{}}
\caption{参数与返回值落在哪个寄存器}\label{tab:apB-reg-args}\\
\toprule
约定 & 整数与指针参数 & 浮点参数 & 整数返回 & 浮点返回 \\
\midrule
\endfirsthead
\toprule
约定 & 整数与指针参数 & 浮点参数 & 整数返回 & 浮点返回 \\
\midrule
\endhead
\bottomrule
\endfoot
\Nix{} \xSixtyFour（系统 V） &
前六个按序 \Tp{rdi}、\Tp{rsi}、\Tp{rdx}、\Tp{rcx}、\Tp{r8}、\Tp{r9}；此后走栈 &
\Tp{xmm0}--\Tp{xmm7} &
\Tp{rax}；需要第二个字时用 \Tp{rdx} &
\Tp{xmm0} \\
\Win{} \xSixtyFour &
前四个按序 \Tp{rcx}、\Tp{rdx}、\Tp{r8}、\Tp{r9}；第五个起走栈 &
同一批参数槽，前四个浮点值进 \Tp{xmm0}--\Tp{xmm3} &
\Tp{rax} &
\Tp{xmm0} \\
\Nix{} 与 \Mac{} 的 \ArmSixtyFour &
前八个按序 \Tp{x0}--\Tp{x7} &
\Tp{v0}--\Tp{v7}，即同一寄存器的 \Tp{d0}--\Tp{d7} 与 \Tp{s0}--\Tp{s7} 视图 &
\Tp{x0}；需要第二个字时用 \Tp{x1} &
\Tp{v0} \\
\Mac{} \xSixtyFour &
沿用系统 V 的寄存器分配 &
沿用系统 V 的那组浮点寄存器 &
\Tp{rax} &
\Tp{xmm0} \\
\Lang{RISC-V} 64 &
前八个按序 \Tp{a0}--\Tp{a7} &
\Tp{fa0}--\Tp{fa7} &
\Tp{a0}；结构型返回时 \Tp{a1} 携带目标地址 &
\Tp{fa0} \\
\end{longtable}

第二张表管被调用方要还回去的东西、栈的要求，以及调用点必须预留的空间。
影子空间是 \Win{} \xSixtyFour{} 的硬性要求，红区是系统 V 的便利，两者不能互相替代。

\begin{longtable}{@{}>{\raggedright\arraybackslash}p{0.12\textwidth} >{\raggedright\arraybackslash}p{0.26\textwidth} >{\raggedright\arraybackslash}p{0.13\textwidth} >{\raggedright\arraybackslash}p{0.21\textwidth}@{}}
\caption{被调用方保存、栈对齐与调用点预留}\label{tab:apB-save-stack}\\
\toprule
约定 & 被调用方保存寄存器 & 栈对齐 & 影子空间与特殊保留区 \\
\midrule
\endfirsthead
\toprule
约定 & 被调用方保存寄存器 & 栈对齐 & 影子空间与特殊保留区 \\
\midrule
\endhead
\bottomrule
\endfoot
\Nix{} \xSixtyFour（系统 V） &
\Tp{rbx}、\Tp{rbp}、\Tp{r12}--\Tp{r15}；其余通用寄存器由调用方负责 &
16 字节 &
无影子空间；调用返回后可用 128 字节红区，叶子函数可不动栈指针 \\
\Win{} \xSixtyFour &
\Tp{rbx}、\Tp{rbp}、\Tp{rsi}、\Tp{rdi}、\Tp{r12}--\Tp{r15}；全部 \Tp{xmm}{} 寄存器由调用方负责 &
16 字节 &
每个跨模块调用前预留 32 字节影子空间；其余保留区以对应文档为准 \\
\Nix{} 与 \Mac{} 的 \ArmSixtyFour &
\Tp{x19}--\Tp{x28}，加 \Tp{x29}（帧指针）与 \Tp{x30}（返回地址） &
16 字节 &
无影子空间这一说；平台各自的保留项以对应文档为准 \\
\Mac{} \xSixtyFour &
与系统 V 同 &
16 字节 &
红区是否可用、系统调用如何陷入：以对应文档为准；不要假设能直接陷入 \\
\Lang{RISC-V} 64 &
以对应 psABI 文档为准 &
以对应 psABI 文档为准 &
无影子空间概念 \\
\end{longtable}

\Mac{} \xSixtyFour{} 那一行值得多说一句：寄存器落点跟系统 V 一致，因此代码生成路径可以共用；
不一致的是系统边界的语义——用户态代码经由系统库进入内核，不存在“把参数摆好然后触发陷入”
这种可移植假设。任何按系统调用编号手写的序列都只属于某一个平台，
这一点属于第 \ref{ch:cl09-dynamic} 章讲的边界问题，而不是调用约定本身。

\begin{note}[title=这张表故意没列的东西]
可变参数如何声明整数与浮点寄存器的使用、同结构体按大小归类后的传递方式、
向量类型与十进制浮点参数的落点、\Win{} \ArmSixtyFour{} 的完整约定、
各平台的返回地址认证与寄存器配对要求，都不在本表范围内。
它们各有确定的文档，凭印象补全比留空更危险。
\end{note}

\begin{lstlisting}[style=asmstyle, caption={同一个调用在两套约定下的差别：参数槽数量不同（结构示意，非某次实测输出）}]
    .text
# 系统 V：前六个整数参数已在 rdi rsi rdx rcx r8 r9 里
    call    six_args
# Microsoft x64：只有前四个在 rcx rdx r8 r9 里
# 调用点还必须先给出 32 字节影子空间，然后第五个才进栈
    call    four_args
# AArch64：前八个整数参数 x0..x7，浮点看 v0..v7
    bl      eight_args
\end{lstlisting}

\section{寄存器角色速记}

读反汇编时真正要问的是“这个寄存器是谁的”。参数与返回只是其中一个视角，
下面这张表按角色把三套约定排在一起；它不改变上一节任何结论，只是换个查法。

\begin{longtable}{@{}>{\raggedright\arraybackslash}p{0.14\textwidth} >{\raggedright\arraybackslash}p{0.22\textwidth} >{\raggedright\arraybackslash}p{0.20\textwidth} >{\raggedright\arraybackslash}p{0.20\textwidth}@{}}
\caption{按角色查寄存器：系统 V、\Win{} \xSixtyFour{} 与 \ArmSixtyFour}\label{tab:apB-roles}\\
\toprule
角色 & \Nix{} 系统 V & \Win{} \xSixtyFour & \ArmSixtyFour \\
\midrule
\endfirsthead
\toprule
角色 & \Nix{} 系统 V & \Win{} \xSixtyFour & \ArmSixtyFour \\
\midrule
\endhead
\bottomrule
\endfoot
参数与返回 & \Tp{rdi}、\Tp{rsi}、\Tp{rdx}、\Tp{rcx}、\Tp{r8}、\Tp{r9}，返回 \Tp{rax} &
\Tp{rcx}、\Tp{rdx}、\Tp{r8}、\Tp{r9}，返回 \Tp{rax} &
\Tp{x0}--\Tp{x7}，返回 \Tp{x0} \\
帧与栈指针 & \Tp{rbp} 可选，\Tp{rsp} 由约定维护 & 同左 & \Tp{x29} 与 \Tp{sp}，按约定配对使用 \\
跨调用必须保住 & \Tp{rbx}、\Tp{rbp}、\Tp{r12}--\Tp{r15} &
\Tp{rbx}、\Tp{rbp}、\Tp{rsi}、\Tp{rdi}、\Tp{r12}--\Tp{r15} &
\Tp{x19}--\Tp{x28}、\Tp{x29}、\Tp{x30} \\
调用即失效 & 其余通用寄存器与全部 \Tp{xmm} 寄存器 &
参数寄存器与全部 \Tp{xmm} 寄存器 &
\Tp{x0}--\Tp{x18} 与浮点侧的对应视图 \\
返回地址放哪 & 由 \Tp{call}{} 压栈，栈上按展开信息找回 & 同左，另需 32 字节影子空间 &
默认进 \Tp{x30}，需要保存时显式入栈 \\
\end{longtable}

可变参数是这套表里少数需要额外记住的一条：系统 V 一侧在调用可变参数函数之前，
要把实际用到的浮点参数寄存器数量放进 \Tp{al}，被调用方靠它决定要不要保存这批寄存器；
\Win{} \xSixtyFour{} 的参数槽不分整数与浮点两组，因此对应文档里没有这一条要求。
差异的根源不是哪个实现更小心，而是参数分类不同（第 \ref{ch:cl10-cc} 章讲过这套分类）。

\begin{warning}[title=手写序列时最容易踩的两处]
其一，把 \Tp{rsi}、\Tp{rdi} 当被调用方保存寄存器：在系统 V 里它们是调用方的，
在 \Win{} \xSixtyFour{} 里不是，同一份内联汇编换平台就得改。
其二，假设红区在 \Win{} 上存在：\Win{} \xSixtyFour{} 反过来要求调用点预留影子空间，
省掉栈指针调整的那套写法在这里会直接踩坏别人的参数槽。
这两处的正确做法都是把序列写进单独的目标文件、由构建脚本按平台选，
而不是在源码里用条件编译拼。
\end{warning}

\section{数据模型}

数据模型决定 \Tp{long} 与指针的宽度，也就决定了一个符号在跨平台重编译后大小是否变化。
三平台的归属很干脆：\Win{} 的 64 位模型是唯一把 \Tp{long} 留在 32 位的主流选择。

\begin{table}[htbp]
\centering
\begin{tabularx}{\textwidth}{@{}>{\raggedright\arraybackslash}p{0.12\textwidth} >{\raggedright\arraybackslash}p{0.11\textwidth} >{\raggedright\arraybackslash}p{0.11\textwidth} >{\raggedright\arraybackslash}p{0.13\textwidth} >{\raggedright\arraybackslash}p{0.12\textwidth} X@{}}
\toprule
模型 & \Tp{int} & \Tp{long} & \Tp{long long} & 指针 & 出现在哪里 \\
\midrule
ILP32 & 32 & 32 & 64 & 32 &
三个平台的 32 位目标；\Nix{} \ArmSixtyFour{} 的 32 位兼容环境属另一类，此处不展开 \\
LP64 & 32 & 64 & 64 & 64 &
\Nix{} \xSixtyFour{} 与 \Mac{} 的两套架构；系统头文件里的 \Tp{size\_t} 与 \Tp{ssize\_t}{} 随此变宽 \\
LLP64 & 32 & 32 & 64 & 64 &
\Win{} \xSixtyFour；把 \Tp{long} 当指针宽度用的旧代码在这里出错 \\
\bottomrule
\end{tabularx}
\caption{三种主流数据模型与平台归属}\label{tab:apB-datamodel}
\end{table}

\Tp{long long} 一列在三行里都是 64，这是语言的下限要求在实践中恰好等于实现值；
不要把这一列当成模型之间的区分点。真正的区分点只有 \Tp{long} 与指针两列。

\begin{warning}[title=跨边界结构里最贵的一列]
同一份结构体定义在 LP64 与 LLP64 下重编译，字段数不变但 \Tp{long} 成员少 4 字节，
后面的偏移整体前移，填充也随之改变。跨平台的磁盘格式与网络协议里不要用 \Tp{long}，
用定宽类型；理由与检查方法见第 \ref{ch:cl11-layout} 章与第 \ref{ch:cl12-endian} 章。
\end{warning}

\section{对齐与布局速查}

对齐要求决定填充放在哪，也就决定同一份定义在不同编译器下是否还是同一个布局。
下表只列标量类型在主流 64 位目标上的自然情况，其余交给实现。

\begin{longtable}{@{}>{\raggedright\arraybackslash}p{0.18\textwidth} >{\raggedright\arraybackslash}p{0.14\textwidth} >{\raggedright\arraybackslash}p{0.16\textwidth} >{\raggedright\arraybackslash}p{0.28\textwidth}@{}}
\caption{常见标量类型的宽度与对齐}\label{tab:apB-align}\\
\toprule
类型 & 宽度（字节） & 对齐（字节） & 备注 \\
\midrule
\endfirsthead
\toprule
类型 & 宽度（字节） & 对齐（字节） & 备注 \\
\midrule
\endhead
\bottomrule
\endfoot
\Tp{char}、\Tp{bool} & 1 & 1 & 无填充需求；位域另算 \\
\Tp{short} & 2 & 2 & \\
\Tp{int}、\Tp{float} & 4 & 4 & \\
\Tp{long} & 4 或 8 & 与自身宽度同 & 数据模型决定，见上一节 \\
\Tp{long long}、\Tp{double} & 8 & 8 & \\
指针、\CPP{} 引用 & 4 或 8 & 与自身宽度同 & 三个平台的 64 位目标上都为 8 \\
\Tp{long double} & 由实现规定 & 由实现规定 & 跨平台结构里不要用这个类型 \\
向量与矩阵类型 & 由实现规定 & 常为 16 或 32 & 具体上限由实现与开关共同决定 \\
结构体与类 & 成员之和加填充 & 取成员中的最大对齐 & 总大小按该对齐向上取整；空基类与虚表指针另算 \\
\Tp{max\_align\_t} & 由实现规定 & 由实现规定 & 这是新类型的对齐上限，超限会被拒绝 \\
\end{longtable}

打包指示的代价，一句话结论：它只删掉填充，成员可能落在非自然边界上，
于是读取需要多次访问、原子操作与向量化路径多半不可用、
某些平台上还会触发对齐检查——把它留给磁盘格式与网络报文，别留给自己内部的数据结构。

\begin{bestpractice}[title=布局假设写进代码而不是写进注释]
任何跨模块、跨版本、跨进程的边界，都用 \Fn{static\_assert} 把尺寸与偏移钉住，
配合 \Fn{alignof} 与 \Fn{offsetof} 一起写；注释会过期，编译期的断言不会。
第 \ref{ch:cl13-cxx-abi} 章有完整的做法。
\end{bestpractice}

\section{三种格式里各类内容落在哪个容器}

最后一张表把 ABI 落到文件层面：同样一段代码，三种格式把它装进不同的容器，
查看工具因此也各不相同（工具名见附录 \ref{ap:clA-cli}）。

\begin{longtable}{@{}>{\raggedright\arraybackslash}p{0.15\textwidth} >{\raggedright\arraybackslash}p{0.22\textwidth} >{\raggedright\arraybackslash}p{0.20\textwidth} >{\raggedright\arraybackslash}p{0.21\textwidth}@{}}
\caption{按内容类别看三种格式的容器}\label{tab:apB-containers}\\
\toprule
内容 & \Elf & \PeFmt & \Macho \\
\midrule
\endfirsthead
\toprule
内容 & \Elf & \PeFmt & \Macho \\
\midrule
\endhead
\bottomrule
\endfoot
机器代码 & \Sect{text} & \Tp{.text} & \Tp{\_\_TEXT}{} 段里的 \Tp{\_\_text} \\
只读数据 & \Sect{rodata} & \Tp{.rdata} & 同一 \Tp{\_\_TEXT}{} 里的常量节 \\
可写已初始化数据 & \Sect{data} & \Tp{.data} & \Tp{\_\_DATA}{} 里的数据节 \\
未初始化数据 & \Sect{bss} & 在可写段里按零长度记录 & \Tp{\_\_DATA}{} 里带符号但没有内容的节 \\
重定位记录 & 带加性的重定位表 & 镜像内的基址重定位表 & 单独的一段链接编辑数据 \\
符号表与字符串 & \Sect{symtab}{} 与配套字符串表 & 导出表与导入表为主 & 符号表加两层级命名空间的引用 \\
展开与异常 & \Sect{eh\_frame}、\Sect{gcc\Bk\_exception\Bk\_table} & \Tp{.pdata}{} 与 \Tp{.xdata} & \Tp{\_\_unwind\_info}，另有函数起始信息 \\
调试信息 & \Sect{debug\_info}{} 一类的独立节 & 可分离到单独的符号文件 & 可分离到 \Tp{dSYM}{} 包 \\
加载承诺 & 程序头一类条目 & 可选头里的数据目录 & 一段有序的加载命令 \\
\end{longtable}

\begin{compare}[title=同一件事的三处位置]
要把异常展开信息搬到只读且不可执行的地方，\Elf{} 里看的是节标志，
\PeFmt{} 里看的是它在数据目录里登记的那一项，\Macho{} 里看的是它属于哪个段。
三处都改对了才叫改完；只改一处的镜像在别的平台上照样按原样被装载。
\end{compare}

\begin{abibox}[title=跨平台 ABI 检查的最小动作]
一个符号声明成 \CL{} 语言链接，它在 \Nix{} 与 \Mac{} 的符号表里就是那个不带装饰的名字，
在 \Win{} 的 \PeFmt{} 导出表里还要看是否带前缀或调用约定装饰；
这一层差异属于第 \ref{ch:cl13-cxx-abi} 章，处理办法也在那里：只让一层薄壳承担差异。
\end{abibox}

% ============================================================
% 附录 C：主流实现的载体语言与自举关系一览
% 起始语言与当前主语言分开写；两者不同时不要合并成一格。
% 记不准的历史细节一律不入表，也不要用传闻补格。
% 回指：第 1 章流水线、第 7 章目标文件、第 9 章动态链接、第 13 章 ABI。
% ============================================================
\chapter{主流实现的载体语言与自举关系一览}\label{ap:clC-selfhost}

这一章回答一个常被问错的问题：一门语言的实现到底是用什么写的。
“用什么写”至少有三层——项目一开始用什么、现在主体用什么、最终生成机器码那一步走什么路径。
表里把它们分成三列，因为这三层经常不一致，而把它们混成一句话正是误读的源头。

\begin{note}[title=读这张表的两个规则]
第一，“起始语言”与“当前实现语言”是两件事：不少项目换过主体语言，也有项目至今仍在生成
别的语言的代码再交给第三方工具。第二，本表只收能确证的条目与路径描述；
早期版本用什么语言写过这种细节，若无法确证就直接删掉整格，不写传闻、不写“大概”。
\end{note}

\section{一览}

“是否自举”一列只说实现本体能不能编译自己，不包含测试套件、构建脚本或工具链外围。
“非本语言的运行时部分”指的是这个实现的运行期里，用别的语言写、且必须一起分发的部分。

\begin{longtable}{@{}>{\raggedright\arraybackslash}p{0.11\textwidth} >{\raggedright\arraybackslash}p{0.17\textwidth} >{\raggedright\arraybackslash}p{0.09\textwidth} >{\raggedright\arraybackslash}p{0.17\textwidth} >{\raggedright\arraybackslash}p{0.14\textwidth} >{\raggedright\arraybackslash}p{0.10\textwidth}@{}}
\caption{主流实现的载体语言、自举关系与代码生成路径}\label{tab:apC-selfhost}\\
\toprule
实现 & 起始与当前语言 & 是否自举 & 代码生成路径 & 非本语言的运行时部分 & 回指 \\
\midrule
\endfirsthead
\toprule
实现 & 起始与当前语言 & 是否自举 & 代码生成路径 & 非本语言的运行时部分 & 回指 \\
\midrule
\endhead
\bottomrule
\endfoot
\Tool{GCC} & 早期 \CL{}，现主要 \CPP{}，另有生成的代码 & 是 &
自有后端直接生成机器码 & 配套运行时含 \CL{} 与汇编 &
第 \ref{ch:cl01-pipeline}、\ref{ch:cl06-codegen} 章 \\
\Tool{Clang} 与 \Tool{LLVM} & 一直是 \CPP{} & 是 &
先生成中级表示，再由 \Tool{LLVM}{} 后端到机器码 &
与系统库的边界仍是 \CL{} 约定 &
第 \ref{ch:cl04-ir}、\ref{ch:cl06-codegen} 章 \\
\Tool{CPython} & \CL{}，现仍 \CL{} & 否 &
译成字节码后解释执行 & 标准库大量 \CL{} 扩展 &
第 \ref{ch:cl01-pipeline} 章 \\
\Tool{Ruby MRI} & \CL{} & 否 &
字节码解释执行 & 扩展接口就是 \CL{} 接口 &
第 \ref{ch:cl01-pipeline} 章 \\
\Tool{V8} & \CPP{} & 否 &
两级即时编译，另带快照 &
线程与内存边界走系统 \CL{} 层 &
第 \ref{ch:cl06-codegen} 章 \\
\Tool{HotSpot} & \CPP{} 与汇编 & 否 &
字节码解释加即时编译的机器码 &
与操作系统交互的那段是汇编 &
第 \ref{ch:cl06-codegen}、\ref{ch:cl09-dynamic} 章 \\
\Tool{OpenJ9} & \CPP{} & 否 &
字节码解释加即时编译 &
与本地代码的边界是 \CL{} 约定 &
第 \ref{ch:cl06-codegen} 章 \\
\Tool{Lua} & \CL{} & 否 &
字节码虚拟机 & 与宿主全部通过 \CL{} &
第 \ref{ch:cl01-pipeline} 章 \\
\Tool{LuaJIT} & \CL{} & 否 &
自带后端直接生成机器码 & 与宿主通过 \CL{} &
第 \ref{ch:cl06-codegen} 章 \\
\Tool{OCaml} & 自身 & 是 &
原生编译到机器码，另有字节码路径 &
运行时里与操作系统交互部分是 \CL{} &
第 \ref{ch:cl01-pipeline}、\ref{ch:cl06-codegen} 章 \\
\Tool{GHC} & \Lang{Haskell} & 是 &
早期经 \CL{} 中转，现行主要走自有后端 &
运行时含 \CL{} 与汇编 &
第 \ref{ch:cl04-ir}、\ref{ch:cl06-codegen} 章 \\
\Tool{Rust} & \Lang{Rust} & 是（必须有种子实现） &
代码生成通常借 \CPP{} 写的 \Tool{LLVM} &
标准运行时与 \CL{} 接口互通 &
第 \ref{ch:cl04-ir}、\ref{ch:cl06-codegen} 章 \\
\Tool{Go} 的 \Tool{gc} & 早期混有 \CL{}，现主体自举 & 是 &
自有后端直接生成机器码 &
运行时的装载器与调度含汇编 &
第 \ref{ch:cl06-codegen}、\ref{ch:cl09-dynamic} 章 \\
\Tool{coreclr} & \CPP{} & 否 &
中间语言经即时编译，另有预编译镜像路径 &
与操作系统的边界为 \CL{} 约定 &
第 \ref{ch:cl06-codegen}、\ref{ch:cl07-object} 章 \\
\Tool{PHP} & \CL{} & 否 &
字节码解释，运行期再编译 & 扩展层是 \CL{} &
第 \ref{ch:cl01-pipeline} 章 \\
\Tool{BEAM} & \CL{} & 否 &
字节码解释执行 & 驱动与调度接口为 \CL{} &
第 \ref{ch:cl01-pipeline} 章 \\
\Tool{SQLite} & \CL{} & 不适用 &
查询在运行期译成虚拟机字节码 &
几乎不需要外部运行时，这是它作为嵌入式实现的范例价值 &
第 \ref{ch:cl01-pipeline}、\ref{ch:cl09-dynamic} 章 \\
\Tool{Redis} & \CL{} & 不适用 &
没有代码生成这一步 & 事件循环与分配器接系统 \CL{} 层 &
第 \ref{ch:cl09-dynamic} 章 \\
\Tool{cl} & \CPP{} & 是 &
直接生成机器码 & 运行库为 \CL{} 与汇编 &
第 \ref{ch:cl06-codegen} 章 \\
\Tool{ld} & \CL{} & 不适用 &
不生成高级代码：解决符号与重定位 & 无 &
第 \ref{ch:cl03-resolution}、\ref{ch:cl08-static-link} 章 \\
\Tool{ld.lld} & \CPP{} & 不适用 &
同上 & 无 &
第 \ref{ch:cl03-resolution}、\ref{ch:cl08-static-link} 章 \\
\end{longtable}

最后一行不属于语言实现，列在这里是为了对照：链接器也有一段代码生成路径，
而它同样由某个高级语言写出、同样要落到目标文件与加载器约定上。
这里并列的两个名字正说明这一点：\Tool{ld} 来自 \Tool{binutils} 那套工具，主体是 \CL{}；
\Tool{ld.lld} 承担同一岗位，主体是 \CPP{}。换的是载体，不是接口形状。
\Tool{ld} 出自 \Tool{binutils} 那套工具，主体是 \CL{}；\Tool{ld.lld} 走的是同一个岗位，
主体是 \CPP{}。同一件事换一种载体语言完成，接口形状不变——这个事实本身就是下一节的证据。

\section{把代码生成路径分成四类}

“编译成什么”这一列可以独立于载体语言来看。同一类路径下，实现用的语言可以完全不同；
反过来，同一门语言的两个实现也可以分属两类。分完之后，很多争论会自动消失：
真正决定性能特征与可移植性的是这一列，不是语言那一列。

\begin{longtable}{@{}>{\raggedright\arraybackslash}p{0.18\textwidth} >{\raggedright\arraybackslash}p{0.30\textwidth} >{\raggedright\arraybackslash}p{0.26\textwidth}@{}}
\caption{代码生成路径的四类与归属}\label{tab:apC-paths}\\
\toprule
路径 & 特征 & 表里属于这一类 \\
\midrule
\endfirsthead
\toprule
路径 & 特征 & 表里属于这一类 \\
\midrule
\endhead
\bottomrule
\endfoot
自有后端到机器码 & 后端与前端同源，指令选择、寄存器分配都由本项目负责 &
\Tool{GCC}、\Tool{Go} 的 \Tool{gc}、\Tool{LuaJIT}、\Tool{GHC} 的现行路径、\Tool{cl} \\
借外部后端 & 前端产出一份中级表示，机器码交给另一个项目 &
\Tool{Clang} 与 \Tool{LLVM}、\Tool{Rust} \\
只到字节码 & 生成的是自家虚拟机的指令，机器码由运行期决定或根本不出现 &
\Tool{CPython}、\Tool{Ruby MRI}、\Tool{Lua}、\Tool{PHP}、\Tool{BEAM}、\Tool{SQLite} \\
分层与即时编译 & 先解释或先出较粗的代码，热了再换更好的机器码 &
\Tool{V8}、\Tool{HotSpot}、\Tool{OpenJ9}、\Tool{coreclr} \\
\end{longtable}

\begin{note}[title=这张表故意没收录的东西]
不少实现同样有名，但要写出“起始语言”与“代码生成路径”两格就必须依赖传闻，
本节因此不列它们：只做过子集或教学用途的实现、以别家后端为主干的实现、
以及把运行时交给第三方虚拟机的实现，都在这里被删掉了。
一行没有写进表，不等于它不重要；只等于这一册给不出可核对的说法。
\end{note}

\section{自举为什么需要种子}

编译器是一种程序，程序要有已经能运行的编译器才能编译。于是任何自举语言的历史上
都有一段“先有鸡还是先有蛋”的时期，而所有项目解决它的办法只有三类。

\begin{table}[htbp]
\centering
\begin{tabularx}{\textwidth}{@{}>{\raggedright\arraybackslash}p{0.16\textwidth} X@{}}
\toprule
形态 & 做法与代价 \\
\midrule
借别的语言起步 &
先用已有实现（通常是 \CL{} 写的解释器或子集编译器）把第一版跑起来，再逐步替换掉种子。
代价是种子版本必须长期保留，否则历史构建无法复现。 \\
子集先行 &
先写一个只支持核心语法、能编译自身的子集，用它编译完整实现。
这是最常见的路线；代价是子集与完整实现之间的差异要靠测试兜住。 \\
永久外部件 &
实现本体自举，但代码生成那一步交给另一个项目。
表里“借 \Tool{LLVM}”那一行就是这类：换掉的只是后端，前端仍由自己维护。 \\
\bottomrule
\end{tabularx}
\caption{自举的三种起步形态}\label{tab:apC-seed}
\end{table}

种子的存在解释了一件工程事实：一门语言的编译器版本不能只发布二进制，
必须能说明“这个版本被哪个版本编译”。流水线阶段本身（第 \ref{ch:cl01-pipeline} 章）
不关心这件事，但构建系统非常关心，因为它决定了一条依赖链能不能在干净机器上重放。
同一份依赖逻辑也决定了运行时为什么必须一起分发（第 \ref{ch:cl07-object} 章）。

\begin{flow}[title=一条种子链通常要走到的四步]
第一步，用别家实现把一个最小子集跑起来；第二步，用这个子集编译完整实现，
此时两个版本互为验证；第三步，把子集编译器扔掉，只保留“能编译当前版本的那个版本”；
第四步，把这条链写成可复现的构建说明，否则第三步之后没人能重新走一遍。
链条上任何一步缺记录，这门语言的实现就只剩二进制可用、无人能改。
\end{flow}

\section{载体为什么几乎都是中级语言}

上表里除了少数几个，实现本体几乎都是 \CL{} 与 \CPP{}。这不是历史偶然，而是三个约束叠加的结果。

其一，加载器与调用约定的接口就是这一层。目标文件的节、符号、重定位类型
（第 \ref{ch:cl07-object} 章）与共享对象的装载方式（第 \ref{ch:cl09-dynamic} 章），
都是按这一层的形状设计的；用这一层写实现，等于顺着平台的默认路径走，
不需要先解决“怎么把运行时塞进去”。

其二，与操作系统线程库、内存分配器、文件描述符的接触面全在这一层。
即时编译要把生成的字节变成可执行内存、装载器要处理地址重排，
这些都必须落到系统接口上，而这些接口的公约数就是本册反复提到的那一个。

其三，工具链在这一层最齐全。消毒器、剖析器、覆盖率统计、内存检查这类工具，
认的是这一层的符号与调试信息形态（附录 \ref{ap:clA-cli} 的三张表就是给它们用的）；
换一种载体语言，这些工具多半要自己重写一遍，成本立即超过换语言带来的好处。

\begin{warning}[title=别把这三条读成“应该用 \CL{} 重写”]
它们解释的是已有实现为什么长这样，不是新项目的选择依据。
一个只需解释执行、不需要即时编译、也不嵌入他人进程的实现，
完全可以把运行时压到最小；表里那一行“几乎不需要外部运行时”的例子走的就是另一条路。
\end{warning}

\section{这带来的实际后果}

\begin{guideline}[title=三条后果]
\begin{enumerate}
\item 跨语言实现的互操作边界仍然是那一层约定：不论宿主与插件各用什么语言写，
最终商定的接口形状是 C 的（附录 \ref{ap:clB-abi} 的寄存器与宽度表就是它的落地细节）。
\item 扩展、插件、嵌入宿主这三件事都靠这一层：宿主提供接口，实现方按接口注册；
谁换实现语言都不改变这个契约，改变契约的是版本与符号可见性策略
（第 \ref{ch:cl09-dynamic} 章、第 \ref{ch:cl10-cc} 章）。
\item 内存分析、性能剖析、崩溃还原这些工具也以这一层为公约数：它们读的是
符号、调试信息与段布局，不是源码语言。工具能不能用，取决于镜像里那些容器在不在，
第 \ref{ch:cl14-ide} 章与此呼应。
\end{enumerate}
\end{guideline}

\CPP{} 的例外值得单独说：一旦边界上放的是类、模板实例与异常，
修饰名与布局规则就成了实现之间的私有约定（第 \ref{ch:cl13-cxx-abi} 章）。
所以表里那些 \CPP{} 写的实现，对外暴露的接口几乎仍退回 \CL{} 形状——
这一层退化不是审美选择，而是“换一个实现方还能链接”这一条要求的直接结果。

\begin{bestpractice}[title=把这三层分别写清楚]
描述一个实现时说清三件事：主体语言、代码生成路径、对外接口形状。
它们分别决定了可维护性、性能特征与互操作范围；三者混在一起谈，
得到的结论多半只能用在某一个实现上。
\end{bestpractice}

\section{怎么核对表里的一行}

这张表可以核对，而且核对成本很低：载体语言看仓库里主体源码的构成，
代码生成路径看它交付的产物，接口形状只看镜像里那些名字。
下面按这三件事给一个最小顺序，工具与附录 \ref{ap:clA-cli} 一致。

\begin{table}[htbp]
\centering
\begin{tabularx}{\textwidth}{@{}>{\raggedright\arraybackslash}p{0.16\textwidth} X@{}}
\toprule
要核对什么 & 从哪里看 \\
\midrule
主体语言 &
仓库里的源码构成与构建入口；只看目录名会误判，因为几乎所有项目都留有一小层接口代码。 \\
代码生成路径 &
产物形态：交付的是可执行镜像、中级表示文件，还是一份带虚拟机的字节码；
再用段与符号确认它是否把代码生成交给了外部后端。 \\
接口形状 &
镜像里的符号名是否带装饰、动态符号是否只在预期的那一层出现；
这决定了它能不能被别人嵌入，也决定了表里那一列怎么写。 \\
\bottomrule
\end{tabularx}
\caption{核对一行所需的最少信息}\label{tab:apC-verify}
\end{table}

\begin{lstlisting}[style=shellstyle, caption={核对接口形状：把实现当黑盒，只看它交付的镜像}]
# 假设某实现已经构建成共享对象
file libembed.so
nm -C libembed.so            > names.txt
readelf -dW libembed.so      > dyn.txt
dumpbin /exports embed.dll   > pe-names.txt
otool -L embed.dylib         > deps.txt
\end{lstlisting}

\begin{note}[title=什么时候该怀疑这一行写错了]
两处信号最直接：一是名字里出现别家实现的修饰风格，说明接口那一层没有退干净；
二是加载期解析到的依赖比构建时声明的多，说明运行时把别的东西也带了进来。
两条都属于第 \ref{ch:cl09-dynamic} 章与第 \ref{ch:cl13-cxx-abi} 章的组合问题，
不是语言选择本身的问题。表里的归属与正文的机制对不上时，先改表，别改说法。
\end{note}


\cleardoublepage


\end{document}
